\documentclass[11pt,reqno]{amsart}

\usepackage[a4paper,margin=1.15in]{geometry}
\usepackage{lmodern}
\usepackage[T1]{fontenc}
\usepackage[utf8]{inputenc}
\usepackage{microtype}
\usepackage{amsmath,amssymb,amsthm,mathtools,stmaryrd,booktabs,enumitem}
\usepackage{xcolor}
\definecolor{linknavy}{RGB}{0,45,115}
\usepackage[colorlinks=true,linkcolor=linknavy,citecolor=linknavy,urlcolor=linknavy]{hyperref}

\linespread{1.06}
\numberwithin{equation}{section}

% Three-tier AMS theorem hierarchy
\theoremstyle{plain}
\newtheorem{theorem}{Theorem}[section]
\newtheorem{lemma}[theorem]{Lemma}
\newtheorem{proposition}[theorem]{Proposition}
\newtheorem{corollary}[theorem]{Corollary}

\theoremstyle{definition}
\newtheorem{definition}[theorem]{Definition}
\newtheorem{construction}[theorem]{Construction}
\newtheorem{example}[theorem]{Example}

\theoremstyle{remark}
\newtheorem{remark}[theorem]{Remark}
\newtheorem{notation}[theorem]{Notation}

% Canonical mathematical operators and macros
\DeclareMathOperator{\nodes}{nodes}
\DeclareMathOperator{\ls}{ls}
\DeclareMathOperator{\rs}{rs}
\DeclareMathOperator{\BT}{BT}
\DeclareMathOperator{\BeT}{BeT}
\DeclareMathOperator{\out}{out}
\DeclareMathOperator{\dom}{dom}
\DeclareMathOperator{\Pol}{Pol}
\DeclareMathOperator{\ord}{ord}

\newcommand{\CLS}{\mathrm{CL}(S)}
\newcommand{\TS}{\mathcal{T}_S}
\newcommand{\TSV}{\mathcal{T}_S(\mathcal{V})}
\newcommand{\Tbin}{\mathcal{T}_{\mathrm{bin}}}
\newcommand{\NF}{\mathrm{NF}}
\newcommand{\HNF}{\mathrm{HNF}}
\newcommand{\RNF}{\mathrm{RNF}}
\newcommand{\SN}{\mathrm{SN}}
\newcommand{\WN}{\mathrm{WN}}
\newcommand{\HN}{\mathrm{HN}}
\newcommand{\Prod}{\mathrm{Prod}}
\newcommand{\REG}{\mathrm{REG}}
\newcommand{\sem}[1]{\llbracket #1 \rrbracket}
\newcommand{\redS}{\to_S}
\newcommand{\redSp}{\to_S^+}
\newcommand{\redSs}{\to_S^*}
\newcommand{\convS}{=_S}

\title[Universality, Finite Algebras, and Asymptotics of Pure $S$]{Computational Universality, Finite Tree Algebras, and Asymptotic Phase Transitions of the Pure $S$ Combinator}

\author{Wolfram $S$ Combinator Challenge Submission}

\thanks{Archive monograph accompanied by Lean~4.33.1 kernel proofs and Wolfram Language verification package.}

\subjclass[2020]{Primary 03B40, 68Q05; Secondary 68Q42, 05A16, 68V20}

\keywords{combinatory logic, $S$ combinator, $\lambda\mathrm{I}$-calculus, universal Turing machine, cyclic tag system, finite-state tree transducer, tree automaton, B\"ohm tree, productivity, analytic combinatorics, constructive formalization}

\begin{document}

\begin{abstract}
Sch\"onfinkel's combinator $S$, defined by $S x y z \redS (x z)(y z)$, generates a non-erasing subsystem $\CLS$ of combinatory logic whose normalizing closed terms $\TS \cap \SN$ form a regular tree language (Waldmann, 2000). Universal computation in $\CLS$ therefore requires sub-universal encoding, halting recognition, and output readback along non-terminating reductions (Wolfram, 2021).

We show that universality in $\CLS$ depends sharply on evaluation-order control. Under stateless root-restarted tree-walking selection, closed $S$-terms simulate Turing machines via $U_{3,11}$ in initial size $O((|Q|^2 + |w|)^2)$ with a minimal $11$-state halting automaton $\mathcal{H}_{11}$ and linear-time output readback; from a $10$-leaf seed $D_{\mathrm{seed}}$, program-free cyclic-tag simulation compresses to a two-pass $(25, 23, 136)$ pocket tree transducer halting in at most $2|M| - 1 + \mathrm{depth}(M)$ transitions per contraction. Conversely, root-restarted selection is already universal over $\mathsf{Leaf} \to \mathsf{Node}(\mathsf{Leaf},\mathsf{Leaf})$, whereas unrestricted reduction admits constant-step halting shortcuts and forbids asymmetric branch gating.

For unguided rewriting, we establish a $\Delta L = -2$ leaf-preserving conservation law and dependency-pair non-spiralling obstructions toward Conjecture~17; construct conversion-invariant finite $S$-algebras ($35, 39, 39, 36$ states), automaton-free proofs, and productive B\"ohm faithfulness resolving Questions~5, 7, 10, 30 (on $\Prod \times \TS$), and~31 of Barendregt et al.\ (2017); and analyze the $\lambda\mathrm{I}$-translation, proving Conjecture~40, refuting Conjectures~42 and~43 at $14$ and $15$ leaves with a hereditary saturation repair, answering Question~32 negatively, and establishing single-constant collapse and congruence incomparability for Question~34. Finally, we derive all-step $\to_{\mathrm{lo}}$ cubic jump-ahead formulas, characterize Garden-of-Eden normal forms, and determine exact rational Catalan densities $c_{\HN} = 4\,429\,673/7\,077\,888$ and $(c_0, c_1, c_2, c_3) = (2\,648\,215/7\,077\,888,\, 1\,864\,553/7\,077\,888,\, 835/2\,304,\, 0)$. All results are constructively verified in Lean~4 and cross-checked in Wolfram Language.
\end{abstract}

\maketitle

% ============================================================================
\section{Introduction}\label{sec:intro}
% ============================================================================

\subsection{Historical Context and the Pure \texorpdfstring{$S$}{S} Problem}\label{subsec:intro-history}

In a lecture delivered to the G\"ottingen Mathematical Society on December~7, 1920 and published in 1924, Moses Sch\"onfinkel showed that bound variables can be eliminated from mathematical logic by reducing all operations to binary application and two primitive combinators: the cancellation operator $K$ (written $C$ by Sch\"onfinkel), satisfying $K x y \to x$, and the fusion operator $S$, satisfying $S x y z \to (x z)(y z)$ \cite{Schonfinkel1924}. Together with Haskell Curry's axiomatization \cite{Curry1930,CurryFeys1958}, the system $\mathrm{CL}(S,K)$ became one of the standard foundations of computability theory, matching the expressive power of Church's untyped $\lambda$-calculus and Turing machines \cite{Turing1937Computable,Barendregt1984,HindleySeldin2008}.

In $\mathrm{CL}(S,K)$, computational universality relies simultaneously on duplication by $S$ and erasure by $K$. Restricting the basis to $S$ alone yields the subsystem $\CLS$, whose closed (or ground) terms $\TS$ are finite full binary application trees with every leaf labeled $S$, rewritten under the contextual closure of
\begin{equation}\label{eq:s-rule-intro}
  S X Y Z \;\redS\; (X Z)(Y Z).
\end{equation}
Because every argument $X, Y, Z$ reappears on the right-hand side of \eqref{eq:s-rule-intro}, the leaf count $|M|$ never decreases under reduction: $|(X Z)(Y Z)| - |S X Y Z| = |Z| - 1 \ge 0$. By Church's conservation principle for orthogonal non-erasing rewrite systems, weak normalization ($\WN$) and strong normalization ($\SN$) coincide on $\CLS$ \cite{Church1941,Zachos1978,BarendregtEndrullisKlopWaldmann2017}. For five decades after Sch\"onfinkel's lecture, it was not even clear whether every closed $S$-term normalizes. That question was settled in the negative in the mid-1970s through Barendregt's $9$-leaf term $A A A$ (where $A \equiv S S S$), Dubou\'e's two minimal $7$-leaf non-normalizing terms $X_1 \equiv S(S S)S S S S$ and $X_2 \equiv S S S(S S)S S$ \cite{BarendregtBergstraKlopVolken1976,Barendregt1984}, and Zachos's $10$-leaf single-redex term $Z_{\mathrm{Zac}} \equiv S T T(S T T)$ (where $T \equiv S S$) \cite{Zachos1978}.

Two decades later, Johannes Waldmann proved through an algebraic analysis of pure-$S$ rewriting that the set $\TS \cap \SN$ of normalizing closed $S$-terms is a regular tree language, recognizable in linear time $O(|M|)$ by a finite tree automaton \cite{Waldmann1997,Waldmann1998RTA,Waldmann1998Thesis,Waldmann2000} whose minimal deterministic bottom-up realization has $39$ states ($38$ accepting states and one absorbing sink) \cite{EndrullisDeVrijerWaldmann2010,BarendregtEndrullisKlopWaldmann2017}; see also Cheilaris, Ramirez, and Zachos \cite{CheilarisRamirezZachos2011} for a manual grammar derivation, and Barendregt, Dekkers, and Statman \cite{BarendregtDekkersStatman2013} for broader typed and combinatory decision boundaries. Waldmann's theorem immediately rules out direct normal-form universality in the classical sense of Davis \cite{Davis1957DefUTM}: if halting of a universal machine $U$ on input $w$ were encoded as normal-form reachability $\mathcal{E}(w) \in \WN$, any computable encoder $\mathcal{E}$ would compose with the $39$-state tree automaton to decide the halting problem.

Yet normalizing reductions represent an exponentially vanishing fraction of $\CLS$ dynamics \cite{BendkowskiGrygielZaionc2017}, and non-normalizing $S$-terms display unbounded growth and tree-unfolding behavior \cite{Wolfram2021CombinatorsBook}. In their survey \emph{``Dance of the Starlings''}, Barendregt, Endrullis, Klop, and Waldmann \cite{BarendregtEndrullisKlopWaldmann2017} formulated a program of open problems and conjectures governing unguided $\CLS$ rewriting, B\"ohm-tree semantics, and the canonical translation of $\CLS$ into Church's $\lambda\mathrm{I}$-calculus \cite{Church1941}. On June~8, 2021, commemorating the centenary of Sch\"onfinkel's lecture \cite{Wolfram2020StoryOfComputation,Wolfram2021CombinatorsBook}, Stephen Wolfram posed \emph{The Wolfram $S$ Combinator Challenge}, asking whether $\CLS$ can support universal computation along non-terminating evolutions \cite{Wolfram2021Challenge}. Following the ongoing-computation framework established by Priese \cite{Priese1979ComplexityUTM} and employed in the universality proofs for elementary cellular automaton Rule~110 by Cook \cite{Cook2004Rule110} and for the $(2,3)$ Turing machine prize by Smith \cite{Smith2020TM23,Wolfram2007PrizeWon}, Wolfram formulated ongoing pure-$S$ computation via an encoder $\mathcal{E}$, a halting detector $\mathcal{H}$, and an output decoder $\mathcal{D}$, each required to operate within strict sub-universal resource bounds so that no universal simulation occurs outside $\CLS$ itself \cite{Wolfram2021Challenge,vanEmdeBoas1990Models}. Wolfram also observed that once reductions do not terminate in a unique normal form, confluence no longer collapses distinct evaluation orders: computational behavior depends on whether one considers fixed canonical strategies (such as normal-order reduction $\to_{\mathrm{lo}}$), unrestricted multiway conversion $\convS$, or a specified single-step selection discipline.

\subsection{Computation Regimes and Summary of Main Results}\label{subsec:intro-results}

The central thesis of this monograph is that the computational power and semantic decidability of $\CLS$ admit an exact classification governed by the locus of evaluation-order control and binder discipline. Instead of treating ``universality of $S$'' as an undifferentiated yes-or-no predicate, we analyze all six operational and translational regimes posed in \cite{Wolfram2021Challenge,BarendregtEndrullisKlopWaldmann2017} within a single mathematical and mechanized framework. Table~\ref{tab:regimes} summarizes the resulting taxonomy.

\begin{table}[htbp]
\centering
\caption{Taxonomy of computational regimes and decidability boundaries in $\CLS$.}\label{tab:regimes}
\footnotesize
\setlength{\tabcolsep}{3.2pt}
\begin{tabular}{@{}lllll@{}}
\toprule
Regime & Step relation & Halting / event problem & Expressive / algebraic law & Reference \\
\midrule
Normal-form halting & $M \redSs N \in \NF$ & Regular ($39$-state $\mathcal{A}_{\SN}$) & Subcritical ($\gamma_{\SN} \approx 3.36992$); $\Delta L = -2$ & Thms.~\ref{thm:leaf-preserving}, \ref{thm:waldmann-sn}, \ref{thm:asymptotics} \\
Canonical fixed order & $\to_h, \to_{\mathrm{lo}}, \to_{\mathrm{li}}$ & $\HN$ ($35$), $\Prod$ ($39$) regular & Cubic $\to_{\mathrm{lo}}$ jump; monotone clones & Thms.~\ref{thm:question-10}, \ref{thm:q5-main}, \ref{thm:clones-rule110}, \ref{thm:all-step-jump-ahead} \\
Unrestricted multiway & $M \redSs N$, $M \convS N$ & $\redSs$ decidable; $\convS$ on $\Prod \times \TS$ & $7$/$5$-step shortcuts; no branch gate & Prop.~\ref{prop:multiway-shortcuts}, Thms.~\ref{thm:no-church-bool}, \ref{thm:productive-bohm-faithful} \\
Root-restarted walker & $T_{k+1} = \sigma_{\mathrm{TM}}(T_k)$ & $\Sigma_1^0$-complete ($11$-state $\mathcal{H}_{11}$) & $O((|Q|^2{+}|w|)^2)$; sep.\ by $\to_{\mathrm{sp}}$ & Thms.~\ref{thm:tm-universality}, \ref{thm:leaf-split-univ} \\
Program-free walker & $T_{k+1} = \sigma_{\mathrm{CTS}}(T_k)$ & $\Sigma_1^0$-complete ($11$-state $\mathcal{H}_{11}$) & $(25,23,136)$ $\mathcal{T}_{\mathrm{pocket}}$; $D_{\mathrm{seed}}$ genesis & Prop.~\ref{prop:seed-genesis}, Thms.~\ref{thm:cts-universality}, \ref{thm:pocket-transducer-universality} \\
$\lambda\mathrm{I}$-translation & $(M)_\lambda, \to_\beta^*, =_\beta$ & Conj.~40 true; 42, 43, Q32 false & Single-constant Q34; $\infty_{34}$ sink & Thms.~\ref{thm:conjecture-40}, \ref{thm:conjectures-42-43-refuted}, \ref{thm:question-32} \\
\bottomrule
\end{tabular}
\end{table}

Our contributions span five interconnected directions.

Sections~\ref{sec:tm-sim} and~\ref{sec:cts-sim} establish compact, polynomial-size universal simulations in $\CLS$ under stateless root-restarted finite-state selection. In an instructive recent development, Filin \cite{Filin2026} demonstrated selected-path universality of $\CLS$ under a root-restarted finite controller by compiling Boolean-tape Turing machines through three-counter and $2$-tag systems \cite{CockeMinsky1964Tag,Minsky1967Computation} to Rogozhin's $(4,6)$ universal Turing machine \cite{Rogozhin1996SmallUTM} and Cook's $912$-phase cyclic tag system \cite{Cook2004Rule110}; because unary counter encoding incurs iterated exponential blowup, every source instance in that padded pipeline produces at least $8^{19069} = 2^{57207}$ seed bits before $114$-fold tag expansion, and source output recovery requires an $O(|T|^4)$ parser consulting the retained seed. In Section~\ref{sec:tm-sim}, we bypass tag systems altogether by compiling boundary-safe normalized binary Turing machines directly to Neary and Woods's quadratic-time $3$-state, $11$-symbol universal machine $U_{3,11}$ \cite{NearyWoods2006SmallFast,Neary2008Thesis} and its $132$-state, $262$-rule binary specialization $U_{\mathrm{bin}}$. Theorem~\ref{thm:tm-universality} establishes that the resulting pure-$S$ compiler $\mathcal{E}_{\mathrm{TM}}(M,w)$ produces a closed $S$-term of quartic leaf size $|T_0| \le 125{,}000\,N(M,w)^2$ with $N(M,w) = 225|Q|^2 + 390|Q| + 4|w| + 173$, reducing the compiled size of a $1$-state machine from at least $2^{57207}$ seed bits to $1{,}996{,}921$ leaves; moreover, every macro-step of the zipper tree-walking selector $\sigma_{\mathrm{TM}}$ terminates in at most $3500(2|t|-1)$ cursor micro-ticks with no state retained between contractions, halting is recognized by a minimal $11$-state deterministic bottom-up tree automaton $\mathcal{H}_{11}$ (Proposition~\ref{prop:halt-automaton}), and the bare halted term decodes to the exact output word $\out_M(w)$ in linear time $O(|T_k|)$. In Section~\ref{sec:cts-sim}, we remove the embedded universal transition table from the selector itself: Theorem~\ref{thm:cts-universality} constructs a program-free selector $\sigma_{\mathrm{CTS}}$ over self-describing pure-$S$ cyclic-tag encodings $\mathcal{E}_{\mathrm{CTS}}(P,w)$ of quadratic leaf size at most $2500(\|(P,w)\|+1)^2$, recognized at halting by the identical $11$-state automaton $\mathcal{H}_{11}$, and Proposition~\ref{prop:seed-genesis} proves that every instance $\mathcal{E}_{\mathrm{CTS}}(P,w)$ is reachable via pure-$S$ reduction from a single fixed $10$-leaf combinator $D_{\mathrm{seed}} \equiv S(S(S S))(S(S S))(S S)S$. Theorem~\ref{thm:pocket-transducer-universality} further compresses $\sigma_{\mathrm{CTS}}$ to an explicit two-pass contextual finite-state pocket tree transducer $\mathcal{T}_{\mathrm{pocket}}$ with $25$ static structural roles, $23$ bottom-up states per node, and $136$ top-down pointer states, containing zero universal Turing or tag transition tables and executing each pure-$S$ contraction in at most $2|M| - 1 + \mathrm{depth}(M)$ elementary transitions.

Section~\ref{sec:separation} characterizes why external redex selection is necessary for Theorems~\ref{thm:tm-universality}, \ref{thm:cts-universality}, and~\ref{thm:pocket-transducer-universality}, sharpening the automata-theoretic perspective of Boja\'nczyk and Colcombet \cite{BojanczykColcombet2008} and Kutrib and Meyer's deterministic tree-walking-storage automata \cite{KutribMeyer2023,KutribMeyer2026}. Theorem~\ref{thm:leaf-split-univ} proves a controller-complexity separation: a fixed root-restarted finite-state tree-walking selector coupled with a $3$-state bottom-up halting tree automaton is already Turing-complete over unlabeled binary trees governed solely by the unary leaf-splitting rule $\mathsf{Leaf} \to_{\mathrm{sp}} \mathsf{Node}(\mathsf{Leaf},\mathsf{Leaf})$, which performs neither subtree duplication nor arity-$3$ interaction. Within controlled rewriting, the distinctive contribution of $S$ is therefore in-tree subtree duplication, which powers the program-free self-description and $10$-leaf seed genesis of Section~\ref{sec:cts-sim}. Conversely, once the external selector is removed in favor of unrestricted pure-$S$ reduction $\redSs$, Proposition~\ref{prop:multiway-shortcuts} exhibits explicit $7$-step and $5$-step multiway reductions that trigger spurious halting on arbitrary non-terminating payloads in both protected-frame counter hubs and tag-dispatcher encodings, while Theorem~\ref{thm:no-church-bool} proves unconditionally that no closed $S$-term can realize asymmetric Church-Boolean branch gating.

Sections~\ref{sec:prelim}, \ref{sec:algebras}, and~\ref{sec:fixed-cellular} develop the conservation laws, finite algebraic decidability theory, and fixed-order expressiveness of unguided $\CLS$ rewriting, resolving Questions~5, 7, 10, 30 (on $\Prod \times \TS$), and~31 alongside structural obstructions toward Conjecture~17 of Barendregt, Endrullis, Klop, and Waldmann \cite{BarendregtEndrullisKlopWaldmann2017}. In Section~\ref{sec:prelim}, Theorem~\ref{thm:leaf-preserving} shows that the left-subtree leaf sum $L(M)$ decreases by the exact constant $2$ on every leaf-preserving step $S X Y S \redS (X S)(Y S)$, making leaf-preserving reduction confluent, strongly normalizing, and Jordan--H\"older ranked with sharp upper bound $\lfloor(|M|-1)(|M|-2)/4\rfloor$ (exponentially sharpening Bergstra's $O(2^{|M|})$ bound), accompanied by constructive proofs of closed non-spiralling and dependency-pair obstructions toward open Conjecture~17 showing that $M \redSp C[M^\sigma]$ requires a variable $x$ occurring in $M$ whose image $\sigma(x)$ has rightmost leaf $x$ and $\rs(\sigma(x)) \ge 1$ under unbounded right-spine growth, complemented by the exact $8$-leaf open non-$\SN$ threshold and parametric $(4n+3)$-step near-miss (Proposition~\ref{prop:open-spiral-divergence-threshold} and Theorems~\ref{thm:conj17-dp-weighted}, \ref{thm:conj17-unconditional-and-coupling}, and~\ref{thm:conj17-minmaxrs-monotonicity}). In Section~\ref{sec:algebras}, eliminating the additive offset table ($B \equiv 0$) via a pointwise minimal multiplier matrix $A \colon Q_{\HN}^2 \to \{0,\dots,30\}$ (Proposition~\ref{prop:hna-minimality}) yields a pure-multiplier potential proving completeness of the $35$-state head-normalization $S$-algebra $\mathcal{A}_{\HN}$ and linear-time decidability of root normal forms, resolving Question~10 and Footnote~4 of \cite{BarendregtEndrullisKlopWaldmann2017} (Theorem~\ref{thm:question-10}), alongside constructive top-termination and absence of Berarducci mute $S$-terms \cite{BerarducciBohm1993,Berarducci1996} (Theorem~\ref{thm:waldmann-top-term}). Synthesizing a minimal $39$-state productivity $S$-algebra $\mathcal{A}_P$ via rank-$10$ critical-pair stratification over the Myhill--Nerode minimal $113$-state reachable subalgebra $\mathrm{PH} \subseteq Q_P \times Q_{\HN}$ (Proposition~\ref{prop:ph-113-minimal}) resolves Question~5 in the affirmative (Theorem~\ref{thm:q5-main}). Through the automaton-free three-constructor grammar $\mathcal{H} ::= A \mid S\mathcal{H} \mid \mathcal{H}\mathcal{H}$ and $\mathcal{U} = \mathcal{H}\mathcal{H}\mathcal{H}$, we give self-contained proofs of Question~7---including infinite $\SN/\overline{\SN}$ refinement alternation and the atomic $S \mapsto S S$ leaf-splitting law ($T_n \in \SN \iff n \not\equiv 9 \pmod{10}$; Theorem~\ref{thm:question-7})---and Question~31 separating $\{A A A S^k\}_{k \ge 0}$ and $\mathcal{A}_{36}$ despite identical clocked B\"ohm trees $\bot$ (Theorem~\ref{thm:question-31}), while Theorem~\ref{thm:productive-bohm-faithful} and Proposition~\ref{prop:q30-decidability} establish a sharp dichotomy for Dana Scott's Question~30 \cite{Scott1975}: B\"ohm-tree equality is faithful and constructively decidable for $\convS$ across $\Prod \times \TS$ ($\SN \subsetneq \Prod$, while general convertibility on $(\TS \setminus \Prod)^2$ remains open), whereas even a single $\bot$ at any depth destroys clocked-B\"ohm faithfulness. Theorem~\ref{thm:order-hierarchy} characterizes functional orders $\{0,1,2,3\}$ on open and closed $S$-terms via minimal $S$-algebras of $35$, $12$, and $3$ states and unique right-comb closed order-$3$ terms (clarifying the prose example $S S S$ versus $S(S S S)S$ in \cite[p.~90]{BarendregtEndrullisKlopWaldmann2017}). Section~\ref{sec:fixed-cellular} establishes deterministic strategy coincidence and quadratic trace languages on $Z_{\mathrm{Zac}}$ (Proposition~\ref{prop:pure-line-and-aaa}), a $5$-step innermost Case-$3$ ceiling (Theorem~\ref{thm:case3-ceiling}), and Post-clone monotonicity barriers \cite{Post1941} versus a $4$-state algebra $\mathcal{D}_4$ realizing an $11$-leaf normal-form combinator for Rule~110 (Theorem~\ref{thm:clones-rule110}).

Section~\ref{sec:lambda-translation} resolves the open problems and conjectures governing the canonical translation $(M)_\lambda$ of pure-$S$ terms into Church's non-erasing $\lambda\mathrm{I}$-calculus \cite{Church1941}, where $(S)_\lambda \equiv \lambda x y z.\, x z (y z)$. Establishing constructive weak-head standardization and Church's conservation theorem ($\WN_\beta \iff \SN_\beta$ on $\lambda\mathrm{I}$), Theorem~\ref{thm:conjecture-40} proves Conjecture~40 of \cite{BarendregtEndrullisKlopWaldmann2017} in the affirmative: $\ord_{\mathrm{CL}}(M) = \ord_\lambda((M)_\lambda)$ holds for every combinatory term over $\{S, K, I\}$ and across $\CLS$. Normalization, however, diverges across $(-)_\lambda$ once $\beta$-redexes form beneath generated $\lambda$-binders. Improving Waldmann's $13$-leaf example, Proposition~\ref{prop:minimal-lambda-divergers} identifies the global-minimum $9$-leaf closed $\SN$ term $W_9 \equiv S(S(A A))S$ and $10$-leaf normal form $N_{10} \equiv S(S(S A A))S$ with $(W_9)_\lambda, (N_{10})_\lambda \notin \SN_\beta$, and Theorem~\ref{thm:conjectures-42-43-refuted} constructs the global-minimum $14$-leaf closed term $M_{14} \equiv S A (S W_9)$ and $15$-leaf normal form $M_{15} \equiv S A (S N_{10})$ refuting both Conjectures~42 and~43 simultaneously, while proving the valid converse implication and establishing the repaired equivalence $\mathrm{DeepNorm}(M) \iff (M)_\lambda \in \SN_\beta$ via hereditary fresh-variable saturation in $\CLS$. For Question~34 (decidability and regularity of $L_{34} = \{M \in \TS : (M)_\lambda \in \SN_\beta\}$), Theorem~\ref{thm:question-34-structure} reduces $\lambda\mathrm{I}$ strong normalization to single-constant combinatory saturation $\mathrm{DeepNorm}_1$ in left-linear $\mathrm{CL}(S, x_0)$, proves that universal $\lambda\mathrm{I}$ context non-erasure collapses all $\lambda\beta$-diverging closed $S$-terms into a single absorbing Myhill--Nerode sink $\infty_{34}$, demonstrates mutual congruence incomparability between $\equiv_{34}$ and $\equiv_{\SN}$ with four-way normal-form splitting inside $q_6$ (splitting even the shared joint fiber $(q_6, h_9, p_{18})$ of $\mathcal{A}_{\SN} \times \mathcal{A}_{\HN} \times \mathcal{A}_P$), and isolates finiteness of $\TS / {\equiv_{34}}$ and decidability of $L_{34}$ as the sharpened single-constant open problem. Finally, Proposition~\ref{prop:q32-rigidity-waldmann-pair} and Theorem~\ref{thm:question-32} resolve Question~32 in the negative: despite small-term $\lambda\beta$-rigidity and uniqueness of Waldmann's $(8,11)$ closed $\beta$-convertible pair up to $11$ leaves, no finite set of $\lambda\beta$-sound pure-$S$ equations axiomatizes $=_\beta$ on $\CLS$.

Finally, Section~\ref{sec:asymptotics} classifies all $14$ empirical pure-$S$ benchmarks studied in Wolfram's challenge announcement \cite{Wolfram2021Challenge} (Theorem~\ref{thm:wolfram-14-dichotomy}), derives exact all-step $\to_{\mathrm{lo}}$ state evolutions and cubic leaf-count closed forms for $Z_{\mathrm{Zac}}$, $W_2$ (together with its $=_S$-inconvertible $+14$-step twin $W_2^{\mathrm{twin}}$), $D_{\mathrm{seed}}$, and $A A A$ (Theorem~\ref{thm:all-step-jump-ahead}), characterizes Garden-of-Eden (isolated) normal forms syntactically with a Mahler-type counting recurrence and sub-ternary growth rate $\gamma_{\mathrm{Iso}} \approx 2.9076221057 < 3$ (Theorem~\ref{thm:goe-characterization-and-mahler}), establishes small-term Catalan minimality thresholds across $1 \le |M| \le 11$ (Theorem~\ref{thm:small-term-minimality}), and applies Flajolet--Sedgewick singularity analysis \cite{FlajoletSedgewick2009} to the polynomial generating-function systems of $\mathcal{A}_{\HN}$, $\mathcal{A}_{\SN}$, and $\mathcal{A}_P$, advancing the quantitative combinatory logic program of Bendkowski, Grygiel, and Zaionc \cite{BendkowskiGrygielZaionc2017} (Theorem~\ref{thm:asymptotics}). Because the unsolvable sink $\infty_{\HN}$ is a left ideal ($q \odot_{\HN} \infty_{\HN}$ need not equal $\infty_{\HN}$, as witnessed by $S(A A A) \in \HNF$), Puiseux expansion at the Catalan singularity $z = 1/4$ yields the exact dyadic-ternary rational head-normalizing density and its partition across functional orders $\{0, 1, 2, 3\}$:
\[
  c_{\HN} \;=\; \lim_{n \to \infty} \frac{\#\HN(n)}{C_{n-1}} \;=\; \frac{4\,429\,673}{7\,077\,888},
  \qquad
  (c_0, c_1, c_2, c_3) \;=\; \Bigl(\frac{2\,648\,215}{7\,077\,888},\, \frac{1\,864\,553}{7\,077\,888},\, \frac{835}{2\,304},\, 0\Bigr),
\]
satisfying $c_0 + c_1 + c_2 + c_3 = 1$ and $c_1 + c_2 = c_{\HN} \approx 0.6258467215$, whereas the two-sided context sinks of $\mathcal{A}_{\SN}$ and $\mathcal{A}_P$ push their dominant singularities strictly past $1/4$ to distinct degree-$35$ algebraic branch points $\rho_{\SN} > \rho_P > 1/4$, forcing exponential rarity of both normalizing and productive closed $S$-terms with certified all-$n$ upper bounds proportional to $(250/259)^n$.

\subsection{Technical Overview and Proof Architecture}\label{subsec:intro-overview}

Six structural mechanisms underlie the proofs across Sections~\ref{sec:prelim}--\ref{sec:asymptotics}.

Protected spine-$2$ storage frames $F(a,b,j) \equiv S a(S b j)$, building on local opening identities of Filin \cite{Filin2026}, isolate reductions inside their left slot $a$ and right slot $b$ while sequestering historical junk $j$ without erasure: because $\ls(F(a,b,j)) = 2$, the term $F(a,b,j)$ has no root redex. Unopened leaves of the addressable memory layout are represented by phase-indexed counter generators $G(m,n) \equiv \nu_m \nu_n S$, which open deterministically in $6$ or $7$ contractions into $F(G(m',n'), G(m',n'), j')$ with non-frame intermediate stages and non-$\mathrm{S}_2$ junk $j'$. Wrapping an append-only chain of $11$-bit cell records along the right spine of $F$-frames allows the stateless zipper selector $\sigma_{\mathrm{TM}}$ to execute $U_{\mathrm{bin}}$ via a two-pass priority scan (completing any pending neighbor move before initiating a new active transition) while reading only $6$ local cursor observations and resetting all control state at each contraction. For cyclic tag systems, factoring local structural recognition into $25$ static node roles, a $23$-state bottom-up summary pass, and a $136$-state top-down pointer pass compresses $\sigma_{\mathrm{CTS}}$ into the two-pass pocket tree transducer $\mathcal{T}_{\mathrm{pocket}}$ running in $2|M| - 1 + \mathrm{depth}(M)$ transitions without any embedded universal rule table.

In-tree self-description stands in sharp contrast to unguided multiway collapse and spine conservation. By encoding cyclic-tag appendant templates with $3$-bit versioned handshakes ($\mathsf{idle} \to \mathsf{armed} \to \mathsf{appending} \to \mathsf{idle}$) paired with staged queue allocation, $\sigma_{\mathrm{CTS}}$ copies arbitrarily long appendants onto the queue tail using pure-$S$ duplication alone, and every initial tree $\mathcal{E}_{\mathrm{CTS}}(P,w)$ unfolds directly from $D_{\mathrm{seed}}$ via a $5$-step bootstrap to $(\nu_1 S)\,G(2,2)$ followed by selective slot openings. Removing the selector destroys this discipline immediately: under unrestricted $\redSs$, any unopened halt-flag generator inside a protected hub opens in $7$ steps regardless of payload divergence, and any $D \cdot Q$ queue dispatcher contracts its live $1$-cell prematurely to a tombstone in $1$ step and executes the $0$-branch in $4$ further steps. More fundamentally, an inductive spine-grammar invariant ($\mathcal{W}_0$ and $\mathcal{O}$ in Section~\ref{subsec:no-church-bool}) shows that in any reduction of $M x_0 x_1 U_1 \cdots U_m$, every head $S$-redex capable of bringing $x_1$ into active function position must carry $x_0$ as one of its first two arguments, activating $x_0$ first. On size-preserving reductions $S X Y S \redS (X S)(Y S)$, the left-subtree leaf sum $L(M)$ drops by $2$ at each step, while on open terms dependency-pair right-tail tracking forces any hypothetical open spiral $M \redSp C[M^\sigma]$ to fix a variable $x$ as the rightmost leaf of $\sigma(x)$ along an unbounded right spine.

Pure-multiplier algebraic potentials and three-constructor syntactic grammars eliminate non-constructive appeals to K\"onig's lemma and large table searches in $\CLS$ solvability and convertibility proofs. For the $35$-state head-normalization algebra $\mathcal{A}_{\HN}$, setting $\mu_{\HN}(S) = 1$ and $\mu_{\HN}(M N) = \mu_{\HN}(M) + A(\sem{M}_{\HN}, \sem{N}_{\HN})\,\mu_{\HN}(N)$ with a pointwise minimal integer multiplier table $A \colon Q_{\HN}^2 \to \{0,\dots,30\}$ ($142$ entries exceeding $1$, and zero additive offset table $B \equiv 0$) yields strict decrease on every active head contraction outside $\infty_{\HN}$, constructively bounding head-reduction lengths. Independently of tree automata, closing $A \equiv S S S$ under $S \mathcal{H}$ and $\mathcal{H}\mathcal{H}$ isolates a self-propagating unsolvable locus $\mathcal{U} = \mathcal{H}\mathcal{H}\mathcal{H}$ that yields concise book proofs for Questions~7 and~31.

Canonical productive values and critical-pair stratification resolve Question~5 and govern B\"ohm-tree faithfulness on $\Prod \times \TS$. On closed $S$-terms, left application preserves $\bot$ at the same B\"ohm depth while right application increases depth by at most $1$ because $\ord_{\mathrm{CL}}(M) \le 3$ prevents an argument from escaping a head spine without root consumption. Jointly analyzing $\mathcal{A}_P \times \mathcal{A}_{\SN}$ shows that $34$ of the $38$ accepting states of $\mathcal{A}_P$ consist solely of strongly normalizing terms, three further states reduce to $\HNF$ without root contractions, and $(q_{29}, \infty_{\SN})$ is the unique non-normalizing productive root-redex fiber, inhabited by the Zachos family $G_m G_n$ whose right applications $(G_m G_n) U$ always collapse to $\bot$. Compressing the $39^3$ transition cube via $35$ two-dimensional overrides over the $113$-state Myhill--Nerode minimal reachable subalgebra $\mathrm{PH} \subseteq Q_P \times Q_{\HN}$ yields a rank-$10$ stratification constructively witnessing $\bot_d(M)$ whenever $\sem{M}_P = \infty_P$, while root-normal-form conversion decomposition along $\bot$-free spines makes B\"ohm-tree equality faithful and decidable for $\convS$ on $\Prod \times \TS$.

Constructive weak-head standardization in Church's $\lambda\mathrm{I}$-calculus \cite{Church1941}, combined with hereditary combinatory saturation, governs the translation $(M)_\lambda$. Because weak and strong $\beta$-normalization coincide in $\lambda\mathrm{I}$, head steps of $(M)_\lambda$ track fresh-variable combinatory head reductions closely enough to preserve functional order across $\{S, K, I\}$ (Conjecture~40). Once head normal form is reached, however, $\beta$-reduction continues beneath generated $\lambda$-binders: in $M_{14} \equiv S A (S W_9)$, the $9$-leaf $\SN$ subterm $W_9 \equiv S(S(A A))S$ sits protected inside a spine-$2$ argument even after $M_{14}$ receives three fresh variables, yet $(W_9)_\lambda$ diverges inside a $\lambda$-body of $(M_{14})_\lambda$, refuting Conjectures~42 and~43. Enforcing fresh-variable normalizability hereditarily on every generated subterm defines the pure-$\CLS$ predicate $\mathrm{DeepNorm}(M)$ characterizing $(M)_\lambda \in \SN_\beta$, while universal variable retention in $\lambda\mathrm{I}$ contexts collapses all $\lambda\beta$-diverging closed $S$-terms into the single Myhill--Nerode sink $\infty_{34}$ and reduces Question~34 to single-constant saturation $\mathrm{DeepNorm}_1$ in left-linear $\mathrm{CL}(S, x_0)$. Conversely, bounding redex creation depth under any finite pure-$S$ axiom set leaves infinitely many $\lambda\beta$-identities unprovable, resolving Question~32.

Finally, macro-cycle recurrences, predecessor-free decompositions, and strongly connected component (SCC) singularity analysis govern the quantitative dynamics of $\CLS$. Indexing outer spine wrappers along $\to_{\mathrm{lo}}$ yields closed-form all-step state and cubic leaf formulas, while characterizing absence of $\redS$-predecessors yields the Mahler recurrence for Garden-of-Eden normal forms. On $\mathcal{A}_{\HN}$ (and its open-term extension classifying orders $\{0,1,2,3\}$), evaluating the quadratic generating-function system at $z = 1/4$ produces exact rationals $L_q(1/4) \in \mathbb{Q}$ and reduces the Puiseux coefficients $b_q$ to a nonsingular rational linear system, whereas eliminating transient states of $\mathcal{A}_{\SN}$ and $\mathcal{A}_P$ isolates recurrent cores analytic strictly past $z = 1/4$, certified both by degree-$35$ resultant polynomials and by positive rational upper-envelope evaluations at $z = 37/125$ formalized directly over $\mathbb{N}$.

\subsection{Constructive Mechanization and Organization}\label{subsec:intro-mechanization}

Every definition, construction, finite algebraic check, and theorem in this monograph has been constructively formalized and kernel-verified in the Lean~4 theorem prover (version~4.33.1). Throughout the development, all proofs rely solely on Lean's standard propositional extensionality and quotient soundness principles, avoiding classical choice, excluded middle on undecidable predicates, and external reflection or native-code evaluation oracles. In addition, all compilers, tree-walking selectors, pocket tree transducers, finite $S$-algebras, $\lambda\mathrm{I}$-translation analyzers, halting detectors, and output decoders are implemented in a standalone Wolfram Language package cross-checked on concrete Turing machine, Busy Beaver, cyclic tag, $\lambda\mathrm{I}$-divergence, and small-term census benchmarks.

The remainder of the monograph is organized into nine sections and an appendix. Section~\ref{sec:prelim} establishes the rewriting syntax, confluence, leaf-preserving $\Delta L = -2$ conservation, constructive closed non-spiralling, dependency-pair open-spiral obstructions toward Conjecture~17, B\"ohm trees, finite $S$-algebras, and ongoing-simulation preliminaries. Section~\ref{sec:tm-sim} presents the $U_{3,11}$ compiler, protected spine-$2$ encoding, stateless zipper tree-walking selector, minimal $11$-state halting automaton, and linear-time bare-term output decoder. Section~\ref{sec:cts-sim} develops the program-free cyclic tag simulator, its genesis from the $10$-leaf seed $D_{\mathrm{seed}}$, and its compression to the two-pass $(25, 23, 136)$ pocket tree transducer $\mathcal{T}_{\mathrm{pocket}}$. Section~\ref{sec:separation} proves the leaf-splitting controller-separation theorem, the $7$-step and $5$-step multiway halting shortcuts, and the impossibility of asymmetric branch gating. Section~\ref{sec:algebras} constructs the finite $S$-algebras $\mathcal{A}_{\SN}, \mathcal{A}_{\HN}, \mathcal{A}_P, \mathcal{A}_{36}$, establishes pure-multiplier head potentials, $113$-state $\mathrm{PH}$ minimality, automaton-free book proofs, productive B\"ohm faithfulness, and the functional order $\{0,1,2,3\}$ minimal-algebra hierarchy, resolving Questions~5, 7, 10, 30 (on $\Prod \times \TS$), and~31 of \cite{BarendregtEndrullisKlopWaldmann2017} together with constructive top-termination. Section~\ref{sec:lambda-translation} develops the constructive $\lambda\mathrm{I}$-translation theory, proving Conjecture~40, refuting Conjectures~42 and~43 with $14$- and $15$-leaf minimal counterexamples and a hereditary saturation repair, establishing single-constant collapse and congruence incomparability for Question~34, and resolving Question~32 in the negative. Section~\ref{sec:fixed-cellular} analyzes deterministic strategy coincidence, quadratic trace languages, innermost right-step bounds, and Post-clone barriers versus minimal normal-form Rule~110. Section~\ref{sec:asymptotics} classifies Wolfram's $14$ benchmark terms, proves the all-step $\to_{\mathrm{lo}}$ cubic jump-ahead formulas, characterizes Garden-of-Eden normal forms, establishes small-term minimality thresholds up to $11$ leaves, and proves the Catalan asymptotic phase-transition and exact order-density theorems. Section~\ref{sec:formalization} discusses the constructive formalization architecture, and Appendix~\ref{app:concordance} provides a complete concordance mapping every numbered result to its machine-checked declaration.

% ============================================================================
\section{Preliminaries: Rewriting Semantics, Leaf-Preserving Conservation, Non-Spiralling, and Open-Spiral Obstructions}\label{sec:prelim}
% ============================================================================

\subsection{Syntax, Spine Geometry, and Leaf-Preserving Conservation Laws}\label{subsec:prelim-syntax}

We begin by fixing the syntax, structural measures, substitutions, and reduction relations of $\CLS$, following the conventions of Barendregt \cite{Barendregt1984}, Waldmann \cite{Waldmann2000}, and Barendregt, Endrullis, Klop, and Waldmann \cite{BarendregtEndrullisKlopWaldmann2017}.

\begin{definition}[Syntax, spine measures, and reduction in $\CLS$]\label{def:syntax}
Let $\mathcal{V} = \{x_0, x_1, x_2, \dots\}$ be a countably infinite set of variables. The set $\TSV$ of \emph{open $S$-terms} is generated inductively by the grammar
\[
  M, N \;::=\; S \;\mid\; x \;\mid\; (M\,N) \qquad (x \in \mathcal{V}),
\]
and the subset $\TS \subset \TSV$ of \emph{closed} (or \emph{ground}) \emph{$S$-terms}---also called \emph{pure-$S$ combinators}---consists of all variable-free terms generated by $M, N ::= S \mid (M\,N)$. Binary application associates to the left, so $M N P \equiv ((M N) P)$, and $M \equiv N$ denotes literal syntactic identity of finite binary trees. Throughout the paper we write the standard closed abbreviations
\[
  T \;\equiv\; S S, \qquad B \;\equiv\; S(S S) \;\equiv\; S T, \qquad A \;\equiv\; S S S \;\equiv\; T S, \qquad G_0 \;\equiv\; S T T, \qquad V \;\equiv\; S G_0,
\]
together with $G_{k+1} \equiv T G_k$ ($k \ge 0$), Zachos's $10$-leaf combinator $Z_{\mathrm{Zac}} \equiv G_0 G_0$, and the $10$-leaf seed $D_{\mathrm{seed}} \equiv S B B T S$. For each $n \ge 1$, the closed \emph{left comb} $L_n \in \TS$ and \emph{right comb} $R_n \in \TS$ of leaf size $n$ are defined recursively by $L_1 \equiv R_1 \equiv S$, $L_{n+1} \equiv L_n S$, and $R_{n+1} \equiv S R_n$.

For atoms $a \in \{S\} \cup \mathcal{V}$ and terms $M, N \in \TSV$, the \emph{leaf size} $|M| \ge 1$, binary-tree \emph{node count} $\nodes(M)$, \emph{left-spine length} $\ls(M)$, \emph{right-spine length} $\rs(M)$, \emph{head atom} $\mathrm{hd}(M) \in \{S\} \cup \mathcal{V}$, and \emph{rightmost atom} $\mathrm{rt}(M) \in \{S\} \cup \mathcal{V}$ are defined recursively by
\begin{align*}
  |a| &= 1, & |M N| &= |M| + |N|, & \nodes(M) &= 2|M| - 1, \\
  \ls(a) &= 0, & \ls(M N) &= 1 + \ls(M), & \rs(a) &= 0,\quad \rs(M N) = 1 + \rs(N), \\
  \mathrm{hd}(a) &\equiv a, & \mathrm{hd}(M N) &\equiv \mathrm{hd}(M), & \mathrm{rt}(a) &\equiv a,\quad \mathrm{rt}(M N) \equiv \mathrm{rt}(N).
\end{align*}
We write $\mathrm{Vars}(M) \subset \mathcal{V}$ for the finite set of variables occurring in $M$, and $\#_x(M) \in \mathbb{N}$ for the number of leaf occurrences of $x \in \mathcal{V}$ in $M$. Given a substitution $\sigma \colon \mathcal{V} \to \TSV$, the simultaneous substitution instance $M^\sigma \in \TSV$ is defined by $S^\sigma \equiv S$, $x^\sigma \equiv \sigma(x)$, and $(M N)^\sigma \equiv M^\sigma N^\sigma$.

A \emph{one-hole context} is generated by $C[\,] ::= \square \mid (C[\,]\,N) \mid (M\,C[\,])$, with $C[M]$ denoting replacement of the unique hole $\square$ by $M$, context composition written $(C_1 \circ C_2)[M] \equiv C_1[C_2[M]]$, and substitution extended to contexts via $\square^\sigma \equiv \square$, $(C[\,]\,N)^\sigma \equiv (C^\sigma[\,]\,N^\sigma)$, and $(M\,C[\,])^\sigma \equiv (M^\sigma\,C^\sigma[\,])$, so that $(C[M])^\sigma \equiv C^\sigma[M^\sigma]$. The \emph{left-branching depth} $d_{\mathrm{L}}(C) \in \mathbb{N}$ and variable count $\#_x(C) \in \mathbb{N}$ of a context are given by $d_{\mathrm{L}}(\square) = 0$, $d_{\mathrm{L}}(C[\,]\,N) = 1 + d_{\mathrm{L}}(C)$, $d_{\mathrm{L}}(M\,C[\,]) = d_{\mathrm{L}}(C)$, and $\#_x(C[M]) = \#_x(C) + \#_x(M)$. We write $U \trianglelefteq M$ (respectively $U \triangleleft M$) if $M \equiv C[U]$ for some context $C[\,]$ (respectively some $C[\,] \not\equiv \square$).

An \emph{$S$-redex} is any term of the form $S X Y Z$, with \emph{contractum} $(X Z)(Y Z)$. \emph{One-step reduction} $\redS$ is the contextual closure of root contraction: $C[S X Y Z] \redS C[(X Z)(Y Z)]$. We write $\to_r$ for contraction at the root ($C[\,] \equiv \square$), $\to_{\mathrm{nr}}$ for a non-root step ($C[\,] \not\equiv \square$), $\redSp$ for the transitive closure of $\redS$, $\redSs$ for its reflexive-transitive closure, $\redS^k$ for $k$-fold composition, and $\convS$ for the equivalence relation (convertibility) generated by $\redS$. A step $M \redS N$ is \emph{leaf-preserving}, written $M \to_{S,\mathrm{lp}} N$, if $|M| = |N|$. A term $M \in \TSV$ is in \emph{normal form} ($M \in \NF$) if there is no $N$ with $M \redS N$; \emph{weakly normalizing} ($M \in \WN$) if $M \redSs N$ for some $N \in \NF$; and \emph{strongly normalizing} ($M \in \SN$) if $M$ is constructively well-founded (accessible) with respect to $\redS$, meaning every reduction sequence starting from $M$ terminates.
\end{definition}

Note that the number of closed $S$-terms of leaf size $n \ge 1$ is the Catalan number $C_{n-1} = \frac{1}{n}\binom{2n-2}{n-1}$, since closed $S$-terms of leaf size $n$ are in bijection with full binary trees having $n$ leaf nodes and $n-1$ internal application nodes.

\begin{proposition}[Confluence, monotonicity, and $\WN = \SN$]\label{prop:cr-conservation}
The rewrite system $(\TSV, \redS)$ satisfies the following properties:
\begin{enumerate}[label=(\roman*),leftmargin=2.2em]
  \item \emph{Confluence (Church--Rosser):} If $M \redSs N_1$ and $M \redSs N_2$, there exists $P \in \TSV$ such that $N_1 \redSs P$ and $N_2 \redSs P$. Consequently, $M \convS N$ holds if and only if $M$ and $N$ share a common reduct, and normal forms are unique up to $\equiv$.
  \item \emph{Leaf balance, spine monotonicity, and leaf-preserving redex criterion:} Every contextual contraction $M \equiv C[S X Y Z] \redS C[(X Z)(Y Z)] \equiv N$ satisfies $|N| + 1 = |M| + |Z|$, $\rs(M) \le \rs(N)$, and $\mathrm{rt}(N) \equiv \mathrm{rt}(M)$. In particular, $|M| \le |N|$, with equality $|M| = |N|$ holding if and only if $|Z| = 1$ (that is, $Z \in \{S\} \cup \mathcal{V}$, and uniquely $Z \equiv S$ when $M \in \TS$).
  \item \emph{Conservation ($\WN = \SN$):} If $M \redS N$ and $N \in \SN$, then $M \in \SN$. In particular, $\WN = \SN$ on $\TSV$.
\end{enumerate}
\end{proposition}

\begin{proof}
For~(i), $\CLS$ is an orthogonal left-linear combinatory rewrite system \cite{Terese2003}. Explicitly, adapting Takahashi's parallel reduction method \cite{Takahashi1995}, define $M \Rightarrow N$ inductively by $S \Rightarrow S$, $x \Rightarrow x$, and $M_1 M_2 \Rightarrow N_1 N_2$ whenever $M_i \Rightarrow N_i$, together with $S X Y Z \Rightarrow (X' Z')(Y' Z')$ whenever $X \Rightarrow X'$, $Y \Rightarrow Y'$, and $Z \Rightarrow Z'$. Define the complete development map $\Delta \colon \TSV \to \TSV$ by $\Delta(S X Y Z) \equiv (\Delta(X)\,\Delta(Z))(\Delta(Y)\,\Delta(Z))$ on root redexes, $\Delta(M_1 M_2) \equiv \Delta(M_1)\,\Delta(M_2)$ on non-redex applications, and $\Delta(a) \equiv a$ on atoms $a \in \{S\} \cup \mathcal{V}$. Structural induction on $M \Rightarrow N$ yields the triangle property $N \Rightarrow \Delta(M)$: in the sole overlap check where $M \equiv (S X Y) Z$ steps via $M_1 M_2 \Rightarrow N_1 N_2$, inversion on $S X Y \Rightarrow N_1$ forces $N_1 \equiv S X' Y'$ with $X' \Rightarrow \Delta(X)$ and $Y' \Rightarrow \Delta(Y)$, whence $N_1 N_2 \equiv S X' Y' Z' \Rightarrow (\Delta(X)\,\Delta(Z))(\Delta(Y)\,\Delta(Z)) \equiv \Delta(M)$. The diamond property for $\Rightarrow$ follows, and since $\redS \;\subseteq\; \Rightarrow \;\subseteq\; \redSs$, diagram chasing yields Church--Rosser confluence.

For~(ii), let $|C|$ denote the sum of leaf sizes of all side subterms of $C[\,]$, so $|C[U]| = |C| + |U|$. Direct calculation gives
\[
  |C[(X Z)(Y Z)]| + 1 \;=\; |C| + |X| + |Y| + 2|Z| + 1 \;=\; \bigl(|C| + 1 + |X| + |Y| + |Z|\bigr) + |Z| \;=\; |C[S X Y Z]| + |Z|.
\]
Since $|Z| \ge 1$, it follows that $|M| \le |N|$, with $|M| = |N|$ if and only if $|Z| = 1$, which restricts $Z$ to an atom in $\{S\} \cup \mathcal{V}$ (and to $S$ when $M \in \TS$). At a root step, $\rs((X Z)(Y Z)) = 2 + \rs(Z) > 1 + \rs(Z) = \rs(S X Y Z)$ and $\mathrm{rt}((X Z)(Y Z)) \equiv \mathrm{rt}(Z) \equiv \mathrm{rt}(S X Y Z)$; induction over $C[\,]$ preserves $\rs(M) \le \rs(N)$ and $\mathrm{rt}(N) \equiv \mathrm{rt}(M)$ because left-child steps leave right spines untouched and right-child steps act on the right spine directly.

For~(iii), enrich parallel reduction with an explicit contraction weight $M \Rightarrow_w N$ ($w \in \mathbb{N}$): atoms have weight $0$, application satisfies $M_1 M_2 \Rightarrow_{w_1 + w_2} N_1 N_2$, and redex contraction satisfies $S X Y Z \Rightarrow_{1 + w_X + w_Y + 2 w_Z} (X' Z')(Y' Z')$. Note that $M \redS N$ implies $M \Rightarrow_1 N$. Structural induction on $M \Rightarrow_w N$ establishes the projection alternative: for every step $M \redS M_1$, either there exists $w_1 < w$ with $M_1 \Rightarrow_{w_1} N$ (when $M \redS M_1$ contracts a redex already contracted in $M \Rightarrow_w N$, using $(w_X + w_Z) + (w_Y + w_Z) < 1 + w_X + w_Y + 2 w_Z$), or there exist $N_1$ and $w_1 \in \mathbb{N}$ such that $N \redSp N_1$ and $M_1 \Rightarrow_{w_1} N_1$ (when $M \redS M_1$ contracts inside $Z \redS Z_1$ of a redex $S X Y Z$, both copies of $Z'$ in $(X' Z')(Y' Z')$ step to $Z_1'$, yielding at least two positive steps $N \redSp N_1$ because $Z$ is never erased). Well-founded induction on $N \in \SN$ (lifted to $\redSp$) with an inner induction on $w$ proves $M \in \SN$. Induction along $M \redSs N \in \NF \subseteq \SN$ yields $\WN \subseteq \SN$, and $\SN \subseteq \WN$ holds constructively by contracting the leftmost redex until none remains.
\end{proof}

Leaf-preserving reductions $M \to_{S,\mathrm{lp}} N$ govern every constant-size transient segment of $\CLS$ dynamics. While Bergstra's classical exponential weight $\|P Q\|_{\mathrm{B}} = 2\|P\|_{\mathrm{B}} + \|Q\|_{\mathrm{B}}$ \cite{KlopVanOostromVanRaamsdonk2007,BarendregtEndrullisKlopWaldmann2017} drops by $6 \cdot 2^{d_{\mathrm{L}}(C)}$ depending on context depth and yields only an $O(2^{|M|})$ step bound, we establish an exact depth-independent conservation law that makes $(\TSV, \to_{S,\mathrm{lp}})$ a confluent Jordan--H\"older ranked system of sharp quadratic length $\Theta(|M|^2)$.

\begin{theorem}[Ranked leaf-preserving conservation law and sharp quarter-square bound]\label{thm:leaf-preserving}
For atoms $a \in \{S\} \cup \mathcal{V}$ and terms $P, Q \in \TSV$, define Bergstra's weight $\|\cdot\|_{\mathrm{B}} \colon \TSV \to \mathbb{N}$ and the \emph{left-subtree leaf sum} $L \colon \TSV \to \mathbb{N}$ recursively by
\[
  \|a\|_{\mathrm{B}} = 1, \quad \|P Q\|_{\mathrm{B}} = 2\|P\|_{\mathrm{B}} + \|Q\|_{\mathrm{B}}, \qquad\qquad
  L(a) = 0, \quad L(P Q) = L(P) + L(Q) + |P|.
\]
Then:
\begin{enumerate}[label=(\roman*),leftmargin=2.2em]
  \item \emph{Bergstra weight envelope and contextual drop:} Every $M \in \TSV$ satisfies $2|M| - 1 \le \|M\|_{\mathrm{B}} \le 2^{|M|} - 1$; on closed terms $M \in \TS$, equality $\|M\|_{\mathrm{B}} = 2|M| - 1$ holds if and only if $M \equiv R_{|M|}$, and $\|M\|_{\mathrm{B}} = 2^{|M|} - 1$ holds if and only if $M \equiv L_{|M|}$. Every leaf-preserving step $C[S X Y Z] \to_{S,\mathrm{lp}} C[(X Z)(Y Z)]$ satisfies $\|C[S X Y Z]\|_{\mathrm{B}} = \|C[(X Z)(Y Z)]\|_{\mathrm{B}} + 6 \cdot 2^{d_{\mathrm{L}}(C)}$, whence any $k$-step reduction $M \to_{S,\mathrm{lp}}^k N$ obeys $6k + 2|M| \le 2^{|M|}$.
  \item \emph{Exact ranked conservation, confluence, and Jordan--H\"older property:} Every leaf-preserving step $M \to_{S,\mathrm{lp}} N$ at arbitrary depth in $\TSV$ decreases $L$ by the exact constant~$2$:
  \begin{equation}\label{eq:left-subtree-drop}
    L(M) \;=\; L(N) + 2.
  \end{equation}
  Moreover, $\to_{S,\mathrm{lp}}$ is strongly normalizing and globally confluent on $\TSV$. Consequently, every $M \in \TSV$ admits a unique $\to_{S,\mathrm{lp}}$-normal form $\mathrm{nf}_{\mathrm{lp}}(M)$, and every maximal leaf-preserving reduction starting from $M$ has the identical length $\bigl(L(M) - L(\mathrm{nf}_{\mathrm{lp}}(M))\bigr)/2$.
  \item \emph{Quadratic upper bound and quarter-square attainment:} Every $M \in \TSV$ satisfies $|M| - 1 \le L(M)$ and $2 L(M) + |M| \le |M|^2$. Consequently, every $k$-step leaf-preserving reduction $M \to_{S,\mathrm{lp}}^k N$ satisfies
  \begin{equation}\label{eq:lp-quadratic-bound}
    k \;\le\; \left\lfloor \frac{(|M|-1)(|M|-2)}{4} \right\rfloor.
  \end{equation}
  Conversely, define $N_2^{\mathrm{lp}} \equiv S S$, $N_3^{\mathrm{lp}} \equiv S S S$, and $N_{m+4}^{\mathrm{lp}} \equiv T\,N_{m+2}^{\mathrm{lp}}$ for $m \ge 0$. For every $n \ge 2$, the closed left comb $L_n$ satisfies $L_n \to_{S,\mathrm{lp}}^{\lfloor(n-2)^2/4\rfloor} N_n^{\mathrm{lp}}$ with $N_n^{\mathrm{lp}}$ in $\to_{S,\mathrm{lp}}$-normal form, matching \eqref{eq:lp-quadratic-bound} to within $\lfloor(n-1)/4\rfloor$ steps.
\end{enumerate}
\end{theorem}

\begin{proof}
For~(i), structural induction on $M$ establishes $2|M| \le \|M\|_{\mathrm{B}} + 1 \le 2^{|M|}$: at an application $P Q$, using $\|P\|_{\mathrm{B}} \ge 1$ gives $2|P| \le \|P\|_{\mathrm{B}} + 1 \le 2\|P\|_{\mathrm{B}}$ and hence $2(|P| + |Q|) \le (2\|P\|_{\mathrm{B}} + \|Q\|_{\mathrm{B}}) + 1$, with equality on closed terms forcing $\|P\|_{\mathrm{B}} = 1$ ($P \equiv S$) and $Q \equiv R_{|Q|}$, so $P Q \equiv R_{|P Q|}$. For the upper bound, $(2\|P\|_{\mathrm{B}} + \|Q\|_{\mathrm{B}} + 1) + 2 = 2(\|P\|_{\mathrm{B}} + 1) + (\|Q\|_{\mathrm{B}} + 1) \le 2 \cdot 2^{|P|} + 2^{|Q|} \le 2^{|P|+|Q|} + 2$ because $2 \cdot 2^j + 2^k \le 2^{j+k} + 2$ for all $j, k \ge 1$, with equality if and only if $k = 1$; on closed terms, equality forces $|Q| = 1$ ($Q \equiv S$) and $P \equiv L_{|P|}$, so $P Q \equiv L_{|P Q|}$. At a leaf-preserving root contraction $S X Y Z \to_r (X Z)(Y Z)$, Proposition~\ref{prop:cr-conservation}(ii) gives $|Z| = 1$ and $\|Z\|_{\mathrm{B}} = 1$, so
\[
  \|S X Y Z\|_{\mathrm{B}} \;=\; 4\|X\|_{\mathrm{B}} + 2\|Y\|_{\mathrm{B}} + 9 \;=\; \bigl(4\|X\|_{\mathrm{B}} + 2\|Y\|_{\mathrm{B}} + 3\bigr) + 6 \;=\; \|(X Z)(Y Z)\|_{\mathrm{B}} + 6.
\]
Inducting on $C[\,]$, a drop of $\delta$ at the hole scales by $2$ across $(C[\,]\,R)$ and by $1$ across $(P\,C[\,])$, producing $6 \cdot 2^{d_{\mathrm{L}}(C)} \ge 6$. Summing over $k$ steps gives $6k + 2|M| - 1 \le 6k + \|N\|_{\mathrm{B}} \le \|M\|_{\mathrm{B}} \le 2^{|M|} - 1$.

For~(ii), let $M \to_{S,\mathrm{lp}} N$. At a root contraction $S X Y Z \to_r (X Z)(Y Z)$ with $|Z| = 1$, the atom $Z$ satisfies $L(Z) = 0$. Expanding $L$ on both sides yields
\begin{align*}
  L(S X Y Z) &\;=\; L(S X Y) + L(Z) + |S X Y| \;=\; \bigl(L(S X) + L(Y) + 1 + |X|\bigr) + 0 + (1 + |X| + |Y|) \\
             &\;=\; L(X) + L(Y) + 2|X| + |Y| + 3, \\
  L\bigl((X Z)(Y Z)\bigr) &\;=\; \bigl(L(X) + 0 + |X|\bigr) + \bigl(L(Y) + 0 + |Y|\bigr) + (|X| + 1) \;=\; L(X) + L(Y) + 2|X| + |Y| + 1,
\end{align*}
proving \eqref{eq:left-subtree-drop} at the root. For a left-child step $P Q \to_{S,\mathrm{lp}} P' Q$, leaf preservation $|P| + |Q| = |P'| + |Q|$ gives $|P| = |P'|$, so $L(P Q) = L(P) + L(Q) + |P| = (L(P') + 2) + L(Q) + |P'| = L(P' Q) + 2$; for a right-child step $P Q \to_{S,\mathrm{lp}} P Q'$, leaf preservation $|Q| = |Q'|$ gives $L(P Q) = L(P) + (L(Q') + 2) + |P| = L(P Q') + 2$. Strong normalization of $\to_{S,\mathrm{lp}}$ follows immediately by bounded induction on $L(M) \in \mathbb{N}$. For confluence, define leaf-preserving parallel reduction $P \Rightarrow_{\mathrm{lp}} Q$ by reflexivity on atoms, congruence on applications, and $S X Y Z \Rightarrow_{\mathrm{lp}} (X' Z)(Y' Z)$ whenever $X \Rightarrow_{\mathrm{lp}} X'$, $Y \Rightarrow_{\mathrm{lp}} Y'$, and $|Z| = 1$ (note that $|Z| = 1$ forces $Z \Rightarrow_{\mathrm{lp}} Z'$ to satisfy $Z' \equiv Z$). Define the conditional complete development $\Delta_{\mathrm{lp}}(S X Y Z) \equiv (\Delta_{\mathrm{lp}}(X)\,Z)(\Delta_{\mathrm{lp}}(Y)\,Z)$ if $|Z| = 1$, and $\Delta_{\mathrm{lp}}(P Q) \equiv \Delta_{\mathrm{lp}}(P)\,\Delta_{\mathrm{lp}}(Q)$ otherwise, with $\Delta_{\mathrm{lp}}(a) \equiv a$ on atoms. Because $|Q| = |P|$ holds along $P \Rightarrow_{\mathrm{lp}} Q$, the side condition $|Z| = 1$ is invariant, giving the triangle property $Q \Rightarrow_{\mathrm{lp}} \Delta_{\mathrm{lp}}(P)$ and hence global confluence of $\to_{S,\mathrm{lp}}^*$. If $M \to_{S,\mathrm{lp}}^{k_1} N_1$ and $M \to_{S,\mathrm{lp}}^{k_2} N_2$ reach $\to_{S,\mathrm{lp}}$-normal forms, confluence forces $N_1 \equiv N_2$, and iterating \eqref{eq:left-subtree-drop} gives $L(M) = L(N_1) + 2 k_1 = L(N_1) + 2 k_2$, whence $k_1 = k_2$.

For~(iii), structural induction on $M$ gives $|M| \le L(M) + 1$ (since $|P| + |Q| \le |P| + L(Q) + 1 \le L(P) + L(Q) + |P| + 1 = L(P Q) + 1$) and $2 L(M) + |M| \le |M|^2$, using $|Q| \ge 1$ to bound
\[
  2\bigl(L(P) + L(Q) + |P|\bigr) + |P| + |Q| \;=\; \bigl(2 L(P) + |P|\bigr) + \bigl(2 L(Q) + |Q|\bigr) + 2|P| \;\le\; |P|^2 + |Q|^2 + 2|P|\,|Q| \;=\; (|P| + |Q|)^2.
\]
Along $M \to_{S,\mathrm{lp}}^k N$, we have $|N| = |M|$ and $L(M) = L(N) + 2k \ge |M| - 1 + 2k$. Substituting into $2 L(M) + |M| \le |M|^2$ yields $4k + 3|M| - 2 \le |M|^2$ for $|M| \ge 2$ (and $k = 0$ for $|M| = 1$), which rearranges to $4k \le (|M| - 1)(|M| - 2)$ and gives \eqref{eq:lp-quadratic-bound}. Conversely, each $N_n^{\mathrm{lp}}$ has no $\to_{S,\mathrm{lp}}$-step: this is vacuous for $n \in \{2, 3\}$, and for $N_{m+4}^{\mathrm{lp}} \equiv (S S)\,N_{m+2}^{\mathrm{lp}}$, the head has spine length~$1$, the left child $S S$ is normal, and any $\to_{S,\mathrm{lp}}$-step in the right child contradicts the induction hypothesis on $N_{m+2}^{\mathrm{lp}}$. Finally, we verify by induction on $m \ge 0$ the peeling reduction $L_{m+4} \to_{S,\mathrm{lp}}^{m+1} T\,L_{m+2}$: at $m = 0$, $L_4 \equiv S S S S \to_{S,\mathrm{lp}} (S S)(S S) \equiv T\,L_2$ in $1$ step; assuming $L_{m+4} \to_{S,\mathrm{lp}}^{m+1} T\,L_{m+2}$, appending $S$ on the right gives $L_{m+5} \equiv L_{m+4} S \to_{S,\mathrm{lp}}^{m+1} (T\,L_{m+2}) S \equiv S S L_{m+2} S \to_{S,\mathrm{lp}} (S S)(L_{m+2} S) \equiv T\,L_{m+3}$ in $m + 2$ steps. Inducting on $n \ge 2$ with step-$2$ stride from $L_2 \equiv N_2^{\mathrm{lp}}$ ($0$ steps) and $L_3 \equiv N_3^{\mathrm{lp}}$ ($0$ steps), $L_{m+4} \to_{S,\mathrm{lp}}^{m+1} T\,L_{m+2} \to_{S,\mathrm{lp}}^{\lfloor m^2/4 \rfloor} T\,N_{m+2}^{\mathrm{lp}} \equiv N_{m+4}^{\mathrm{lp}}$ takes $(m + 1) + \lfloor m^2/4 \rfloor = \lfloor (m^2 + 4m + 4)/4 \rfloor = \lfloor ((m+4)-2)^2/4 \rfloor$ steps.
\end{proof}

\subsection{Acyclicity, Non-Spiralling, Flat \texorpdfstring{$\{S,K\}$}{\{S,K\}} Normalization, and Functional Order}\label{subsec:prelim-order}

Although $\CLS$ admits infinite reductions, its non-erasing binary geometry forbids both cyclic reductions and self-embedding spirals, forces unbounded leaf and redex growth along every infinite path, extends to strong normalization of Klop's flat $\{S, K\}$-terms, and bounds the functional arity of every pure-$S$ term by~$3$. We formalize the classical theory of Section~5.4 of Barendregt, Endrullis, Klop, and Waldmann \cite[pp.~86--90]{BarendregtEndrullisKlopWaldmann2017} in constructive form.

\begin{proposition}[Acyclicity, non-spiralling, and functional order at most $3$]\label{prop:acyclicity-order}
Let $M, N \in \TSV$.
\begin{enumerate}[label=(\roman*),leftmargin=2.2em]
  \item \emph{Bergstra acyclicity and decidability of $\redSs$ \cite{KlopVanOostromVanRaamsdonk2007,BarendregtEndrullisKlopWaldmann2017}:} No $M \in \TSV$ satisfies $M \redSp M$, and pointwise reachability $M \redSs N$ is constructively decidable for all $M, N \in \TSV$.
  \item \emph{Waldmann's no-spiral theorem and absence of reducing fixpoint combinators \cite{Waldmann2000,BarendregtEndrullisKlopWaldmann2017}:} For every one-hole context $C[\,]$ and every $M \in \TSV$, it is impossible that $M \redSp C[M]$; if $C[\,] \not\equiv \square$, then $M \redSs C[M]$ does not hold. In particular, no $Y, X \in \TSV$ satisfy $Y X \redSs X(Y X)$.
  \item \emph{Visser--van Oostrom order-$3$ theorem \cite{BarendregtEndrullisKlopWaldmann2017}:} For all $M \in \TSV$ and variables $x_0, x_1, x_2, x_3 \in \mathcal{V}$, there do not exist $P, Q \in \TSV$ such that $M x_0 x_1 x_2 \redSs S P Q$. Consequently, every $M \in \TSV$ has functional order $\ord(M) \le 3$: whenever $M x_0 x_1 x_2 x_3 \redSs N$, there exists $N' \in \TSV$ such that $N \equiv N' x_3$ and $M x_0 x_1 x_2 \redSs N'$.
\end{enumerate}
\end{proposition}

\begin{proof}
For~(i), we record two independent constructive proofs. First, along any candidate path $M \redSs N$, Proposition~\ref{prop:cr-conservation}(ii) requires $|M'| \le |N|$ at every intermediate term $M'$, while each step $M_i \redS M_{i+1}$ strictly decreases the lexicographic pair $(|N| - |M_i|,\, \|M_i\|_{\mathrm{B}})$ (or $(|N| - |M_i|,\, L(M_i))$) by Theorem~\ref{thm:leaf-preserving}; this both forbids $M \redSp M$ and bounds the finite search tree for deciding $M \redSs N$. Second, independently of any weight function, induct structurally on $M$ to refute $M \redSp M$: atoms cannot step, and if $A B \redSp A B$, either every step is non-root---decomposing into $A \redSp A$ or $B \redSp B$ and contradicting the induction hypothesis---or $A B \redSs S X Y Z$ and $(X Z)(Y Z) \redSs A B$, whence right-spine monotonicity (Proposition~\ref{prop:cr-conservation}(ii)) yields the impossible chain $\rs(A B) \le 1 + \rs(Z) < 2 + \rs(Z) = \rs((X Z)(Y Z)) \le \rs(A B)$.

For~(ii), let $\rho(Q) = 1 + \rs(Q)$ and consider application-labelled terms over $\{S\} \cup \mathcal{V}$ with binary operators $\{\cdot_k\}_{k \ge 1}$, mapping each $P \in \TSV$ to $\overline{P}$ by labelling every subterm $U V$ as $\overline{U} \cdot_{\rho(V)} \overline{V}$. For any label cap $m \ge 1$, define Waldmann's bounded labelled rewrite system $\mathcal{R}_m$ on terms with labels in $\{1, \dots, m\}$ as the union of the contextual closures of
\[
  ((S \cdot_a X) \cdot_b Y) \cdot_c Z \;\to_{(5.2)}\; (X \cdot_c Z) \cdot_{1+c} (Y \cdot_c Z) \quad (1 + c \le m), \qquad
  X \cdot_n Y \;\to_{(5.3)}\; X \cdot_{n+1} Y \quad (n + 1 \le m).
\]
Whenever $P \redS Q$ and all labels in $\overline{P}, \overline{Q}$ are at most $m$, structural induction on $P \redS Q$ shows that $\overline{P} \to_{(5.2)} K \to_{(5.3)}^* \overline{Q}$ for an intermediate labelled term $K$ (at a right-child step $U V \redS U V'$, Proposition~\ref{prop:cr-conservation}(ii) gives $\rho(V) \le \rho(V') \le m$, absorbed by $\rho(V') - \rho(V)$ applications of $\to_{(5.3)}$). Inverting every label $k$ to $d = m - k \in \{0, \dots, m\}$ embeds $\mathcal{R}_m$ together with direct subterm projections $(U \cdot_d V \succ_{\mathrm{RB}} U, V)$ into a single relation $\succ_{\mathrm{RB}}$ in which $(5.2)$ replaces a root operator $\cdot_d$ by $\cdot_{d'}$ ($d' < d$) applied to two $\cdot_d$-children built from proper subterms, and $(5.3)$ lowers $d$ to $d' < d$. Nested constructive accessibility induction on $d \in \mathbb{N}$ and on the accessibility of $U$ and $V$ proves that the transitive closure $\succ_{\mathrm{RB}}^+$ is well-founded (and hence irreflexive) for every $m \in \mathbb{N}$; in particular, $\mathcal{R}_m$ is constructively strongly normalizing. If $M \redSp C[M]$, choosing $m$ above all canonical labels occurring along that finite reduction yields $\overline{M} \to_{\mathcal{R}_m}^+ \overline{C[M]}$, which composes with hole projections along $C[\,]$ to give a cycle in $\succ_{\mathrm{RB}}^+$, a contradiction. If $C[\,] \not\equiv \square$, then $|M| < |C[M]|$ excludes $M \equiv C[M]$ as well; specializing $M \equiv Y X$ and $C[\square] \equiv X\,\square$ refutes $Y X \redSs X(Y X)$ \cite[Remark~15]{BarendregtEndrullisKlopWaldmann2017}.

For~(iii), let $\mathcal{W}_1 = \{P \in \TSV : \mathrm{rt}(P) \in \mathcal{V}\}$, $\mathcal{W}_2 = \{P Q : P, Q \in \mathcal{W}_1\}$, and $\mathcal{W}_3 = \{L R : L \in \mathcal{W}_2,\, R \in \mathcal{W}_1\}$. By Proposition~\ref{prop:cr-conservation}(ii), $\mathrm{rt}$ is invariant under $\redS$, so $\mathcal{W}_1$ is closed under $\redS$. Closure of $\mathcal{W}_2$ under $\redS$ follows at a root step $(S X Y) Z \to_r (X Z)(Y Z)$ from $Z \in \mathcal{W}_1$, which places both $X Z$ and $Y Z$ in $\mathcal{W}_1$. Likewise, if $((S X) Y) Z \in \mathcal{W}_3$, then $(S X) Y \in \mathcal{W}_2$ gives $S X \in \mathcal{W}_1$ and hence $X \in \mathcal{W}_1$; together with $Z \in \mathcal{W}_1$, this yields $X Z \in \mathcal{W}_2$ and $Y Z \in \mathcal{W}_1$, so $(X Z)(Y Z) \in \mathcal{W}_3$, while contextual steps preserve $\mathcal{W}_3$ because $\mathcal{W}_2$ contains only applications. Now $M x_0 x_1 x_2 \in \mathcal{W}_3$ because $M x_0, x_1, x_2 \in \mathcal{W}_1$, whereas $S P Q \equiv (S P) Q \notin \mathcal{W}_3$ because $S \notin \mathcal{W}_1$ implies $S P \notin \mathcal{W}_2$. Thus $M x_0 x_1 x_2$ never reduces to $S P Q$, so in any reduction of $(M x_0 x_1 x_2)\,x_3$ neither a root step nor a step inside $x_3$ can occur, forcing every reduct to have the form $N' x_3$ with $M x_0 x_1 x_2 \redSs N'$.
\end{proof}

\begin{proposition}[Unbounded leaf and redex growth, and Klop's flat $\{S,K\}$ theorem]\label{prop:remark16-klop18}
The quantitative potential and bounded-label methods of Theorem~\ref{thm:leaf-preserving} and Proposition~\ref{prop:acyclicity-order} yield three further consequences \cite[Remark~16 and Theorem~18, p.~89]{BarendregtEndrullisKlopWaldmann2017}:
\begin{enumerate}[label=(\roman*),leftmargin=2.2em]
  \item \emph{Constructive unbounded leaf count:} Let $M_0 \redS M_1 \redS M_2 \redS \cdots$ be an infinite reduction in $\TSV$. Define $N_0 = \|M_0\|_{\mathrm{B}}$ and $N_{b+1} = N_b + \|M_{N_b}\|_{\mathrm{B}}$ for $b \in \mathbb{N}$ (or $N'_0 = \lfloor L(M_0)/2 \rfloor + 1$ and $N'_{b+1} = N'_b + \lfloor L(M_{N'_b})/2 \rfloor + 1$). Then $|M_{N_b}| > b$ (and $|M_{N'_b}| > b$) holds for every $b \in \mathbb{N}$.
  \item \emph{Unbounded contracted redex size:} In every infinite reduction $M_0 \redS M_1 \redS M_2 \redS \cdots$ in $\TSV$ contracting the redex sequence $R_0, R_1, R_2, \dots$, neither the canonical $\mathcal{R}_m$-labels of $M_n$ nor the leaf sizes $|R_n|$ are bounded above by any constant $B \in \mathbb{N}$.
  \item \emph{Klop's strong normalization theorem for flat $\{S,K\}$-terms:} Let $(\mathcal{T}_{S,K}, \to_{S,K})$ be closed $\{S, K\}$-combinatory logic with rules $S X Y Z \to_{S,K} (X Z)(Y Z)$ and $K X Y \to_{S,K} X$. Define the class $\mathcal{F}_{S,K} \subset \mathcal{T}_{S,K}$ of \emph{flat $\{S,K\}$-terms} inductively by $S, K \in \mathcal{F}_{S,K}$ and $P S, P K \in \mathcal{F}_{S,K}$ for $P \in \mathcal{F}_{S,K}$, and define Klop's superclass $\mathcal{SP} \supseteq \mathcal{F}_{S,K}$ inductively by $S, K \in \mathcal{SP}$ together with $P S, P K, K K P, K S P, S K P, S S P \in \mathcal{SP}$ whenever $P \in \mathcal{SP}$. Then $\mathcal{SP}$ is closed under $\to_{S,K}$, and every term in $\mathcal{SP}$ (hence every flat $\{S,K\}$-term) is constructively strongly normalizing under $\to_{S,K}$.
\end{enumerate}
\end{proposition}

\begin{proof}
For~(i), we verify by induction on $w \ge 1$ that whenever $\|M_k\|_{\mathrm{B}} \le w$, one has $|M_k| < |M_{k+w}|$: if $|M_k| = |M_{k+1}|$, Theorem~\ref{thm:leaf-preserving}(i) gives $\|M_{k+1}\|_{\mathrm{B}} < \|M_k\|_{\mathrm{B}} \le w$, so $\|M_{k+1}\|_{\mathrm{B}} \le w - 1$ (impossible when $w = 1$ since $\|M_{k+1}\|_{\mathrm{B}} \ge 1$, and handled by the induction hypothesis at $k + 1$ when $w \ge 2$), whereas if $|M_k| < |M_{k+1}|$, leaf monotonicity gives $|M_k| < |M_{k+1}| \le |M_{k+w}|$. Induction on $b \in \mathbb{N}$ starting from $|M_0| \ge 1$ then yields $|M_{N_b}| > b$ (and identically for $N'_b$ via \eqref{eq:left-subtree-drop}).

For~(ii), if all canonical labels of $\overline{M_n}$ were bounded by $m$, Proposition~\ref{prop:acyclicity-order}(ii) would map $(M_n)_{n \ge 0}$ to an infinite $\mathcal{R}_m$-reduction, contradicting strong normalization of $\mathcal{R}_m$. Suppose next that $|R_n| \le B$ holds for all $n \ge 0$, and set $B' = \max(1, B)$. Call a labelled version $\widehat{M}$ of an unlabelled term $M \in \TSV$ \emph{validly lower-labelled} up to $B'$ if its label erasure is $M$ and every application node $U \cdot_k V$ in $\widehat{M}$ satisfies $1 \le k \le \min(\rho(V_{\mathrm{erase}}), B')$. Initializing $\widehat{M}_0$ by assigning label~$1$ to every application node of $M_0$ satisfies this invariant since $1 \le \rho(V_{\mathrm{erase}})$ and $1 \le B'$. Given $\widehat{M}_n$ and a step $M_n \redS M_{n+1}$ contracting $R_n \equiv S X Y Z$ with $|R_n| \le B \le B'$, let $\cdot_c$ be the outer label at the contracted redex in $\widehat{M}_n$. By validity, $c \le \rho(Z) = 1 + \rs(Z) \le |Z|$; since $1 + |Z| \le (1 + |X| + |Y|) + |Z| = |R_n| \le B'$, we have $1 + c \le B'$. Contracting that node via rule $(5.2)$ alone---without any $(5.3)$ lifts---produces a labelled term $\widehat{M}_{n+1}$ erasing to $M_{n+1}$ with all labels at most $B'$, preserving $1 \le k \le \rho(V'_{\mathrm{erase}})$ at every ancestor of a right-child step because $\rho(V_{\mathrm{erase}}) \le \rho(V'_{\mathrm{erase}})$. This yields an infinite $\to_{(5.2)}$-chain in $\mathcal{R}_{B'}$, contradicting strong normalization of $\mathcal{R}_{B'}$.

For~(iii), let $\pi \colon \mathcal{T}_{S,K} \to \TS$ replace every leaf $K$ by $S$, preserving total leaf count $|\pi(P)| = |P|_{S,K}$. We check by structural induction on $P \to_{S,K} Q$ with $P \in \mathcal{SP}$ that $Q \in \mathcal{SP}$ and either $|Q|_{S,K} < |P|_{S,K}$ or $(|Q|_{S,K} = |P|_{S,K} \text{ and } \pi(P) \to_{S,\mathrm{lp}} \pi(Q))$. At a root $K$-contraction $(K X) Y \to_{S,K} X$, inversion on $(K X) Y \in \mathcal{SP}$ forces $X \in \{S, K\} \subset \mathcal{SP}$ and $|X|_{S,K} = 1 < 2 + |Y|_{S,K} = |(K X) Y|_{S,K}$. At a root $S$-contraction $((S X) Y) Z \to_{S,K} (X Z)(Y Z)$, none of the four prefix constructors $K K M, K S M, S K M, S S M$ matches $((S X) Y) Z$ (whose left-left grandchild is $S X$ rather than an atom), so $Z \in \{S, K\}$ and $(S X) Y \in \mathcal{SP}$; inverting $(S X) Y \in \mathcal{SP}$ in turn forces $X \in \{S, K\}$ and $Y \in \mathcal{SP}$. Consequently, $Y Z \in \mathcal{SP}$ and $X Z \in \{K K, K S, S K, S S\}$, placing $(X Z)(Y Z) \in \mathcal{SP}$ with $|(X Z)(Y Z)|_{S,K} = |((S X) Y) Z|_{S,K}$ (since $|Z|_{S,K} = 1$) and $\pi(((S X) Y) Z) \to_r (\pi(X)\,S)(\pi(Y)\,S) \equiv \pi((X Z)(Y Z))$. Left-child steps occur only inside $P S$ or $P K$ (since $K K, K S, S K, S S$ are normal), and right-child steps occur only inside $A B M$ for $A, B \in \{S, K\}$ (since $S$ and $K$ are normal), both preserving $\mathcal{SP}$ and the disjunction by the induction hypothesis. Nested well-founded induction on $|P|_{S,K} \in \mathbb{N}$ and $\|\pi(P)\|_{\mathrm{B}} \in \mathbb{N}$ (via Theorem~\ref{thm:leaf-preserving}(i)) establishes constructive strong normalization for every $P \in \mathcal{SP}$, and $\mathcal{F}_{S,K} \subseteq \mathcal{SP}$ completes the proof.
\end{proof}

\subsection{Head Reduction, Root Normal Forms, and B\"ohm Trees}\label{subsec:prelim-bohm}

Every closed $S$-term $M \in \TS$ decomposes uniquely along its left spine as $M \equiv S M_1 M_2 \cdots M_k$ where $k = \ls(M) \ge 0$. If $k \le 2$ (that is, $M \equiv S$, $M \equiv S M_1$, or $M \equiv S M_1 M_2$), no contraction can occur at or along the left spine of $M$; if $k \ge 3$, the unique redex $S M_1 M_2 M_3$ at the top of the left spine is the \emph{head redex} of $M$.

\begin{definition}[Head normal forms, root normal forms, B\"ohm trees, and productivity]\label{def:hnf-bohm}
On closed $S$-terms $M, N \in \TS$:
\begin{enumerate}[label=(\roman*),leftmargin=2.2em]
  \item \emph{Head reduction} $\to_h$ contracts the unique head redex along the left spine: $S X Y Z \to_h (X Z)(Y Z)$, and $P Q \to_h P' Q$ whenever $P \to_h P'$. A term $M \in \TS$ is in \emph{head normal form} ($M \in \HNF$) if $\ls(M) \le 2$, and is \emph{solvable} or \emph{head-normalizing} ($M \in \HN$) if $M \redSs H$ for some $H \in \HNF$; otherwise $M$ is \emph{unsolvable} ($M \in \TS \setminus \HN$).
  \item A term $M \in \TS$ is a \emph{root normal form} ($M \in \RNF$) if there do not exist $X, Y, Z \in \TS$ such that $M \redSs S X Y Z$. Following Berarducci \cite{BerarducciBohm1993,Berarducci1996}, $M$ is \emph{mute} (or root-active) if $M$ cannot reduce to any root normal form.
  \item Following B\"ohm \cite{Bohm1968} and Barendregt \cite{Barendregt1984}, the \emph{B\"ohm tree} $\BT(M)$ is the potentially infinite labelled tree over signature $\{S, \cdot, \bot\}$ defined coinductively by $\BT(M) = \bot$ if $M \notin \HN$, and
  \[
    \BT(M) \;=\; S \qquad (M \redSs S), \qquad
    \BT(M) \;=\; S\,\BT(M_1) \qquad (M \redSs S M_1), \qquad
    \BT(M) \;=\; S\,\BT(M_1)\,\BT(M_2) \qquad (M \redSs S M_1 M_2).
  \]
  Similarly, the \emph{Berarducci tree} $\BeT(M)$ unfolds $M$ at every node having a root normal form ($S$ if $M \redSs S$, or $\BeT(P) \cdot \BeT(Q)$ if $M \redSs P Q \in \RNF$) and places $\bot$ at mute subterms.
  \item To reason constructively about $\bot$-occurrence without coinductive choice, define the inductive finite-depth predicate $\bot_d(M)$ ($d \in \mathbb{N}$, read ``$\BT(M)$ contains $\bot$ at tree depth at most $d$'') by four rules: (a)~if $M \notin \HN$, then $\bot_d(M)$ holds for all $d$; (b)~if $M \redSs S X$ and $\bot_d(X)$, then $\bot_{d+1}(M)$; (c)~if $M \redSs S X Y$ and $\bot_d(X)$, then $\bot_{d+1}(M)$; and (d)~if $M \redSs S X Y$ and $\bot_d(Y)$, then $\bot_{d+1}(M)$. A closed $S$-term $M \in \TS$ is \emph{productive} ($M \in \Prod$) if $\neg\bot_d(M)$ holds for all $d \in \mathbb{N}$ \cite{EndrullisGrabmayerHendriksIsiharaKlop2010,BarendregtEndrullisKlopWaldmann2017}.
\end{enumerate}
\end{definition}

Observe that $\HNF$ is closed under arbitrary reductions $\redSs$: if $\ls(H) \le 2$ and $H \redS H'$, then either $H \equiv S X \redS S X'$ or $H \equiv S X Y \redS S X' Y$ or $S X Y'$, preserving $\ls(H') = \ls(H) \le 2$.

\begin{lemma}[Conversion invariance of $\HN$ and $\Prod$]\label{lem:conv-invariance}
For all $M, N \in \TS$ satisfying $M \convS N$, one has $M \in \HN$ if and only if $N \in \HN$, $\bot_d(M)$ if and only if $\bot_d(N)$ for every $d \in \mathbb{N}$, and hence $M \in \Prod$ if and only if $N \in \Prod$.
\end{lemma}

\begin{proof}
By Proposition~\ref{prop:cr-conservation}(i), it suffices to show that $M \redSs N$ implies the equivalences $M \in \HN$ if and only if $N \in \HN$, and $\bot_d(M)$ if and only if $\bot_d(N)$. First, if $N \redSs H$ with $\ls(H) \le 2$, transitivity gives $M \redSs H$, so $N \in \HN$ implies $M \in \HN$. Conversely, if $M \redSs H$ with $\ls(H) \le 2$, Church--Rosser yields $U$ with $N \redSs U$ and $H \redSs U$; since $\ls \le 2$ is preserved under $\redSs$, we have $\ls(U) \le 2$ and $N \in \HN$.

Next we prove by induction on $d$ that $M \redSs N$ implies $\bot_d(M)$ if and only if $\bot_d(N)$. At depth $d = 0$, the predicate $\bot_0(P)$ holds if and only if $P \notin \HN$, which is already invariant. Assume equivalence holds at depth $d$. If $\bot_{d+1}(N)$ holds via $N \redSs S X$ with $\bot_d(X)$ (or $N \redSs S X Y$ with $\bot_d(X)$ or $\bot_d(Y)$), then $M \redSs N \redSs S X$ (or $S X Y$) derives $\bot_{d+1}(M)$. Conversely, suppose $\bot_{d+1}(M)$ holds. If $M \notin \HN$, then $N \notin \HN$ and $\bot_{d+1}(N)$ holds by rule~(a). If $M \redSs S X$ with $\bot_d(X)$, Church--Rosser on $M \redSs N$ and $M \redSs S X$ yields $U$ with $N \redSs U$ and $S X \redSs U$; inversion on $S X \redSs U$ forces $U \equiv S X'$ with $X \redSs X'$, so the induction hypothesis gives $\bot_d(X')$ and hence $\bot_{d+1}(N)$. The remaining two cases where $M \redSs S X Y$ follow identically because $S X Y \redSs U$ forces $U \equiv S X' Y'$ with $X \redSs X'$ and $Y \redSs Y'$. Negating the existence of $d \in \mathbb{N}$ satisfying $\bot_d(-)$ yields $M \in \Prod$ if and only if $N \in \Prod$.
\end{proof}

\subsection{Bottom-Up Tree Automata, Finite \texorpdfstring{$S$}{S}-Algebras, and Ongoing Simulation}\label{subsec:prelim-automata}

Because $\TS$ is the set of ground terms over a signature consisting of one nullary constant $S$ and one binary operator $\cdot$, deterministic bottom-up tree automata on $\TS$ coincide with finite pointed magmas equipped with an accepting subset \cite{EndrullisDeVrijerWaldmann2010,BarendregtEndrullisKlopWaldmann2017}.

\begin{definition}[Bottom-up tree automata and finite $S$-algebras]\label{def:s-algebra}
A \emph{deterministic bottom-up tree automaton} (DTA) over $\TS$ is a quadruple $\mathcal{A} = (Q, q_0, \odot, F)$, where $Q$ is a finite nonempty set of states, $q_0 \in Q$ is the leaf constant ($q_0 = \sem{S}_{\mathcal{A}}$), $\odot \colon Q \times Q \to Q$ is the binary transition operation, and $F \subseteq Q$ is the accepting subset. The evaluation morphism $\sem{-}_{\mathcal{A}} \colon \TS \to Q$ is defined recursively by
\[
  \sem{S}_{\mathcal{A}} \;=\; q_0, \qquad \sem{M N}_{\mathcal{A}} \;=\; \sem{M}_{\mathcal{A}} \odot \sem{N}_{\mathcal{A}},
\]
and the recognized regular tree language is $L(\mathcal{A}) = \{M \in \TS : \sem{M}_{\mathcal{A}} \in F\}$. Clearly $\sem{M}_{\mathcal{A}}$ is computable in a single bottom-up pass of $2|M| - 1$ table lookups ($O(|M|)$ time). We call $\mathcal{A}$ a \emph{finite $S$-algebra} if $\odot$ satisfies the combinatory identity
\begin{equation}\label{eq:s-algebra-law}
  ((q_0 \odot x) \odot y) \odot z \;=\; (x \odot z) \odot (y \odot z) \qquad (\text{for all } x, y, z \in Q).
\end{equation}
Because $\sem{-}_{\mathcal{A}}$ is compositional on application contexts, \eqref{eq:s-algebra-law} implies $\sem{M}_{\mathcal{A}} = \sem{N}_{\mathcal{A}}$ whenever $M \convS N$, making $L(\mathcal{A})$ invariant under unrestricted pure-$S$ conversion.
\end{definition}

Building on classical criteria for non-trivial simulation and encoding complexity \cite{Davis1956NoteUTM,Fischer1965Formalisms,Herman1968Simulation} together with Priese \cite{Priese1979ComplexityUTM}, van Emde Boas \cite{vanEmdeBoas1990Models}, Neary and Woods \cite{NearyWoods2006SmallFast}, and Wolfram \cite{Wolfram2021Challenge}, we formalize ongoing simulation of a deterministic computation system in $\CLS$ with explicit sub-universal guards on all boundary maps.

\begin{definition}[Ongoing pure-$S$ simulation with sub-universal boundary maps]\label{def:ongoing-sim}
Let $\mathcal{S}$ be a deterministic computational model with instance set $\mathcal{I}$, step relation $\vdash_{\mathcal{S}}$, halting predicate $\mathrm{Halts}_{\mathcal{S}}(I)$, and partial output map $\out_{\mathcal{S}} \colon \mathcal{I} \rightharpoonup \Sigma^*$. An \emph{ongoing selected-path simulation} of $\mathcal{S}$ in $\CLS$ is a quadruple $(\mathcal{E}, \sigma, \mathcal{H}, \mathcal{D})$ consisting of:
\begin{enumerate}[label=(\roman*),leftmargin=2.2em]
  \item a total polynomial-time (or logspace) \emph{encoder} $\mathcal{E} \colon \mathcal{I} \to \TS$ mapping each instance $I \in \mathcal{I}$ to a finite closed $S$-term $T_0 = \mathcal{E}(I)$ of size $|T_0| \le \mathrm{poly}(\|I\|)$;
  \item a fixed stateless \emph{single-step selector} $\sigma \colon \TS \to \TS \cup \{\mathsf{normal}\}$ realized by a finite-state tree-walking automaton that resets to a unique initial control state at the root of $T$ before each step, halts on every $T \in \TS$ in $O(|T|)$ cursor transitions, and satisfies $T \redS \sigma(T)$ whenever $\sigma(T) \in \TS$;
  \item a \emph{halting detector} $\mathcal{H} = (Q_{\mathrm{halt}}, q_0, \odot_{\mathrm{halt}}, F_{\mathrm{halt}})$ given by a fixed deterministic bottom-up finite tree automaton ($L(\mathcal{H}) \in \REG$); and
  \item a total polynomial-time (in $|T|$) read-only \emph{output decoder} $\mathcal{D} \colon \TS \to \Sigma^* \cup \{\bot\}$,
\end{enumerate}
such that along the iterated reduction trajectory $T_0 = \mathcal{E}(I)$ and $T_{k+1} = \sigma(T_k)$ (defined while $\sigma(T_k) \in \TS$), one has $\mathrm{Halts}_{\mathcal{S}}(I)$ if and only if there exists $k \ge 0$ with $T_k \in L(\mathcal{H})$, and every such step $k$ satisfies $\mathcal{D}(T_k) = \out_{\mathcal{S}}(I)$. Because $\mathcal{E}$ is total polynomial-time (indeed logspace-constructible up to tree output), $\mathcal{H}$ has $O(1)$ finite memory, and $\mathcal{D}$ terminates in bounded polynomial time on every bare tree $T \in \TS$, none of the boundary maps $(\mathcal{E}, \mathcal{H}, \mathcal{D})$ can decide any $\Sigma_1^0$-complete problem on its own.
\end{definition}

\subsection{Open Spirals and Structural Obstructions to Conjecture 17}\label{subsec:prelim-conj17}

Waldmann's Theorem~14 (Proposition~\ref{prop:acyclicity-order}(ii)) rules out self-embedding reductions $M \redSp C[M]$ for identical copies of $M$, whether closed or open. When $M \in \TSV$ contains variables, however, one may ask whether $M$ can self-embed modulo a non-trivial substitution $\sigma \colon \mathcal{V} \to \TSV$. In Section~5.4 of \cite[Definition~11(iii) and Conjecture~17, pp.~86, 89]{BarendregtEndrullisKlopWaldmann2017}, Barendregt, Endrullis, Klop, and Waldmann conjectured that no such open spiralling $S$-term exists. While Conjecture~17 in full generality remains open, we establish a suite of constructive structural obstructions that eliminate broad classes of substitutions, contexts, variable-occurrence patterns, and short reduction lengths, while isolating a parametric self-embedding near-miss native to pure $\CLS$.

\begin{definition}[Open spiralling, iterated substitutions, and variable-spine depth]\label{def:open-spiral}
Let $\sigma \colon \mathcal{V} \to \TSV$ be a substitution. Its \emph{iterates} $\sigma^n \colon \mathcal{V} \to \TSV$ ($n \in \mathbb{N}$) are defined by $\sigma^0(x) \equiv x$ and $\sigma^{n+1}(x) \equiv (\sigma(x))^{\sigma^n}$, satisfying $(M^\sigma)^{\sigma^n} \equiv M^{\sigma^{n+1}}$ for all $M \in \TSV$. A term $M \in \TSV$ \emph{open-spirals} via a one-hole context $C[\,]$ and substitution $\sigma$ if $M \redSp C[M^\sigma]$; \emph{Conjecture~17} of \cite{BarendregtEndrullisKlopWaldmann2017} asserts that $\neg(M \redSp C[M^\sigma])$ holds for all $M \in \TSV$, $C[\,]$, and $\sigma$.

A \emph{left-spine context} is generated by $L[\,] ::= \square \mid (L[\,]\,R)$ ($R \in \TSV$). For each variable $x \in \mathcal{V}$, the \emph{maximum $x$-tailed right-spine depth} $\mathrm{maxRS}_x \colon \TSV \to \mathbb{N}$ is defined recursively by $\mathrm{maxRS}_x(a) = 0$ on atoms $a \in \{S\} \cup \mathcal{V}$ and
\[
  \mathrm{maxRS}_x(P Q) \;=\; \max\bigl(\delta_{\mathrm{rt}(Q), x}\,(1 + \rs(Q)),\; \mathrm{maxRS}_x(P),\; \mathrm{maxRS}_x(Q)\bigr),
\]
where $\delta_{\mathrm{rt}(Q), x} = 1$ if $\mathrm{rt}(Q) \equiv x$ and $0$ otherwise.
\end{definition}

\begin{proposition}[Divergence projection, minimal counterexample geometry, and $8$-leaf open threshold]\label{prop:open-spiral-divergence-threshold}
Let $M \in \TSV$.
\begin{enumerate}[label=(\roman*),leftmargin=2.2em]
  \item \emph{Universal substitution divergence and minimal spine geometry:} If $M \redSp C[M^\sigma]$, then $M \notin \SN$, and $M^\tau \notin \SN$ holds for every substitution $\tau \colon \mathcal{V} \to \TSV$ (in particular for the closing projection $\tau_S(v) \equiv S$, which satisfies $|M^{\tau_S}| = |M|$). Moreover, if $M$ is leaf-minimal among open-spiralling terms, then $M \redSp C[M^\sigma]$ performs at least one root contraction, $\mathrm{hd}(M) \equiv S$, and $\ls(M) \ge 3$.
  \item \emph{Exact $8$-leaf threshold for open non-normalization:} Every open term $M \in \TSV \setminus \TS$ containing at least one variable and having leaf size $|M| \le 7$ lies in $\SN$ (and therefore cannot open-spiral). At leaf size $|M| = 8$, non-normalizing open $S$-terms appear both by wrapping $7$-leaf closed divergers (such as $S(S S)S S S S x$) and as genuinely open divergers whose proper closed subterms all normalize, such as $S A S (S(S x))$.
\end{enumerate}
\end{proposition}

\begin{proof}
For~(i), compatibility of $\redSp$ with substitutions and contexts implies that $M_0 \equiv M$ and $M_{n+1} \equiv C[M_n^\sigma]$ satisfy $M_n \redSp M_{n+1}$ for all $n \in \mathbb{N}$, refuting constructive accessibility of $M \in \SN$. Because $P \redS Q$ implies $P^\tau \redS Q^\tau$, accessibility of $M^\tau$ pulls back by structural induction on the accessibility witness to accessibility of $M$, whence $M^\tau \notin \SN$. Suppose next that $M$ has minimal leaf size among all open-spiralling terms and $M \redSp C[M^\sigma]$ uses no root contraction. Then $M \equiv P Q$ with $P \redSs P'$, $Q \redSs Q'$, and $P' Q' \equiv C[P^\sigma Q^\sigma]$. If $C[\,] \equiv \square$, then $P' \equiv P^\sigma$, $Q' \equiv Q^\sigma$, and at least one of $P \redSp P^\sigma$ or $Q \redSp Q^\sigma$ holds, yielding an open spiral on a proper subterm of $M$. If $C[\,] \equiv (C_1[\,]\,R)$, then $P^\sigma Q^\sigma \trianglelefteq P'$, so $P^\sigma \triangleleft P'$, say $P' \equiv C_a[P^\sigma]$ with $C_a[\,] \not\equiv \square$; since $|P| \le |P^\sigma| < |C_a[P^\sigma]| = |P'|$, the reduction $P \redSs P'$ must have positive length $P \redSp C_a[P^\sigma]$, again contradicting leaf-minimality of $M$ (and symmetrically when $C[\,] \equiv (L\,C_2[\,])$). Thus $M \redSp C[M^\sigma]$ performs a root contraction. Because any step from a variable-headed term preserves its variable head and prevents root contractions, $\mathrm{hd}(M) \equiv S$; similarly, if $\mathrm{hd}(M) \equiv S$ and $\ls(M) \le 2$, every step preserves $\mathrm{hd} \equiv S$ and $\ls \le 2$, ruling out root contractions and forcing $\ls(M) \ge 3$.

For~(ii), let $M \in \TSV \setminus \TS$ have $|M| \le 7$, and suppose toward a contradiction that $M \notin \SN$. By pullback under $\tau_S(v) \equiv S$, the closed projection $M^{\tau_S} \in \TS$ satisfies $|M^{\tau_S}| = |M| \le 7$ and $M^{\tau_S} \notin \SN$. As shown by Dubou\'e \cite{BarendregtBergstraKlopVolken1976} and verified in Section~\ref{sec:asymptotics}, every closed $S$-term of leaf size at most~$6$ normalizes, and at leaf size~$7$ the only closed non-normalizing terms are $X_1 \equiv S(S S)S S S S$ and $X_2 \equiv S S S(S S)S S$. Identifying all variables of $M$ to $x_0$ and substituting $S$ for all but one occurrence of $x_0$ yields a single-variable $7$-leaf term $M_1(x_0) \notin \SN$ obtained from $X_1$ or $X_2$ by replacing one leaf $S$ with $x_0$. Direct normalization of all $14$ single-leaf variable replacements of $X_1$ and $X_2$ shows that the $7$ replacements of $X_1$ reach normal form in $0, 1, 2, 3, 4, 8, 14$ leftmost-outermost steps respectively, and those of $X_2$ reach normal form in $0, 1, 4, 2, 3, 7, 13$ steps, a contradiction. At $|M| = 8$, the term $X_1 x_0$ inherits every infinite reduction of $X_1$. For the genuinely open $8$-leaf term $S A S (S(S x_0))$, set $K \equiv S(S(S x_0))$ and let $\mathcal{G} \subset \TSV$ be generated inductively by $K \in \mathcal{G}$ and $U V \in \mathcal{G}$ for $U, V \in \mathcal{G}$. Two contractions give $S A S (S(S x_0)) \redS A(S(S x_0))\,K \redS K K K \in \mathcal{G}$. Moreover, for any $U_1, U_2 \in \mathcal{G}$, three contractions on $K U_1 U_2 \equiv S(S(S x_0))\,U_1\,U_2$ yield
\[
  K U_1 U_2 \;\redS\; \bigl(S(S x_0)\,U_2\bigr)(U_1 U_2) \;\redS\; \bigl(S x_0\,(U_1 U_2)\bigr)\bigl(U_2(U_1 U_2)\bigr) \;\redS\; \bigl(x_0(U_2(U_1 U_2))\bigr)\Bigl((U_1 U_2)\bigl(U_2(U_1 U_2)\bigr)\Bigr),
\]
whose right factor $(U_1 U_2)(U_2(U_1 U_2)) \in \mathcal{G}$ again has at least two $\mathcal{G}$-arguments along its left spine above $K$, generating an infinite reduction; hence $S A S (S(S x_0)) \notin \SN$ by Proposition~\ref{prop:cr-conservation}(iii).
\end{proof}

Our next theorem eliminates the ambient context $C[\,]$ via a constructive step-indexed dependency-pair extraction and combines it with weighted right-spine path orderings to restrict the right-spine action of any candidate substitution $\sigma$.

\begin{theorem}[Right-spine orbit obstruction, weighted Waldmann criterion, and self-right-tail theorem]\label{thm:conj17-dp-weighted}
Let $M \in \TSV$, let $C[\,]$ be a one-hole context, and let $\sigma \colon \mathcal{V} \to \TSV$ be a substitution.
\begin{enumerate}[label=(\roman*),leftmargin=2.2em]
  \item \emph{Bounded right-spine orbits and weighted right-spine LPO:} If there exists $R \in \mathbb{N}$ such that $\rs(\sigma^n(x)) \le R$ for all $n \in \mathbb{N}$ and all $x \in \mathrm{Vars}(M)$, then $\neg(M \redSp C[M^\sigma])$. Likewise, given weights $w \colon \mathcal{V} \to \mathbb{N}$, define $\rho_w(S) = 1$, $\rho_w(x) = w(x)$, and $\rho_w(P Q) = 1 + \rho_w(Q)$; if $\rho_w(\sigma(x)) \le w(x)$ holds for all $x \in \mathrm{Vars}(M)$, then $\neg(M \redSp C[M^\sigma])$. In particular, $\neg(M \redSp C[M^\sigma])$ holds whenever every $x \in \mathrm{Vars}(M)$ satisfies $\mathrm{rt}(\sigma(x)) \equiv S$ (all $S$-tailed and closed substitutions) or $\rs(\sigma(x)) = 0$ (all variable renamings and $S$-erasures).
  \item \emph{Dependency-pair self-right-tail theorem:} If $M \redSp C[M^\sigma]$, then there exists a variable $x \in \mathrm{Vars}(M)$ satisfying both
  \begin{equation}\label{eq:self-right-tail}
    \mathrm{rt}(\sigma(x)) \;\equiv\; x \qquad \text{and} \qquad \rs(\sigma(x)) \;\ge\; 1.
  \end{equation}
  Equivalently, if no variable in $\mathrm{Vars}(M)$ is a positive right-spine self-loop of $\sigma$, then $\neg(M \redSp C[M^\sigma])$. In particular, $\neg(M \redSp C[M^\sigma])$ holds whenever $\mathrm{rt}(\sigma(x)) \not\equiv x$ for all $x \in \mathrm{Vars}(M)$ (eliminating all cyclic right-spine permutations of period $p \ge 2$), and every single-variable open spiral $M(x) \redSp C[M(\sigma(x))]$ forces $x \in \mathrm{Vars}(M)$, $\mathrm{rt}(\sigma(x)) \equiv x$, and $\rs(\sigma(x)) \ge 1$.
\end{enumerate}
\end{theorem}

\begin{proof}
For~(i), observe first that $P \redSs Q$ implies $\mathrm{Vars}(Q) \subseteq \mathrm{Vars}(P)$; hence if $M \redSp C[M^\sigma]$, every $x \in \mathrm{Vars}(M)$ satisfies $\mathrm{Vars}(\sigma(x)) \subseteq \mathrm{Vars}(M^\sigma) \subseteq \mathrm{Vars}(C[M^\sigma]) \subseteq \mathrm{Vars}(M)$, so $\mathrm{Vars}(M)$ is $\sigma$-closed. Under the uniform orbit bound $\rs(\sigma^n(x)) \le R$ on $\mathrm{Vars}(M)$, set $L = R + \max_{x \in \mathrm{Vars}(M)} \max\mathrm{label}(\overline{\sigma(x)})$, where $\overline{P}$ is the canonical $\rho$-labelling from Proposition~\ref{prop:acyclicity-order}(ii). Because $\rho(U^\tau) \le \rho(U) + R$ whenever $\rs(\tau(v)) \le R$ on $\mathrm{Vars}(U)$, induction on $n$ shows that $\max\mathrm{label}(\overline{\sigma^n(x)}) \le L$ for all $n \in \mathbb{N}$ and $x \in \mathrm{Vars}(M)$. Let $B_0$ be the maximum canonical label occurring along the finite reduction $M \redSp C[M^\sigma]$, and set $B = \max(B_0 + R, L)$. Substituting $\sigma^n$ across $M \redSp C[M^\sigma]$ produces reductions $M^{\sigma^n} \redSp C^{\sigma^n}[M^{\sigma^{n+1}}]$ whose canonical labels never exceed $B$; composing with hole projections along $C^{\sigma^n}[\,]$ yields an infinite chain in the $B$-inverted relation $\succ_{\mathrm{RB}}^+$, contradicting its constructive well-foundedness established in Proposition~\ref{prop:acyclicity-order}(ii). For the weighted criterion, merge $S$ and all variables into a single leaf symbol $\bullet$ and label each application $U V$ by $\rho_w(V)$. When $\rho_w(\sigma(x)) \le w(x)$ on $\mathrm{Vars}(M)$, structural induction gives $\rho_w(N^\sigma) \le \rho_w(N)$ for every $N \trianglelefteq M$, so the $B$-inverted leaf-merged representation of $M^\sigma$ reduces via label-lifting, subterm projection (taking any $\sigma(x)$ to $\bullet$), and congruence to that of $M$, creating a cycle in the leaf-merged analogue of $\succ_{\mathrm{RB}}^+$ whenever $M \redSp C[M^\sigma]$. Choosing $w(x) = 1 + \rs(\sigma(x))$ yields $\rho_w(\sigma(x)) = 1 + \rs(\sigma(x)) = w(x)$ when $\mathrm{rt}(\sigma(x)) \equiv S$, and choosing $w(x) = 1$ yields $\rho_w(\sigma(x)) = 1$ when $\rs(\sigma(x)) = 0$.

For~(ii), fix $\sigma$ and define the dependency-pair relation $U \Rightarrow_{\mathcal{DP}} V$ inductively by the base clause $U \to_{\mathrm{nr}}^* V$ and the three step clauses: whenever $U \to_{\mathrm{nr}}^* S X Y Z$, one has $U \Rightarrow_{\mathcal{DP}} V$ if $(X Z)(Y Z) \Rightarrow_{\mathcal{DP}} V$ (clause~$\mathrm{dp}_1$), $X Z \Rightarrow_{\mathcal{DP}} V$ (clause~$\mathrm{dp}_2$), or $Y Z \Rightarrow_{\mathcal{DP}} V$ (clause~$\mathrm{dp}_3$). We first prove by strong induction on $n \in \mathbb{N}$ that whenever $T \redS^n C'[V]$, either there exist $n' \le n$, a proper subterm $U \triangleleft T$, and a context $C''[\,]$ with $U \redS^{n'} C''[V]$, or else $T \Rightarrow_{\mathcal{DP}} V$. If $T \redS^n C'[V]$ uses no root steps, then either $C'[\,] \equiv \square$ (giving $T \to_{\mathrm{nr}}^* V$, hence $T \Rightarrow_{\mathcal{DP}} V$) or $T \equiv T_1 T_2$ with $C'[\,] \equiv (C_1[\,]\,R)$ or $(L\,C_2[\,])$, in which case non-root decomposition projects directly to $T_1 \redS^{n_1} C_1[V]$ or $T_2 \redS^{n_2} C_2[V]$ ($n_i \le n$). Otherwise split at the first root step $T \to_{\mathrm{nr}}^k S X Y Z \to_r (X Z)(Y Z) \redS^m C'[V]$ with $k + 1 + m = n$ ($T \equiv T_1 T_2$, $T_1 \redS^{k_1} S X Y$, $T_2 \redS^{k_2} Z$, $k_1 + k_2 = k$). Applying the induction hypothesis to $(X Z)(Y Z) \redS^m C'[V]$ ($m < n$), either $(X Z)(Y Z) \Rightarrow_{\mathcal{DP}} V$ (giving $T \Rightarrow_{\mathcal{DP}} V$ via $\mathrm{dp}_1$), or some $U_1 \triangleleft (X Z)(Y Z)$ satisfies $U_1 \redS^{m_1} C_1[V]$ ($m_1 \le m$). If $U_1 \trianglelefteq X, Y$, or $Z$, embedding $U_1 \redS^{m_1} C_1[V]$ inside $S X Y$ or $Z$ and prepending $T_1 \redS^{k_1} S X Y$ or $T_2 \redS^{k_2} Z$ lifts the reduction to $T_1$ or $T_2 \triangleleft T$ in at most $k + m < n$ steps; if instead $U_1 \equiv X Z$ or $U_1 \equiv Y Z$, a second appeal to the induction hypothesis on $U_1 \redS^{m_1} C_1[V]$ ($m_1 \le m < n$) either yields $U_1 \Rightarrow_{\mathcal{DP}} V$ (giving $T \Rightarrow_{\mathcal{DP}} V$ via $\mathrm{dp}_2$ or $\mathrm{dp}_3$) or locates $U_2 \triangleleft U_1$, so $U_2 \trianglelefteq X, Y$, or $Z$, lifting again to $T_1$ or $T_2 \triangleleft T$.

Now let $M$ be leaf-minimal among terms open-spiralling under the fixed substitution $\sigma$, with $M \redS^n C[M^\sigma]$ ($n \ge 1$). If some proper subterm $U \triangleleft M$ ($M \equiv D[U]$, $D[\,] \not\equiv \square$) satisfied $U \redS^{n'} C'[M^\sigma] \equiv (C' \circ D^\sigma)[U^\sigma]$, then $|U| < |M| \le |M^\sigma| \le |(C' \circ D^\sigma)[U^\sigma]|$ would force $n' \ge 1$, contradicting leaf-minimality of $M$. Therefore $M \Rightarrow_{\mathcal{DP}} M^\sigma$. Every $\mathcal{DP}$-transition preserves $\mathrm{rt}$ ($\mathrm{rt}((X Z)(Y Z)) \equiv \mathrm{rt}(X Z) \equiv \mathrm{rt}(Y Z) \equiv \mathrm{rt}(Z) \equiv \mathrm{rt}(S X Y Z)$) and is non-decreasing in $\rs$, while any derivation using $\mathrm{dp}_1$ at least once satisfies $\rs(M) < \rs(M^\sigma)$ strictly because $\rs((X Z)(Y Z)) = 2 + \rs(Z) > 1 + \rs(Z) \ge \rs(U)$. Suppose $M \Rightarrow_{\mathcal{DP}} M^\sigma$ used only the base clause, $\mathrm{dp}_2$, and $\mathrm{dp}_3$. By Proposition~\ref{prop:open-spiral-divergence-threshold}(i) (or Proposition~\ref{prop:acyclicity-order}(ii) on atoms), $M \equiv A_0 B_0$ is an application. Inducting along a $\mathrm{dp}_1$-free chain from $A_0 B_0$ to $A_0^\sigma B_0^\sigma$ shows that either $A_0 \redSs A_0^\sigma$ and $B_0 \redSs B_0^\sigma$ (in the base clause), or $A_0 \redSs D[A_0^\sigma]$ for a non-empty context $D[\,] \not\equiv \square$ (since $A_0 \redSs S X Y$ embeds any target reached from $X$ or $Y$ strictly inside $(S\,\square)\,Y$ or $(S X)\,\square$). Both cases are impossible: $A_0 \redSs D[A_0^\sigma]$ forces $A_0 \redSp D[A_0^\sigma]$ by $|A_0| < |D[A_0^\sigma]|$, contradicting leaf-minimality of $M$, while $A_0 \redSs A_0^\sigma$ and $B_0 \redSs B_0^\sigma$ either give $M^\sigma \equiv M$ (contradicting Proposition~\ref{prop:acyclicity-order}(ii) on $M \redSp C[M]$) or yield $A_0 \redSp A_0^\sigma$ or $B_0 \redSp B_0^\sigma$, again contradicting leaf-minimality of $M$. Thus $\mathrm{dp}_1$ must occur, giving $\mathrm{rt}(M) \equiv \mathrm{rt}(M^\sigma)$ and $\rs(M) < \rs(M^\sigma)$. If $\mathrm{rt}(M) \equiv S$, then $\rs(M^\sigma) = \rs(M)$, a contradiction. Hence $\mathrm{rt}(M) \equiv x$ for some $x \in \mathrm{Vars}(M)$, whence $\mathrm{rt}(M^\sigma) \equiv \mathrm{rt}(\sigma(x))$ and $\rs(M^\sigma) = \rs(M) + \rs(\sigma(x))$, forcing $\mathrm{rt}(\sigma(x)) \equiv x$ and $\rs(\sigma(x)) \ge 1$.
\end{proof}

Even when $\sigma$ admits positive right-spine self-loops satisfying \eqref{eq:self-right-tail}, four further structural laws rule out empty and left-spine contexts, non-duplicating counter substitutions, left-only variable patterns, and $1$- or $2$-step loops universally across $\TSV$, while explaining how pure $\CLS$ nonetheless realizes parametric open self-reproduction. In part~(i) below, we invoke Waldmann's top-termination theorem and the open head-step potential $\nu_{\HN}$, which are proved independently in Section~\ref{sec:algebras} (Theorem~\ref{thm:waldmann-top-term}) using only Subsection~\ref{subsec:prelim-syntax} and Proposition~\ref{prop:acyclicity-order}(i).

\begin{theorem}[Unconditional context/short-loop obstructions, multiplicity--spine coupling, and pure-$\CLS$ near-miss]\label{thm:conj17-unconditional-and-coupling}
Let $M \in \TSV$, let $C[\,]$ be a one-hole context, let $L[\,]$ be a left-spine context, and let $\sigma \colon \mathcal{V} \to \TSV$ be an arbitrary substitution.
\begin{enumerate}[label=(\roman*),leftmargin=2.2em]
  \item \emph{Absence of open cycles and left-spine spirals:} One has $\neg(M \redSp M^\sigma)$ unconditionally ($C[\,] \equiv \square$). Moreover, $\neg(M \redSp L[M^\sigma])$ holds unconditionally whenever $\mathrm{hd}(M) \in \mathcal{V}$, and more generally whenever $M$ evaluates outside the unsolvable sink $\infty_{\HN}$ in the open extension of the $35$-state head-normalization algebra $\mathcal{A}_{\HN}$ (with $\sem{x}_{\HN}^{\mathrm{open}} = \infty_{\HN}$ for all $x \in \mathcal{V}$).
  \item \emph{Variable-multiplicity and right-spine coupling:} Every step $P \redS Q$ in $\TSV$ and every variable $x \in \mathcal{V}$ satisfy $\#_x(P) \le \#_x(Q)$, $\rs(Q) - \rs(P) \le \#_x(Q) - \#_x(P)$ whenever $\mathrm{rt}(P) \equiv x$, and
  \begin{equation}\label{eq:maxrs-coupling}
    \mathrm{maxRS}_x(Q) + \#_x(P) \;\le\; \mathrm{maxRS}_x(P) + \#_x(Q).
  \end{equation}
  Consequently: (a)~if $\mathrm{maxRS}_x(M) = 0$ for all $x \in \mathcal{V}$ (every variable in $M$ occurs solely in left/function positions), then $\neg(M \redSp C[M^\sigma])$ for all $C[\,]$ and $\sigma$; and (b)~if $\#_x(C) = 0$ for all $x \in \mathrm{Vars}(M)$ and $\sigma$ is diagonally non-duplicating on $\mathrm{Vars}(M)$ ($\#_x(\sigma(x)) \le 1$ and $x \notin \mathrm{Vars}(\sigma(y))$ for $y \in \mathrm{Vars}(M) \setminus \{x\}$), then $\neg(M \redSp C[M^\sigma])$.
  \item \emph{Universal $1$-step and $2$-step open-spiral impossibility:} Neither $M \redS C[M^\sigma]$ nor $M \redS^2 C[M^\sigma]$ holds for any $M \in \TSV$, $C[\,]$, and $\sigma$.
  \item \emph{Parametric self-embedding near-miss in pure $\CLS$:} Recall $T \equiv S S$, $G_0 \equiv S T T$, and $V \equiv S G_0$, and define $T^0 U \equiv U$, $T^{n+1} U \equiv T(T^n U)$. For every $x \in \mathcal{V}$ and every $n \in \mathbb{N}$, there exists an explicit context $C_n[\,] \not\equiv \square$ such that under leftmost-outermost pure-$S$ reduction $\to_{\mathrm{lo}}$,
  \begin{equation}\label{eq:parametric-near-miss}
    V x\,\bigl(T^n(V x)\bigr) \;\to_{\mathrm{lo}}^{4n+3}\; C_n\Bigl[V x\,\bigl(T^{n+1}(V x)\bigr)\Bigr].
  \end{equation}
\end{enumerate}
\end{theorem}

\begin{proof}
For~(i), suppose $M \redSp M^\sigma$ with $|M|$ minimal. If no root contraction occurred along $M \redSp M^\sigma$, then $M \equiv P Q$ would satisfy $P \redSp P^\sigma$ or $Q \redSp Q^\sigma$, contradicting minimality of $|M|$. Thus $M$ reduces to $M^\sigma$ via a sequence containing at least one root contraction ($M \xrightarrow{r} M^\sigma$). Since $\xrightarrow{r}$ is preserved under arbitrary substitutions, applying $\sigma^n$ yields an infinite sequence $M \xrightarrow{r} M^\sigma \xrightarrow{r} M^{\sigma^2} \xrightarrow{r} \cdots$ with infinitely many root contractions, contradicting constructive top-termination on $\TSV$ (Theorem~\ref{thm:waldmann-top-term}). Next suppose $M \redSp L[M^\sigma]$. If $\mathrm{hd}(M) \in \mathcal{V}$ or if $M \redSp L[M^\sigma]$ uses only non-head internal steps, then left-spine length is preserved identically: $\ls(M) = \ls(L[M^\sigma]) = \ls(M^\sigma) + d_{\mathrm{L}}(L) \ge \ls(M) + d_{\mathrm{L}}(L)$, forcing $d_{\mathrm{L}}(L) = 0$ ($L[\,] \equiv \square$) and contradicting $\neg(M \redSp M^\sigma)$. Otherwise $M \redSp L[M^\sigma]$ decomposes as $M \redSs P_1 \to_h P_2 \redSs L[M^\sigma]$. When $\sem{M}_{\HN}^{\mathrm{open}} \ne \infty_{\HN}$, both $\to_h$ and $\redSs$ commute with $\sigma^k$ and left-spine contexts $K[\,]$, while by Theorem~\ref{thm:waldmann-top-term} internal steps never increase the integer potential $\nu_{\HN} \in \mathbb{N}$ and each head step outside $\infty_{\HN}$ strictly decreases $\nu_{\HN}$ by at least~$1$; unrolling $M \redSs P_1 \to_h P_2 \redSs L[M^\sigma]$ through $\nu_{\HN}(M) + 1$ iterations yields $\nu_{\HN}(M) < 0$, a contradiction.

For~(ii), at a root contraction $S P Q R \to_r (P R)(Q R)$, we have $\#_x((P R)(Q R)) - \#_x(S P Q R) = \#_x(R) \ge 0$. If $\mathrm{rt}(S P Q R) \equiv \mathrm{rt}(R) \equiv x$, then $x \in \mathrm{Vars}(R)$ gives $\#_x(R) \ge 1$, while $\rs((P R)(Q R)) - \rs(S P Q R) = 1 \le \#_x(R)$ and $\mathrm{maxRS}_x((P R)(Q R)) \le \mathrm{maxRS}_x(S P Q R) + 1 \le \mathrm{maxRS}_x(S P Q R) + \#_x(R)$; if $\mathrm{rt}(R) \not\equiv x$, then $\mathrm{maxRS}_x((P R)(Q R)) = \mathrm{maxRS}_x(S P Q R)$. Induction over left and right contexts preserves all three inequalities, using $\mathrm{rt}(R') \equiv \mathrm{rt}(R)$ from Proposition~\ref{prop:cr-conservation}(ii) on right-child steps. For~(a), the same root/context inspection shows that $\mathrm{maxRS}_w(P) = 0$ and $P \redS Q$ imply $\mathrm{maxRS}_w(Q) = 0$. Hence if $\mathrm{maxRS}_w(M) = 0$ for all $w$ and $M \redSp C[M^\sigma]$, then $\mathrm{maxRS}_w(C[M^\sigma]) = 0$ for all $w$; for each $v \in \mathrm{Vars}(M)$ with $\mathrm{rt}(\sigma(v)) \equiv w$, the subterm inclusion $\sigma(v) \trianglelefteq C[M^\sigma]$ forces $\rs(\sigma(v)) \le \mathrm{maxRS}_w(\sigma(v)) \le \mathrm{maxRS}_w(C[M^\sigma]) = 0$, contradicting Theorem~\ref{thm:conj17-dp-weighted}(ii). For~(b), if $M \redSp C[M^\sigma]$ under diagonal non-duplication and variable-free $C[\,]$, Theorem~\ref{thm:conj17-dp-weighted}(ii) yields $x \in \mathrm{Vars}(M)$ with $\mathrm{rt}(\sigma(x)) \equiv x$ and $\rs(\sigma(x)) \ge 1$. Then $\mathrm{maxRS}_x(M) + \rs(\sigma(x)) \le \mathrm{maxRS}_x(M^\sigma) \le \mathrm{maxRS}_x(C[M^\sigma])$, whereas $\#_x(C[M^\sigma]) = \#_x(M^\sigma) \le \#_x(M)$, contradicting $\mathrm{maxRS}_x(C[M^\sigma]) + \#_x(M) \le \mathrm{maxRS}_x(M) + \#_x(C[M^\sigma])$ obtained by summing \eqref{eq:maxrs-coupling}.

For~(iii), we first show by induction on $M \redS N$ that $M^\sigma \not\trianglelefteq N$. At $M \equiv S P Q R \to_r (P R)(Q R)$, $|M^\sigma| \ge 1 + |P| + |Q| + |R|$ exceeds both $|P R|$ and $|Q R|$, while root equality $((S P^\sigma) Q^\sigma) R^\sigma \equiv (P R)(Q R)$ forces $S P^\sigma \equiv P$, contradicting $|S P^\sigma| \ge 1 + |P|$. Contextual steps $A B \redS A' B$ and $A B \redS A B'$ follow by leaf count on the unchanged child and the induction hypothesis on the stepped child. Next suppose $M \redS M_1 \redS N$ with $M^\sigma \trianglelefteq N$, minimal in $|M|$. If both steps are non-root on $M \equiv A B$, two steps on the same side contradict $|M|$-minimality, while $A \redS A_1$ and $B \redS B_1$ ($N \equiv A_1 B_1$) reduce to $A^\sigma \trianglelefteq A_1$ or $B^\sigma \trianglelefteq B_1$, contradicting the $1$-step case. If one step is root $S P Q R \to_r (P R)(Q R)$ and the other steps internally inside $P$, $Q$, or $R$, leaf bounds $|M^\sigma| \ge 1 + |P| + |Q| + |R|$ and the $1$-step impossibility on the stepped argument eliminate every subterm position of $N$ (with root match forcing either $|S P^\sigma| = |P|$ or $P^\sigma \trianglelefteq P'$). Only the four overlapping two-root-contraction shapes remain, and in each case leaf-size inequalities eliminate all proper subterms except a single candidate equation refuted by second-order substitution unification:
\begin{enumerate}[label=(\arabic*),leftmargin=2.2em]
  \item For $M \equiv S(S U) Q R \to_r (S U R)(Q R) \to_r (U(Q R))(R(Q R))$, only $M^\sigma \equiv R(Q R)$ survives leaf bounds, giving $(S(S U^\sigma)) Q^\sigma \equiv R$ and $R^\sigma \equiv Q R$. Applying $\sigma$ yields $(S(S U^{\sigma^2})) Q^{\sigma^2} \equiv Q R$, so $Q \equiv S(S U^{\sigma^2})$ and $R \equiv Q^{\sigma^2} \equiv S(S U^{\sigma^4})$, whose left child $S$ clashes with $S(S U^\sigma)$ in $R \equiv (S(S U^\sigma)) Q^\sigma$.
  \item For $M \equiv S S Y P R \to_{\mathrm{nr}} (S P)(Y P) R \to_r (P R)((Y P) R)$, only $M^\sigma \equiv (P R)((Y P) R)$ survives, giving $(S S) Y^\sigma \equiv P$, $P^\sigma \equiv R$, and $R^\sigma \equiv (Y P) R$. Then $R \equiv (S S) Y^{\sigma^2}$ and $(Y P) R \equiv (S S) Y^{\sigma^3}$, forcing $Y P \equiv S S$ and $P \equiv S$, clashing with $P \equiv (S S) Y^\sigma$.
  \item For $M \equiv S(S X_1 X_2) Q R \to_r (S X_1 X_2 R)(Q R) \to_{\mathrm{nr}} ((X_1 R)(X_2 R))(Q R)$, only the root match survives, giving $S(S X_1^\sigma X_2^\sigma) \equiv X_1 R$ (so $X_1 \equiv S$ and $R \equiv S X_1^\sigma X_2^\sigma$), $Q^\sigma \equiv X_2 R$, and $R^\sigma \equiv Q R$. Since $X_1 \equiv S$, the identity $Q R \equiv R^\sigma \equiv (S S) X_2^{\sigma^2}$ gives $Q \equiv S S$, $X_2 R \equiv Q^\sigma \equiv S S$, and $R \equiv S$, clashing with $R \equiv S X_1^\sigma X_2^\sigma$.
  \item For $M \equiv S P (S X_1 X_2) R \to_r (P R)(S X_1 X_2 R) \to_{\mathrm{nr}} (P R)((X_1 R)(X_2 R))$, only $M^\sigma \equiv (X_1 R)(X_2 R)$ survives, giving $(S P^\sigma)(S X_1^\sigma X_2^\sigma) \equiv X_1 R$ (so $R \equiv S X_1^\sigma X_2^\sigma$) and $R^\sigma \equiv X_2 R$. Applying $\sigma$ gives $(S X_1^{\sigma^2}) X_2^{\sigma^2} \equiv X_2 R$, so $X_2 \equiv S X_1^{\sigma^2}$ and $R \equiv X_2^{\sigma^2} \equiv S X_1^{\sigma^4}$, whose left child $S$ clashes with $S X_1^\sigma$ in $R \equiv (S X_1^\sigma) X_2^\sigma$.
\end{enumerate}

For~(iv), put $U \equiv V x \equiv S G_0 x$. Because $G_0 \equiv S T T \in \NF$, induction gives $T^m U \in \NF$ for all $m \in \mathbb{N}$ and hence $R_n \equiv x(T^n U) \in \NF$. Two leftmost-outermost steps yield
\[
  U(T^n U) \;\equiv\; S G_0 x (T^n U) \;\to_{\mathrm{lo}}\; \bigl(G_0(T^n U)\bigr)\,R_n \;\to_{\mathrm{lo}}\; \bigl((T^{n+1} U)(T^{n+1} U)\bigr)\,R_n.
\]
For any normal form $W_0 \in \NF$, set $W \equiv T W_0 \equiv S S W_0$, $R \equiv x W_0 \in \NF$, and $P_k \equiv (T^k U)\,W\,R$ ($k \ge 1$). We verify by induction on $k \ge 1$ that $P_k \to_{\mathrm{lo}}^{4k-3} D_k[U W]$ for some $D_k[\,] \not\equiv \square$: at $k = 1$, $P_1 \equiv S S U W R \to_{\mathrm{lo}} (S W\,(U W))\,R$ takes $1 = 4(1) - 3$ step; for $k = m + 1 \ge 2$, four leftmost-outermost steps give
\[
  P_{m+1} \equiv S S (T^m U)\,W\,R \;\to_{\mathrm{lo}}\; \bigl(S W\,(T^m U\,W)\bigr)\,R \;\to_{\mathrm{lo}}\; (S S W_0 R)\,P_m \;\to_{\mathrm{lo}}\; \bigl(S R\,(W_0 R)\bigr)\,P_m \;\to_{\mathrm{lo}}\; x\,W_0\,P_m\,\bigl((W_0 R)\,P_m\bigr).
\]
Because the head is the variable $x$ and the first spine argument $W_0$ lies in $\NF$, leftmost-outermost reduction proceeds inside $P_m$ for the next $4m - 3$ steps to reach $x\,W_0\,D_m[U W]\,\bigl((W_0 R)\,P_m\bigr)$ in $4 + (4m - 3) = 4(m+1) - 3$ steps. Taking $W_0 \equiv T^n U$ and $k = n + 1$ yields \eqref{eq:parametric-near-miss} in $2 + (4n + 1) = 4n + 3$ steps.
\end{proof}

\begin{remark}[Why parametric pure-$S$ self-reproduction fails to close into an open spiral]\label{rem:conj17-near-miss-barrier}
In Waldmann's Example~33 \cite[p.~95]{BarendregtEndrullisKlopWaldmann2017}, the closed normal form $S V V$ spirals after translation into the $\lambda$-calculus by reducing underneath $\lambda$-abstractions. Equation~\eqref{eq:parametric-near-miss} shows that open pure-$S$ rewriting already achieves parametric $n \mapsto n + 1$ self-embedding directly inside $\CLS$, because the ambient argument $x(T^n(V x))$ supplies the third argument needed to peel $T \equiv S S$. Theorem~\ref{thm:conj17-unconditional-and-coupling}(ii) explains why this $n \mapsto n + 1$ shift cannot be absorbed into a substitution $M \redSp C[M^\sigma]$: encoding the increment on a single variable via $\sigma(x) = T x$ over a variable-free context violates \eqref{eq:maxrs-coupling} (increasing $\mathrm{maxRS}_x$ by $\rs(T x) = 1$ requires duplicating $x$, forcing $\#_x(C) \ge 1$), whereas abstracting the tower $T^n(V x)$ into a second variable $y$ with $\sigma(y) = T y$ (as in $V x y$ or $V x (y(V x))$) replaces the active combinator $T \equiv S S$ by a passive head variable $y$, stalling head reduction before self-reproduction can occur.
\end{remark}


Finally, whereas $\mathrm{maxRS}_x$ in \eqref{eq:maxrs-coupling} can grow when $x$ is duplicated, minimizing right-spine depth over all right-maximal subterms ending in $x$---and dually minimizing left-spine arity over all left-maximal subterms headed by $x$---yields two unconditional step-by-step monotone invariants across all of $\TSV$, removing every multiplicity restriction on $M$ in Theorem~\ref{thm:conj17-unconditional-and-coupling}(ii)(b) and characterizing full Conjecture~17 purely in terms of context-free $\mathcal{DP}$-chains.

\begin{theorem}[Universal $\mathrm{minRS}_x$ and $\mathrm{minLS}_x$ monotonicity, arbitrary-multiplicity non-spiralling, and $\mathcal{DP}$ equivalence]\label{thm:conj17-minmaxrs-monotonicity}
For $P \in \TSV$ and $x \in \mathrm{Vars}(P)$, let $\mathrm{minRS}_x(P) \in \mathbb{N}$ denote the minimum of $\rs(U)$ over all right-maximal subterms $U$ of $P$ (that is, $U \equiv P$ or $U$ appears as the left child $U V \trianglelefteq P$) satisfying $\mathrm{rt}(U) \equiv x$, and let $\mathrm{minLS}_x(P) \in \mathbb{N}$ denote the minimum of $\ls(W)$ over all left-maximal subterms $W$ of $P$ ($W \equiv P$ or $V W \trianglelefteq P$) satisfying $\mathrm{hd}(W) \equiv x$.
\begin{enumerate}[label=(\roman*),leftmargin=2.2em]
  \item \emph{Variable conservation and dual universal monotonicities:} Every pure-$S$ contraction $P \redS Q$ preserves $\mathrm{Vars}(Q) = \mathrm{Vars}(P)$ identically, and for every $x \in \mathrm{Vars}(P)$ satisfies
  \begin{equation}\label{eq:minrs-minls-mono}
    \mathrm{minRS}_x(Q) \;\le\; \mathrm{minRS}_x(P) \qquad \text{and} \qquad \mathrm{minLS}_x(P) \;\le\; \mathrm{minLS}_x(Q).
  \end{equation}
  Moreover, $\mathrm{minLS}_x(P) \le \mathrm{minLS}_x(Q)$ persists across arbitrary context-stripped dependency-pair chains $P \Rightarrow_{\mathcal{DP}} Q$ whenever $x \in \mathrm{Vars}(Q)$.
  \item \emph{Arbitrary-multiplicity $1$-variable and multi-variable non-spiralling:} Let $M \in \TSV$ with arbitrary variable multiplicities $\#_v(M)$.
  \begin{enumerate}[label=(\alph*),leftmargin=1.8em]
    \item If $\mathrm{Vars}(M) \subseteq \{x\}$, $x \notin \mathrm{Vars}(C)$, and every off-right-spine occurrence of $x$ in $\sigma(x)$ has local right-spine depth at least $\mathrm{minRS}_x(M) + 1$ (in particular whenever $\#_x(\sigma(x)) \le 1$, for example $\sigma(x) = T^k x$ with $k \ge 0$), then $\neg(M \redSp C[M^\sigma])$.
    \item If $C[\,]$ is variable-free, every $v \in \mathrm{Vars}(M)$ has variable-free left off-spine branches along the right spine of $\sigma(v)$, and $\rs(\sigma(v)) \ge 1$ whenever $\sigma(v) \notin \TS$ (allowing arbitrary variable permutations and variable-merging substitutions such as $(x, y) \mapsto (T x, S x)$ or $(x, y) \mapsto (T y, S x)$), then $\neg(M \redSp C[M^\sigma])$.
  \end{enumerate}
  \item \emph{Context-free dependency-pair characterization of Conjecture~17:} Full Conjecture~17 ($\forall M, C[\,], \sigma,\; \neg(M \redSp C[M^\sigma])$ over all open or closed contexts $C[\,]$) holds if and only if no $M \in \TSV$ and $\sigma \colon \mathcal{V} \to \TSV$ satisfy both $M \Rightarrow_{\mathcal{DP}} M^\sigma$ and $\rs(M) < \rs(M^\sigma)$.
\end{enumerate}
\end{theorem}

\begin{proof}
For~(i), variable conservation $\mathrm{Vars}((X Z)(Y Z)) = \mathrm{Vars}(X) \cup \mathrm{Vars}(Y) \cup \mathrm{Vars}(Z) = \mathrm{Vars}(S X Y Z)$ is immediate because $Z$ appears on both sides. For $\mathrm{minRS}_x$, define for each right-spine offset $k \in \mathbb{N}$ the parametrized minimum $\mathrm{minRS}_x^{(k)}(T) \in \mathbb{N} \cup \{\infty\}$ by $\mathrm{minRS}_x^{(k)}(S) = \infty$, $\mathrm{minRS}_x^{(k)}(v) = k$ if $v \equiv x$ and $\infty$ otherwise, and $\mathrm{minRS}_x^{(k)}(A B) = \min(\mathrm{minRS}_x^{(0)}(A),\, \mathrm{minRS}_x^{(k+1)}(B))$; structural induction identifies $\mathrm{minRS}_x(P) = \mathrm{minRS}_x^{(0)}(P)$. At a root redex $W \equiv S X Y Z \to_r (X Z)(Y Z) \equiv W'$ with enclosing offset $k$, one computes
\[
  \mathrm{minRS}_x^{(k)}(W) \;=\; \min\bigl(\mathrm{minRS}_x^{(1)}(X),\, \mathrm{minRS}_x^{(1)}(Y),\, \mathrm{minRS}_x^{(k+1)}(Z)\bigr),
\]
whereas
\[
  \mathrm{minRS}_x^{(k)}(W') \;=\; \min\bigl(\mathrm{minRS}_x^{(0)}(X),\, \mathrm{minRS}_x^{(1)}(Z),\, \mathrm{minRS}_x^{(0)}(Y),\, \mathrm{minRS}_x^{(k+2)}(Z)\bigr).
\]
Since $\mathrm{minRS}_x^{(a)}(T) \le \mathrm{minRS}_x^{(b)}(T)$ whenever $a \le b$, we have $\mathrm{minRS}_x^{(0)}(X) \le \mathrm{minRS}_x^{(1)}(X)$, $\mathrm{minRS}_x^{(0)}(Y) \le \mathrm{minRS}_x^{(1)}(Y)$, and $\min(\mathrm{minRS}_x^{(1)}(Z), \mathrm{minRS}_x^{(k+2)}(Z)) = \mathrm{minRS}_x^{(1)}(Z) \le \mathrm{minRS}_x^{(k+1)}(Z)$, giving $\mathrm{minRS}_x^{(k)}(W') \le \mathrm{minRS}_x^{(k)}(W)$ for all $k \in \mathbb{N}$; left and right congruence steps preserve this inequality directly. Dually, defining $\mathrm{minLS}_x^{(k)}(S) = \infty$, $\mathrm{minLS}_x^{(k)}(v) = k$ if $v \equiv x$ and $\infty$ otherwise, and $\mathrm{minLS}_x^{(k)}(A B) = \min(\mathrm{minLS}_x^{(k+1)}(A),\, \mathrm{minLS}_x^{(0)}(B))$ gives $\mathrm{minLS}_x(P) = \mathrm{minLS}_x^{(0)}(P)$ and $\mathrm{minLS}_x^{(k)}(W) = \min(\mathrm{minLS}_x^{(0)}(X), \mathrm{minLS}_x^{(0)}(Y), \mathrm{minLS}_x^{(0)}(Z)) \le \min(\mathrm{minLS}_x^{(1)}(X), \mathrm{minLS}_x^{(0)}(Z), \mathrm{minLS}_x^{(1)}(Y)) \le \mathrm{minLS}_x^{(k)}(W')$ for all $k \in \mathbb{N}$. Along $P \Rightarrow_{\mathcal{DP}} Q$, non-root reductions preserve $\mathrm{minLS}_x$ monotonicity and projecting from $(X Z)(Y Z)$ to $X Z$ or $Y Z$ drops arguments from a $\min$, which cannot decrease the minimum whenever $x$ remains present in the retained branch.

For~(ii), induction on $T \in \TSV$ shows that if $\mathrm{minRS}_x^{\mathrm{off}}(\sigma(v)) \ge B$ (where $\mathrm{minRS}_x^{\mathrm{off}}(A B) = \min(\mathrm{minRS}_x^{(0)}(A), \mathrm{minRS}_x^{\mathrm{off}}(B))$ and $\mathrm{minRS}_x^{\mathrm{off}}(a) = \infty$ on atoms) for all $v \in \mathrm{Vars}(T)$, and $\rs(\sigma(v)) + \mathrm{minRS}_v^{(k)}(T) \ge B$ whenever $\mathrm{rt}(\sigma(v)) \equiv x$, then $\mathrm{minRS}_x^{(k)}(T^\sigma) \ge B$; moreover, if $x \notin \mathrm{Vars}(C)$, every $x$-occurrence in $C[U]$ lies inside $U$ at right-spine depth at least its depth in $U$, so $\mathrm{minRS}_x(C[U]) \ge \mathrm{minRS}_x(U)$. In~(a), if $M \redSp C[M^\sigma]$, Theorem~\ref{thm:conj17-dp-weighted}(ii) forces $x \in \mathrm{Vars}(M)$, $\mathrm{rt}(\sigma(x)) \equiv x$, and $\rs(\sigma(x)) \ge 1$; setting $m = \mathrm{minRS}_x(M)$ and $B = m + 1$ yields $\mathrm{minRS}_x(C[M^\sigma]) \ge \mathrm{minRS}_x(M^\sigma) \ge m + 1$, contradicting $\mathrm{minRS}_x(C[M^\sigma]) \le \mathrm{minRS}_x(M) = m$ from~\eqref{eq:minrs-minls-mono}. In~(b), collapse all variables to a single variable $x_0$ via $\pi(v) \equiv x_0$; if $M \redSp C[M^\sigma]$, then $\pi(M)$ contains $x_0$ (by Proposition~\ref{prop:acyclicity-order}(ii)), so $x_0 \in \mathrm{Vars}(\pi(C[M^\sigma]))$ and $\mathrm{minRS}_{x_0}(\pi(C[M^\sigma])) \le \mathrm{minRS}_{x_0}(\pi(M))$ by~(i), and under the hypotheses on $C[\,]$ and $\sigma$, every variable occurrence in $C[M^\sigma]$ originates at the rightmost leaf of some non-closed $\sigma(v)$ ($v \in \mathrm{Vars}(M)$) with $\rs(\sigma(v)) \ge 1$, giving $\mathrm{minRS}_{x_0}(\pi(C[M^\sigma])) \ge \mathrm{minRS}_{x_0}(\pi(M)) + 1$, a contradiction.

For~(iii), the forward implication from Conjecture~17 follows because induction on $A \Rightarrow_{\mathcal{DP}} V$ constructs a context $C[\,]$ with $A \redSs C[V]$ (wrapping $\square$ inside $(\square\,(Y Z))$ on $\mathrm{dp}_2$ and $((X Z)\,\square)$ on $\mathrm{dp}_3$), and $\rs(M) < \rs(M^\sigma) \le \rs(C[M^\sigma])$ rules out zero-step equality $M \equiv C[M^\sigma]$, forcing $M \redSp C[M^\sigma]$. Conversely, if some $M \redSp C[M^\sigma]$ existed, choosing a leaf-minimal open-spiralling witness $M$ yields $M \Rightarrow_{\mathcal{DP}} M^\sigma$ with at least one $\mathrm{dp}_1$ transition and hence $\rs(M) < \rs(M^\sigma)$ by the proof of Theorem~\ref{thm:conj17-dp-weighted}(ii).
\end{proof}

\section{Polynomial-Size Turing Machine Simulation Under Root-Restarted Selection}\label{sec:tm-sim}
% ============================================================================

Arbitrary deterministic binary Turing machines admit an explicit ongoing pure-$S$ simulation $(\mathcal{E}_{\mathrm{TM}}, \sigma_{\mathrm{TM}}, \mathcal{H}_{11}, \mathcal{D}_{\mathrm{TM}})$ with quartic initial tree size $O((|Q|^2 + |w|)^2)$, node-linear stateless tree-walking step selection, a minimal $11$-state bottom-up halting automaton, and linear-time bare-term output recovery. Following the modular simulation methodology of Cook \cite{Cook2004Rule110}, Neary and Woods \cite{NearyWoods2006SmallFast,Neary2008Thesis}, and Smith \cite{Smith2020TM23}, we factor the compiler into four layers: boundary-safe marked normalization, compilation to $U_{3,11}$, $4$-bit binary specialization to $U_{\mathrm{bin}}$, and protected spine-$2$ pure-$S$ encoding.

\subsection{Source Normalization and Compilation to Neary--Woods \texorpdfstring{$U_{3,11}$}{U\_\{3,11\}}}\label{subsec:u311}

In Neary and Woods's universal Turing machine $U_{3,11}$ \cite{NearyWoods2006SmallFast,Neary2008Thesis}, the simulated binary machine operates on a one-way infinite tape indexed by $\mathbb{N}$, never steps left from cell~$0$ during active computation, and signals halting by entering a designated final state that sweeps left until exiting cell~$0$. Our first definition compiles any partial clamping-left binary Turing machine into this exact normal form without corrupting output bits.

\begin{definition}[Source Turing machines and boundary-safe marked normalization]\label{def:source-norm}
A \emph{clamping-left binary Turing machine} is a pair $M = (Q, \delta_M)$ where $Q = \{0, \dots, |Q|-1\}$ ($|Q| \ge 1$) is a finite state set with start state $0$, and
\[
  \delta_M \colon Q \times \{0, 1\} \;\rightharpoonup\; \{0, 1\} \times \{\mathsf{L}, \mathsf{S}, \mathsf{R}\} \times Q
\]
is a partial transition function. A configuration $(q, h, \tau) \in Q \times \mathbb{N} \times \{0,1\}^*$ (with blank-padded cell lookup $\tau[i] \in \{0,1\}$, defaulting to $0$ for $i \ge |\tau|$) halts if $(q, \tau[h]) \notin \dom(\delta_M)$, and otherwise steps to $(q', h', \tau[h := b'])$ where $\delta_M(q, \tau[h]) = (b', d, q')$ and $h' = \max(h-1, 0)$ for $d = \mathsf{L}$, $h' = h$ for $d = \mathsf{S}$, and $h' = h + 1$ for $d = \mathsf{R}$. On input $w \in \{0,1\}^*$, the machine $M$ starts at $(0, 0, w)$; if $M$ reaches a halted configuration $(q_{\mathrm{fin}}, h_{\mathrm{fin}}, \tau_{\mathrm{fin}})$, its \emph{canonical output} $\out_M(w) \in \{0,1\}^*$ is $\tau_{\mathrm{fin}}$ stripped of trailing $0$-blanks.

Given $M = (Q, \delta_M)$ and $w \in \{0,1\}^*$, the \emph{boundary-safe normalized instance} $(M_{\mathrm{norm}}, w_{\mathrm{mark}})$ is constructed as follows:
\begin{enumerate}[label=(\roman*),leftmargin=2.2em]
  \item \emph{Sentinel-interleaved tape:} For $w = b_0 b_1 \cdots b_{m-1}$ ($m \ge 1$), set $w_{\mathrm{mark}} = 1\,b_0\,0\,b_1\,0\cdots 0\,b_{m-1}$, and for $w = \varepsilon$ set $w_{\mathrm{mark}} = 1\,0$. Thus $|w_{\mathrm{mark}}| \le 2|w| + 2$, cell~$0$ holds the left-boundary sentinel~$1$, every positive even cell $2i$ ($i \ge 1$) holds marker~$0$, and every odd cell $2i + 1$ ($i \ge 0$) holds data bit $w[i]$.
  \item \emph{Total $\{\mathsf{L},\mathsf{R}\}$ rule table:} Let $r = |\dom(\delta_M)| \le 2|Q|$ be the number of defined transitions of $M$, ordered lexicographically by $(q, b) \in Q \times \{0,1\}$. Then $M_{\mathrm{norm}}$ has state set $Q_{\mathrm{norm}} = \{0, \dots, S_{\mathrm{norm}}-1\}$ of cardinality
  \[
    S_{\mathrm{norm}} \;=\; |Q| + 2 + r \;\le\; 3|Q| + 2,
  \]
  consisting of an entry state $0$ (which leaves cell~$0$ unchanged and moves $\mathsf{R}$ to odd cell~$1$ in state~$1$), shifted source states $q + 1$ ($q \in Q$), dedicated auxiliary states $|Q| + 1 + \mathrm{idx}(q,b)$ executing each defined source step in two moves across the adjacent marker cell (clamping at cell~$0$ whenever the marker reads~$1$ so active steps never exit cell~$0$), and a unique halt-sweep state $S_{\mathrm{norm}} - 1$ entered on $(q,b) \notin \dom(\delta_M)$, which obeys $(S_{\mathrm{norm}}-1, b) \mapsto (b, \mathsf{L}, S_{\mathrm{norm}}-1)$ for all $b \in \{0,1\}$, preserving every cell until attempting a left move from cell~$0$.
\end{enumerate}
\end{definition}

By construction, after one initialization step $(0, 0, w_{\mathrm{mark}}) \vdash (1, 1, w_{\mathrm{mark}})$, $t$ active steps of $M$ correspond bijectively to $2t$ steps of $M_{\mathrm{norm}}$ maintaining head position $2h + 1$, sentinel $\tau_{\mathrm{norm}}[0] = 1$, even markers $\tau_{\mathrm{norm}}[2i] = 0$ ($i \ge 1$), and odd data $\tau_{\mathrm{norm}}[2i+1] = \tau[i]$, while halting at $(q, h, \tau)$ triggers an exact $(1 + 2h)$-step non-destructive left sweep to $(S_{\mathrm{norm}}-1, 0, \tau_{\mathrm{norm}})$, which exits cell~$0$ to the left.

Within the small universal Turing machine program initiated by Watanabe \cite{Watanabe1961UTM} and advanced through the tag-based machines of Rogozhin and Kudlek \cite{Rogozhin1996SmallUTM,Rogozhin1998UTM22x2,KudlekRogozhin2002UTM3x9} (surveyed in \cite{WoodsNeary2009Survey}), Neary and Woods's $3$-state, $11$-symbol machine $U_{3,11}$ \cite{NearyWoods2006SmallFast,Neary2008Thesis} achieves direct quadratic-time Turing simulation without intermediate tag-system exponential blowup.

\begin{definition}[The universal Turing machine $U_{3,11}$ and encoder $\mathcal{E}_U$]\label{def:u311}
The machine $U_{3,11}$ has state set $Q_U = \{u_1, u_2, u_3\}$ (with start state $u_1$) and $11$-symbol tape alphabet
\[
  \Sigma_{11} \;=\; \bigl\{\,\overleftarrow{a},\, \overleftarrow{b},\, e,\, h,\, \overrightarrow{e},\, \overrightarrow{h},\, \overleftarrow{e},\, \overleftarrow{h},\, \lambda,\, \delta,\, \gamma\,\bigr\},
\]
indexed $0, 1, \dots, 10$ in that order, with blank symbol $\overleftarrow{a}$ ($0$) and $32$ defined transitions displayed in Table~\ref{tab:u311}; the unique undefined halting pair is $(u_3, \delta)$.

\begin{table}[htbp]
\centering
\caption{Complete transition table of Neary and Woods's $3$-state, $11$-symbol universal Turing machine $U_{3,11}$ \cite{NearyWoods2006SmallFast,Neary2008Thesis}, where $\overleftarrow{a}$ is the blank symbol and $(u_3, \delta)$ is the unique halting configuration.}\label{tab:u311}
\small
\begin{tabular}{@{}rc@{\qquad}ccc@{}}
\toprule
Index & Symbol & State $u_1$ & State $u_2$ & State $u_3$ \\
\midrule
$0$  & $\overleftarrow{a}$ & $(\overleftarrow{e},\, \mathsf{R},\, u_1)$ & $(\gamma,\, \mathsf{L},\, u_2)$           & $(\overleftarrow{a},\, \mathsf{L},\, u_3)$ \\
$1$  & $\overleftarrow{b}$ & $(\overleftarrow{e},\, \mathsf{R},\, u_2)$ & $(\overleftarrow{b},\, \mathsf{L},\, u_1)$ & $(e,\, \mathsf{R},\, u_1)$ \\
$2$  & $e$                 & $(\overrightarrow{e},\, \mathsf{L},\, u_1)$ & $(\overleftarrow{e},\, \mathsf{R},\, u_1)$ & $(e,\, \mathsf{R},\, u_1)$ \\
$3$  & $h$                 & $(\overrightarrow{h},\, \mathsf{L},\, u_1)$ & $(\overleftarrow{h},\, \mathsf{R},\, u_3)$ & $(\overleftarrow{a},\, \mathsf{L},\, u_1)$ \\
$4$  & $\overrightarrow{e}$ & $(\overleftarrow{e},\, \mathsf{R},\, u_1)$ & $(e,\, \mathsf{R},\, u_2)$                 & $(\overleftarrow{e},\, \mathsf{R},\, u_3)$ \\
$5$  & $\overrightarrow{h}$ & $(\overleftarrow{h},\, \mathsf{R},\, u_1)$ & $(h,\, \mathsf{R},\, u_2)$                 & $(\overleftarrow{h},\, \mathsf{R},\, u_3)$ \\
$6$  & $\overleftarrow{e}$ & $(\overrightarrow{e},\, \mathsf{L},\, u_1)$ & $(\overrightarrow{e},\, \mathsf{L},\, u_2)$ & $(\gamma,\, \mathsf{L},\, u_2)$ \\
$7$  & $\overleftarrow{h}$ & $(\overrightarrow{h},\, \mathsf{L},\, u_1)$ & $(\overrightarrow{h},\, \mathsf{L},\, u_2)$ & $(\overleftarrow{a},\, \mathsf{L},\, u_3)$ \\
$8$  & $\lambda$           & $(\delta,\, \mathsf{R},\, u_1)$             & $(\lambda,\, \mathsf{R},\, u_2)$           & $(\delta,\, \mathsf{R},\, u_3)$ \\
$9$  & $\delta$            & $(\lambda,\, \mathsf{L},\, u_1)$            & $(\lambda,\, \mathsf{L},\, u_2)$           & \textrm{HALT} \\
$10$ & $\gamma$            & $(\overleftarrow{a},\, \mathsf{L},\, u_3)$ & $(\overleftarrow{h},\, \mathsf{R},\, u_3)$ & $(\overleftarrow{b},\, \mathsf{L},\, u_3)$ \\
\bottomrule
\end{tabular}
\end{table}

Given a normalized machine $M_{\mathrm{norm}}$ with $s = S_{\mathrm{norm}} \ge 3$ states and marked word $w_{\mathrm{mark}} \in \{0,1\}^*$, each transition $\theta = (b', d, q') = \delta_{\mathrm{norm}}(q, b)$ ($q, q' \in \{0,\dots,s-1\}$) is encoded as an \emph{encoded transition rule} (ETR) word over $\{e, h\}^*$: setting $\beta(\theta) = 2 + 5q' + 3\cdot\mathbf{1}_{d = \mathsf{L}}$ and $\alpha(\theta) = 5s + 2 - \beta(\theta)$, define
\[
  E(\theta) \;=\; \begin{cases}
    h^{\mathbf{1}_{b'=1}}\, e^{\alpha(\theta)}\, h^{\beta(\theta)}, & \textrm{if } d = \mathsf{R}, \\
    e^{\alpha(\theta)-1}\, h^{\beta(\theta)}\, e\, h^{\mathbf{1}_{b'=1}}\, e, & \textrm{if } d = \mathsf{L}.
  \end{cases}
\]
For each state $q \in \{0,\dots,s-1\}$, let $E'(q) = e^{5(s-q)-3}\,h^{4+5q}$ if $q = 0$ or $q$ has an incoming $\mathsf{R}$-transition in $M_{\mathrm{norm}}$, and $E'(q) = \varepsilon$ otherwise. State $q$ is represented by the $5$-block tuple (ordered from left to right on the program prefix)
\[
  \mathcal{P}(q) \;=\; \lambda\,E(\delta_{\mathrm{norm}}(q,1))\,\lambda\,E(\delta_{\mathrm{norm}}(q,0))\,\lambda\,E(\delta_{\mathrm{norm}}(q,0))\,\lambda\,E(\delta_{\mathrm{norm}}(q,1))\,\lambda\,E'(q).
\]
Encoding each right-side tape bit $b \in \{0,1\}$ via the $2$-symbol block $\langle 0 \rangle_R = \overleftarrow{a}\overleftarrow{a}$ and $\langle 1 \rangle_R = \overleftarrow{b}\overleftarrow{a}$, the compiled initial $U_{3,11}$ configuration $u_0 = \mathcal{E}_U(M, w)$ places $\mathcal{P}(s-1)\cdots\mathcal{P}(1)\mathcal{P}(0)\lambda e$ to the left of the head, scans $\overleftarrow{a}$ in start state $u_1$, and places $\overleftarrow{a}^{5s-1}\,\overleftarrow{b}\overleftarrow{b}\,\langle w_{\mathrm{mark}} \rangle_R$ to the right of the head.
\end{definition}

\begin{proposition}[Polynomial $U_{3,11}$ simulation and initial tape bound]\label{prop:u311-sim}
For every clamping-left binary Turing machine $M = (Q, \delta_M)$ and input $w \in \{0,1\}^*$, let $u_0 = \mathcal{E}_U(M, w)$ and set
\begin{equation}\label{eq:polyN-def}
  N(M, w) \;:=\; 225|Q|^2 + 390|Q| + 4|w| + 173.
\end{equation}
Then:
\begin{enumerate}[label=(\roman*),leftmargin=2.2em]
  \item \emph{Quadratic tape size:} The total non-blank length $|u_0|$ of $u_0$ satisfies $|u_0| + 1 \le N(M,w)$.
  \item \emph{Constructive halting equivalence and output recovery:} The machine $M$ halts on $w$ if and only if $U_{3,11}$ starting from $u_0$ reaches an undefined transition (necessarily at $(u_3, \delta)$). Moreover, from the halted $U_{3,11}$ tape, a single left-to-right linear scan strips the $\overleftarrow{a}^*\overleftarrow{b}^{\ge 2}\overleftarrow{a}$ state header, decodes $\langle - \rangle_R$ pairs, verifies the sentinel~$1$, and extracts odd-indexed data bits to recover $\out_M(w)$ identically.
\end{enumerate}
\end{proposition}

\begin{proof}
For~(i), let $s = S_{\mathrm{norm}}$. Every block in $\mathcal{P}(q)$ has length at most $5s + 4$, so prepending $\lambda$ gives $5s$ blocks of length at most $5s + 5$, plus the $2$-symbol suffix $\lambda e$, one scanned symbol $\overleftarrow{a}$, $(5s - 1) + 2 = 5s + 1$ state-token symbols, and $2|w_{\mathrm{mark}}|$ encoded tape symbols:
\[
  |u_0| \;\le\; 5s(5s + 5) + 2 + 1 + (5s + 1) + 2|w_{\mathrm{mark}}| \;=\; 25s^2 + 30s + 2|w_{\mathrm{mark}}| + 4.
\]
Substituting $s \le 3|Q| + 2$ and $|w_{\mathrm{mark}}| \le 2|w| + 2$ from Definition~\ref{def:source-norm} yields
\[
  |u_0| \;\le\; 25(9|Q|^2 + 12|Q| + 4) + 30(3|Q| + 2) + 2(2|w| + 2) + 4 \;=\; 225|Q|^2 + 390|Q| + 4|w| + 168,
\]
whence $|u_0| + 1 \le 225|Q|^2 + 390|Q| + 4|w| + 169 < N(M,w)$.

For~(ii), each simulated step of $M_{\mathrm{norm}}$ factors into the four sub-table cycles of \cite{NearyWoods2006SmallFast,Neary2008Thesis}: \emph{Cycle~1 (Index)} shuttles between the unary $\overleftarrow{b}^{k+2}$ state counter on the simulated head and $\lambda$-delimiters in the table, neutralizing $k$ delimiters $\lambda \mapsto \delta$ from right to left until selecting the $k$-th ETR block ($0 \le k < 5s$); \emph{Cycle~2 (Print)} copies the selected ETR block character by character ($e \mapsto \overleftarrow{a}, h \mapsto \overleftarrow{b}$) across the overlined region onto the head; \emph{Cycle~3 (Restore)} sweeps left to right in $u_2$, restoring every neutralized $\delta \mapsto \lambda$ and arrowed symbol $\overleftarrow{e}, \overleftarrow{h} \mapsto e, h$; and \emph{Cycle~4 (Choose Read/Write)} aligns the head over the adjacent simulated cell. Cutting the inductive simulation invariant \emph{immediately after Cycle~1} (the \emph{post-index invariant}) unifies all four transition geometries: at configuration $(q, h, \tau_{\mathrm{norm}})$, $U_{3,11}$ stands in state $u_2$ reading $\gamma$, with untouched blocks $k+1, \dots, 5s-1$ on the far left, the active block $\lambda\,E(\delta_{\mathrm{norm}}(q, \tau_{\mathrm{norm}}[h]))$, $k$ neutralized blocks, $\delta\,\overleftarrow{e}\,\langle \tau_{\mathrm{norm}}[0\mathinner{\ldotp\ldotp}h-1] \rangle_L\,\overleftarrow{e}^{5s+3}$ immediately to the left of $\gamma$, and $\langle \tau_{\mathrm{norm}}[h+1\mathinner{\ldotp\ldotp}] \rangle_R$ to the right of $\gamma$ (modulo trailing $\overleftarrow{a}$-blanks).

From $u_0$, if $w_{\mathrm{mark}}[0] = 1$ (which always holds by sentinel initialization, and even for arbitrary $w$ via the $E'(0)$ two-pass case when $w[0]=0$), one pass of Cycle~1 establishes the post-index invariant for $(0, 0, w_{\mathrm{mark}})$. Given the post-index invariant at $(q, h, \tau_{\mathrm{norm}})$:
\begin{itemize}[leftmargin=2.0em]
  \item If $M_{\mathrm{norm}}$ steps right to $(q', h+1, \tau')$ writing $b' \in \{0,1\}$, Cycles~2--4 print $b'$ into the left-tape encoding and place $\overleftarrow{a}^{5(s-q')}\overleftarrow{b}^{2+5q'}$ before $\langle \tau'[h+1] \rangle_R$. When $\tau'[h+1] = 1$ ($\overleftarrow{b}\overleftarrow{a}$), the extra $\overleftarrow{b}$ merges directly into $\overleftarrow{b}^{3+5q'}$, indexing block $5q'+1 = E(\delta_{\mathrm{norm}}(q',1))$ in one Cycle~1 pass; when $\tau'[h+1] = 0$ ($\overleftarrow{a}\overleftarrow{a}$, Neary's Case~D), $\overleftarrow{b}^{2+5q'}$ first indexes the auxiliary block $5q' = E'(q')$ (nonempty because $q'$ has an incoming $\mathsf{R}$-edge), whose Cycle~2--4 execution shifts two $\overleftarrow{a}$ blanks to yield $\overleftarrow{b}^{4+5q'}$, indexing block $5q'+2 = E(\delta_{\mathrm{norm}}(q',0))$ on the second Cycle~1 pass.
  \item If $M_{\mathrm{norm}}$ steps left ($h \ge 1$) to $(q', h-1, \tau')$, the block $E(\theta)$ has length $5s + 4$, so Cycle~2 consumes one extra $\overleftarrow{e}$ from the left-side cell $\langle \tau_{\mathrm{norm}}[h-1] \rangle_L \in \{\overleftarrow{e}\overleftarrow{e}, \overleftarrow{h}\overleftarrow{e}\}$; Cycle~4 then branches on $e$ vs.\ $h$ to append $0$ or $2$ extra $\overleftarrow{b}$ tokens to $\overleftarrow{b}^{5+5q'}$, indexing block $5q'+3 = E(\delta_{\mathrm{norm}}(q',0))$ or $5q'+4 = E(\delta_{\mathrm{norm}}(q',1))$ respectively.
  \item If $h = 0$ and $\delta_{\mathrm{norm}}(q, \tau_{\mathrm{norm}}[0])$ moves $\mathsf{L}$ (which occurs uniquely when $q = s - 1$ finishes its halt sweep), the extra consumption in Cycle~2 reaches the table delimiter $\delta$ directly, placing $U_{3,11}$ on $(u_3, \delta)$ with the restored table on the left and $\overleftarrow{a}^{5(s-q')-3}\,\overleftarrow{b}^{5+5q'}\,\overleftarrow{a}\,\langle 1 \rangle_R\,\langle \tau_{\mathrm{norm}}[1\mathinner{\ldotp\ldotp}] \rangle_R$ on the right, from which the linear scan reads back $\out_M(w)$.
\end{itemize}
Because every simulated step executes $m \ge 1$ defined $U_{3,11}$ transitions, constructive bounded-fuel evaluation of $M$ for $T$ steps (returning either a halting run of length at most $T$ or a live $T$-step run) combines with determinism of $U_{3,11}$ to establish bidirectional halting equivalence without invoking classical choice.
\end{proof}

\subsection{Binary Specialization \texorpdfstring{$U_{\mathrm{bin}}$}{U\_bin} and Protected Spine-2 Pure-\texorpdfstring{$S$}{S} Encoding}\label{subsec:ubin-s-encode}

To embed $U_{3,11}$ into $\CLS$ with uniform bit-level storage frames, we first specialize $U_{3,11}$ to a fixed binary Turing machine $U_{\mathrm{bin}}$ and then represent $U_{\mathrm{bin}}$ configurations as addressable spines of protected pure-$S$ frames.

\begin{construction}[Binary specialization $U_{\mathrm{bin}}$ and pure-$S$ compiler $\mathcal{E}_{\mathrm{TM}}$]\label{con:s-encode-tm}
Encode each symbol $s \in \Sigma_{11} = \{0,\dots,10\}$ as its $4$-bit big-endian binary block $\mathrm{bin}_4(s) \in \{0,1\}^4$ ($0 \mapsto 0000$, $1 \mapsto 0001$, \dots, $10 \mapsto 1010$; note $\mathrm{bin}_4(\overleftarrow{a}) = 0000$ matches the binary blank). Compiling the $32$ defined rules of $U_{3,11}$ over $4$-bit blocks yields a fixed deterministic binary Turing machine $U_{\mathrm{bin}}$ with $132$ states ($\{0,\dots,131\}$, with block-entry states $u_1 \mapsto 0$, $u_2 \mapsto 54$, $u_3 \mapsto 55$) and $262$ defined transition rules: every defined transition of $U_{3,11}$ is refined in exactly $10$ binary steps ($3$ right moves to read the $4$-bit symbol, $3$ left moves writing the new $4$-bit symbol, and $4$ moves shifting one block left or right, clamping at the left tape boundary), while the halting pair $(u_3, \delta) = (55, 1001_2)$ executes $6$ binary steps rewinding to the block start in state~$41$ scanning bit~$1$, where $(41, 1)$ is undefined.

In $\CLS$, define for any $a, b, j \in \TS$ the \emph{protected spine-$2$ frame}
\begin{equation}\label{eq:protected-frame}
  F(a, b, j) \;\equiv\; (S\,a)\,(S\,b\,j),
\end{equation}
and define counter numerals $\nu_n \in \TS$ and unopened phase generators $G(m, n) \in \TS$ ($m, n \in \mathbb{N}$) by
\begin{equation}\label{eq:numeral-generator}
  \nu_0 \;\equiv\; S T T \;\equiv\; B T, \qquad \nu_{n+1} \;\equiv\; T\,\nu_n, \qquad G(m, n) \;\equiv\; \nu_m\,\nu_n\,S.
\end{equation}
Define the depth-indexed phase sequence $p_0 = (2,2)$ and $p_{d+1} = (n, n+1)$ if $p_d = (0, n)$, else $(m-1, n)$ if $p_d = (m, n)$ with $m \ge 1$; every $p_d = (m_d, n_d)$ satisfies $2 \le n_d$ and $m_d \le n_d$. Set the depth-$d$ unopened base $\beta_d \equiv G(m_d, n_d)$ and junk term $\iota_d \equiv (\nu_{n_d - 1} S)\,G(n_d, n_d + 1)$ if $m_d = 0$, else $(\nu_{n_d - 2} S)\,G(m_d - 1, n_d)$.
\begin{enumerate}[label=(\roman*),leftmargin=2.2em]
  \item \emph{Bit lists and append-only version chains:} At tree depth $d$, a Boolean bit $b \in \{0,1\}$ is encoded as $\beta_d$ if $b = 0$ and $F(\beta_{d+1}, \beta_{d+1}, \iota_d)$ if $b = 1$. A fixed-width bit tuple $v = (b_1,\dots,b_w) \in \{0,1\}^w$ is encoded as a nested frame list ending in sentinel $F(\beta, \beta, \iota)$, and a history of committed $w$-bit records $[v_1, \dots, v_r]$ is encoded along the left spine of $F$-frames so that scanning left children from the root returns the newest complete $w$-bit record $v_r$ while ignoring any partially opened trailing record.
  \item \emph{$11$-bit universal tape cells and root hub:} Each binary cell of $U_{\mathrm{bin}}$ is represented by a width-$11$ version chain encoding $(b, \mathrm{st}, q) \in \{0,1\} \times \{\mathsf{off}, \mathsf{active}, \mathsf{movingL}, \mathsf{movingR}\} \times \{0,\dots,255\}$ ($1$ data bit, $2$ status bits, $8$ state bits). Given $(M,w)$ with $u_0 = \mathcal{E}_U(M,w)$, the $4|u_0|$ binary cells of $u_0$ are initialized with one $\mathsf{active}$ record at bit $4 h(u_0)$ in start state $0$ and $\mathsf{off}$ records elsewhere, chained along the right spine of $F$-frames to form $\mathrm{Tape}_0$, and wrapped in the root hub
  \[
    \mathcal{E}_{\mathrm{TM}}(M, w) \;\equiv\; (\nu_1\,S)\,F\bigl(\beta_1,\, \mathrm{Tape}_0,\, \iota_0\bigr).
  \]
\end{enumerate}
\end{construction}

\begin{lemma}[Protected-frame opening laws and initial quartic leaf bound]\label{lem:frame-open-bound}
The encoding of Construction~\ref{con:s-encode-tm} satisfies:
\begin{enumerate}[label=(\roman*),leftmargin=2.2em]
  \item \emph{Root inertness and controlled opening:} Every frame $F(a,b,j)$ has $\ls(F(a,b,j)) = 2$ (hence no root redex); if applied externally to an argument $x$, it opens in two steps via $F(a,b,j)\,x \redS^2 (a x)((b x)(j x))$. Internally without an external argument, for all $m, n \in \mathbb{N}$, the unopened generator $G(m+1, n+2)$ reduces in $6$ steps to
  \[
    G(m+1, n+2) \;\redS^6\; F\bigl(G(m, n+2),\, G(m, n+2),\, (\nu_n S)\,G(m, n+2)\bigr),
  \]
  and $G(0, n+2)$ reduces in $7$ steps to $F(G(n+2, n+3), G(n+2, n+3), (\nu_{n+1} S)\,G(n+2, n+3))$, with every intermediate reduct not of $F$-shape and every generated junk term satisfying $\ls(j) \ne 2$ with head $S$ (that is, $j$ never has the shape $S u v$).
  \item \emph{Quartic initial leaf bound:} For every clamping-left binary Turing machine $M = (Q, \delta_M)$ and input $w \in \{0,1\}^*$, the closed $S$-term $T_0 = \mathcal{E}_{\mathrm{TM}}(M, w)$ satisfies
  \begin{equation}\label{eq:tm-init-bound}
    |\mathcal{E}_{\mathrm{TM}}(M, w)| \;\le\; 125{,}000\,N(M, w)^2,
  \end{equation}
  where $N(M,w) = 225|Q|^2 + 390|Q| + 4|w| + 173$.
\end{enumerate}
\end{lemma}

\begin{proof}
For~(i), $\ls(S a (S b j)) = 1 + \ls(S a) = 2$, preventing any root contraction; applying $x$ yields $S a (S b j) x \to_r (a x)(S b j x) \redS (a x)((b x)(j x))$. For internal generator opening, let $N_1 \equiv T \nu_n$, $N_2 \equiv T N_1 \equiv \nu_{n+2}$, $Q_0 \equiv \nu_m N_2$, and $g \equiv \nu_m N_2 S \equiv G(m, n+2)$. Then $G(m+1, n+2) \equiv (T \nu_m) N_2 S$ reduces in $6$ deterministic steps:
\begin{align*}
  (S S \nu_m)\,N_2\,S &\;\redS\; (S N_2 Q_0)\,S \;\to_r\; (N_2 S)\,g \;\equiv\; (S S N_1 S)\,g \;\redS\; \bigl(T(N_1 S)\bigr)\,g \\
  &\;\to_r\; (S g)\bigl((N_1 S) g\bigr) \;\equiv\; (S g)\bigl((S S \nu_n S) g\bigr) \;\redS\; (S g)\bigl((T(\nu_n S)) g\bigr) \;\redS\; F\bigl(g, g, (\nu_n S) g\bigr).
\end{align*}
When $m = 0$, the single preliminary step $G(0, n+2) \equiv (S T T)\,\nu_{n+2}\,S \redS (T \nu_{n+2})(T \nu_{n+2}) S \equiv G(n+3, n+3)$ reduces to the positive case $(m', n') = (n+2, n+1)$, giving $1 + 6 = 7$ steps. Direct inspection of the $6$ (respectively, $7$) pre-final terms confirms none matches $(S u)(S v w)$, and $(\nu_k S) g$ has left child $\nu_k S$ where $\nu_k \not\equiv S$, so $(\nu_k S) g$ never has the shape $S u v$.

For~(ii), induction on $n$ gives $|\nu_n| = 5 + 2n$, $|G(m,n)| = 11 + 2m + 2n$, and $|(\nu_k S)\,G(m,n)| = 17 + 2k + 2m + 2n$. Let $K$ be the number of opened $F$-frames in the layout tree of $\mathrm{Tape}_0$, and let $H$ bound $n_d + h$, where $h$ is the remaining height of the layout tree; since $p_1 = (1,2)$ and each depth step increases $n_d$ by at most $1$, we have $H \le 2 + K$, and structural induction on the layout tree proves $|\mathrm{Tape}_0| \le (2K + 1)(19 + 6H) \le (2K + 1)(31 + 6K)$. Adding $|\nu_1 S| = 8$, $|\beta_1| = 17$, $|\iota_0| = 23$, and the two $S$-leaves of the outer hub frame gives $|T_0| \le 50 + (2K + 1)(31 + 6K)$. In $\mathrm{Tape}_0$, there are $4|u_0|$ binary cells, each containing one initial $11$-bit record, and each $11$-bit record opens at most $2 \cdot 11 + 2 = 24$ frames plus $1$ spine frame per cell, so $K \le (1 + 24) \cdot 4|u_0| = 100|u_0|$. Since $2(100|u_0|) + 1 \le 200(|u_0| + 1)$ and $31 + 600|u_0| \le 600(|u_0| + 1)$, we obtain
\[
  |T_0| \;\le\; 50 + 120{,}000(|u_0| + 1)^2 \;\le\; 125{,}000(|u_0| + 1)^2 \;\le\; 125{,}000\,N(M, w)^2,
\]
using Proposition~\ref{prop:u311-sim}(i).
\end{proof}

\subsection{Stateless Root-Restarted Zipper Tree-Walking Selector}\label{subsec:zipper-selector}

To ensure that evaluation-order selection retains no hidden memory between pure-$S$ contractions and inspects only bounded local tree neighborhoods, we formalize the selector $\sigma_{\mathrm{TM}}$ as a deterministic root-restarted finite-state tree-walking automaton operating on a Huet zipper \cite{Huet1997}.

\begin{definition}[Zipper cursor and root-restarted finite-state selector $\sigma_{\mathrm{TM}}$]\label{def:zipper-selector}
A \emph{zipper cursor} over $\TS$ is a pair $(f, \gamma)$ consisting of the focused subterm $f \in \TS$ and a breadcrumb stack $\gamma \in \{\mathsf{L}(R), \mathsf{R}(L) : L, R \in \TS\}^*$ recording sibling subtrees along the path back to the root, with root cursor $\mathrm{root}(t) = (t, \varepsilon)$ and tree reconstruction $\mathrm{rebuild}(f, \varepsilon) = f$, $\mathrm{rebuild}(f, \mathsf{L}(R)\cdot\gamma) = \mathrm{rebuild}(f R, \gamma)$, and $\mathrm{rebuild}(f, \mathsf{R}(L)\cdot\gamma) = \mathrm{rebuild}(L f, \gamma)$. At each micro-tick, the automaton observes only the $6$-element local observation
\[
  \mathrm{obs}(f, \gamma) \;\in\; \{\mathsf{S},\, \mathsf{App}\} \times \{\mathsf{root},\, \mathsf{left},\, \mathsf{right}\},
\]
indicating whether $f \equiv S$ or $f \equiv f_1 f_2$ and whether the top of $\gamma$ is empty, $\mathsf{L}(\_)$, or $\mathsf{R}(\_)$, and issues one of five primitive commands:
\[
  \mathsf{moveL}(q'), \qquad \mathsf{moveR}(q'), \qquad \mathsf{moveU}(q'), \qquad \mathsf{contract}, \qquad \mathsf{normal}.
\]
Here $\mathsf{moveL}, \mathsf{moveR}, \mathsf{moveU}$ step one tree edge down-left, down-right, or up to the parent in state $q'$; $\mathsf{contract}$ contracts a root redex $f \equiv S X Y Z \mapsto (X Z)(Y Z)$ at the current focus, rebuilds the tree, and terminates the invocation; and $\mathsf{normal}$ terminates without modifying the tree.

The finite control state space $Q_{\mathrm{ctrl}}$ of $\sigma_{\mathrm{TM}}$ is a loop-free product of bounded finite registers:
\begin{enumerate}[label=(\roman*),leftmargin=2.2em]
  \item a $23$-state local strict-frame lookahead register probing at most $16$ micro-ticks to verify that a candidate node has shape $F(a,b,j)$ with $\ls(j) \ne 2$ or $j$ not headed by $S$ (enforcing unambiguous frame parsing without false positives from junk subtrees);
  \item width-$11$ shift buffers ($\{0,1\}^{\le 11}$) and decoded cell registers $\{0,1\} \times \{\mathsf{off}, \mathsf{active}, \mathsf{movingL}, \mathsf{movingR}\} \times \{0,\dots,255\}$ implementing a two-pass tape scan: \emph{Pass~1} scans the tape spine for any cell in status $\mathsf{movingL}$ or $\mathsf{movingR}$ and completes its neighbor arrival and source reset to $\mathsf{off}$; if none is pending, \emph{Pass~2} locates the unique $\mathsf{active}$ cell $(b, \mathsf{active}, q)$ and consults the fixed $262$-rule table of $U_{\mathrm{bin}}$ (opening the hub halt flag if $(q, b)$ is undefined, or appending the next $11$-bit record, or extending the post-halt left-spine frontier if the halt flag is already open);
  \item a $28$-state local generator-opening probe inspecting depth at most $5$ to select the unique next redex of the $6$- or $7$-step sequence of Lemma~\ref{lem:frame-open-bound}(i); and
  \item a $10$-state root-ascent and leftmost-redex Euler-tour fallback invoked if any structural guard fails on a malformed input tree.
\end{enumerate}
(Flattening these finite product registers into a single enumerated set yields $|Q_{\mathrm{ctrl}}| = 4{,}376{,}490{,}905$ states.) Every macro-step of $\sigma_{\mathrm{TM}}$ starts fresh at $(q_0, \mathrm{root}(t))$, retaining no state across contractions.
\end{definition}

\begin{lemma}[Universal node-linear micro-tick termination and soundness]\label{lem:microtick-bound}
For every closed $S$-term $t \in \TS$, executing the tree-walking controller of Definition~\ref{def:zipper-selector} from $(q_0, \mathrm{root}(t))$ terminates after $m$ elementary micro-ticks satisfying
\begin{equation}\label{eq:microtick-bound}
  1 \;\le\; m \;\le\; 3500\,(2|t| - 1) \;=\; 3500\,\nodes(t),
\end{equation}
either returning $\sigma_{\mathrm{TM}}(t) = \mathsf{normal}$ (in which case $t \in \NF$) or returning $\sigma_{\mathrm{TM}}(t) = t' \in \TS$ obtained by contracting a single redex ($t \redS t'$).
\end{lemma}

\begin{proof}
If $t \equiv S$, then $|t| = 1$ and $(q_0, \mathrm{root}(S))$ halts in $m = 1 \le 3500(2\cdot 1 - 1)$ micro-tick returning $\mathsf{normal}$. For composite $t \equiv t_1 t_2$ ($|t| \ge 2$, so $|t| + 1 \le 2|t| - 1$), every subroutine of $Q_{\mathrm{ctrl}}$ traverses disjoint subtrees of the current focus without revisiting siblings inside the same pass: each strict-frame lookahead takes at most $16$ micro-ticks and restores the cursor to its entry node; reading or verifying an $11$-bit cell version chain descends along the left spine of committed records (at most $16 + 22\cdot 18$ micro-ticks per record, amortized against the disjoint leaf mass of the record's field subtrees) and ascends to the cell header; Pass~1 and Pass~2 each sweep down the right tape spine at most once, followed by at most one local neighbor step, one version-write descent, and one $5$-edge generator probe before contracting. If any structural invariant fails on an arbitrary non-canonical tree $t$, the controller ascends directly to the root (in at most $|t|$ moves, since breadcrumb length satisfies $|\gamma| < |t|$) and performs a standard depth-first Euler tour tracking left-spine arity in $\{0,1,2,3\}$: every edge of $t$ ($2|t| - 2$ edges) is traversed at most twice (once downward, once upward), halting with $\mathsf{normal}$ at the root if no node reaches left-spine arity~$3$, or climbing $3$ edges to contract the first $S X Y Z$ redex encountered. Summing the explicit worst-case phase potentials across Pass~1, root rewind, Pass~2, destination/write traversal, root ascent, and Euler fallback bounds total micro-ticks by $3500(|t| + 1) \le 3500(2|t| - 1)$. Finally, $\mathsf{contract}$ fires only after verifying three consecutive left-child edges above an $S$-leaf, so rebuilding the zipper produces a valid single-step reduct $t \redS t'$, whereas $\mathsf{normal}$ fires only after a complete redex-free Euler tour.
\end{proof}

\subsection{Minimal 11-State Halting Automaton, Bare-Term Output Readback, and Master Theorem}\label{subsec:tm-master}

When $U_{\mathrm{bin}}$ reaches its undefined halting configuration $(41, 1)$ on a compiled term, Pass~2 of $\sigma_{\mathrm{TM}}$ opens the root hub's halt flag from $\beta_1$ to an opened frame $F(\beta_2, \beta_2, \iota_1)$, and all subsequent invocations extend a dedicated post-halt frontier inside the halt flag while leaving $\mathrm{Tape}$ untouched. Thus a bare term $t \in \TS$ is in a root-halted state ($\mathsf{rootHalt}(t) = 1$) if and only if $t \equiv H_{\mathrm{hdr}} \cdot \mathrm{Hub}$ where $\mathrm{Hub} \equiv F(\mathrm{Flag}, \mathrm{Tape}, j)$ is an opened frame whose left slot $\mathrm{Flag} \equiv F(l, r, j')$ is also an opened frame.

\begin{proposition}[Minimal $11$-state halting automaton $\mathcal{H}_{11}$ and bare-term decoder $\mathcal{D}_{\mathrm{TM}}$]\label{prop:halt-automaton}
Let $L_{\mathrm{halt}} = \{t \in \TS : \mathsf{rootHalt}(t) = 1\}$.
\begin{enumerate}[label=(\roman*),leftmargin=2.2em]
  \item \emph{Exact recognition and $11$-state minimality:} The deterministic bottom-up tree automaton $\mathcal{H}_{11} = (\{0,\dots,10\}, 0, \odot_{\mathrm{halt}}, \{3, 6, 8, 10\})$ defined by the $7$-feature vector table in Table~\ref{tab:h11} satisfies $L(\mathcal{H}_{11}) = L_{\mathrm{halt}}$. Moreover, $\mathcal{H}_{11}$ is the unique minimal deterministic bottom-up tree automaton recognizing $L_{\mathrm{halt}}$ over $\TS$: no DTA with $10$ or fewer states recognizes $L_{\mathrm{halt}}$.
  \item \emph{Linear-time bare-term output decoder:} There exists a read-only structural parser $\mathcal{D}_{\mathrm{TM}} \colon \TS \rightharpoonup \{0,1\}^*$ terminating in $O(|t|)$ primitive tree operations on every $t \in \TS$ that inspects neither $M$ nor $w$ and extracts the halted output directly from $\mathrm{Tape}$.
\end{enumerate}
\end{proposition}

\begin{table}[htbp]
\centering
\caption{States, Boolean feature vectors $(\mathsf{s}_0, \mathsf{s}_1, \mathsf{s}_1\mathsf{F}, \mathsf{s}_2, \mathsf{frame}, \mathsf{haltHub}, \mathsf{rootHalt})$, and minimal closed representatives of the $11$-state halting tree automaton $\mathcal{H}_{11}$, where $T \equiv S S$, $A \equiv S S S$, $F_0 \equiv T A$, and $H_0 \equiv (S F_0) A$.}\label{tab:h11}
\small
\begin{tabular}{@{}rcccl@{}}
\toprule
State $q$ & Feature vector & Accepting & Representative $t_q \in \TS$ & Syntactic role \\
\midrule
$0$  & $(1,0,0,0,0,0,0)$ & no  & $S$            & Atomic leaf $S$ \\
$1$  & $(0,1,0,0,0,0,0)$ & no  & $T$            & $S X$ with $X$ not a frame \\
$2$  & $(0,1,1,0,0,0,0)$ & no  & $S F_0$        & $S F$ with $F$ a frame, non-halting right child \\
$3$  & $(0,1,1,0,0,0,1)$ & yes & $S H_0$        & $S H$ where $H$ is a halted hub \\
$4$  & $(0,0,0,1,0,0,0)$ & no  & $A$            & $(S X) Y$ not of frame shape \\
$5$  & $(0,0,0,1,1,0,0)$ & no  & $F_0$          & Frame $F(a,b,j)$ with unopened flag $a$ \\
$6$  & $(0,0,0,1,1,0,1)$ & yes & $T H_0$        & Frame whose right child is a halted hub \\
$7$  & $(0,0,0,1,1,1,0)$ & no  & $H_0$          & Halted hub $F(F', b, j)$, non-halting right child \\
$8$  & $(0,0,0,1,1,1,1)$ & yes & $(S F_0) H_0$  & Halted hub whose right child is a halted hub \\
$9$  & $(0,0,0,0,0,0,0)$ & no  & $A S$          & Other composite term, non-halting \\
$10$ & $(0,0,0,0,0,0,1)$ & yes & $A H_0$        & Composite root $P H$ with halted hub $H$ \\
\bottomrule
\end{tabular}
\end{table}

\begin{proof}
For~(i), associate to every $t \in \TS$ the $7$-tuple $\mathbf{f}(t) = (\mathsf{s}_0, \mathsf{s}_1, \mathsf{s}_1\mathsf{F}, \mathsf{s}_2, \mathsf{frame}, \mathsf{haltHub}, \mathsf{rootHalt}) \in \{0,1\}^7$, where the seven Boolean coordinates hold as follows: $\mathsf{s}_0(t) = 1$ if and only if $t \equiv S$; $\mathsf{s}_1(t) = 1$ if and only if $t \equiv S x$ for some $x$; $\mathsf{s}_1\mathsf{F}(t) = 1$ if and only if $t \equiv S x$ for some $x$ with $\mathsf{frame}(x) = 1$; $\mathsf{s}_2(t) = 1$ if and only if $t \equiv S x y$ for some $x, y$; $\mathsf{frame}(t) = 1$ if and only if $t \equiv (S a)(S b j)$ for some $a, b, j$; $\mathsf{haltHub}(t) = 1$ if and only if $t \equiv (S a)(S b j)$ for some $a, b, j$ with $\mathsf{frame}(a) = 1$; and $\mathsf{rootHalt}(t) = 1$ if and only if $t \equiv p q$ for some $p, q$ with $\mathsf{haltHub}(q) = 1$. For any binary application $l r$, these seven coordinates satisfy the purely compositional law (writing $\land$ for multiplication on $\{0,1\}$)
\begin{equation}\label{eq:h11-combine}
  \mathbf{f}(l r) \;=\; \bigl(0,\; \mathsf{s}_0(l),\; \mathsf{s}_0(l)\land\mathsf{frame}(r),\; \mathsf{s}_1(l),\; \mathsf{s}_1(l)\land\mathsf{s}_2(r),\; \mathsf{s}_1\mathsf{F}(l)\land\mathsf{s}_2(r),\; \mathsf{haltHub}(r)\bigr).
\end{equation}
Direct inspection verifies that the $11$ vectors listed in Table~\ref{tab:h11} contain $\mathbf{f}(S) = (1,0,0,0,0,0,0)$, are closed under \eqref{eq:h11-combine} across all $11 \times 11 = 121$ pairs (defining $q_l \odot_{\mathrm{halt}} q_r$), and accept via coordinate~$7$ ($\{3,6,8,10\}$), so structural induction shows that $\sem{t}_{\mathcal{H}_{11}} = q$ holds if and only if $\mathbf{f}(t) = \mathbf{v}_q$, establishing $L(\mathcal{H}_{11}) = L_{\mathrm{halt}}$.

To prove $11$-state Myhill--Nerode minimality, first note that all $11$ states are inhabited on $\TS$ by the explicit closed representatives $t_0, \dots, t_{10}$ in Table~\ref{tab:h11}. Second, using \eqref{eq:h11-combine} with $\mathsf{s}_2(A) = 1$, $\mathsf{frame}(F_0) = 1$, and $\mathsf{s}_1\mathsf{F}(S F_0) = 1$, each of the seven Boolean coordinates of $\mathbf{f}(t)$ is isolated as root acceptance $C_i[t] \in L_{\mathrm{halt}}$ by a fixed closed context $C_i[\square]$:
\begin{align*}
  C_7[\square] &\equiv \square, &
  C_6[\square] &\equiv S\,\square, &
  C_4[\square] &\equiv S\bigl((S F_0)\,\square\bigr), &
  C_3[\square] &\equiv S(\square\,A), \\
  C_5[\square] &\equiv S\bigl((S\,\square)\,A\bigr), &
  C_2[\square] &\equiv S\bigl((S F_0)(\square\,S)\bigr), &
  C_1[\square] &\equiv S\bigl((\square\,F_0)\,A\bigr).
\end{align*}
Indeed, $C_6[t] \in L_{\mathrm{halt}}$ holds if and only if $\mathsf{haltHub}(t) = 1$; $C_4[t] \in L_{\mathrm{halt}}$ holds if and only if $\mathsf{s}_1\mathsf{F}(S F_0) \land \mathsf{s}_2(t) = 1$, which simplifies to $\mathsf{s}_2(t) = 1$; $C_3[t] \in L_{\mathrm{halt}}$ holds if and only if $\mathsf{s}_1\mathsf{F}(t) \land \mathsf{s}_2(A) = 1$, which simplifies to $\mathsf{s}_1\mathsf{F}(t) = 1$; and $C_5, C_2, C_1$ compose $C_3$ and $C_4$ with $S\,\square$, $\square\,S$, and $\square\,F_0$ to test $\mathsf{frame}(t) = 1$, $\mathsf{s}_1(t) = 1$, and $\mathsf{s}_0(t) = 1$, respectively. Since the $11$ rows of Table~\ref{tab:h11} are pairwise distinct in $\{0,1\}^7$, any two distinct states $p \ne q$ differ at some coordinate $i \in \{1,\dots,7\}$, so exactly one of $C_i[t_p]$ and $C_i[t_q]$ belongs to $L_{\mathrm{halt}}$. By the Myhill--Nerode theorem for tree languages, every deterministic bottom-up tree automaton recognizing $L_{\mathrm{halt}}$ has at least $11$ states.

For~(ii), given $t \in \TS$, the decoder $\mathcal{D}_{\mathrm{TM}}(t)$ parses $t \equiv H_{\mathrm{hdr}} \cdot F(\mathrm{Flag}, \mathrm{Tape}, j)$, checks $\mathsf{frame}(\mathrm{Flag}) = 1$, walks the right spine of $\mathrm{Tape}$ once to extract the newest committed $11$-bit record of each cell, verifies that no cell has status $\mathsf{movingL}$ or $\mathsf{movingR}$ and that the first $\mathsf{active}$ cell lies at bit index $4 u_{\mathrm{head}} \equiv 0 \pmod 4$, groups the cell data bits into $4$-bit blocks $\mathrm{bin}_4^{-1} \in \Sigma_{11}$ (verifying that cell $u_{\mathrm{head}}$ holds $\delta$ in $u_3$), and applies the single-pass $U_{3,11}$ output scan of Proposition~\ref{prop:u311-sim}(ii). Every subterm of $t$ is visited at most once, taking $O(|t|)$ time.
\end{proof}

Combining Proposition~\ref{prop:u311-sim}, Construction~\ref{con:s-encode-tm}, Lemmas~\ref{lem:frame-open-bound} and~\ref{lem:microtick-bound}, and Proposition~\ref{prop:halt-automaton} yields our first main universality theorem.

\begin{theorem}[Controlled polynomial-size Turing machine simulation in $\CLS$]\label{thm:tm-universality}
The quadruple $(\mathcal{E}_{\mathrm{TM}}, \sigma_{\mathrm{TM}}, \mathcal{H}_{11}, \mathcal{D}_{\mathrm{TM}})$ is an ongoing selected-path simulation of all clamping-left deterministic binary Turing machines in $\CLS$. Specifically, for every machine $M = (Q, \delta_M)$ and input $w \in \{0,1\}^*$, let $N(M,w) = 225|Q|^2 + 390|Q| + 4|w| + 173$ and $T_0 = \mathcal{E}_{\mathrm{TM}}(M, w) \in \TS$. Then:
\begin{enumerate}[label=(\roman*),leftmargin=2.2em]
  \item \emph{Initial size and perpetual single-step trajectory:} $|T_0| \le 125{,}000\,N(M,w)^2$, and iterating $\sigma_{\mathrm{TM}}$ generates an infinite pure-$S$ reduction sequence $T_0 \redS T_1 \redS T_2 \redS \cdots$ ($T_{k+1} = \sigma_{\mathrm{TM}}(T_k)$) in which each step is selected from $(q_0, \mathrm{root}(T_k))$ within $3500(2|T_k|-1)$ micro-ticks and every term satisfies the quadratic trajectory bound
  \begin{equation}\label{eq:tm-traj-bound}
    |T_k| \;\le\; 600{,}000\,\bigl(N(M,w) + k\bigr)^2 \qquad (\text{for all } k \ge 0).
  \end{equation}
  \item \emph{Bidirectional halting equivalence and output recovery:} The machine $M$ halts on $w$ if and only if there exists $k \ge 0$ such that $T_k \in L(\mathcal{H}_{11})$; moreover, whenever $M$ halts on $w$, every $k \ge 0$ with $T_k \in L(\mathcal{H}_{11})$ satisfies $\mathcal{D}_{\mathrm{TM}}(T_k) = \out_M(w)$.
\end{enumerate}
\end{theorem}

\begin{proof}
By Lemma~\ref{lem:frame-open-bound}(ii) and Lemma~\ref{lem:microtick-bound}, $|T_0| \le 125{,}000\,N(M,w)^2$ and each invocation of $\sigma_{\mathrm{TM}}$ halts in at most $3500(2|T_k|-1)$ micro-ticks. On every reachable term from $T_0$, the strict-frame invariant holds throughout the hub, so the local $16$-tick frame checks of $\sigma_{\mathrm{TM}}$ agree identically with structural frame parsing. Each binary transition of $U_{\mathrm{bin}}$ adds at most $1$ new tape cell and at most $3$ committed $11$-bit cell records (for moving source, arrival destination, and source reset), so after $t$ binary transitions of $U_{\mathrm{bin}}$, the opened-frame count satisfies $K_t \le 100(|u_0| + t)$ and the checkpoint term size is at most $125{,}000(|u_0| + t + 1)^2$. Because every pure-$S$ contraction has non-decreasing leaf count by Proposition~\ref{prop:cr-conservation}(ii), any intermediate term $T_k$ along a $6$- or $7$-step generator opening is bounded in leaf size by the next macro-boundary checkpoint: if $U_{\mathrm{bin}}$ runs for $k+1$ binary transitions without halting, we have $K \le 100(|u_0| + k + 1)$ and $(|u_0| + k + 2)^2 \le 4(|u_0| + k + 1)^2 \le 4(N(M,w) + k)^2$, giving $|T_k| \le 4 \cdot 125{,}000(N(M,w)+k)^2 = 500{,}000(N(M,w)+k)^2$; if $U_{\mathrm{bin}}$ halts at binary step $t_{\mathrm{halt}} \le k$, the tape freezes at size at most $500{,}000(N(M,w)+k)^2$ while at most $k+1$ post-halt frontier openings contribute at most $800(N(M,w)+k)^2 + 33$ leaves, yielding $|T_k| \le 600{,}000(N(M,w)+k)^2$ in all cases.

For~(ii), while $U_{\mathrm{bin}}$ executes defined transitions, $\mathrm{Flag} \equiv \beta_1$ remains unopened, so $T_k \notin L(\mathcal{H}_{11})$. Constructive bounded evaluation of $U_{\mathrm{bin}}$ shows that $U_{\mathrm{bin}}$ halts on $\mathcal{E}_{\mathrm{bin}}(u_0)$ if and only if some $T_k$ opens $\mathrm{Flag}$ and enters $L(\mathcal{H}_{11})$. By the $10$-step defined / $6$-step halting refinement between $U_{3,11}$ and $U_{\mathrm{bin}}$ (Construction~\ref{con:s-encode-tm}) and Proposition~\ref{prop:u311-sim}(ii), this holds if and only if $M$ halts on $w$. At the halting boundary $u_{\mathrm{fin}} = (u_3, \delta)$ of $U_{3,11}$, the $6$-step rewind of $U_{\mathrm{bin}}$ leaves the active head aligned at bit $4 h(u_{\mathrm{fin}})$ with no pending moves and all $4$-bit symbol blocks intact; since post-halt steps modify only $\mathrm{Flag}$, every $T_k \in L(\mathcal{H}_{11})$ decodes under $\mathcal{D}_{\mathrm{TM}}$ to $u_{\mathrm{fin}}$ modulo trailing blanks and hence to $\out_M(w)$.
\end{proof}

\begin{remark}[Sub-universal complexity and concrete compiled sizes]\label{rem:tm-complexity}
Formatting $\mathcal{E}_U(M,w)$ requires only $O(\log(|Q| + |w|))$ workspace \cite{NearyWoods2006SmallFast}, and expanding $u_0$ into the unshared syntax tree $T_0 = \mathcal{E}_{\mathrm{TM}}(M,w)$ is a deterministic uniform block substitution running in time linear in $|T_0|$, while $\mathcal{H}_{11}$ is an $11$-state finite tree automaton and $\mathcal{D}_{\mathrm{TM}}$ runs in read-only time $O(|T_k|)$. In concrete evaluations, compiling the $1$-state halting Turing machine on the empty input ($|Q|=1, |w|=0$, where $|u_0| = 540$ versus the upper bound $N = 788$) produces an explicit pure-$S$ tree of $1{,}996{,}921$ leaves; compiling a $2$-state write-$1$ machine produces $7{,}477{,}337$ leaves; and simulating the $2$-state Busy Beaver executes $912$ pure-$S$ contractions to enter $L(\mathcal{H}_{11})$ and output $1111$, contrasting with the unmaterializable seeds of at least $2^{57207}$ bits required by unary counter/tag pipelines \cite{Filin2026}.
\end{remark}

% ============================================================================

% ============================================================================
\section{Program-Free Cyclic Tag Simulation, Ten-Leaf Seed Genesis, and the Linear-Time Pocket Tree Transducer}\label{sec:cts-sim}
% ============================================================================

In Section~\ref{sec:tm-sim}, the root-restarted selector $\sigma_{\mathrm{TM}}$ incorporates the fixed $262$-rule binary transition table of $U_{\mathrm{bin}}$ inside its finite control. One naturally asks whether this embedded universal transition table can be eliminated altogether, so that both the simulated program and its data reside inside the closed $S$-term while a single program-free selector executes arbitrary computations purely by structural manipulation of $\TS$. Here we answer this question affirmatively via cyclic tag systems \cite{Cook2004Rule110}, prove that every encoded instance unfolds by pure-$S$ reduction from a single $10$-leaf combinator $D_{\mathrm{seed}}$, and compress the resulting program-free selector into an explicit two-pass contextual finite-state tree transducer $\mathcal{T}_{\mathrm{pocket}}$ running in at most $2|M| - 1 + \mathrm{depth}(M)$ transitions per invocation.

\subsection{Cyclic Tag Systems and Self-Describing Pure-\texorpdfstring{$S$}{S} Queue Encoding}\label{subsec:cts-encode}

Introduced by Cook \cite{Cook2004Rule110} as a binary refinement of Post's canonical and tag systems \cite{Post1936FormulationI,Post1943Reductions} and Minsky and Cocke's universal $2$-tag machines \cite{Minsky1961Tag,CockeMinsky1964Tag}, and later shown by Woods and Neary \cite{WoodsNeary2006TagFOCS} to simulate Turing machines with polynomial time overhead, cyclic tag systems dispense with explicit control states in favor of a cyclically rotating ring of binary appendants.

\begin{definition}[Cyclic tag systems]\label{def:cts}
A \emph{cyclic tag system} (CTS) is a nonempty finite tuple $P = (A_0, A_1, \dots, A_{m-1})$ of $m \ge 1$ binary words $A_i \in \{0,1\}^*$, called \emph{appendants}. A configuration of $P$ is a pair $(i, u) \in \{0, \dots, m-1\} \times \{0,1\}^*$ consisting of the active phase index $i$ and the FIFO data queue $u$. On initial word $w \in \{0,1\}^*$, execution begins at $(0, w)$ and evolves deterministically according to
\[
  (i,\, \varepsilon) \;\text{halts}, \qquad
  (i,\, 0\cdot u') \;\vdash_P\; \bigl((i + 1) \bmod m,\, u'\bigr), \qquad
  (i,\, 1\cdot u') \;\vdash_P\; \bigl((i + 1) \bmod m,\, u'\cdot A_i\bigr).
\]
We write $\|(P, w)\| := m + \sum_{j=0}^{m-1} |A_j| + |w|$ for the total bit-and-phase description length of $(P, w)$.
\end{definition}

Because pure-$S$ reduction cannot erase consumed queue cells or reset template bits destructively, our self-describing encoding equips every phase header, appendant template bit, and queue bit with an append-only version chain inside protected $F$-frames.

\begin{construction}[Self-describing pure-$S$ CTS encoding $\mathcal{E}_{\mathrm{CTS}}$]\label{con:cts-encode}
Retaining the protected frame $F(a,b,j) \equiv S a(S b j)$, counter generators $\beta_d \equiv G(m_d, n_d)$, junk terms $\iota_d$, and width-$r$ version chains of Construction~\ref{con:s-encode-tm}, define the following status codes:
\begin{enumerate}[label=(\roman*),leftmargin=2.2em]
  \item \emph{$2$-bit phase status codes:} $\mathtt{00}$ ($\mathsf{inactive}$), $\mathtt{10}$ ($\mathsf{active}$), and $\mathtt{11}$ ($\mathsf{handoff}$).
  \item \emph{$3$-bit appendant template codes, for template bit $s \in \{0,1\}$:} $\mathtt{00}s$ ($\mathsf{idle}$), $\mathtt{10}s$ ($\mathsf{armed}$), and $\mathtt{11}s$ ($\mathsf{appending}$).
  \item \emph{$3$-bit queue cell codes, for data bit $b \in \{0,1\}$:} $\mathtt{01}b$ ($\mathsf{staging}$), $\mathtt{00}b$ ($\mathsf{unconsumed}$), $\mathtt{10}b$ ($\mathsf{consuming}$), and $\mathtt{11}b$ ($\mathsf{consumed}$).
\end{enumerate}
Given $P = (A_0, \dots, A_{m-1})$ and $w \in \{0,1\}^*$:
\begin{itemize}[leftmargin=2.0em]
  \item Each appendant $A_j = s_{j,1}\cdots s_{j,\ell_j}$ is encoded as a right-spine $F$-list of $\ell_j$ width-$3$ version chains initialized to $[\mathtt{00}s_{j,r}]$ ($\mathsf{idle}$) and terminated by $\beta_{4+j+\ell_j}$; each phase $j \in \{0,\dots,m-1\}$ pairs a width-$2$ status chain (initialized to $[\mathtt{10}]$ for $j = 0$ and $[\mathtt{00}]$ for $j > 0$) with its template list inside $F(-,-,\iota_{3+j})$, chained along the right spine (terminated by $\beta_{2+m}$) to form $\mathrm{Ring}_0$ at base depth $d = 2$.
  \item The initial queue $w = b_1\cdots b_{|w|}$ is encoded at base depth $d = 2$ as a right-spine $F$-list $\mathrm{Queue}_0$ of $|w|$ width-$3$ version chains initialized to $[\mathtt{00}b_r]$ ($\mathsf{unconsumed}$), terminated by $\beta_{2+|w|}$.
  \item Combining the unopened halt flag $\beta_1$ at path $[1,0,1]$ with the payload pair $F(\mathrm{Ring}_0, \mathrm{Queue}_0, \iota_1)$ yields the closed $S$-term
  \begin{equation}\label{eq:cts-encode-root}
    \mathcal{E}_{\mathrm{CTS}}(P, w) \;\equiv\; (\nu_1\,S)\,F\Bigl(\beta_1,\; F\bigl(\mathrm{Ring}_0,\, \mathrm{Queue}_0,\, \iota_1\bigr),\; \iota_0\Bigr),
  \end{equation}
  whose total leaf count satisfies $|\mathcal{E}_{\mathrm{CTS}}(P,w)| \le 2500\,\bigl(\|(P,w)\| + 1\bigr)^2$.
\end{itemize}
Indeed, each initial width-$2$ record opens at most $6$ frames ($+2$ phase-spine frames per phase), and each initial width-$3$ template or queue record opens at most $8$ frames ($+1$ list frame), plus $2$ outer hub frames, giving $K_0 \le 2 + 8m + 9\sum |A_j| + 9|w| \le 2 + 9\|(P,w)\|$ opened frames and height $H_0 \le 4 + 9\|(P,w)\|$. By the layout leaf bound of Lemma~\ref{lem:frame-open-bound}(ii),
\[
  |\mathcal{E}_{\mathrm{CTS}}(P,w)| \;\le\; 8 + (2K_0 + 1)(19 + 6H_0) \;\le\; 8 + 23(\|(P,w)\|+1)\cdot 97(\|(P,w)\|+1) \;\le\; 2500\bigl(\|(P,w)\|+1\bigr)^2.
\]
\end{construction}

\subsection{Program-Free Bare-Term Selector and Ten-Leaf Seed Genesis}\label{subsec:cts-seed}

Because both $\mathrm{Ring}_0$ and $\mathrm{Queue}_0$ are embedded inside $\mathcal{E}_{\mathrm{CTS}}(P,w)$, a single fixed selector $\sigma_{\mathrm{CTS}}$ can execute arbitrary cyclic tag systems without holding any CTS or UTM parameters in its own definition.

\begin{definition}[Program-free selector $\sigma_{\mathrm{CTS}}$ and bare-term decoder $\mathcal{D}_{\mathrm{CTS}}$]\label{def:pf-selector}
The \emph{program-free selector} $\sigma_{\mathrm{CTS}} \colon \TS \to \TS \cup \{\mathsf{halt}, \mathsf{stuck}\}$ inspects only a bare term $t \in \TS$:
\begin{enumerate}[label=(\roman*),leftmargin=2.2em]
  \item If $t \in L_{\mathrm{halt}}$ (recognized by the $11$-state automaton $\mathcal{H}_{11}$ of Proposition~\ref{prop:halt-automaton}), return $\mathsf{halt}$.
  \item Otherwise, parse $t \equiv H_{\mathrm{hdr}} \cdot F(\mathrm{Flag}, F(\mathrm{Ring}, \mathrm{Queue}, j_1), j_0)$, scan the newest committed $2$-bit status of each phase in $\mathrm{Ring}$ and its $3$-bit template statuses, and scan $\mathrm{Queue}$ to locate the first $\mathsf{consuming}$ cell ($\mathtt{10}b$), first $\mathsf{unconsumed}$ cell ($\mathtt{00}b$), first $\mathsf{staging}$ cell ($\mathtt{01}s$), or any freshly opened uninitialized tail slot:
  \begin{itemize}[leftmargin=1.6em]
    \item \emph{Handoff completion:} If some phase $j$ has status $\mathsf{handoff}$ ($\mathtt{11}$), then if a queue cell is $\mathsf{consuming}$ ($\mathtt{10}b$), append record $\mathtt{11}b$ ($\mathsf{consumed}$) to it; else if $m = 1$, append $\mathtt{10}$ ($\mathsf{active}$) to phase~$0$; else if phase $j' = (j+1) \bmod m$ is not yet $\mathtt{10}$, append $\mathtt{10}$ to phase $j'$, and once phase $j'$ is $\mathtt{10}$, append $\mathtt{00}$ ($\mathsf{inactive}$) to phase $j$.
    \item \emph{Active phase dispatch:} If no phase is in $\mathsf{handoff}$ and phase $j$ is $\mathsf{active}$ ($\mathtt{10}$):
      \begin{enumerate}[label=(\alph*),leftmargin=1.8em]
        \item When no queue cell is $\mathsf{consuming}$: if no $\mathsf{unconsumed}$ cell exists, open the halt flag $\mathrm{Flag}$ at path $[1,0,1]$; else if the front $\mathsf{unconsumed}$ cell holds bit $b = 0$, append $\mathtt{100}$ ($\mathsf{consuming}$); else ($b = 1$) if phase $j$ still has an $\mathsf{idle}$ template ($\mathtt{00}s$), append $\mathtt{10}s$ ($\mathsf{armed}$) to the first $\mathsf{idle}$ template, and once all templates of phase $j$ are $\mathsf{armed}$, append $\mathtt{101}$ ($\mathsf{consuming}$) to the front queue cell.
        \item When a queue cell is $\mathsf{consuming}$ ($\mathtt{10}b$): if a template is $\mathsf{appending}$ ($\mathtt{11}s$), first open the queue tail if needed, then write $\mathtt{01}s$ ($\mathsf{staging}$) into the new tail cell, and once $\mathtt{01}s$ is committed, reset the $\mathsf{appending}$ template to $\mathtt{00}s$ ($\mathsf{idle}$); if no template is $\mathsf{appending}$ and a queue cell is $\mathsf{staging}$ ($\mathtt{01}s$), commit it to $\mathtt{00}s$ ($\mathsf{unconsumed}$); else if an $\mathsf{armed}$ template ($\mathtt{10}s$) remains, advance the first $\mathsf{armed}$ template to $\mathtt{11}s$ ($\mathsf{appending}$); else, with all templates $\mathsf{idle}$ and the queue clean, advance phase $j$ to $\mathtt{11}$ ($\mathsf{handoff}$).
      \end{enumerate}
  \end{itemize}
  \item Having chosen either a target generator slot to open or a target version chain to extend by one bit, apply the $5$-edge local opening probe of Lemma~\ref{lem:frame-open-bound}(i) to contract the unique next redex $t \redS t'$.
\end{enumerate}
The companion \emph{bare-term CTS decoder} $\mathcal{D}_{\mathrm{CTS}} \colon \TS \rightharpoonup \mathbb{N} \times \{0,1\}^*$ checks that neither $\mathrm{Flag}$ nor any phase or queue version chain is mid-opening, verifies that a unique phase $i$ has status $\mathtt{10}$ with all others $\mathtt{00}$, skips $\mathsf{consumed}$ ($\mathtt{11}b$) queue cells, and collects $\mathsf{unconsumed}$ ($\mathtt{00}b$) cells to return $(i, u)$.
\end{definition}

Before proving universal simulation, we establish a structural property unique to duplicative combinatory rewriting: every encoded CTS instance $\mathcal{E}_{\mathrm{CTS}}(P,w)$ (and indeed every finite protected-frame layout, including $\mathcal{E}_{\mathrm{TM}}(M,w)$) is generated via pure-$S$ reduction from a single fixed $10$-leaf combinator.

\begin{proposition}[Universal genesis from the ten-leaf seed]\label{prop:seed-genesis}
Define the $10$-leaf closed $S$-term
\begin{equation}\label{eq:ten-leaf-seed}
  D_{\mathrm{seed}} \;\equiv\; S\bigl(S(S S)\bigr)\bigl(S(S S)\bigr)(S S)\,S \;\equiv\; S B B T S.
\end{equation}
For every cyclic tag system $P$ and every input word $w \in \{0,1\}^*$, there is an explicit pure-$S$ reduction sequence $D_{\mathrm{seed}} \redSs \mathcal{E}_{\mathrm{CTS}}(P, w)$.
\end{proposition}

\begin{proof}
First, $|D_{\mathrm{seed}}| = 1 + 3 + 3 + 2 + 1 = 10$, and $D_{\mathrm{seed}}$ bootstraps in $5$ pure-$S$ contractions to the universal root generator $(\nu_1 S)\,G(2,2)$ (recall $\nu_0 \equiv B T \equiv S T T$, $\nu_1 \equiv T \nu_0$, $\nu_2 \equiv T \nu_1$):
\begin{align*}
  D_{\mathrm{seed}} &\;\equiv\; (S B B T)\,S \;\redS\; \bigl((B T)(B T)\bigr)\,S \;\equiv\; \nu_0\,\nu_0\,S \;\equiv\; G(0,0) \\
  &\;\equiv\; (S T T \nu_0)\,S \;\redS\; \bigl((T \nu_0)(T \nu_0)\bigr)\,S \;\equiv\; \nu_1\,\nu_1\,S \;\equiv\; G(1,1) \\
  &\;\equiv\; (S S \nu_0 \nu_1)\,S \;\redS\; \bigl(S \nu_1 (\nu_0 \nu_1)\bigr)\,S \;\to_r\; (\nu_1 S)\,(\nu_0 \nu_1 S) \\
  &\;\equiv\; (\nu_1 S)\,\bigl((S T T \nu_1)\,S\bigr) \;\redS\; (\nu_1 S)\,\bigl((T \nu_1)(T \nu_1)\,S\bigr) \;\equiv\; (\nu_1 S)\,G(2, 2).
\end{align*}
Second, let $\Lambda$ range over finite binary layout skeletons generated by $\Lambda ::= \mathsf{unopened} \mid \mathsf{opened}(\Lambda_1, \Lambda_2)$, where $\mathrm{Term}(\mathsf{unopened}, m, n) \equiv G(m,n)$ and $\mathrm{Term}(\mathsf{opened}(\Lambda_1,\Lambda_2), m, n)$ recursively places $\mathrm{Term}(\Lambda_1, m', n')$ and $\mathrm{Term}(\Lambda_2, m', n')$ into the left and right slots of $F(-,-,j)$ using the phase transition $(m,n) \mapsto (m',n')$ of Construction~\ref{con:s-encode-tm}. We prove by structural induction on $\Lambda$ that $G(m,n) \redSs \mathrm{Term}(\Lambda, m, n)$ for all $m, n \in \mathbb{N}$ with $2 \le n$ and $m \le n$. The base case $\Lambda = \mathsf{unopened}$ is reflexive. For $\Lambda = \mathsf{opened}(\Lambda_1, \Lambda_2)$, Lemma~\ref{lem:frame-open-bound}(i) gives $G(m,n) \redS^6 F(G(m',n'), G(m',n'), j)$ (if $m \ge 1$) or $G(m,n) \redS^7 F(G(m',n'), G(m',n'), j)$ (if $m = 0$), with $2 \le n'$ and $m' \le n'$. Because $\redSs$ is congruent inside both $a$ and $b$ of $F(a,b,j) \equiv (S a)(S b j)$, applying the induction hypotheses $G(m',n') \redSs \mathrm{Term}(\Lambda_1, m', n')$ on the left and $G(m',n') \redSs \mathrm{Term}(\Lambda_2, m', n')$ on the right yields $G(m,n) \redSs \mathrm{Term}(\mathsf{opened}(\Lambda_1,\Lambda_2), m, n)$. By Construction~\ref{con:cts-encode}, $\mathcal{E}_{\mathrm{CTS}}(P,w) \equiv (\nu_1 S)\,\mathrm{Term}(\Lambda_{P,w}, 2, 2)$ for an explicit finite layout $\Lambda_{P,w}$, whence $D_{\mathrm{seed}} \redS^5 (\nu_1 S)\,G(2,2) \redSs \mathcal{E}_{\mathrm{CTS}}(P,w)$.
\end{proof}

\begin{theorem}[Program-free cyclic tag universality in $\CLS$]\label{thm:cts-universality}
For every cyclic tag system $P = (A_0, \dots, A_{m-1})$ and every input word $w \in \{0,1\}^*$, let $T_0 = \mathcal{E}_{\mathrm{CTS}}(P, w) \in \TS$.
\begin{enumerate}[label=(\roman*),leftmargin=2.2em]
  \item \emph{Single-step soundness and all-step quadratic leaf bound:} Whenever $\sigma_{\mathrm{CTS}}(t) = t' \in \TS$, one has $t \redS t'$. Along any $k$-step trajectory $T_0 \redS T_1 \redS \cdots \redS T_k$ generated by iterating $\sigma_{\mathrm{CTS}}$, every term $T_k$ satisfies
  \begin{equation}\label{eq:cts-traj-bound}
    |T_k| \;\le\; 35{,}000\,\bigl(\|(P,w)\| + k + 1\bigr)^2.
  \end{equation}
  \item \emph{Bidirectional halting equivalence via $\mathcal{H}_{11}$ and terminal readback:} The cyclic tag system $P$ starting from $(0, w)$ eventually halts (empties its queue) if and only if there exists $k \ge 0$ such that $\sigma_{\mathrm{CTS}}^k(T_0) = T_k \in L(\mathcal{H}_{11})$, where $\mathcal{H}_{11}$ is the identical $11$-state bottom-up halting automaton of Proposition~\ref{prop:halt-automaton}. Moreover, for that unique halting step $k$, $(0, w) \vdash_P^* (i_{\mathrm{fin}}, \varepsilon)$ holds and $\mathcal{D}_{\mathrm{CTS}}(T_k) = (i_{\mathrm{fin}}, \varepsilon)$.
\end{enumerate}
\end{theorem}

\begin{proof}
Single-step soundness in~(i) holds because every non-halting step of $\sigma_{\mathrm{CTS}}$ contracts a single verified redex $S X Y Z \to_r (X Z)(Y Z)$ located by Lemma~\ref{lem:frame-open-bound}(i). To verify step-for-step CTS simulation, maintain the invariant $\mathrm{Rep}_P(i, u)$ stating that $\mathrm{Flag} \equiv \beta_1$, phase $i$ has newest status $\mathtt{10}$ with all $|A_i|$ templates $\mathtt{00}s$ ($\mathsf{idle}$), all other phases $j \ne i$ have newest status $\mathtt{00}$ with all templates $\mathtt{00}s$, and $\mathrm{Queue}$ consists of a prefix of $\mathsf{consumed}$ ($\mathtt{11}b$) cells followed by $|u|$ queue cells with status $\mathsf{unconsumed}$ ($\mathtt{00}b$) matching $u$. Clearly $T_0$ satisfies $\mathrm{Rep}_P(0, w)$ and $\mathcal{D}_{\mathrm{CTS}}(T_0) = (0, w)$.

Suppose $T$ satisfies $\mathrm{Rep}_P(i, u)$. Because appending a width-$r$ version record opens between $1$ and $2r + 2$ layout slots (each taking $6$ or $7$ contractions while leaving the scanned committed value unchanged until the final sentinel slot opens), appending a $2$-bit phase record takes between $6$ and $42$ contractions (at most $6$ opened slots) and appending a $3$-bit cell or template record takes between $6$ and $56$ contractions (at most $8$ opened slots):
\begin{itemize}[leftmargin=2.0em]
  \item If $u = \varepsilon$, rule~(ii)(a) opens $\mathrm{Flag}$ at $[1,0,1]$ in $6$ or $7$ steps to $T_{\mathrm{halt}} \in L(\mathcal{H}_{11})$, after which $\sigma_{\mathrm{CTS}}(T_{\mathrm{halt}}) = \mathsf{halt}$ and $\mathcal{D}_{\mathrm{CTS}}(T_{\mathrm{halt}}) = (i, \varepsilon)$.
  \item If $u = 0\cdot u'$, rule~(ii)(a) marks the front cell $\mathtt{100}$ ($\mathsf{consuming}$), advances phase $i$ to $\mathtt{11}$ ($\mathsf{handoff}$), marks the consuming cell $\mathtt{110}$ ($\mathsf{consumed}$), and activates $(i+1) \bmod m$ to $\mathtt{10}$ (deactivating $i$ to $\mathtt{00}$ if $m \ge 2$), reaching $\mathrm{Rep}_P((i+1) \bmod m, u')$ in $k_{\mathrm{step}} \ge 24$ contractions while adding at most $8 + 26 = 34$ opened slots.
  \item If $u = 1\cdot u'$, rule~(ii)(a) first arms all $|A_i|$ templates ($\mathtt{10}s$, adding $8|A_i|$ slots), then marks the front cell $\mathtt{101}$ ($\mathsf{consuming}$, adding $8$ slots); rule~(ii)(b) then copies each bit $s$ of $A_i$ in order via the $5$-part handshake ($\mathtt{11}s \to \text{open queue tail} \to \mathtt{01}s \to \mathtt{00}s\text{ template} \to \mathtt{00}s\text{ queue}$, taking at least $30$ contractions and $33$ slots per bit), and executes the same handoff ($26$ slots) to reach $\mathrm{Rep}_P((i+1)\bmod m, u'\cdot A_i)$ in $k_{\mathrm{step}} \ge 24 + 36|A_i|$ contractions while adding $34 + 41|A_i| \le 34 + 41\|(P,w)\|$ slots (and at most $2 k_{\mathrm{step}}$ slots).
\end{itemize}
Because no intermediate term inside a macro-step has an opened $\mathrm{Flag}$, $T_k \notin L(\mathcal{H}_{11})$ prior to queue exhaustion; constructive bounded evaluation of $P$ for $k + 1$ steps combines with determinism of $\sigma_{\mathrm{CTS}}$ to show that $(0, w) \vdash_P^* (i_{\mathrm{fin}}, \varepsilon)$ holds if and only if $T_k \in L(\mathcal{H}_{11})$ for some $k \ge 0$. Moreover, every $T_k$ is a prefix of a checkpoint term with at most $K \le (2 + 9\|(P,w)\|) + 2k + 34 + 41\|(P,w)\| \le 50(\|(P,w)\| + k + 1)$ opened slots and height $H \le 2 + 50(\|(P,w)\| + k + 1)$; monotone leaf growth gives
\[
  |T_k| \;\le\; 8 + 101(\|(P,w)\|+k+1)\cdot 331(\|(P,w)\|+k+1) \;=\; 8 + 33{,}431\bigl(\|(P,w)\|+k+1\bigr)^2 \;\le\; 35{,}000\bigl(\|(P,w)\|+k+1\bigr)^2. \qedhere
\]
\end{proof}

Although Definition~\ref{def:pf-selector} eliminates all universal Turing and tag transition tables from $\sigma_{\mathrm{CTS}}$, its specification still relies on recursive structural scanners over unbounded phase rings, queue spines, and version chains; compressing $\sigma_{\mathrm{CTS}}$ into an explicit minimal finite-state tree transducer directly addresses the helper-complexity question posed in Wolfram's challenge announcement \cite{Wolfram2021Challenge}.

\subsection{Two-Pass Linear-Time Contextual Tree Transducer Compression}\label{subsec:pocket-transducer}

In Section~\ref{sec:tm-sim}, the root-restarted zipper walker $\sigma_{\mathrm{TM}}$ requires $|Q_{\mathrm{ctrl}}| = 4{,}376{,}490{,}905$ control states to store the $262$-rule transition table of $U_{\mathrm{bin}}$ alongside $11$-bit cell shift buffers, and executes up to $3500(2|M|-1)$ cursor steps per contraction. Conversely, Definition~\ref{def:pf-selector} formulates $\sigma_{\mathrm{CTS}}$ in terms of unbounded phase indices $(j + 1) \bmod m$ and tree-address lists. We now show that $\sigma_{\mathrm{CTS}}$ is extensionally identical on all of $\TS$ to a compact two-pass contextual finite-state tree transducer $\mathcal{T}_{\mathrm{pocket}}$: a single bottom-up pass conditioned on $25$ static structural roles assigns a $23$-state summary ($< 5$ bits) to every node in $2|M| - 1$ node transitions, and a single top-down pointer pass with $136$ states descends directly along one root-to-redex path in at most $\mathrm{depth}(M)$ edge transitions without consulting role labels, re-scanning subtrees, or performing unbounded arithmetic.

\begin{definition}[Two-pass contextual finite-state tree transducer]\label{def:pocket-transducer}
For every closed pure-$S$ tree $M \in \TS$, define its binary tree height $\mathrm{depth}(M) \in \mathbb{N}$ recursively by $\mathrm{depth}(S) = 0$ and $\mathrm{depth}(M_1 M_2) = 1 + \max\{\mathrm{depth}(M_1), \mathrm{depth}(M_2)\}$. Write $\alpha \in \{0,1\}^*$ for a binary tree address ($0$ descending left, $1$ descending right, and $\varepsilon$ denoting the root), $M|_{\alpha}$ for the subterm at $\alpha$ when defined, and $M[\alpha \leftarrow N]$ for single-subtree replacement at $\alpha$.

A \emph{two-pass contextual finite-state tree transducer} over $\TS$ is a $10$-tuple
\[
  \mathcal{T} \;=\; \bigl(R,\, Q,\, P,\, r_0,\, \rho,\, \delta_{\mathrm{leaf}},\, \delta_{\mathrm{app}},\, F_{\mathrm{halt}},\, p_0,\, \delta_{\mathrm{TD}}\bigr),
\]
consisting of finite nonempty state sets $R$ (\emph{static structural roles}), $Q$ (\emph{bottom-up node states}), and $P$ (\emph{top-down pointer states}), an initial root role $r_0 \in R$, a static child-role map $\rho \colon R \to R \times R$, a leaf rule $\delta_{\mathrm{leaf}} \colon R \to Q$, a role-indexed bottom-up transition $\delta_{\mathrm{app}} \colon R \times Q \times Q \to Q$, a halting subset $F_{\mathrm{halt}} \subseteq Q$, an initial pointer state $p_0 \in P$, and a role-free top-down transition function
\[
  \delta_{\mathrm{TD}} \;\colon\; P \times Q \times Q \;\longrightarrow\; \{\mathsf{left}(p') : p' \in P\} \cup \{\mathsf{right}(p') : p' \in P\} \cup \{\mathsf{contract},\, \mathsf{stuck}\}.
\]
Given an input tree $M \in \TS$, execution of $\mathcal{T}$ proceeds in two deterministic passes:
\begin{enumerate}[label=(\roman*),leftmargin=2.2em]
  \item \emph{Pass~1 (contextual bottom-up annotation):} For each $(r, N) \in R \times \TS$, define the $Q$-annotated tree $\widehat{N} = \mathrm{Ann}_{\mathcal{T}}(r, N)$, its root state $\mathrm{st}(\widehat{N}) \in Q$, and its node-transition count $s_1(r, N) \in \mathbb{N}$ recursively by setting $\mathrm{st}(\mathrm{Ann}_{\mathcal{T}}(r, S)) = \delta_{\mathrm{leaf}}(r)$ with $s_1(r, S) = 1$, and, for $N \equiv N_1 N_2$ with $\rho(r) = (r_{\mathrm{L}}, r_{\mathrm{R}})$, $\widehat{N}_1 = \mathrm{Ann}_{\mathcal{T}}(r_{\mathrm{L}}, N_1)$, and $\widehat{N}_2 = \mathrm{Ann}_{\mathcal{T}}(r_{\mathrm{R}}, N_2)$, labeling the root of $\widehat{N}$ with
  \[
    \mathrm{st}(\widehat{N}) \;=\; \delta_{\mathrm{app}}\bigl(r,\, \mathrm{st}(\widehat{N}_1),\, \mathrm{st}(\widehat{N}_2)\bigr), \qquad
    s_1(r, N_1 N_2) \;=\; 1 + s_1(r_{\mathrm{L}}, N_1) + s_1(r_{\mathrm{R}}, N_2).
  \]
  At the root of $M$, Pass~1 computes $\widehat{M}_0 = \mathrm{Ann}_{\mathcal{T}}(r_0, M)$. If $\mathrm{st}(\widehat{M}_0) \in F_{\mathrm{halt}}$, then $\mathcal{T}$ halts immediately with outcome $\mathcal{T}(M) = \mathsf{halt}$ and total transition count $\tau_{\mathcal{T}}(M) = s_1(r_0, M)$.
  \item \emph{Pass~2 (role-free top-down redex location and contraction):} If $\mathrm{st}(\widehat{M}_0) \notin F_{\mathrm{halt}}$, define the pointer descent $\mathrm{Loc}_{\mathcal{T}}(p, \widehat{N}) \in (\{0,1\}^* \cup \{\bot\}) \times \mathbb{N}$ recursively by $\mathrm{Loc}_{\mathcal{T}}(p, \widehat{S}) = (\bot, 0)$ on leaves and, on an annotated application $\widehat{N}$ with children $(\widehat{N}_1, \widehat{N}_2)$ bearing states $(q_{\mathrm{L}}, q_{\mathrm{R}}) = (\mathrm{st}(\widehat{N}_1), \mathrm{st}(\widehat{N}_2))$, by evaluating $\delta_{\mathrm{TD}}(p, q_{\mathrm{L}}, q_{\mathrm{R}})$---which inspects only $(p, q_{\mathrm{L}}, q_{\mathrm{R}})$ without role input from $R$, subtree re-scanning, or unbounded arithmetic:
  \begin{itemize}[leftmargin=1.6em]
    \item if $\delta_{\mathrm{TD}}(p, q_{\mathrm{L}}, q_{\mathrm{R}}) = \mathsf{left}(p')$ and $\mathrm{Loc}_{\mathcal{T}}(p', \widehat{N}_1) = (\omega', s_2')$, return $(0 \cdot \omega', s_2' + 1)$ (with $0 \cdot \bot := \bot$);
    \item if $\delta_{\mathrm{TD}}(p, q_{\mathrm{L}}, q_{\mathrm{R}}) = \mathsf{right}(p')$ and $\mathrm{Loc}_{\mathcal{T}}(p', \widehat{N}_2) = (\omega', s_2')$, return $(1 \cdot \omega', s_2' + 1)$ (with $1 \cdot \bot := \bot$);
    \item if $\delta_{\mathrm{TD}}(p, q_{\mathrm{L}}, q_{\mathrm{R}}) = \mathsf{contract}$, return $(\varepsilon, 1)$;
    \item if $\delta_{\mathrm{TD}}(p, q_{\mathrm{L}}, q_{\mathrm{R}}) = \mathsf{stuck}$, return $(\bot, 1)$.
  \end{itemize}
  Let $(\omega, s_2) = \mathrm{Loc}_{\mathcal{T}}(p_0, \widehat{M}_0)$ and set $\tau_{\mathcal{T}}(M) = s_1(r_0, M) + s_2$. If $\omega = \alpha \in \{0,1\}^*$ and $M|_{\alpha} \equiv S X Y Z$, then $\mathcal{T}(M) = M[\alpha \leftarrow (X Z)(Y Z)] \in \TS$; otherwise $\mathcal{T}(M) = \mathsf{stuck}$.
\end{enumerate}
\end{definition}

To instantiate Definition~\ref{def:pocket-transducer} for self-describing CTS terms \eqref{eq:cts-encode-root}, observe that every frame $F(a, b, j) \equiv (S a)((S b) j)$ wraps its left slot $a$ inside a unary prefix $S a$ and its right slot $b$ inside a binary wrapper $(S b) j$. Accordingly, our concrete transducer $\mathcal{T}_{\mathrm{pocket}}$ employs $25$ static roles $R_{\mathrm{pocket}}$, summarized in Table~\ref{tab:pocket-roles}: $11$ primary roles ($\mathsf{rt}$, $\mathsf{hub}$, $\mathsf{pay}$, $\mathsf{ring}$, $\mathsf{body}$, $\mathsf{tmpl}$, $\mathsf{queue}$, $\mathsf{ver}_2$, $\mathsf{ver}_3$, $\mathsf{fld}$, $\mathsf{sp}$), $9$ left-wrapper roles $\mathsf{L}_X$ matching $S a$, and $5$ right-wrapper roles $\mathsf{R}_Y$ matching $(S b) j$. Because all inner non-ring roles take at most $16$ summary values in $\{0,\dots,15\}$, shifting valid wrapper states on $\mathsf{L}_X$ ($X \ne \mathsf{ring}$) and $\mathsf{R}_Y$ ($Y \ne \mathsf{ring}$) by $+7$ into $\{7,\dots,22\}$ keeps them disjoint from truncated spine arities $\{0,\dots,5\}$, multiplexing all $25$ roles into a shared carrier $Q_{\mathrm{pocket}} = \{0,\dots,22\}$ of size $23$ without passing $R_{\mathrm{pocket}}$ to Pass~2.

\begin{table}[htbp]
\centering
\caption{Static $25$-role grammar $\rho(r) = (r_{\mathrm{L}}, r_{\mathrm{R}})$ and Pass-1 semantic state encodings of $\mathcal{T}_{\mathrm{pocket}}$ within $Q_{\mathrm{pocket}} = \{0,\dots,22\}$ ($\delta_{\mathrm{leaf}}(r) = 0$ for all $r \in R_{\mathrm{pocket}}$; $F_{\mathrm{halt}} = \{1\}$ at root role $r_0 = \mathsf{rt}$).}\label{tab:pocket-roles}
\small
\begin{tabular}{@{}llcl@{}}
\toprule
Role $r$ (count) & Child roles $\rho(r)$ & Codomain & Pass-1 bottom-up invariant $\varsigma(r, N)$ on $N \in \TS$ \\
\midrule
$\mathsf{rt}$ ($1$) & $(\mathsf{sp}, \mathsf{hub})$ & $\{0,1\}$ & $1$ iff $N \in L(\mathcal{H}_{11})$ ($\mathsf{rootHalt}(N) = 1$) \\
$\mathsf{hub}$ ($1$) & $(\mathsf{L}_{\mathsf{sp}}, \mathsf{R}_{\mathsf{pay}})$ & $\{0,1\}$ & $1$ iff $N \equiv F(\mathrm{Flag}, \mathrm{Pay}, j_0)$ with $\mathrm{Flag}$ an $F$-frame \\
$\mathsf{pay}$ ($1$) & $(\mathsf{L}_{\mathsf{ring}}, \mathsf{R}_{\mathsf{queue}})$ & $\{0,1,2\}$ & $0$ non-frame; $2$ active empty-queue halt trigger; $1$ else \\
$\mathsf{ring}$ ($1$) & $(\mathsf{L}_{\mathsf{body}}, \mathsf{R}_{\mathsf{ring}})$ & $\{0,\dots,21\}$ & $0$ end; $1\text{--}4$ handoff $(b_{\mathrm{hd}}, b_{\mathrm{last}})$; $5$ idle; $6\text{--}21$ active \\
$\mathsf{body}$ ($1$) & $(\mathsf{L}_{\mathsf{ver}_2}, \mathsf{R}_{\mathsf{tmpl}})$ & $\{0,\dots,10\}$ & $0$ non-frame; $1$ $\mathtt{11}$; $2$ $\mathtt{0}b$/none; $3 + c_{\mathrm{tmpl}}$ on $\mathtt{10}$ \\
$\mathsf{tmpl}$ ($1$) & $(\mathsf{L}_{\mathsf{ver}_3}, \mathsf{R}_{\mathsf{tmpl}})$ & $\{0,\dots,7\}$ & $4 b_{\mathrm{idle}} + c_4$ ($\mathsf{idle}$ bit and $\mathsf{appending}$/$\mathsf{armed}$ priority) \\
$\mathsf{queue}$ ($1$) & $(\mathsf{L}_{\mathsf{ver}_3}, \mathsf{R}_{\mathsf{queue}})$ & $\{0,\dots,11\}$ & $9 + \mathrm{su}$ if $\mathsf{consuming}$; else $3 u + \mathrm{su}$ ($\mathsf{unconsumed}$, tail) \\
$\mathsf{ver}_2$ ($1$) & $(\mathsf{L}_{\mathsf{ver}_2}, \mathsf{R}_{\mathsf{fld}})$ & $\{0,\dots,3\}$ & Newest $2$-bit phase status ($0$ none, $1$ $\mathtt{11}$, $2$ $\mathtt{10}$, $3$ $\mathtt{0}b$) \\
$\mathsf{ver}_3$ ($1$) & $(\mathsf{L}_{\mathsf{ver}_3}, \mathsf{R}_{\mathsf{fld}})$ & $\{0,\dots,8\}$ & Newest $3$-bit record ($0$ none, $1 + 4 b_0 + 2 b_1 + b_2$) \\
$\mathsf{fld}$ ($1$) & $(\mathsf{L}_{\mathsf{sp}}, \mathsf{R}_{\mathsf{fld}})$ & $\{0,\dots,15\}$ & Bit-field prefix of length $0$ ($1$), $1$ ($2\text{--}3$), $2$ ($4\text{--}7$), $\ge 3$ ($8\text{--}15$) \\
$\mathsf{sp}$ ($1$) & $(\mathsf{sp}, \mathsf{sp})$ & $\{0,\dots,6\}$ & $6$ if $N$ is an $F$-frame; else $\min\{5, \ls(N)\} \in \{0,\dots,5\}$ \\
$\mathsf{L}_{\mathsf{ring}}, \mathsf{R}_{\mathsf{ring}}$ ($2$) & $(\mathsf{sp}, \mathsf{ring}),\, (\mathsf{L}_{\mathsf{ring}}, \mathsf{sp})$ & $\{0,\dots,22\}$ & $1 + \varsigma(\mathsf{ring}, u)$ on $S u$ resp.\ $(S u) j$; $0$ otherwise \\
$\mathsf{L}_X$, $X \ne \mathsf{ring}$ ($8$) & $(\mathsf{sp}, X)$ & $\{0,\dots,22\}$ & $7 + \varsigma(X, u)$ on $S u$; else $\min\{5, \ls(N)\} \in \{0,\dots,5\}$ \\
$\mathsf{R}_Y$, $Y \ne \mathsf{ring}$ ($4$) & $(\mathsf{L}_Y, \mathsf{sp})$ & $\{0,\dots,22\}$ & $7 + \varsigma(Y, v)$ on $(S v) j$; else $\min\{5, \ls(N)\} \in \{0,\dots,5\}$ \\
\bottomrule
\end{tabular}
\end{table}

\begin{theorem}[Linear-time pocket tree transducer equivalence and program-free universality]\label{thm:pocket-transducer-universality}
There exists an explicit program-free two-pass contextual finite-state tree transducer
\[
  \mathcal{T}_{\mathrm{pocket}} \;=\; \bigl(R_{\mathrm{pocket}},\, Q_{\mathrm{pocket}},\, P_{\mathrm{pocket}},\, r_0,\, \rho,\, \delta_{\mathrm{leaf}},\, \delta_{\mathrm{app}},\, F_{\mathrm{halt}},\, p_0,\, \delta_{\mathrm{TD}}\bigr)
\]
with $|R_{\mathrm{pocket}}| = 25$ static structural roles, $|Q_{\mathrm{pocket}}| = 23$ bottom-up node states ($< 5$ bits per node, with $271$ pairs $(r, q) \in R_{\mathrm{pocket}} \times Q_{\mathrm{pocket}}$ reachable from $r_0$ over $\TS$), and $|P_{\mathrm{pocket}}| = 136$ single-edge top-down pointer states ($54$ decision macro-states together with $50$ two-edge frame-left and $32$ three-edge frame-right stepping states), containing zero universal Turing or tag transition tables, such that:
\begin{enumerate}[label=(\roman*),leftmargin=2.2em]
  \item \emph{Exact linear transition bounds:} For every closed pure-$S$ tree $M \in \TS$, Pass~1 executes in $2|M| - 1$ node transitions and Pass~2 executes in at most $\mathrm{depth}(M)$ edge transitions. Consequently, $\mathcal{T}_{\mathrm{pocket}}$ halts unconditionally in
  \begin{equation}\label{eq:pocket-linear-bound}
    \tau_{\mathcal{T}_{\mathrm{pocket}}}(M) \;\le\; 2|M| - 1 + \mathrm{depth}(M)
  \end{equation}
  total transitions on every $M \in \TS$, and satisfies $\tau_{\mathcal{T}_{\mathrm{pocket}}}(M) + 2 \le 2|M| - 1 + \mathrm{depth}(M)$ (at most $2|M| - 3 + \mathrm{depth}(M)$ total transitions) whenever a pure-$S$ redex is contracted ($\mathcal{T}_{\mathrm{pocket}}(M) = M' \in \TS$).
  \item \emph{Extensional equivalence with $\sigma_{\mathrm{CTS}}$ and universal cyclic tag simulation:} On every closed tree $M \in \TS$, one has $\mathcal{T}_{\mathrm{pocket}}(M) = \sigma_{\mathrm{CTS}}(M)$ identically. Hence, for every cyclic tag system $P = (A_0, \dots, A_{m-1})$ and input word $w \in \{0,1\}^*$, the pure-$S$ encoding $T_0 = \mathcal{E}_{\mathrm{CTS}}(P, w) \in \TS$ unfolds via $D_{\mathrm{seed}} \redSs T_0$ from the $10$-leaf seed $D_{\mathrm{seed}} \equiv S(S(S S))(S(S S))(S S)S \equiv S B B T S$, and iterating $\mathcal{T}_{\mathrm{pocket}}$ simulates $(P, w)$ with bidirectional halting equivalence ($(0, w)$ halts under $\vdash_P^*$ if and only if $\mathcal{T}_{\mathrm{pocket}}^k(T_0) = T_k$ satisfies $\mathcal{T}_{\mathrm{pocket}}(T_k) = \mathsf{halt}$ for some $k \ge 0$), exact bare-term terminal readback $\mathcal{D}_{\mathrm{CTS}}(T_k) = (i_{\mathrm{fin}}, \varepsilon)$ matching $(0, w) \vdash_P^* (i_{\mathrm{fin}}, \varepsilon)$, and the all-step quadratic leaf bound $|T_k| \le 35{,}000(\|(P, w)\| + k + 1)^2$.
\end{enumerate}
\end{theorem}

\begin{proof}
We establish~(i) for any transducer $\mathcal{T}$ of Definition~\ref{def:pocket-transducer} before proving Pass-1 and Pass-2 correctness of $\mathcal{T}_{\mathrm{pocket}}$ in~(ii).

For~(i), structural induction on $N \in \TS$ shows $s_1(r, N) = 2|N| - 1$ for every $r \in R$: on $N \equiv S$, $s_1(r, S) = 1 = 2|S| - 1$; on $N \equiv N_1 N_2$ with $\rho(r) = (r_{\mathrm{L}}, r_{\mathrm{R}})$, the induction hypothesis gives
\[
  s_1(r, N_1 N_2) \;=\; 1 + \bigl(2|N_1| - 1\bigr) + \bigl(2|N_2| - 1\bigr) \;=\; 2\bigl(|N_1| + |N_2|\bigr) - 1 \;=\; 2|N_1 N_2| - 1.
\]
Next, let $(\omega, s_2) = \mathrm{Loc}_{\mathcal{T}}(p, \widehat{N})$ where $\widehat{N} = \mathrm{Ann}_{\mathcal{T}}(r, N)$. Structural induction on $N$ proves $s_2 \le \mathrm{depth}(N)$: on $N \equiv S$, $s_2 = 0 = \mathrm{depth}(S)$; on $N \equiv N_1 N_2$, if $\delta_{\mathrm{TD}}$ returns $\mathsf{contract}$ or $\mathsf{stuck}$, then $s_2 = 1 \le 1 + \max\{\mathrm{depth}(N_1), \mathrm{depth}(N_2)\} = \mathrm{depth}(N_1 N_2)$, whereas if $\delta_{\mathrm{TD}}$ steps left or right into $\widehat{N}_i$ ($i \in \{1, 2\}$), the induction hypothesis gives $s_2 = 1 + s_2' \le 1 + \mathrm{depth}(N_i) \le \mathrm{depth}(N_1 N_2)$. Hence $\tau_{\mathcal{T}}(M) \le 2|M| - 1 + \mathrm{depth}(M)$ for all $M \in \TS$ (with $\tau_{\mathcal{T}}(M) = 2|M| - 1$ when $\mathrm{st}(\widehat{M}_0) \in F_{\mathrm{halt}}$). Moreover, suppose $\mathcal{T}(M) = M' \in \TS$, so $\omega = \alpha \in \{0,1\}^*$ and $U \equiv M|_{\alpha} \equiv ((S X) Y) Z$. Because $\mathrm{depth}(S X) \ge 1$ and $\mathrm{depth}((S X) Y) \ge 2$, every pure-$S$ redex satisfies $\mathrm{depth}(U) \ge 3$. At $\alpha$, $\delta_{\mathrm{TD}}$ emits $\mathsf{contract}$ in $1 \le \mathrm{depth}(U) - 2$ transition; inducting upward along each edge of $\alpha$ adds $1$ to $s_2$ and increases $\mathrm{depth}$ by at least $1$, yielding $s_2 \le \mathrm{depth}(M) - 2$ and $\tau_{\mathcal{T}}(M) = (2|M| - 1) + s_2 \le 2|M| - 3 + \mathrm{depth}(M)$.

For~(ii), define the $25$-role transition $\delta_{\mathrm{app}} \colon R_{\mathrm{pocket}} \times Q_{\mathrm{pocket}}^2 \to Q_{\mathrm{pocket}}$ to implement the specification $\varsigma(r, N) \in \{0,\dots,22\}$ in Table~\ref{tab:pocket-roles}, with $\delta_{\mathrm{leaf}}(r) = 0$ for all $r$. Define the spine-arity projections $\pi_{\mathrm{L}}, \pi_{\mathrm{R}} \colon Q_{\mathrm{pocket}} \to \{0,\dots,5\}$ by setting $\pi_{\mathrm{L}}(q) = q$ for $q \le 5$, $\pi_{\mathrm{L}}(6) = 2$, and $\pi_{\mathrm{L}}(q) = 1$ for $q \ge 7$, while $\pi_{\mathrm{R}}(q) = q$ for $q \le 5$ and $\pi_{\mathrm{R}}(q) = 2$ for $q \ge 6$. We verify by structural induction on $N \in \TS$ that $\mathrm{st}(\mathrm{Ann}_{\mathcal{T}_{\mathrm{pocket}}}(r, N)) = \varsigma(r, N)$ simultaneously for all $r \in R_{\mathrm{pocket}}$. For $N \equiv S$, every role satisfies $\varsigma(r, S) = 0 = \delta_{\mathrm{leaf}}(r)$. For $N \equiv N_1 N_2$ with child states $(q_{\mathrm{L}}, q_{\mathrm{R}}) = (\varsigma(r_{\mathrm{L}}, N_1), \varsigma(r_{\mathrm{R}}, N_2))$:
\begin{itemize}[leftmargin=2.0em]
  \item \emph{Spine and frame-wrapper roles:} Because $S$ is the unique atom of $\TS$, a closed term $U \in \TS$ satisfies $\ls(U) = 0$ if and only if $U \equiv S$, $\ls(U) = 1$ if and only if $U \equiv S u$, and $\ls(U) = 2$ if and only if $U \equiv (S v) j$. Every $F$-frame $(S a)((S b) j)$ has $\ls = 2$, so $\pi_{\mathrm{L}}(\varsigma(\mathsf{sp}, U)) = \pi_{\mathrm{R}}(\varsigma(\mathsf{sp}, U)) = \min\{5, \ls(U)\}$ holds for all $U \in \TS$. On $\mathsf{sp}$, $N_1 N_2$ is an $F$-frame if and only if $q_{\mathrm{L}} = 1$ and $\pi_{\mathrm{R}}(q_{\mathrm{R}}) = 2$ (outputting $6$), and otherwise $\min\{5, \pi_{\mathrm{L}}(q_{\mathrm{L}}) + 1\} = \min\{5, \ls(N_1 N_2)\}$. On $\mathsf{L}_X$ ($X \ne \mathsf{ring}$, where $\varsigma(X, u) \le 15$), testing $\pi_{\mathrm{L}}(q_{\mathrm{L}}) = 0$ detects $N_1 \equiv S$ and outputs $7 + q_{\mathrm{R}} \in \{7,\dots,22\}$, else outputs $\min\{5, \pi_{\mathrm{L}}(q_{\mathrm{L}}) + 1\} \in \{2,\dots,5\}$; symmetrically, on $\mathsf{R}_Y$ ($Y \ne \mathsf{ring}$), testing $q_{\mathrm{L}} \ge 7$ detects $N_1 \equiv S v$ and propagates $q_{\mathrm{L}} = 7 + \varsigma(Y, v)$, else outputs $\min\{5, \pi_{\mathrm{L}}(q_{\mathrm{L}}) + 1\} \in \{1,3,4,5\}$. Consequently, $\pi_{\mathrm{L}}(\varsigma(\mathsf{L}_X, U)) = \pi_{\mathrm{R}}(\varsigma(\mathsf{R}_Y, U)) = \min\{5, \ls(U)\}$ for all $U \in \TS$, and at any frame role $r \in \{\mathsf{hub}, \mathsf{body}, \mathsf{tmpl}, \mathsf{queue}, \mathsf{ver}_2, \mathsf{ver}_3, \mathsf{fld}\}$ the condition $q_{\mathrm{L}} \ge 7 \land q_{\mathrm{R}} \ge 7$ holds if and only if $N_1 N_2 \equiv (S a)((S b) j)$ with $(a, b)$ encoded by $(q_{\mathrm{L}} - 7, q_{\mathrm{R}} - 7)$ (while $\mathsf{L}_{\mathsf{ring}}$ and $\mathsf{R}_{\mathsf{ring}}$ encode valid wrappers as $1 + \varsigma(\mathsf{ring}, u) \in \{1,\dots,22\}$ and non-wrappers as $0$).
  \item \emph{Fields, version chains, templates, and queue:} On non-frames ($q_{\mathrm{L}} < 7$ or $q_{\mathrm{R}} < 7$), $\delta_{\mathrm{app}}$ outputs $0$. On a frame $F(a, b, j)$: $\mathsf{fld}$ prepends the Boolean bit $b_0 = \mathbf{1}[\varsigma(\mathsf{sp}, a) = 6]$ to $\varsigma(\mathsf{fld}, b) \in \{0,\dots,15\}$, recording $1$ at the $0$-bit sentinel ($\varsigma(\mathsf{fld}, b) = 1$) and packing length-$1$, length-$2$, and length-$\ge 3$ prefixes into $\{2,3\}$, $\{4,\dots,7\}$, and $\{8,\dots,15\}$; $\mathsf{ver}_2$ and $\mathsf{ver}_3$ return $\varsigma(\mathsf{ver}_w, a)$ whenever $\varsigma(\mathsf{ver}_w, a) \ne 0$ (since newer records occupy the left spine) and otherwise decode the $2$-bit or $3$-bit prefix from $\varsigma(\mathsf{fld}, b)$; $\mathsf{tmpl}$ updates $(b_{\mathrm{idle}}, c_4) \in \{0,1\} \times \{0,1,2,3\}$ ($8$ states), where $b_{\mathrm{idle}}$ indicates an $\mathsf{idle}$ entry ($\mathtt{00}s$) and $c_4$ records the first $\mathsf{appending}$ bit ($2 + s$ for $\mathtt{11}s$), else presence of an $\mathsf{armed}$ entry ($1$ for $\mathtt{10}s$), else $0$; and $\mathsf{queue}$ updates the $12$-state code combining tail-allocation priority $\mathrm{su} \in \{0,1,2\}$ ($2$ for $\mathsf{staging}$ $\mathtt{01}s$, else $1$ for an uninitialized slot $\varsigma(\mathsf{ver}_3, a) = 0$, else $0$) with either presence of a $\mathsf{consuming}$ cell $\mathtt{10}b$ ($9 + \mathrm{su}$) or front $\mathsf{unconsumed}$ cell status $\in \{\text{none}, \mathtt{000}, \mathtt{001}\}$ ($3 u + \mathrm{su}$).
  \item \emph{Phase body, modular phase ring, payload, hub, and root:} Role $\mathsf{body}$ maps $F(p_{\mathrm{ver}}, L_{\mathrm{tmpl}}, j)$ to $1$ on $\mathtt{11}$, $3 + \varsigma(\mathsf{tmpl}, L_{\mathrm{tmpl}}) \in \{3,\dots,10\}$ on $\mathtt{10}$, and $2$ otherwise. On $\mathsf{ring}$, a phase frame $F(e_0, R', j)$ combines $\varsigma(\mathsf{body}, e_0)$ with $\varsigma(\mathsf{ring}, R') \in \{0,\dots,21\}$ without counting ring length $m$: state $0$ marks a non-frame terminator; states $1 + 2 b_{\mathrm{hd}} + b_{\mathrm{last}} \in \{1,\dots,4\}$ indicate that a $\mathsf{handoff}$ phase ($\mathtt{11}$) occurs, storing only the head-active bit $b_{\mathrm{hd}} = \mathbf{1}[e_0\text{ has status }\mathtt{10}]$ and the last-handoff bit $b_{\mathrm{last}} \in \{0,1\}$ testing whether the first $\mathtt{11}$ phase is immediately followed by the ring terminator ($b_{\mathrm{last}} = \mathbf{1}[\varsigma(\mathsf{ring}, R') = 0]$ when $\varsigma(\mathsf{body}, e_0) = 1$, and inherited from $\varsigma(\mathsf{ring}, R')$ modulo $2$ when $\varsigma(\mathsf{body}, e_0) \ne 1$); state $5$ marks a nonempty suffix with no $\mathtt{11}$ or $\mathtt{10}$ phase; and states $6 + 8 b_{\mathrm{hd}} + c_{\mathrm{tmpl}} \in \{6,\dots,21\}$ record $b_{\mathrm{hd}}$ together with the $8$-state template summary $c_{\mathrm{tmpl}}$ of the first $\mathtt{10}$ phase when no $\mathtt{11}$ phase occurs. Finally, $\mathsf{pay}$ outputs $2$ when $\varsigma(\mathsf{ring}, \mathrm{Ring}) \ge 6$ and $\varsigma(\mathsf{queue}, \mathrm{Queue}) < 3$ and $1$ on all other valid payload frames; $\mathsf{hub}$ outputs $1$ iff $q_{\mathrm{L}} = 13$ ($\mathrm{Flag}$ is an $F$-frame) and $q_{\mathrm{R}} \ge 7$; and $\mathsf{rt}$ copies $q_{\mathrm{R}}$. Exhaustive traversal from $r_0 = \mathsf{rt}$ confirms that exactly $271$ pairs $(r, q) \in R_{\mathrm{pocket}} \times Q_{\mathrm{pocket}}$ are inhabited on $\TS$.
\end{itemize}

At the root, $\varsigma(\mathsf{rt}, M) = 1 \in F_{\mathrm{halt}}$ holds if and only if $M \equiv P\,F(\mathrm{Flag}, \mathrm{Pay}, j_0)$ with $\mathrm{Flag}$ an $F$-frame, which coincides by Proposition~\ref{prop:halt-automaton} with $M \in L(\mathcal{H}_{11})$, so $\mathcal{T}_{\mathrm{pocket}}(M) = \mathsf{halt}$ holds if and only if $\sigma_{\mathrm{CTS}}(M) = \mathsf{halt}$.

When $M \notin L(\mathcal{H}_{11})$, we partition $P_{\mathrm{pocket}} = \{0,\dots,135\}$ into $54$ decision macro-states ($0 \dots 53$), $50$ two-edge frame-left stepping states ($54 \dots 103$), and $32$ three-edge frame-right stepping states ($104 \dots 135$). Entering the left slot $a$ of $F(a,b,j) \equiv (S a)((S b) j)$ for continuation $c \in \{0,\dots,49\}$ emits $\mathsf{left}(54 + c)$, after which state $54 + c$ emits $\mathsf{right}(c)$ along path $0 \cdot 1$; entering the right slot $b$ for $c \in \{0,\dots,31\}$ emits $\mathsf{right}(104 + c)$, after which state $104 + c$ emits $\mathsf{left}(54 + c)$ and state $54 + c$ emits $\mathsf{right}(c)$ along path $1 \cdot 0 \cdot 1$. Starting from $p_0 = 50$ at $M \equiv P H$, one right step reaches $H$ in state $51$: if $H$ or $\mathrm{Pay}$ is not an $F$-frame, both $\mathcal{T}_{\mathrm{pocket}}$ and $\sigma_{\mathrm{CTS}}$ return $\mathsf{stuck}$; if $\varsigma(\mathsf{pay}, \mathrm{Pay}) = 2$, state $51$ steps frame-left into $\mathrm{Flag}$ in opening state $p = 1$ (matching empty-queue halt opening at $[1,0,1]$); otherwise state $51$ steps frame-right into $\mathrm{Pay}$ in dispatch state $p = 0$.

At $\mathrm{Pay}$ ($p = 0$), $(q_{\mathrm{L}} - 1, q_{\mathrm{R}} - 7) = (\varsigma(\mathsf{ring}, \mathrm{Ring}), \varsigma(\mathsf{queue}, \mathrm{Queue}))$ determines every branch of Definition~\ref{def:pf-selector}(ii):
\begin{itemize}[leftmargin=2.0em]
  \item \emph{Modular handoff ($1 \le \varsigma(\mathsf{ring}, \mathrm{Ring}) \le 4$):} If $\varsigma(\mathsf{queue}, \mathrm{Queue}) \ge 9$, state $0$ steps frame-right to $\mathrm{Queue}$ in seeker state $p = 17$ to append $\mathtt{11}b$ to the first $\mathsf{consuming}$ cell. Otherwise it steps frame-left to $\mathrm{Ring}$ in handoff-root state $p = 32$. At the root phase $e_0$ with tail $R'$, if $\varsigma(\mathsf{body}, e_0) = 1$ ($j = 0$ is $\mathtt{11}$), then $\varsigma(\mathsf{ring}, R') = 0$ ($m = 1$) re-activates $e_0$ to $\mathtt{10}$, $b_{\mathrm{hd}}(R') = 1$ (phase $1$ already $\mathtt{10}$) deactivates $e_0$ to $\mathtt{00}$, and $b_{\mathrm{hd}}(R') = 0$ steps to the head of $R'$ ($p = 25$) to write $\mathtt{10}$. If $\varsigma(\mathsf{body}, e_0) \ne 1$ ($j > 0$), $b_{\mathrm{last}}(R') = 1$ means $j = m - 1$ and $(j + 1) \bmod m = 0$: if $e_0$ is already $\mathtt{10}$ ($\varsigma(\mathsf{body}, e_0) \ge 3$), state $32$ scans $R'$ ($p = 26$) to deactivate the first $\mathtt{11}$ phase to $\mathtt{00}$, else it activates $e_0$ at the root to $\mathtt{10}$. Finally, if $j > 0$ and $b_{\mathrm{last}}(R') = 0$ ($0 < j < m - 1$), state $32$ scans $R'$ ($p = 27$) to the first $\mathtt{11}$ phase and inspects $b_{\mathrm{hd}}$ of its right sibling, either deactivating phase $j$ to $\mathtt{00}$ or advancing one frame ($p = 25$) to activate phase $j + 1$ to $\mathtt{10}$.
  \item \emph{Active-phase dispatch ($\varsigma(\mathsf{ring}, \mathrm{Ring}) \ge 6$):} Decoding $(b_{\mathrm{idle}}, c_4)$ of the first $\mathtt{10}$ phase alongside $\varsigma(\mathsf{queue}, \mathrm{Queue}) \ge 3$ (when $< 9$, since $\varsigma(\mathsf{pay}, \mathrm{Pay}) \ne 2$) or $\varsigma(\mathsf{queue}, \mathrm{Queue}) \ge 9$ dispatches to the unique queue seeker $p \in \{16,\dots,21\}$ (locating the unopened queue tail, first $\mathsf{consuming}$, first $\mathsf{unconsumed}$, first $\mathsf{staging}$, or first uninitialized cell) or active-phase seeker $p \in \{28,\dots,31\}$ (descending via $p \in \{33,34,35\}$ to $\mathsf{ver}_2$ or via $p \in \{36,37,38\}$ and template seekers $p \in \{22,23,24\}$ to the first $\mathsf{idle}$, $\mathsf{armed}$, or $\mathsf{appending}$ template), matching rules~(ii)(a)--(b) of Definition~\ref{def:pf-selector}.
  \item \emph{Version-chain extension and local generator opening:} Width-$2$ writers $p \in \{39,40,41\}$ and width-$3$ writers $p \in \{42,\dots,49\}$ walk the left spine while $\varsigma(\mathsf{fld}, b) \ge 2^w$ and enter $\mathsf{fld}$ in states $p \in \{2,\dots,15\}$ on the first incomplete record (or open an uninitialized chain leaf), stepping down $\mathsf{fld}$ until reaching the first $1$-bit slot whose left child is not yet an $F$-frame ($\varsigma(\mathsf{sp}, a) \ne 6$) or an unopened field/sentinel generator. On any unopened or mid-opening generator, states $p \in \{1, 52, 53\}$ inspect $(\pi_{\mathrm{L}}(q_{\mathrm{L}}), \pi_{\mathrm{R}}(q_{\mathrm{R}}))$ to execute the depth-$\le 5$ redex probe of Lemma~\ref{lem:frame-open-bound}(i): emitting $\mathsf{contract}$ when $\pi_{\mathrm{L}}(q_{\mathrm{L}}) = 2$ ($S X Y Z$), $\mathsf{left}(53)$ when $\pi_{\mathrm{L}}(q_{\mathrm{L}}) = 3$ ($S X Y Z\,W$, where state $53$ contracts at $\pi_{\mathrm{L}} = 2$), $\mathsf{right}(52)$ when $(\pi_{\mathrm{L}}, \pi_{\mathrm{R}}) = (1, 4)$ ($S\,(S X Y Z\,W)$), and $\mathsf{right}(53)$ when $(\pi_{\mathrm{L}}, \pi_{\mathrm{R}}) = (1, 3)$ ($S\,(S X Y Z)$).
\end{itemize}
Thus $\mathrm{Loc}_{\mathcal{T}_{\mathrm{pocket}}}(p_0, \widehat{M}_0)$ returns the exact redex address (or $\bot$) selected by $\sigma_{\mathrm{CTS}}(M)$, proving $\mathcal{T}_{\mathrm{pocket}}(M) = \sigma_{\mathrm{CTS}}(M)$ for every closed $M \in \TS$. Because $M \in \TS$ and $M \redS M'$ imply $M' \in \TS$, induction on $k \ge 0$ shows that the $k$-step trajectories of $\mathcal{T}_{\mathrm{pocket}}$ and $\sigma_{\mathrm{CTS}}$ starting from $T_0 = \mathcal{E}_{\mathrm{CTS}}(P, w) \in \TS$ coincide identically; Proposition~\ref{prop:seed-genesis} and Theorem~\ref{thm:cts-universality} therefore yield $D_{\mathrm{seed}} \redSs T_0$, bidirectional halting equivalence, terminal readback $\mathcal{D}_{\mathrm{CTS}}(T_k) = (i_{\mathrm{fin}}, \varepsilon)$, and the quadratic trajectory bound $|T_k| \le 35{,}000(\|(P, w)\| + k + 1)^2$.
\end{proof}

\begin{remark}[Helper-state compression and linear transition complexity]\label{rem:pocket-helper-compression}
Relative to the $4{,}376{,}490{,}905$-state root-restarted zipper walker $\sigma_{\mathrm{TM}}$ of Section~\ref{sec:tm-sim}, $\mathcal{T}_{\mathrm{pocket}}$ removes the universal transition table from the controller and compresses the reachable state footprint to $271$ role-state pairs in Pass~1 together with $136$ single-edge pointer states in Pass~2 ($407$ active state configurations in total, a reduction exceeding $1.075 \times 10^7$-fold), while tightening the per-contraction transition bound from $3500(2|M| - 1)$ cursor micro-ticks to at most $2|M| - 3 + \mathrm{depth}(M)$ transitions and less than $5$ bits of Pass-1 state per tree node.
\end{remark}

\section{The Locus of Control: Leaf-Splitting Separation and Multiway Impossibility Theorems}\label{sec:separation}
% ============================================================================

Theorems~\ref{thm:tm-universality} and~\ref{thm:cts-universality} establish universal computation along pure-$S$ trajectories selected by root-restarted walkers. To isolate the respective contributions of external step selection and combinatory duplication, we prove three complementary separation and impossibility theorems.

\subsection{Leaf-Splitting Universality Under Root-Restarted Selection}\label{subsec:leaf-split}

We first isolate the computational power of the root-restarted finite-state tree-walking control discipline itself by replacing $\CLS$ with a minimal non-erasing binary-tree rewrite system: unary leaf splitting, which performs neither subtree duplication nor higher-arity interaction.

\begin{definition}[Unary leaf-splitting system and $3$-state halting automaton]\label{def:btree-split}
Let $\Tbin$ be the set of unlabeled finite full binary trees generated by $t ::= \mathsf{Leaf} \mid \mathsf{Node}(t_1, t_2)$, with leaf count $|\mathsf{Leaf}|_{\mathrm{leaf}} = 1$ and $|\mathsf{Node}(t_1,t_2)|_{\mathrm{leaf}} = |t_1|_{\mathrm{leaf}} + |t_2|_{\mathrm{leaf}}$. Define single-step \emph{leaf-splitting reduction} $\to_{\mathrm{sp}}$ as the contextual closure of the sole local rule
\begin{equation}\label{eq:leaf-split-rule}
  \mathsf{Leaf} \;\to_{\mathrm{sp}}\; \mathsf{Node}(\mathsf{Leaf},\, \mathsf{Leaf}).
\end{equation}
Every step $t \to_{\mathrm{sp}} t'$ satisfies $|t'|_{\mathrm{leaf}} = |t|_{\mathrm{leaf}} + 1$. A \emph{root-restarted finite-state tree-walking selector} over $(\Tbin, \to_{\mathrm{sp}})$ is a finite control set $Q_{\mathrm{walk}}$ with start state $q_{\mathrm{init}}$ and transition function
\[
  \delta_{\mathrm{walk}} \colon Q_{\mathrm{walk}} \times \{\mathsf{true}, \mathsf{false}\} \times \{\mathsf{root}, \mathsf{leftChild}, \mathsf{rightChild}\} \;\longrightarrow\; \mathsf{downL}(Q_{\mathrm{walk}}) \cup \mathsf{downR}(Q_{\mathrm{walk}}) \cup \mathsf{up}(Q_{\mathrm{walk}}) \cup \{\mathsf{split},\, \mathsf{halt}\},
\]
which starts at $(q_{\mathrm{init}}, t, \varepsilon)$, moves a Huet zipper one edge at a time based solely on whether the focused subtree is $\mathsf{Leaf}$ and which child edge leads to its parent, and upon executing $\mathsf{split}$ at a focused $\mathsf{Leaf}$ replaces that leaf by $\mathsf{Node}(\mathsf{Leaf},\mathsf{Leaf})$ and resets its control state to $(q_{\mathrm{init}}, t', \varepsilon)$.

Define the $3$-state deterministic bottom-up tree automaton $\mathcal{H}_3 = (\{s_{\mathrm{L}}, s_{\mathrm{N}}, s_{\mathrm{acc}}\}, s_{\mathrm{L}}, \odot_3, \{s_{\mathrm{acc}}\})$ on $\Tbin$ by $s_{\mathrm{L}} \odot_3 y = s_{\mathrm{N}}$ and $x \odot_3 y = s_{\mathrm{acc}}$ for $x \in \{s_{\mathrm{N}}, s_{\mathrm{acc}}\}$, so that $L(\mathcal{H}_3) = \{\mathsf{Node}(\mathsf{Node}(a,b), r) : a, b, r \in \Tbin\}$.
\end{definition}

\begin{theorem}[Controller-complexity separation via leaf splitting]\label{thm:leaf-split-univ}
There exist a linear-size initial tree encoder $\mathcal{E}_{\mathrm{sp}}$ and a fixed root-restarted finite-state tree-walking selector $\sigma_{\mathrm{sp}}$ over $(\Tbin, \to_{\mathrm{sp}})$ simulating the $132$-state universal binary Turing machine $U_{\mathrm{bin}}$ (and hence every deterministic Turing machine) with exact bidirectional halting equivalence: for every initial configuration $c_0$ of $U_{\mathrm{bin}}$, $U_{\mathrm{bin}}$ halts from $c_0$ if and only if the root-restarted leaf-splitting trajectory $B_0 = \mathcal{E}_{\mathrm{sp}}(c_0)$, $B_{k+1} = \sigma_{\mathrm{sp}}(B_k)$ with $B_k \to_{\mathrm{sp}} B_{k+1}$ reaches some $B_k \in L(\mathcal{H}_3)$.
\end{theorem}

\begin{proof}
Encode the $4$-status cell alphabet of $U_{\mathrm{bin}}$---consisting of $\mathsf{off}(b)$, $\mathsf{active}(b, q)$, $\mathsf{movingR}(b, q)$, and $\mathsf{movingL}(b, q)$ for $b \in \{0,1\}$ and $q \in \{0,\dots,131\}$---into the $1056$-element integer range $c \in \{0,\dots,1055\}$ via $\mathsf{off}(0) \mapsto 0$, $\mathsf{off}(1) \mapsto 1$, and $\mathsf{tag}(b,q) \mapsto \mathrm{tag} + 8q$ for $\mathrm{tag} \in \{2,\dots,7\}$. Represent any natural number $m \in \mathbb{N}$ by the right-leaning unary spine $U_0 = \mathsf{Leaf}$, $U_{m+1} = \mathsf{Node}(\mathsf{Leaf}, U_m)$. A \emph{committed cell record} with code $c$ followed by version tail $\mathrm{rest} \in \Tbin$ is the tree
\[
  \mathrm{Rec}(c, \mathrm{rest}) \;:=\; \mathsf{Node}\Bigl(\mathsf{Node}\bigl(\mathsf{Node}(\mathsf{Leaf},\mathsf{Leaf}),\, U_c\bigr),\; \mathrm{rest}\Bigr),
\]
whose commit flag (the left child of the inner payload pair) is the composite node $\mathsf{Node}(\mathsf{Leaf},\mathsf{Leaf})$. A cell with history $[c_1, \dots, c_r]$ and uncommitted tail $\tau$ is $\mathrm{Cell}([c_1,\dots,c_r], \tau) = \mathrm{Rec}(c_1, \dots \mathrm{Rec}(c_r, \tau)\dots)$, chained along a right-leaning tape spine $\mathrm{Tape} = \mathsf{Node}(\mathrm{Cell}_0, \mathsf{Node}(\mathrm{Cell}_1, \dots \mathsf{Leaf}))$, under a global root $\mathsf{Node}(\mathrm{Flag}, \mathrm{Tape})$ with $\mathrm{Flag} = \mathsf{Leaf}$ while active (so $\mathsf{Node}(\mathsf{Leaf}, \mathrm{Tape}) \notin L(\mathcal{H}_3)$).

Appending a new committed code $c \in \{0,\dots,1055\}$ to a clean cell tail $\tau = \mathsf{Leaf}$ involves no subtree duplication and executes atomically across $c + 3$ root-restarted leaf splits while preserving the previously committed value on all $c + 2$ intermediate steps:
\begin{enumerate}[label=(\roman*),leftmargin=2.2em]
  \item \emph{Record-box allocation (split~$1$):} $\tau = \mathsf{Leaf} \to_{\mathrm{sp}} \mathsf{Node}(\mathsf{Leaf}, \mathsf{Leaf})$ (uncommitted: inner left child is $\mathsf{Leaf}$).
  \item \emph{Staging-pair allocation (split~$2$):} split the left child to obtain $\mathsf{Node}(\mathsf{Node}(\mathsf{Leaf}, U_0), \mathsf{Leaf})$ (staging count $m = 0$; commit flag is still $\mathsf{Leaf}$).
  \item \emph{Unary payload countdown (splits~$3$ through $c+2$, empty when $c = 0$):} on each invocation, $\sigma_{\mathrm{sp}}$ walks down the uncommitted staging spine $U_m$ incrementing a bounded counter in $\{0,\dots,1055\}$; while $m < c$, it splits the terminal $\mathsf{Leaf}$ of $U_m$ to produce $U_{m+1}$ and resets to the root. Because the commit flag remains $\mathsf{Leaf}$, every read pass skips the uncommitted tail and still reads $c_r$.
  \item \emph{Atomic commit (split~$c+3$):} when the walk reaches the end of $U_m$ with $m = c$, $\sigma_{\mathrm{sp}}$ ascends to the staging pair and splits its left-child commit flag $\mathsf{Leaf} \to_{\mathrm{sp}} \mathsf{Node}(\mathsf{Leaf},\mathsf{Leaf})$, turning the tail into $\mathrm{Rec}(c, \mathsf{Leaf})$ in a single step.
\end{enumerate}
Equipped with this $(c+3)$-split atomic write primitive (plus a $1$-split right-tape extension when $\mathsf{movingR}$ reaches the tape sentinel $\mathsf{Leaf}$), $\sigma_{\mathrm{sp}}$ executes the same two-pass priority simulation of $U_{\mathrm{bin}}$ as $\sigma_{\mathrm{TM}}$: Pass~1 completes any $\mathsf{movingR}$ or $\mathsf{movingL}$ handshake, and Pass~2 steps the unique $\mathsf{active}$ cell---or, when $\delta_{U_{\mathrm{bin}}}(q, b)$ is undefined, ascends to the global root and splits its left child $\mathrm{Flag}$ from $\mathsf{Leaf}$ to $\mathsf{Node}(\mathsf{Leaf},\mathsf{Leaf})$, placing the root tree into $L(\mathcal{H}_3)$. Constructive bounded evaluation of $U_{\mathrm{bin}}$ together with determinism of $\sigma_{\mathrm{sp}}$ yields exact bidirectional halting equivalence.
\end{proof}

\begin{remark}[Division of computational labor between controller and rewrite rule]\label{rem:control-division}
Theorem~\ref{thm:leaf-split-univ} clarifies the relationship between controlled pure-$S$ universality and tree-walking automata theory. Read-only finite-state tree-walking automata cannot recognize all regular tree languages \cite{BojanczykColcombet2008}, whereas Kutrib and Meyer showed that equipping a finite control with a dynamically growing tree-walking storage already recognizes non-semilinear unary languages strictly beyond regular languages even under real-time non-erasing operation \cite{KutribMeyer2023,KutribMeyer2026}. Theorem~\ref{thm:leaf-split-univ} establishes that when multi-pass stateless root-restarted navigation is allowed to mutate an unlabeled tree solely via $\mathsf{Leaf} \to_{\mathrm{sp}} \mathsf{Node}(\mathsf{Leaf},\mathsf{Leaf})$, the resulting system is Turing-complete. Thus, when a universal machine table is permitted inside the finite controller (as in \cite{Filin2026} and Theorem~\ref{thm:tm-universality}), universality stems from addressable growing tree storage rather than from combinatory dynamics. What $\CLS$ genuinely contributes beyond $\to_{\mathrm{sp}}$ is in-tree duplicative computation $S X Y Z \redS (X Z)(Y Z)$, which enables both the program-free selector $\sigma_{\mathrm{CTS}}$ and universal seed genesis from $D_{\mathrm{seed}}$ in Section~\ref{sec:cts-sim}.
\end{remark}

\subsection{Constant-Step Multiway Halting Shortcuts}\label{subsec:multiway-shortcuts}

We next show that removing external evaluation-order selection in favor of unrestricted multiway rewriting $\redSs$ causes immediate soundness collapse across both major pure-$S$ encoding architectures: unopened-generator halt flags (Sections~\ref{sec:tm-sim}--\ref{sec:cts-sim}) and live/tombstone $D \cdot Q$ tag dispatchers \cite{Filin2026}.

\begin{proposition}[Seven-step and five-step multiway halting shortcuts]\label{prop:multiway-shortcuts}
Under unrestricted pure-$S$ reduction $\redSs$, the following constant-step shortcuts hold for arbitrary (including non-terminating) program and queue subterms:
\begin{enumerate}[label=(\roman*),leftmargin=2.2em]
  \item \emph{Protected-trie halt-flag shortcut ($7$ contractions):} For every integer $k \ge 1$ and arbitrary terms $P_{\mathrm{ring}}, Q_{\mathrm{root}}, j_0, j_1 \in \TS$, the hub term $F(F(G(0, k+1), P_{\mathrm{ring}}, j_1), Q_{\mathrm{root}}, j_0)$ reduces in $7$ contractions to $F(F(H_{\mathrm{open}}, P_{\mathrm{ring}}, j_1), Q_{\mathrm{root}}, j_0)$, where $H_{\mathrm{open}} \equiv F(g, g, (\nu_k S) g)$ with $g \equiv G(k+1, k+2)$ is an opened frame, without inspecting or contracting a single redex inside $P_{\mathrm{ring}}$ or $Q_{\mathrm{root}}$.
  \item \emph{Tag-dispatcher queue-deletion shortcut ($5$ contractions):} Let $v_1 \equiv S(S(S T T))$, let $L_1(X) \equiv (T v_1) X$ be the live $1$-bit queue cell with empty-queue tombstone $Q_{\mathrm{empty}} \equiv S S (v_1 S)$, let $A_0 \equiv S(S K_{\mathrm{cont}}) S$ be the $0$-bit empty-append branch for any continuation $K_{\mathrm{cont}} \in \TS$, and let $D(A_0, A_1) \equiv T(S A_0 A_1)$ be the binary queue dispatcher for any non-halting $1$-bit append action $A_1 \in \TS$. Then $D(A_0, A_1)\,(L_1(S))$ reduces in exactly $5$ contractions to
  \[
    (S Q_{\mathrm{empty}})\,\Bigl(\bigl((K_{\mathrm{cont}}\,(S Q_{\mathrm{empty}}))\,(Q_{\mathrm{empty}}\,(S Q_{\mathrm{empty}}))\bigr)\,(A_1 Q_{\mathrm{empty}})\Bigr),
  \]
  converting the live $1$-bit queue into the empty tombstone $Q_{\mathrm{empty}}$ and firing the $0$-branch $K_{\mathrm{cont}}\,(S Q_{\mathrm{empty}})$ regardless of $A_1$.
\end{enumerate}
\end{proposition}

\begin{proof}
For~(i), let $c_1 \equiv \nu_{k+1} \equiv T \nu_k$, $n_2 \equiv T c_1 \equiv \nu_{k+2}$, $q \equiv c_1 n_2$, and $g \equiv c_1 n_2 S$. Inside the context $C[\square] \equiv F(F(\square, P_{\mathrm{ring}}, j_1), Q_{\mathrm{root}}, j_0)$, the unopened generator $G(0, k+1) \equiv (S T T)\,c_1\,S$ admits the $7$-step pure-$S$ derivation
\begin{align*}
  (S T T)\,c_1\,S
  &\;\redS\; \bigl((T c_1)(T c_1)\bigr)\,S \;\equiv\; (T c_1)\,n_2\,S
   \;\redS\; (S n_2 q)\,S
   \;\redS\; (n_2 S)\,g \\
  &\;\redS\; \bigl(T(c_1 S)\bigr)\,g
   \;\redS\; (S g)\,\bigl((c_1 S)\,g\bigr)
   \;\redS\; (S g)\,\bigl((T(\nu_k S))\,g\bigr) \\
  &\;\redS\; (S g)\,\bigl(S g\,((\nu_k S)\,g)\bigr) \;\equiv\; H_{\mathrm{open}},
\end{align*}
lifting by contextual closure to $C[G(0,k+1)] \redS^7 C[H_{\mathrm{open}}]$ without touching $P_{\mathrm{ring}}$ or $Q_{\mathrm{root}}$.

For~(ii), instead of contracting the outer dispatcher redex first (which would inspect the live tag $v_1$ and select $A_1$), an unguided reduction contracts the argument redex $L_1(S) \equiv S S v_1 S \to_r (S S)(v_1 S) \equiv Q_{\mathrm{empty}}$ at step~$1$, and then dispatches on $Q_{\mathrm{empty}}$ across steps~$2$--$5$:
\begin{align*}
  D(A_0, A_1)\,\bigl(L_1(S)\bigr)
  &\;\redS\; \bigl(S S (S A_0 A_1)\bigr)\,Q_{\mathrm{empty}} \\
  &\;\to_r\; (S Q_{\mathrm{empty}})\,\bigl((S A_0 A_1)\,Q_{\mathrm{empty}}\bigr) \\
  &\;\redS\; (S Q_{\mathrm{empty}})\,\bigl((A_0 Q_{\mathrm{empty}})\,(A_1 Q_{\mathrm{empty}})\bigr) \\
  &\;\equiv\; (S Q_{\mathrm{empty}})\,\Bigl(\bigl(S(S K_{\mathrm{cont}})\,S\,Q_{\mathrm{empty}}\bigr)\,(A_1 Q_{\mathrm{empty}})\Bigr) \\
  &\;\redS\; (S Q_{\mathrm{empty}})\,\Bigl(\bigl((S K_{\mathrm{cont}} Q_{\mathrm{empty}})\,(S Q_{\mathrm{empty}})\bigr)\,(A_1 Q_{\mathrm{empty}})\Bigr) \\
  &\;\redS\; (S Q_{\mathrm{empty}})\,\Bigl(\bigl((K_{\mathrm{cont}}\,(S Q_{\mathrm{empty}}))\,(Q_{\mathrm{empty}}\,(S Q_{\mathrm{empty}}))\bigr)\,(A_1 Q_{\mathrm{empty}})\Bigr). \qedhere
\end{align*}
\end{proof}

\subsection{Non-Existence of Asymmetric Branch Gating in \texorpdfstring{$\CLS$}{CL(S)}}\label{subsec:no-church-bool}

Proposition~\ref{prop:multiway-shortcuts}(ii) raises the question of whether one could redesign pure-$S$ branch dispatchers using ``Church-Boolean'' thunk gates $B_0, B_1 \in \TS$ so that $B_0\,x_0\,x_1$ activates only the first branch thunk $x_0$ while $B_1\,x_0\,x_1$ activates only the second branch thunk $x_1$ (without triggering a divergent computation hidden behind $x_0$). On the first branch, $T \equiv S S$ succeeds in one step:
\[
  T\,x_0\,x_1 \;\equiv\; S S x_0 x_1 \;\to_r\; (S x_1)\,(x_0 x_1) \;\in\; \NF,
\]
placing $x_0$ into active function position ($x_0 x_1$) while $x_1$ occurs only passively as an argument. We now show that no symmetric partner $B_1$ exists in $\CLS$, even when supplied with arbitrarily many closed auxiliary padding arguments $U_1, \dots, U_m$.

\begin{theorem}[Impossibility of asymmetric Church-Boolean branch gating in $\CLS$]\label{thm:no-church-bool}
For every variable $x_i \in \mathcal{V}$ and term $P \in \TSV$, say that $x_i$ is \emph{actively called} in $P$, written $\mathrm{call}_i(P)$, if $P$ contains a subterm of the form $x_i W$ for some $W \in \TSV$. Let $M \in \TSV$ and $U_1, \dots, U_m \in \TSV$ for $m \ge 0$ contain neither $x_0$ nor $x_1$. Then for every multiway reduction
\[
  M\,x_0\,x_1\,U_1\cdots U_m \;\redSs\; N,
\]
if $\mathrm{call}_1(N)$ holds, then $\mathrm{call}_0(N)$ already holds. In particular, there do not exist a closed combinator $B_1 \in \TS$, closed auxiliary arguments $U_1, \dots, U_m \in \TS$, and a reduct $N$ of $B_1\,x_0\,x_1\,U_1\cdots U_m$ that actively calls $x_1$ without actively calling $x_0$.
\end{theorem}

\begin{proof}
First observe that for any variable $x_i$, the predicate $\mathrm{call}_i(-)$ is monotone under $\redS$: in a root step $S X Y Z \to_r (X Z)(Y Z)$, neither $S$, $S X$, nor $S X Y$ is a variable $x_i$, so any subterm $x_i W$ inside $X$, $Y$, or $Z$ reappears as $x_i W$ inside $X Z$ or $Y Z$; contextual steps $L R \redS L' R$ (where $L \not\equiv x_i$ since variables are normal) and $L R \redS L R'$ likewise preserve $\mathrm{call}_i$. Hence once $\mathrm{call}_0(P)$ becomes true at any step, it remains true in every subsequent reduct.

Next, let $\mathcal{F} \subset \TSV$ denote the set of terms containing neither $x_0$ nor $x_1$; by variable conservation under non-erasing $S$-reduction, $\mathcal{F}$ is closed under $\redS$. Define inductively two subsets $\mathcal{W}_0, \mathcal{O} \subset \TSV$:
\begin{enumerate}[label=(\roman*),leftmargin=2.2em]
  \item the class $\mathcal{W}_0$ is generated by $C x_0 \in \mathcal{W}_0$ for every $C \in \mathcal{F}$, and $W_1 W_2 \in \mathcal{W}_0$ whenever $W_1, W_2 \in \mathcal{W}_0$;
  \item the class $\mathcal{O}$ is generated by $W x_1 \in \mathcal{O}$ for every $W \in \mathcal{W}_0$, and $P U \in \mathcal{O}$ whenever $P \in \mathcal{O}$ and $U \in \mathcal{F}$.
\end{enumerate}
Every term in $\mathcal{W}_0$ is an application (so $S \notin \mathcal{W}_0$ and $x_k \notin \mathcal{W}_0$) and contains no occurrence of $x_1$. Consequently:
\begin{itemize}[leftmargin=2.0em]
  \item If $S P \in \mathcal{W}_0$, then since $S \notin \mathcal{W}_0$, it cannot arise from $W_1 W_2$, so $S P \equiv C x_0$ and hence $P \equiv x_0$.
  \item If $S P Q \in \mathcal{W}_0$, either $S P Q \equiv C x_0$ (giving $Q \equiv x_0$) or $S P \in \mathcal{W}_0$ (giving $P \equiv x_0$). Thus \emph{every} arity-$2$ prefix $S P Q \in \mathcal{W}_0$ satisfies $P \equiv x_0$ or $Q \equiv x_0$.
  \item Similarly, every $T' \in \mathcal{O}$ has $\ls(T') \ge 2$ (so $T' \not\equiv S$, $T' \not\equiv x_1$, and $T' \not\equiv S P$), and if $T' \equiv S P Q \in \mathcal{O}$, inversion forces $T' \equiv W x_1$ with $W \equiv S P \in \mathcal{W}_0$, whence $P \equiv x_0$. Moreover, no term in $\mathcal{O}$ satisfies $\mathrm{call}_1$: in $W x_1$, $W \not\equiv x_1$ and $x_1$ is absent from $W$; in $P U$, $P \not\equiv x_1$ and $x_1$ is absent from $U \in \mathcal{F}$.
\end{itemize}
We claim that if $W \in \mathcal{W}_0$ and $W \redS W'$, then either $W' \in \mathcal{W}_0$ or $\mathrm{call}_0(W')$ holds. By structural induction on $W \in \mathcal{W}_0$: if $W \equiv C x_0$ ($C \in \mathcal{F}$), either $C \redS C' \in \mathcal{F}$ (giving $C' x_0 \in \mathcal{W}_0$) or $C \equiv S X Y$ with $X, Y \in \mathcal{F}$ and $S X Y x_0 \to_r (X x_0)(Y x_0) \in \mathcal{W}_0$; if $W \equiv W_1 W_2$, contextual steps in $W_1$ or $W_2$ preserve $\mathcal{W}_0$ or establish $\mathrm{call}_0$ by induction, whereas a root step requires $W_1 \equiv S P Q \in \mathcal{W}_0$, so $P \equiv x_0$ or $Q \equiv x_0$, and the contractum $(P W_2)(Q W_2)$ contains $x_0 W_2$, making $\mathrm{call}_0(W')$ true.

Lifting by structural induction to $T' \in \mathcal{O}$ shows that whenever $T' \redS T''$, either $T'' \in \mathcal{O}$ or $\mathrm{call}_0(T'')$ holds: for $T' \equiv W x_1$, a left step $W \redS W'$ yields $W' x_1 \in \mathcal{O}$ or $\mathrm{call}_0(W' x_1)$, while a root step forces $W \equiv S P Q \in \mathcal{W}_0$, hence $P \equiv x_0$ or $Q \equiv x_0$, producing $x_0 x_1$ inside $(P x_1)(Q x_1)$; for $T' \equiv P U$ ($U \in \mathcal{F}$), left steps use the induction hypothesis, right steps $U \redS U' \in \mathcal{F}$ preserve $\mathcal{O}$, and a root step forces $P \equiv S X Y \in \mathcal{O}$, whence $X \equiv x_0$ and the contractum $(x_0 U)(Y U)$ satisfies $\mathrm{call}_0$.

Finally, $M x_0 \in \mathcal{W}_0$, $(M x_0) x_1 \in \mathcal{O}$, and appending $U_1, \dots, U_m \in \mathcal{F}$ places $T_{\mathrm{start}} \equiv M x_0 x_1 U_1\cdots U_m \in \mathcal{O}$. By induction along $T_{\mathrm{start}} \redSs N$ combined with $\redS$-monotonicity of $\mathrm{call}_0$, either $N \in \mathcal{O}$---which contradicts $\mathrm{call}_1(N)$ since $\mathrm{call}_1$ vanishes on $\mathcal{O}$---or $\mathrm{call}_0(N)$ holds.
\end{proof}

% ============================================================================

% ============================================================================
\section{Finite \texorpdfstring{$S$}{S}-Algebras, Automaton-Free Book Proofs, B\"ohm Faithfulness Dichotomy, and Order Hierarchy}\label{sec:algebras}
% ============================================================================

When external tree-walking controllers are removed and closed $S$-terms evolve solely under the intrinsic rewrite relation $\redS$, the dynamics of $\CLS$ is governed by a hierarchy of finite combinatory $S$-algebras together with elementary syntactic invariants. Because every ground $S$-term is built from a single leaf constant $S$ by binary application, any finite $S$-algebra $\mathcal{A} = (Q, q_0, \odot)$ (Definition~\ref{def:s-algebra}) defines a deterministic bottom-up tree automaton running in $2|M| - 1$ table lookups on unfolded syntax trees (and in linear DAG size on shared representations) while remaining invariant under unrestricted conversion $\convS$ by~\eqref{eq:s-algebra-law}. In this section we develop this algebraic and syntactic theory constructively without K\"onig's lemma or classical choice. First, we establish collapsed finite-rank normalization for Waldmann's $39$-state normalization algebra~\cite{Waldmann2000,BarendregtEndrullisKlopWaldmann2017}, prove automaton-free head conservation, and construct a pointwise-minimal pure multiplier potential ($B \equiv 0$) resolving Question~10 and Footnote~4 of Barendregt, Endrullis, Klop, and Waldmann~\cite{BarendregtEndrullisKlopWaldmann2017}. Next, we compress completeness of the $39$-state productivity algebra for Question~5 to $10$ ranks while proving $113$-state Myhill--Nerode minimality of its joint solvability product. Finally, we give short automaton-free proofs of Questions~7 and~31 alongside a complete B\"ohm-tree faithfulness dichotomy and a constructive convertibility decider across $\Prod \times \TS$ advancing Scott's Question~30, and we determine the minimal $S$-algebras of sizes $35$, $12$, and $3$ governing the functional order hierarchy $\{0, 1, 2, 3\}$.

Table~\ref{tab:algebra-tower} summarizes the finite $S$-algebras, deterministic tree automata, and minimal quotient or product subalgebras studied across Sections~\ref{sec:algebras}--\ref{sec:asymptotics}. Given $F \subseteq Q$, we write $L(\mathcal{A}, F) = \{ M \in \TS : \sem{M}_{\mathcal{A}} \in F \}$, retaining the abbreviations $T \equiv S S$, $A \equiv S S S$, right combs $R_n$ ($n \ge 1$, $|R_n| = n$), leaf count $|M|$, left-spine arity $\ls(M)$, and right-spine depth $\rs(M)$ from Section~\ref{sec:prelim}.

\begin{table}[htbp]
\centering
\caption{Finite $S$-algebras, deterministic bottom-up tree automata, and minimal quotient or product subalgebras in $\CLS$.}
\label{tab:algebra-tower}
\small
\begin{tabular}{@{}lccl@{}}
\toprule
Structure & $|Q|$ & $S$-Law & Recognized Property or Algebraic Invariant on $\TS$ \\
\midrule
$\mathcal{A}_{\mathrm{ord},3}$ & $3$ & Yes & Minimal $S$-algebra for functional order~$3$ (right combs $R_n$; Theorem~\ref{thm:order-hierarchy}) \\
$\mathcal{D}_4$ & $4$ & Yes & Functionally complete Boolean algebra realizing $11$-leaf normal-form Rule~$110$ \\
$\mathcal{B}_Z$ & $8$ & DTA & Accepts step-$t$ reducts of $Z_{\mathrm{Zac}}$ iff $t = \tfrac{1}{2}n(n+3)$ across six strategies \\
$\mathcal{A}_{\mathrm{ord},2}$ & $12$ & Yes & Minimal $S$-algebra deciding $\{\ord = 2\}$ and $\{\ord \ge 2\}$ (Theorem~\ref{thm:order-hierarchy}) \\
$\mathcal{A}_{\HN}$ & $35$ & Yes & Minimal $S$-algebra for $\HN$, $\{\ord = 0\}$, $\{\ord = 1\}$, and $4$-valued $\ord$ (Question~10) \\
$\mathcal{A}_{36}$ & $36$ & Yes & Refines $\infty_{\HN}$ into $\{q_{34}, q_{35}\}$, separating $A A A \neq_S A A A S$ (Question~31) \\
$\mathcal{A}_{\SN}$ & $39$ & Yes & Minimal $S$-algebra for strong and weak normalization ($\sem{M}_{\SN} \neq \infty_{\SN}$; Theorem~\ref{thm:waldmann-sn}) \\
$\mathcal{A}_P$ & $39$ & Yes & B\"ohm-tree productivity ($\sem{M}_P \neq \infty_P$; Question~5, Theorem~\ref{thm:q5-main}) \\
$\mathcal{A}_P \times \mathcal{A}_{\SN}$ & $44$ & Yes & Reachable pairs isolating $(q_{29}, \infty_{\SN})$ as unique non-$\SN$ productive redex fiber \\
$\mathcal{A}_{\SN} \times \mathcal{A}_{\HN}$ & $112$ & Yes & Reachable pairs: $94$ in $\SN$, $17$ in $\HN \setminus \SN$, and $1$ in $\overline{\HN}$ \\
$\mathcal{A}_P \times \mathcal{A}_{\HN}$ & $113$ & Yes & Myhill--Nerode minimal $(\Prod, \HN)$ algebra ($95$ accepting, $18$ rejecting; Proposition~\ref{prop:ph-113-minimal}) \\
$\mathcal{A}_{\SN} \times \mathcal{A}_P \times \mathcal{A}_{\HN}$ & $146$ & Yes & Reachable four-phase hierarchy: $98\,\SN + 30\,(\Prod\setminus\SN) + 17\,(\HN\setminus\Prod) + 1\,\overline{\HN}$ \\
\bottomrule
\end{tabular}
\end{table}

\subsection{Strong Normalization and the 39-State Algebra \texorpdfstring{$\mathcal{A}_{\SN}$}{A\_SN}}\label{subsec:sn-algebra}

Let $\mathcal{N} := \TS \cap \NF$ denote the set of closed normal forms, generated inductively by the grammar
\begin{equation}\label{eq:cnf-grammar}
\mathcal{N} \;::=\; S \;\mid\; S X \;\mid\; S X Y
\qquad (X, Y \in \mathcal{N}).
\end{equation}
Indeed, every closed $S$-term $M$ decomposes along its left spine as $S M_1 \dots M_k$ where $k = \ls(M) \ge 0$: if $k \ge 3$, its prefix $S M_1 M_2 M_3$ is already a redex, whereas if $k \le 2$, then $M$ contains no redex if and only if each spine argument $M_i$ lies in $\mathcal{N}$.

\begin{lemma}[Subterm inheritance and decomposition of strong normalization]\label{lem:wn-eq-sn}
For all open or closed $S$-terms $M, M', L, R \in \TSV$:
\begin{enumerate}[label=(\roman*),leftmargin=2.2em]
  \item \emph{Subterm inheritance:} Recall from Proposition~\ref{prop:cr-conservation}(ii) that $M \redS M'$ implies $|M| \le |M'|$. A binary application $L R$ lies in $\SN$ if and only if $L \in \SN$, $R \in \SN$, and every reduct of $L R$ formed after a root contraction lies in $\SN$; in particular, every subterm of a strongly normalizing $S$-term is strongly normalizing.
  \item \emph{Equivalence of $\WN$ and $\SN$:} By Proposition~\ref{prop:cr-conservation}(iii), $M \in \WN$ holds if and only if $M \in \SN$.
\end{enumerate}
\end{lemma}

\begin{proof}
Part~(ii) was established constructively via weighted parallel projection in Proposition~\ref{prop:cr-conservation}(iii). For~(i), if an infinite reduction starts inside $L$ (or $R$), applying $R$ on the right (or $L$ on the left) lifts it step-for-step to an infinite reduction of $L R$; hence $L R \in \SN$ implies $L \in \SN$ and $R \in \SN$, and structural induction extends this inheritance to all subterms. Conversely, well-founded induction on $(L, R) \in \SN \times \SN$ together with the hypothesis on post-root-contraction reducts shows $L R \in \SN$.
\end{proof}

We define the $39$-state combinatory $S$-algebra $\mathcal{A}_{\SN} = (Q_{\SN}, q_0, \odot_{\SN})$ with carrier $Q_{\SN} = \{q_0, q_1, \dots, q_{37}, \infty_{\SN}\}$ (identifying $\infty_{\SN} = q_{38}$), accepting subset $F_{\SN} = Q_{\SN} \setminus \{\infty_{\SN}\}$ of cardinality $38$, and transition operation $\odot_{\SN}$ obtained by transposing Figure~5.11 of Barendregt et al.~\cite[p.~82]{BarendregtEndrullisKlopWaldmann2017} and adjoining $\infty_{\SN}$ as a two-sided absorbing sink ($\infty_{\SN} \odot_{\SN} y = x \odot_{\SN} \infty_{\SN} = \infty_{\SN}$ for all $x, y \in Q_{\SN}$). Direct evaluation over all $39^3 = 59{,}319$ triples in $Q_{\SN}^3$ confirms~\eqref{eq:s-algebra-law}.

\begin{theorem}[Regularity of strong normalization {\cite{Waldmann2000}; \cite[Theorem~9]{BarendregtEndrullisKlopWaldmann2017}}]\label{thm:waldmann-sn}
For every closed $S$-term $M \in \TS$,
\[
M \in \SN \;\Longleftrightarrow\; M \in \WN \;\Longleftrightarrow\; \sem{M}_{\SN} \neq \infty_{\SN},
\]
and membership $M \in \SN$ is decidable in $2|M| - 1$ deterministic table lookups.
\end{theorem}

\begin{proof}
Equivalence of $M \in \SN$ and $M \in \WN$ holds by Lemma~\ref{lem:wn-eq-sn}(ii). For soundness (that $M \in \WN$ implies $\sem{M}_{\SN} \neq \infty_{\SN}$), inspection of row $q_0$ and of all $38^2 = 1{,}444$ pairs $(x, y) \in F_{\SN}^2$ in $\mathcal{A}_{\SN}$ shows that $q_0 \neq \infty_{\SN}$, that $x \neq \infty_{\SN}$ implies $q_0 \odot_{\SN} x \neq \infty_{\SN}$, and that $x, y \neq \infty_{\SN}$ implies $(q_0 \odot_{\SN} x) \odot_{\SN} y \neq \infty_{\SN}$. Structural induction on the grammar~\eqref{eq:cnf-grammar} therefore yields $\sem{V}_{\SN} \neq \infty_{\SN}$ for every $V \in \mathcal{N}$. If $M \redSs V \in \mathcal{N}$, conversion invariance of $\mathcal{A}_{\SN}$ gives $\sem{M}_{\SN} = \sem{V}_{\SN} \neq \infty_{\SN}$.

For constructive completeness (that $\sem{M}_{\SN} \neq \infty_{\SN}$ implies $M \redSs V$ for some $V \in \mathcal{N}$), we avoid K\"onig's lemma by constructing a total recursive normalization operator $\mathrm{app}_{\mathcal{N}}(F, Z) \in \mathcal{N}$ satisfying $F Z \redSs \mathrm{app}_{\mathcal{N}}(F, Z)$ for all canonical normal forms $F, Z \in \mathcal{N}$ with $\sem{F}_{\SN} \odot_{\SN} \sem{Z}_{\SN} \neq \infty_{\SN}$. Once $\mathrm{app}_{\mathcal{N}}$ is established, structural induction on $M$ normalizes every accepted term: $\mathrm{norm}(S) = S$, and for $M \equiv L R$ with $\sem{L R}_{\SN} \neq \infty_{\SN}$, the two-sided sink property of $\infty_{\SN}$ forces $\sem{L}_{\SN} \neq \infty_{\SN}$ and $\sem{R}_{\SN} \neq \infty_{\SN}$, whence $L R \redSs \mathrm{norm}(L)\,\mathrm{norm}(R) \redSs \mathrm{app}_{\mathcal{N}}(\mathrm{norm}(L), \mathrm{norm}(R)) \in \mathcal{N}$.

Rather than storing an unstructured $59{,}319$-entry three-dimensional array, the rank map $\tau_{\SN} \colon Q_{\SN}^3 \to \{0, 1, \dots, 6\}$ collapses to a two-dimensional sub-rank matrix $\hat{\tau}_{\SN} \colon Q_{\SN}^2 \to \{0, 1, \dots, 6\}$ modified on $35$ explicit overrides ($34$ entries at $x = q_1$ assigned rank~$2$, plus the single triple $(q_9, q_0, q_1)$ assigned rank~$3$), setting $\tau_{\SN}(x, y, z) = \hat{\tau}_{\SN}((q_0 \odot_{\SN} x) \odot_{\SN} y,\, z)$ everywhere else. Let $I_2 \subseteq F_{\SN}$ denote the set of states inhabited by spine-$2$ normal forms $S A_1 A_2$. Finite verification across $F_{\SN}^3$ confirms $\tau_{\SN}(x, y, z) \le \hat{\tau}_{\SN}((q_0 \odot_{\SN} x) \odot_{\SN} y,\, z)$ and establishes two strict decrease laws whenever $((q_0 \odot_{\SN} x) \odot_{\SN} y) \odot_{\SN} z \neq \infty_{\SN}$, alongside $\hat{\tau}_{\SN}(x, z) \le \tau_{\SN}(x, y, z)$ and $\hat{\tau}_{\SN}(y, z) \le \tau_{\SN}(x, y, z)$:
\begin{align}
\tau_{\SN}(u, z, y \odot_{\SN} z) &\;<\; \tau_{\SN}(q_0 \odot_{\SN} u, y, z)
&&\textrm{whenever } x = q_0 \odot_{\SN} u, \label{eq:sn-rank-1}\\
\hat{\tau}_{\SN}(x \odot_{\SN} z, y \odot_{\SN} z) &\;<\; \tau_{\SN}(x, y, z)
&&\textrm{whenever } x \in I_2. \label{eq:sn-rank-2}
\end{align}
Assign to each call $(F, Z)$ its rank $r(F, Z) = 0$ if $F \equiv S$ or $F \equiv S X$, and $r(S X Y, Z) = \tau_{\SN}(\sem{X}_{\SN}, \sem{Y}_{\SN}, \sem{Z}_{\SN}) \ge 1$ otherwise. We proceed by well-founded lexicographic induction on $(r(F, Z), |F|) \in \mathbb{N}^2$. If $F \equiv S$ or $F \equiv S X$, then $F Z \in \mathcal{N}$ already. If $F \equiv S X Y$ with $X, Y \in \mathcal{N}$, then $(S X Y) Z \redS (X Z)(Y Z)$, and conversion invariance plus the two-sided sink property ensure $\sem{X Z}_{\SN} \neq \infty_{\SN}$, $\sem{Y Z}_{\SN} \neq \infty_{\SN}$, and $\sem{(X Z)(Y Z)}_{\SN} \neq \infty_{\SN}$. Since $|Y| < |S X Y|$ and $r(Y, Z) \le \hat{\tau}_{\SN}(\sem{Y}_{\SN}, \sem{Z}_{\SN}) \le r(S X Y, Z)$, the induction hypothesis normalizes $Y Z \redSs W \in \mathcal{N}$ with $\sem{W}_{\SN} = \sem{Y}_{\SN} \odot_{\SN} \sem{Z}_{\SN}$. Similarly, $|X| < |S X Y|$ and $r(X, Z) \le r(S X Y, Z)$ normalize $X Z \redSs V \in \mathcal{N}$ with $\sem{V}_{\SN} = \sem{X}_{\SN} \odot_{\SN} \sem{Z}_{\SN}$. We now branch on $X \in \mathcal{N}$:
\begin{enumerate}[label=(\roman*),leftmargin=2.2em]
  \item If $X \equiv S$, then $V \equiv S Z$ and $(X Z)(Y Z) \redSs S Z W \in \mathcal{N}$, completing normalization.
  \item If $X \equiv S U$ with $U \in \mathcal{N}$, then $V \equiv S U Z$, and by~\eqref{eq:sn-rank-1} the outer call $(S U Z, W)$ has strictly smaller first component $\tau_{\SN}(\sem{U}_{\SN}, \sem{Z}_{\SN}, \sem{W}_{\SN}) < r(S(S U)Y, Z)$.
  \item If $X \equiv S A_1 A_2$ with $A_1, A_2 \in \mathcal{N}$, then $\sem{X}_{\SN} \in I_2$, and by~\eqref{eq:sn-rank-2} we have $r(V, W) \le \hat{\tau}_{\SN}(\sem{V}_{\SN}, \sem{W}_{\SN}) < r(S X Y, Z)$, so $\mathrm{app}_{\mathcal{N}}(V, W) \in \mathcal{N}$ terminates. \qedhere
\end{enumerate}
\end{proof}

\begin{remark}[Automaton minimization and table orientation]\label{rem:sn-history}
Two bibliographical details merit record when comparing implementations of $\mathcal{A}_{\SN}$ across the literature. First, Waldmann's original top-down grammar construction~\cite{Waldmann1998Thesis,Waldmann2000} yielded $43$ non-terminals; Endrullis, de~Vrijer, and Waldmann~\cite[Section~5, p.~15]{EndrullisDeVrijerWaldmann2010} subsequently applied local termination and Myhill--Nerode minimization to obtain the minimal $39$-state totalized deterministic bottom-up tree automaton ($38$ accepting states and $1$ undefined sink). Second, in Figure~5.11 of Barendregt, Endrullis, Klop, and Waldmann~\cite[p.~82]{BarendregtEndrullisKlopWaldmann2017}, the printed $38 \times 38$ array indexes the left-applicand state $x = \sem{M}_{\SN}$ horizontally across columns and the right-applicand state $y = \sem{N}_{\SN}$ vertically down rows, transposed relative to the figure caption; reading rows as left applicands and columns as right applicands violates~\eqref{eq:s-algebra-law} already on $T T T \equiv (S S)(S S)(S S)$.
\end{remark}

\subsection{Automaton-Free Head Conservation, Optimal Multiplier Potential, and Question 10}\label{subsec:hn-q10}

Recall from Definition~\ref{def:hnf-bohm}(i)--(ii) that a closed $S$-term $M \in \TS$ lies in $\HNF$ when $\ls(M) \le 2$, is solvable ($M \in \HN$) when $M \redSs H$ for some $H \in \HNF$, and lies in $\RNF$ when no reduction from $M$ ever contracts at the root. In classical treatments~\cite{Barendregt1984}, equivalence between multiway solvability ($M \in \HN$) and termination of deterministic head reduction $\to_h$ invokes full standardization. In $\CLS$, however, a direct weighted-parallel pullback establishes head conservation constructively without standardization or finite automata. Write $\mathrm{SHN} \subseteq \TSV$ for the set of terms $M$ that are \emph{strongly head-normalizing}, meaning that deterministic head reduction $\to_h$ from $M$ is well-founded.

\begin{lemma}[Automaton-free head conservation]\label{lem:head-conservation}
For all open or closed $S$-terms $M, N \in \TSV$, if $M \redS N$ and $N \in \mathrm{SHN}$, then $M \in \mathrm{SHN}$. Consequently, for every closed $S$-term $M \in \TS$, one has $M \in \HN$ if and only if $M \in \mathrm{SHN}$.
\end{lemma}

\begin{proof}
Recall from the proof of Proposition~\ref{prop:cr-conservation}(iii) the weighted parallel reduction relation $M \Rightarrow_w N$ ($w \in \mathbb{N}$), where atoms have weight $0$, application satisfies $A_1 A_2 \Rightarrow_{w_1 + w_2} B_1 B_2$, redex contraction satisfies $S X Y Z \Rightarrow_{1 + w_X + w_Y + 2 w_Z} (X' Z')(Y' Z')$, and $M \redS N$ implies $M \Rightarrow_1 N$. Structural induction on $M \Rightarrow_w N$ yields a sharp head-projection alternative for any head step $M \to_h M_1$:
\begin{enumerate}[label=(\roman*),leftmargin=2.2em]
  \item if $M \equiv S X Y Z$ is contracted at the root in $M \Rightarrow_w N$, then $M_1 \equiv (X Z)(Y Z)$ satisfies $M_1 \Rightarrow_{w_1} N$ with $w_1 = (w_X + w_Z) + (w_Y + w_Z) < 1 + w_X + w_Y + 2 w_Z = w$ (no left-child head step can occur inside $S X Y$, whose left-spine arity is $2$);
  \item if $M \equiv A_1 A_2 \Rightarrow_{w_1 + w_2} B_1 B_2$ and $M \to_h M_1$ contracts at the root ($A_1 \equiv S X Y$, forcing $B_1 \equiv S X' Y'$), then $N \equiv S X' Y' B_2 \to_h (X' B_2)(Y' B_2) \equiv N_1$ with $M_1 \Rightarrow_{w'} N_1$; and if $A_1 \to_h A_1'$, the induction hypothesis on $A_1 \Rightarrow_{w_1} B_1$ either decreases $w_1$ while keeping $B_1 B_2$ fixed or steps $B_1 \to_h B_1'$ and hence $B_1 B_2 \to_h B_1' B_2$.
\end{enumerate}
Well-founded outer induction on $N \in \mathrm{SHN}$ combined with inner induction on $w \in \mathbb{N}$ therefore shows that $N \in \mathrm{SHN}$ and $M \Rightarrow_w N$ imply $M \in \mathrm{SHN}$. Since every $H \in \TSV$ with $\ls(H) \le 2$ admits no head step and thus lies in $\mathrm{SHN}$, induction along $M \redSs H$ proves that $M \in \HN$ implies $M \in \mathrm{SHN}$. Conversely, on closed $M \in \TS$, iterating the unique partial head-step function along $\mathrm{SHN}$ terminates at a closed head-irreducible term $H$, which necessarily satisfies $\ls(H) \le 2$ and witnesses $M \in \HN$.
\end{proof}

Before introducing the $35$-state head-normalization algebra, we record three syntactic rigidity properties of multi-argument applications in $\CLS$ refining Visser's variable-tail method and Proposition~\ref{prop:acyclicity-order}(iii).

\begin{lemma}[Spine-$2$ persistence, Fibonacci $2$-HNFs, and third-argument head exclusion]\label{lem:visser-fibonacci}
Let $U \in \TS$ be a closed $S$-term and let $x, y, z \in \mathcal{V}$ be distinct variables.
\begin{enumerate}[label=(\roman*),leftmargin=2.2em]
  \item \emph{Persistence of spine arity at least $2$:} For all open or closed terms $A_1, A_2, A_3, W \in \TSV$, neither $(A_1 A_2) A_3 \redSs S$ nor $(A_1 A_2) A_3 \redSs S W$ holds. Consequently, $F X \redSs S W$ implies $F \equiv S$ and $X \redSs W$.
  \item \emph{Fibonacci classification of parametric $2$-HNFs:} Recall from Definition~\ref{def:syntax} the closed right combs $R_1 \equiv S$ and $R_{n+1} \equiv S R_n$ ($n \ge 1$), and define the Fibonacci word trees $W_0(x, y) \equiv x$, $W_1(x, y) \equiv y$, and $W_{m+2}(x, y) \equiv W_m(x, y)\,W_{m+1}(x, y)$. Then $U x y \redSs S P Q$ holds if and only if there exists $k \in \mathbb{N}$ such that $U \redSs R_{k+1}$, $P \equiv W_k(x, y)$, and $Q \equiv W_{k+1}(x, y)$. Moreover, if $U x y \redSs S P Q$ with $U \not\equiv S$, then both right spines of $P$ and $Q$ terminate in $y$.
  \item \emph{Third-argument head exclusion:} Whenever $U x y z \redSs H$ with $H$ in open head normal form (either $\ls(H) \le 2$ with head $S$, or head in $\mathcal{V}$), the leftmost leaf of $H$ belongs to $\{x, y\}$ and is never $z$ or $S$.
\end{enumerate}
\end{lemma}

\begin{proof}
For~(i), define the predicate $\Phi_2(K)$ asserting $\ls(K) \ge 2$, meaning $K \equiv (K_1 K_2) K_3$. A root contraction $S X Y Z \redS (X Z)(Y Z)$ produces a term of left-spine arity $\ls(X Z) + 1 \ge 2$; a left step $(K_1 K_2) K_3 \redS L' K_3$ satisfies $\ls(L') \ge 1$ by inspection of root, left, and right subcases on $K_1 K_2$; and a right step leaves $\ls$ unchanged. Induction along $\redSs$ shows that $\Phi_2$ is closed forward under reduction, excluding $S$ ($\ls = 0$) and $S W$ ($\ls = 1$). If $F X \redSs S W$, then $F$ cannot be a variable (head variable $v$ is preserved) or an application $F_1 F_2$ (by $\Phi_2$), forcing $F \equiv S$ and $X \redSs W$.

For~(ii), forward verification $R_{k+1} W_m W_{m+1} \redSs S W_{m+k} W_{m+k+1}$ holds by induction on $k \in \mathbb{N}$, since $(S R_{k+1}) W_m W_{m+1} \redS R_{k+1} W_{m+1} (W_m W_{m+1}) \equiv R_{k+1} W_{m+1} W_{m+2}$. Conversely, we prove by strong induction on step count $n$ that if $U W_m W_{m+1} \redS^n S P Q$ with $U$ closed, then $U \redSs R_{k+1}$, $P \equiv W_{m+k}$, and $Q \equiv W_{m+k+1}$ for some $k \ge 0$. Because every $W_j(x, y)$ contains no $S$-leaves, it is irreducible. If no root contraction occurs at the outer application $(U W_m) W_{m+1}$, then $U W_m \redS^n S P$ and $Q \equiv W_{m+1}$; by~(i), $U \equiv S \equiv R_1$ and $W_m \redSs P$, so $P \equiv W_m$ ($k = 0$). Otherwise, $(U W_m) \redS^{n_1} S A_1 A_2$ and $(A_1 W_{m+1})(A_2 W_{m+1}) \redS^{n_2} S P Q$ with $n_1 + n_2 + 1 = n$. Recall from the proof of Proposition~\ref{prop:acyclicity-order}(iii) that the classes $\mathcal{W}_1$, $\mathcal{W}_2 = \mathcal{W}_1 \cdot \mathcal{W}_1$, and $\mathcal{W}_3 = \mathcal{W}_2 \cdot \mathcal{W}_1$ of terms ending in variables are closed forward under $\redS$. If $U W_m \redS^{n_1} S A_1 A_2$ used a root contraction $(S X_1 X_2) W_m \redS (X_1 W_m)(X_2 W_m) \redSs S A_1 A_2$, then $(X_1 W_m)(X_2 W_m) \in \mathcal{W}_2$ would force $S A_1 \in \mathcal{W}_1$ and $A_1 \in \mathcal{W}_1$, placing $(A_1 W_{m+1})(A_2 W_{m+1}) \in \mathcal{W}_3$, which never reduces to $S P Q$ since $S \notin \mathcal{W}_1$. Thus $U W_m \redS^{n_1} S A_1 A_2$ performs no outer root contraction, giving $U \redS^{n_1} S A_1$, $A_2 \equiv W_m$, and $A_1 W_{m+1} W_{m+2} \redS^{n_2} S P Q$ with $n_2 < n$; the induction hypothesis on $A_1$ yields $A_1 \redSs R_{k'+1}$ and $U \redSs S R_{k'+1} \equiv R_{k'+2}$. When $k \ge 1$, both $W_k(x, y)$ and $W_{k+1}(x, y)$ end on their right spine in $y$.

For~(iii), suppose $(U x y) z \redSs H \in \HNF$. Since $z$ is irreducible, either $H \equiv F' z$ with $U x y \redSs F'$ (in which case $\ls(F') \ge 2$ by~(i), so $\ls(F' z) \ge 3$, forcing the head of $H$ to be a variable occurring in $U x y$, namely $x$ or $y$), or else $U x y \redSs S P Q$ and $(P z)(Q z) \redSs H$. By~(ii), $P \equiv W_k(x, y)$ and $Q \equiv W_{k+1}(x, y)$, so $(P z)(Q z)$ contains no $S$-leaves and is already in normal form with head leaf equal to that of $W_k(x, y)$, which belongs to $\{x, y\}$.
\end{proof}

In Section~5.3 of Barendregt, Endrullis, Klop, and Waldmann~\cite[Figure~5.14, p.~85; Question~10 and Footnote~4, p.~86]{BarendregtEndrullisKlopWaldmann2017}, a candidate $35$-state transition table was synthesized via local termination heuristics to recognize $\HN$, and its proof of correctness---together with linear-time decidability of $\RNF$---was posed as Question~10 and Footnote~4. Let $\mathcal{A}_{\HN} = (Q_{\HN}, q_0, \odot_{\HN})$ denote this $35$-state structure with carrier $Q_{\HN} = \{q_0, q_1, \dots, q_{33}, \infty_{\HN}\}$ ($\infty_{\HN} = q_{34}$) and accepting set $F_{\HN} = Q_{\HN} \setminus \{\infty_{\HN}\}$. Direct evaluation over all $35^3 = 42{,}875$ triples verifies~\eqref{eq:s-algebra-law}. Unlike $\infty_{\SN} \in Q_{\SN}$, the state $\infty_{\HN}$ exhibits a directional asymmetry:
\begin{equation}\label{eq:hn-asymmetry}
\infty_{\HN} \odot_{\HN} y \;=\; \infty_{\HN} \quad\textrm{for all } y \in Q_{\HN},
\qquad\textrm{whereas}\qquad
q_0 \odot_{\HN} \infty_{\HN} \;=\; q_{10} \;\neq\; \infty_{\HN},
\end{equation}
and $(q_0 \odot_{\HN} x) \odot_{\HN} \infty_{\HN} \neq \infty_{\HN}$ for every $x \in Q_{\HN}$. Algebraically, $\infty_{\HN}$ is a left ideal element because applying an unsolvable head $M$ to $N$ yields an unsolvable term $M N$, whereas $S M$ and $S X M$ already lie in $\HNF$ regardless of whether $M$ has a head normal form.

To prove completeness (that $\sem{M}_{\HN} \neq \infty_{\HN}$ implies $M \in \HN$) constructively without K\"onig's lemma, we assign unit weight $\mu_{\HN}(S) = 1$ to the head combinator leaf itself. Because each head contraction $S X Y Z \to_h (X Z)(Y Z)$ destroys that leading $S$-leaf, no additive offset table $B$ is needed: a single $35 \times 35$ multiplier matrix $A$ suffices, and we show below that $A$ is pointwise minimal among all valid multiplier matrices.

\begin{lemma}[Pure multiplier head-step potential]\label{lem:hn-potential}
There exists an explicit non-negative integer matrix $A \colon Q_{\HN}^2 \to \{0, 1, \dots, 30\}$ (with $881$ zero entries, $202$ entries equal to $1$, and $142$ entries strictly greater than $1$, attaining its maximum $A(q_{10}, q_0) = 30$ and satisfying $A(w, \infty_{\HN}) = 0$ for all $w \in Q_{\HN}$) such that across all $8{,}329$ active head-redex triples $(x, y, z) \in Q_{\HN}^3$ satisfying $((q_0 \odot_{\HN} x) \odot_{\HN} y) \odot_{\HN} z \neq \infty_{\HN}$, writing $s_1 = q_0 \odot_{\HN} x$, $s_2 = s_1 \odot_{\HN} y$, $u = x \odot_{\HN} z$, and $v = y \odot_{\HN} z$, the following four conditions hold simultaneously:
\begin{align}
u &\;\neq\; \infty_{\HN}, \label{eq:pot-1}\\
1 &\;\le\; A(q_0, x), \label{eq:pot-2}\\
A(u, v) &\;\le\; A(s_1, y), \label{eq:pot-3}\\
A(x, z) + A(u, v)\,A(y, z) &\;\le\; A(s_2, z). \label{eq:pot-4}
\end{align}
Define the pure multiplier potential $\mu_{\HN} \colon \TSV \to \mathbb{N}$ recursively by $\mu_{\HN}(S) = \mu_{\HN}(v) = 1$ ($v \in \mathcal{V}$) and
\begin{equation}\label{eq:mu-hn-def}
\mu_{\HN}(L R) \;=\; \mu_{\HN}(L) + A(\sem{L}_{\HN}, \sem{R}_{\HN})\,\mu_{\HN}(R).
\end{equation}
Then for every closed $S$-term $M \in \TS$ with $\sem{M}_{\HN} \neq \infty_{\HN}$ and every head contraction $M \to_h M'$, we have $\mu_{\HN}(M') + 1 \le \mu_{\HN}(M)$.
\end{lemma}

\begin{proof}
Verification of~\eqref{eq:pot-1}--\eqref{eq:pot-4} over $\{ (x, y, z) \in Q_{\HN}^3 : ((q_0 \odot_{\HN} x) \odot_{\HN} y) \odot_{\HN} z \neq \infty_{\HN} \}$ ($8{,}329$ triples out of $35^3 = 42{,}875$, where~\eqref{eq:pot-1} follows automatically from left-absorption of $\infty_{\HN}$ and~\eqref{eq:s-algebra-law}) is a finite integer arithmetic check. We prove $\mu_{\HN}(M') + 1 \le \mu_{\HN}(M)$ by induction on the derivation of $M \to_h M'$.

In the base case, $M \equiv ((S X) Y) Z$ contracts at the root to $M' \equiv (X Z)(Y Z)$, with $x = \sem{X}_{\HN}$, $y = \sem{Y}_{\HN}$, $z = \sem{Z}_{\HN}$, and $((q_0 \odot_{\HN} x) \odot_{\HN} y) \odot_{\HN} z = \sem{M}_{\HN} \neq \infty_{\HN}$. Expanding $\mu_{\HN}(M)$ via~\eqref{eq:mu-hn-def} and $\mu_{\HN}(S) = 1$ yields
\[
\mu_{\HN}(((S X) Y) Z)
\;=\; 1 + A(q_0, x)\,\mu_{\HN}(X) + A(s_1, y)\,\mu_{\HN}(Y) + A(s_2, z)\,\mu_{\HN}(Z),
\]
whereas expanding $\mu_{\HN}((X Z)(Y Z))$ with $\sem{X Z}_{\HN} = u$ and $\sem{Y Z}_{\HN} = v$ yields
\[
\mu_{\HN}((X Z)(Y Z))
\;=\; \mu_{\HN}(X) + A(u, v)\,\mu_{\HN}(Y) + \bigl(A(x, z) + A(u, v)\,A(y, z)\bigr)\mu_{\HN}(Z).
\]
Since $\mu_{\HN}(X), \mu_{\HN}(Y), \mu_{\HN}(Z) \ge 0$, comparing the coefficients of $\mu_{\HN}(X)$, $\mu_{\HN}(Y)$, and $\mu_{\HN}(Z)$ via~\eqref{eq:pot-2}--\eqref{eq:pot-4} while retaining the leading constant~$1$ from $\mu_{\HN}(S) = 1$ gives $\mu_{\HN}((X Z)(Y Z)) + 1 \le \mu_{\HN}(((S X) Y) Z)$.

In the inductive step, $M \equiv L R \to_h L' R$ where $L \to_h L'$. Because $\sem{L}_{\HN} \odot_{\HN} \sem{R}_{\HN} = \sem{L R}_{\HN} \neq \infty_{\HN}$ and $\infty_{\HN}$ is a left sink by~\eqref{eq:hn-asymmetry}, we have $\sem{L}_{\HN} \neq \infty_{\HN}$. The induction hypothesis gives $\mu_{\HN}(L') + 1 \le \mu_{\HN}(L)$, and conversion invariance gives $\sem{L'}_{\HN} = \sem{L}_{\HN}$, so $A(\sem{L'}_{\HN}, \sem{R}_{\HN}) = A(\sem{L}_{\HN}, \sem{R}_{\HN})$ is preserved, whence $\mu_{\HN}(L R) - \mu_{\HN}(L' R) = \mu_{\HN}(L) - \mu_{\HN}(L') \ge 1$.
\end{proof}

\begin{proposition}[Pointwise minimality of $A$ and impossibility of rank-$1$ multipliers]\label{prop:hna-minimality}
Call a matrix $A' \colon Q_{\HN}^2 \to \mathbb{N}$ \emph{head-valid} if it satisfies~\eqref{eq:pot-2}--\eqref{eq:pot-4} across all $8{,}329$ active triples $(x, y, z) \in Q_{\HN}^3$ with $((q_0 \odot_{\HN} x) \odot_{\HN} y) \odot_{\HN} z \neq \infty_{\HN}$.
\begin{enumerate}[label=(\roman*),leftmargin=2.2em]
  \item \emph{Global pointwise minimality:} The matrix $A$ of Lemma~\ref{lem:hn-potential} satisfies $A(u, v) \le A'(u, v)$ for every head-valid matrix $A'$ and all $u, v \in Q_{\HN}$; in particular, every head-valid matrix satisfies $\max A' \ge A'(q_{10}, q_0) \ge 30$.
  \item \emph{Impossibility of masked rank-$1$ and $1$D potentials:} There exist no Boolean mask $M \colon Q_{\HN}^2 \to \{0, 1\}$ and one-dimensional weight vectors $w_L, w_R \colon Q_{\HN} \to \mathbb{N}$ such that $A'(u, v) = M(u, v)\,w_L(u)\,w_R(v)$ is head-valid. In particular, no left-only weight $A'(u, v) = w_L(u)$ or right-only weight $A'(u, v) = w_R(v)$ satisfies~\eqref{eq:pot-2}--\eqref{eq:pot-4}.
\end{enumerate}
\end{proposition}

\begin{proof}
For~(i), assign each pair $(u, v) \in Q_{\HN}^2$ a derivation level $\ell(u, v) \in \{0, 1, \dots, 13\}$ ($\ell(u, v) = 0$ precisely where $A(u, v) = 0$). A finite certificate table associates to every pair $(u, v)$ with $\ell(u, v) \ge 1$ an active triple $(x, y, z)$ whose lower-level premises strictly precede $\ell(u, v)$ and force $A(u, v) \le A'(u, v)$ via one of~\eqref{eq:pot-2} (when $u = q_0, v = x$), \eqref{eq:pot-3} (when $u = q_0 \odot_{\HN} x, v = y$), or~\eqref{eq:pot-4} (when $u = (q_0 \odot_{\HN} x) \odot_{\HN} y, v = z$). Induction over $\ell = 0, 1, \dots, 13$ yields $A(u, v) \le A'(u, v)$ for all $u, v \in Q_{\HN}$, and evaluating at $(q_{10}, q_0)$ gives $30 = A(q_{10}, q_0) \le A'(q_{10}, q_0)$.

For~(ii), suppose toward a contradiction that $A'(u, v) = M(u, v)\,w_L(u)\,w_R(v)$ is head-valid. Applying~(i) at four entries of $A$ gives
\[
A'(q_5, q_0) \ge 15, \qquad A'(q_{10}, q_1) \ge 15, \qquad A'(q_0, q_0) \ge 1, \qquad A'(q_{10}, q_0) \ge 30.
\]
Hence $M$ equals $1$ at all four pairs, and $w_L(q_5), w_R(q_0), w_L(q_{10}), w_R(q_1), w_L(q_0) \ge 1$. Now instantiate~\eqref{eq:pot-4} at the triple $(x, y, z) = (q_5, q_0, q_0)$: in $\mathcal{A}_{\HN}$ one has $q_0 \odot_{\HN} q_5 = q_5$, $q_5 \odot_{\HN} q_0 = q_{10}$, $q_0 \odot_{\HN} q_0 = q_1$, and $((q_0 \odot_{\HN} q_5) \odot_{\HN} q_0) \odot_{\HN} q_0 = q_{10} \odot_{\HN} q_0 = q_{23} \neq \infty_{\HN}$, so~\eqref{eq:pot-4} requires
\[
w_L(q_5)\,w_R(q_0) + \bigl(w_L(q_{10})\,w_R(q_1)\bigr)\bigl(w_L(q_0)\,w_R(q_0)\bigr) \;\le\; w_L(q_{10})\,w_R(q_0).
\]
Regrouping the second summand as $\bigl(w_R(q_1)\,w_L(q_0)\bigr)\bigl(w_L(q_{10})\,w_R(q_0)\bigr) \ge 1 \cdot w_L(q_{10})\,w_R(q_0)$ and noting $w_L(q_5)\,w_R(q_0) \ge 1$ makes the left-hand side strictly exceed $w_L(q_{10})\,w_R(q_0)$, a contradiction. Setting $M \equiv 1$ and $w_R \equiv 1$ (or $w_L \equiv 1$) rules out one-dimensional weights.
\end{proof}

\begin{theorem}[Resolution of Question~10 and Footnote~4 of Barendregt et al.\ {\cite[p.~86]{BarendregtEndrullisKlopWaldmann2017}}]\label{thm:question-10}
For every closed $S$-term $M \in \TS$, the $35$-state $S$-algebra $\mathcal{A}_{\HN}$ satisfies the following properties:
\begin{enumerate}[label=(\roman*),leftmargin=2.2em]
  \item \emph{Regularity and constructive bound for $\HN$:} One has $M \in \HN$ if and only if $\sem{M}_{\HN} \neq \infty_{\HN}$, and $\sem{M}_{\HN} = \infty_{\HN}$ if and only if $M$ admits an infinite head-reduction sequence $M = M_0 \to_h M_1 \to_h \cdots$. Specifically, if $\sem{M}_{\HN} \neq \infty_{\HN}$, deterministic head reduction from $M$ reaches $\HNF$ in at most $\mu_{\HN}(M) - 1$ steps.
  \item \emph{Linear-time decidability of $\RNF$ (Footnote~4):} A closed $S$-term $M$ lies in $\RNF$ if and only if
  \begin{equation}\label{eq:footnote4-rnf}
  \ls(M) \le 2 \quad\textrm{or}\quad \bigl(M \equiv P Q \;\textrm{ and }\; \sem{P}_{\HN} = \infty_{\HN}\bigr),
  \end{equation}
  which is decidable in $O(|M|)$ deterministic operations.
  \item \emph{Five-composite-argument $\RNF$ ceiling:} In the spine transition graph of $\mathcal{A}_{\HN}$, the transition $q_{13} \xrightarrow{q_0} q_{13}$ (entered via $q_0 \xrightarrow{q_0} q_1 \xrightarrow{y} q_{13}$) is the unique accepting self-loop, and consuming composite arguments ($y \neq q_0$) contracts the reachable state sets strictly along cardinalities $35 \to 26 \to 14 \to 4 \to 2 \to 1 = \{\infty_{\HN}\}$. Consequently, for every closed $M \in \TS$ and all closed $U_1, \dots, U_5 \in \TS \setminus \{S\}$, we have $\sem{M U_1 U_2 U_3 U_4 U_5}_{\HN} = \infty_{\HN}$, and $M U_1 U_2 U_3 U_4 U_5 U_6 \in \RNF$ for every $U_6 \in \TS$.
  \item \emph{Joint $(\SN, \HN)$ subalgebra:} The reachable subalgebra of $\mathcal{A}_{\SN} \times \mathcal{A}_{\HN}$ generated by $(q_0, q_0)$ has exact cardinality $112$, partitioning into $94$ states in $F_{\SN} \times F_{\HN}$ ($\SN$), $17$ states in $\{\infty_{\SN}\} \times F_{\HN}$ ($\HN \setminus \SN$), and the single state $(\infty_{\SN}, \infty_{\HN})$ ($\overline{\HN}$).
\end{enumerate}
\end{theorem}

\begin{proof}
For~(i), every closed $H \in \HNF$ has the form $S$, $S A_1$, or $S A_1 A_2$, and inspection of $\mathcal{A}_{\HN}$ confirms $q_0 \neq \infty_{\HN}$, $q_0 \odot_{\HN} x \neq \infty_{\HN}$, and $(q_0 \odot_{\HN} x) \odot_{\HN} y \neq \infty_{\HN}$ for all $x, y \in Q_{\HN}$; conversion invariance therefore shows that $M \in \HN$ implies $\sem{M}_{\HN} \neq \infty_{\HN}$. Conversely, every closed $M \notin \HNF$ satisfies $\ls(M) \ge 3$ and admits a unique head step $M \to_h M'$. If $\sem{M}_{\HN} \neq \infty_{\HN}$, conversion invariance preserves $\sem{M'}_{\HN} \neq \infty_{\HN}$ and Lemma~\ref{lem:hn-potential} decreases $\mu_{\HN} \ge 1$ by at least~$1$ at each step, terminating in $\HNF$ within $\mu_{\HN}(M) - 1$ steps; if instead $\sem{M}_{\HN} = \infty_{\HN}$, then every head reduct $M_k$ preserves $\sem{M_k}_{\HN} = \infty_{\HN}$ and therefore satisfies $\ls(M_k) \ge 3$, generating an infinite deterministic head-reduction chain.

For~(ii), if $\ls(M) \le 2$, then $M \in \HNF \subseteq \RNF$. Suppose now $\ls(M) \ge 3$, so $M \equiv P Q$ with $\ls(P) \ge 2$. A reduction $P Q \redSs N$ contracts at the root if and only if $P \redSs S X Y$ for some $X, Y \in \TS$. By Lemma~\ref{lem:visser-fibonacci}(i), since $\ls(P) \ge 2$, no reduct of $P$ can be $S$ or $S W$; hence $P$ reduces to some $S X Y$ if and only if $P \in \HN$, which by~(i) holds if and only if $\sem{P}_{\HN} \neq \infty_{\HN}$. Negating this equivalence yields $P Q \in \RNF$ if and only if $\sem{P}_{\HN} = \infty_{\HN}$, establishing~\eqref{eq:footnote4-rnf}.

For~(iii), inspection of $\mathcal{A}_{\HN}$ shows that $q_0$ never occurs in the image of $\odot_{\HN}$; hence a closed $S$-term $U$ evaluates to $q_0$ in $\mathcal{A}_{\HN}$ if and only if $U \equiv S$. Computing $S_{k+1} = \{ u \odot_{\HN} y : u \in S_k,\, y \in Q_{\HN} \setminus \{q_0\} \}$ starting from $S_0 = Q_{\HN}$ yields $|S_0| = 35$, $|S_1| = 26$, $|S_2| = 14$, $|S_3| = 4$ ($\{q_{23}, q_{24}, q_{25}, \infty_{\HN}\}$), $|S_4| = 2$ ($\{q_{23}, \infty_{\HN}\}$), and $S_5 = \{\infty_{\HN}\}$. Thus $\sem{M U_1 \dots U_5}_{\HN} = \infty_{\HN}$, and~\eqref{eq:footnote4-rnf} implies $M U_1 \dots U_5 U_6 \in \RNF$; five composite arguments is sharp, since $S S U_1 U_2 U_3 U_4$ has $\HN$ state $q_{23} \neq \infty_{\HN}$ for $U_i \equiv S S$.

For~(iv), breadth-first closure of $\{(q_0, q_0)\} \subseteq Q_{\SN} \times Q_{\HN}$ under componentwise $\odot$ stabilizes at $112$ pairs, none of which lies in $F_{\SN} \times \{\infty_{\HN}\}$, reflecting $\SN \subsetneq \HN$ algebraically.
\end{proof}

\subsection{Global Top-Termination, Root Normal Forms, and Absence of Mute Terms}\label{subsec:top-termination}

In untyped $\lambda$-calculus and full $\mathrm{CL}(S, K)$, unsolvable terms divide into those reducing to an application of unsolvable terms and \emph{mute} (or root-active) terms in the sense of Berarducci~\cite{BerarducciBohm1993,Berarducci1996}, such as $\Omega \equiv (\lambda x.\,x x)(\lambda x.\,x x) \to_r \Omega$. Waldmann~\cite{Waldmann2000} (see also~\cite[Theorem~29]{BarendregtEndrullisKlopWaldmann2017}) established via semantic labelling~\cite{Zantema1995} that $\CLS$ is globally \emph{top-terminating}: no infinite reduction sequence contracts at the root infinitely often. Using $\mathcal{A}_{\HN}$ and Lemma~\ref{lem:hn-potential}, we give a direct constructive proof yielding a total root-normalizing function $\mathrm{rnf}$.

\begin{theorem}[Constructive global top-termination and root normal forms {\cite{Waldmann2000}; \cite[Theorem~29]{BarendregtEndrullisKlopWaldmann2017}}]\label{thm:waldmann-top-term}
\begin{enumerate}[label=(\roman*),leftmargin=2.2em]
  \item \emph{Global top-termination:} Every open or closed $S$-term $M \in \TSV$ is top-terminating: there exists no infinite sequence $M = M_0 \to_{\mathrm{nr}}^* M'_0 \to_r M_1 \to_{\mathrm{nr}}^* M'_1 \to_r \cdots$ containing infinitely many root contractions $\to_r$ interspersed with non-root steps $\to_{\mathrm{nr}}$.
  \item \emph{Total root normalizer $\mathrm{rnf}$ and absence of Berarducci mute terms:} There is a total constructive procedure $\mathrm{rnf} \colon \TS \to \RNF$ satisfying $M \redSs \mathrm{rnf}(M)$ for every $M \in \TS$. Consequently, $\CLS$ contains no Berarducci mute terms~\cite{Berarducci1996}: every unsolvable closed $M \in \TS \setminus \HN$ reduces to an applicative root normal form $P Q \in \RNF$ with $P \in \TS \setminus \HN$, and every closed $S$-term possesses a $\bot$-free Berarducci tree $\BeT(M)$.
  \item \emph{Root-step bound on $A A A$:} By Proposition~\ref{prop:acyclicity-order}(i), $\redSp$ is irreflexive on $\TSV$; moreover, every multiway reduction starting from $A A A \equiv (S S S)(S S S)(S S S)$ executes at most $3$ root contractions, reaching its root normal form in $6$ deterministic head steps.
\end{enumerate}
\end{theorem}

\begin{proof}
For~(i), consider first closed $M \in \TS$ with $\sem{M}_{\HN} \neq \infty_{\HN}$. Define the truncated pure-multiplier potential $\nu_{\HN} \colon \TS \to \mathbb{N}$ by $\nu_{\HN}(U) = 0$ if $\sem{U}_{\HN} = \infty_{\HN}$, $\nu_{\HN}(S) = 1$, and $\nu_{\HN}(L R) = \nu_{\HN}(L) + A(\sem{L}_{\HN}, \sem{R}_{\HN})\,\nu_{\HN}(R)$ whenever $\sem{L R}_{\HN} \neq \infty_{\HN}$. Because $A(u, \infty_{\HN}) = 0$ for all $u \in Q_{\HN}$ and $\infty_{\HN}$ is a left sink, $\nu_{\HN}$ agrees with $\mu_{\HN}$ modulo zeroing out $\infty_{\HN}$-subterms, which carry multiplier $0$ whenever they appear on the right of an accepted parent. By Lemma~\ref{lem:hn-potential} and conversion invariance of $\sem{\cdot}_{\HN}$, every non-root step $N \to_{\mathrm{nr}} N'$ anywhere inside $N$ satisfies $\nu_{\HN}(N') \le \nu_{\HN}(N)$, while every root contraction $N \to_r N'$ on $\sem{N}_{\HN} \neq \infty_{\HN}$ satisfies $\nu_{\HN}(N') + 1 \le \nu_{\HN}(N)$. Thus $M$ executes at most $\nu_{\HN}(M) - 1$ root contractions.

Next consider closed $M \in \TS$ with $\sem{M}_{\HN} = \infty_{\HN}$. By Proposition~\ref{prop:cr-conservation}(ii), right-spine depth $\rs$ satisfies $\rs(N) \le \rs(N')$ for every non-root step $N \to_{\mathrm{nr}} N'$, whereas every root contraction $S X Y Z \to_r (X Z)(Y Z)$ yields $\rs((X Z)(Y Z)) = \rs(S X Y Z) + 1$. Stratifying $Q_{\HN}$ into nested root-contraction capacity classes $G_{\ge 0} \supset G_{\ge 1} \supset \dots \supset G_{\ge 5}$ of cardinalities $(35, 34, 26, 18, 17, 16)$ and tracking the maximal strongly connected spine prefixes on $\infty_{\HN}$ once $\rs(N) \ge 6$ shows that at most finitely many root steps can occur before the left child of $N \equiv P Q$ locks permanently into $\infty_{\HN}$, where Theorem~\ref{thm:question-10}(ii) forces $N \in \RNF$. For an arbitrary open term $M \in \TSV$, the variable-erasing homomorphism $\sigma_S \colon \TSV \to \TS$ sending every $v \in \mathcal{V}$ to $S$ maps root steps to root steps and non-root steps to non-root steps, lifting well-foundedness of $\to_{\mathrm{nr}}^* \cdot \to_r$ from $\TS$ to $\TSV$.

For~(ii), well-founded recursion on the relative root relation $\to_{\mathrm{nr}}^* \cdot \to_r$ defines $\mathrm{rnf}(M)$: check whether $M \in \RNF$ in $O(|M|)$ time via Theorem~\ref{thm:question-10}(ii); if so, return $M$; otherwise $M \equiv P Q$ has $\ls(P) \ge 2$ and $\sem{P}_{\HN} \neq \infty_{\HN}$, so head-normalizing $P \to_h^* S X Y$ within $\mu_{\HN}(P) - 1$ steps and contracting $(S X Y) Q \to_r (X Q)(Y Q)$ performs one relative root step, allowing recursive invocation $\mathrm{rnf}((X Q)(Y Q))$. If $M \in \TS \setminus \HN$, then $\mathrm{rnf}(M) \in \RNF \setminus \HN$ cannot satisfy $\ls(\mathrm{rnf}(M)) \le 2$, so by~\eqref{eq:footnote4-rnf} it has the form $P Q$ with $\sem{P}_{\HN} = \infty_{\HN}$. Coinductively unfolding $\mathrm{rnf}$ on $P$ and $Q$ generates a total $\bot$-free Berarducci tree $\BeT(M)$, proving the absence of mute terms.

For~(iii), for $A A A \equiv (S S S)(S S S)(S S S)$ we have $\rs(A A A) = 2$, and finite enumeration of the root-contracting reachable cone confirms that $A A A$ admits at most $3$ root contractions along any multiway path, reaching the root normal form $\bigl((S A A)\,(S A (S A A))\bigr)\bigl(S(S A A)\,(S A (S A A))\bigr)$ in $6$ deterministic head steps.
\end{proof}

\subsection{Regularity of B\"ohm-Tree Productivity, Rank-10 Completeness, and 113-State Minimality}\label{subsec:q5-productivity}

Recall from Definition~\ref{def:hnf-bohm}(iii)--(iv) and Lemma~\ref{lem:conv-invariance} that a closed $S$-term $M \in \TS$ is \emph{productive} ($M \in \Prod$) when $\BT(M)$ contains no occurrence of $\bot$ at any finite depth---formally, when $\neg\bot_d(M)$ holds for all $d \in \mathbb{N}$---and that $\bot_d$ is conversion-invariant and monotone in $d$ ($\bot_d(M)$ implies $\bot_{d+1}(M)$). Every strongly normalizing term is productive because its normal form is a finite $\bot$-free B\"ohm tree. As discovered by Zachos~\cite{Zachos1978} (see~\cite[p.~80]{BarendregtEndrullisKlopWaldmann2017}), $\Prod$ strictly contains $\SN$: defining the normal-form tower
\begin{equation}\label{eq:zachos-tower}
G_0 \;\equiv\; S T T \;\equiv\; S(S S)(S S),
\qquad
G_{k+1} \;\equiv\; T G_k \;\equiv\; (S S)\,G_k
\qquad (k \in \mathbb{N}),
\end{equation}
we have $G_0 G_n \equiv S T T G_n \to_h (T G_n)(T G_n) \equiv G_{n+1} G_{n+1}$ and $G_{k+1} G_m \equiv S S G_k G_m \to_h S G_m (G_k G_m) \in \HNF$ (whose left child $G_m$ lies in $\mathcal{N}$), so $Z_{\mathrm{Zac}} \equiv G_0 G_0$ ($10$ leaves) is non-normalizing yet generates an infinite $\bot$-free B\"ohm tree. Question~5 of Barendregt et al.~\cite[p.~80]{BarendregtEndrullisKlopWaldmann2017} asks how to characterize $\Prod$, whether $\Prod$ is regular, and whether productivity is decidable.

Our resolution of Question~5 begins with two structural lemmas showing via potential-free head conservation how non-erasure in $\CLS$ forces $\bot$ to propagate across both sides of binary application, and why Zachos pairs $G_m G_n$ lose productivity upon receiving any further right argument.

\begin{lemma}[Left and right $\bot$-propagation on closed $S$-terms]\label{lem:bot-propagation}
For every depth $d \in \mathbb{N}$ and all closed $S$-terms $Y, Z, F, W \in \TS$:
\begin{equation}\label{eq:both-sink-laws}
\bot_d(Y) \;\Longrightarrow\; \bot_d(Y Z)
\qquad\textrm{and}\qquad
\bot_d(W) \;\Longrightarrow\; \bot_{d+1}(F W).
\end{equation}
Consequently, $M N \in \Prod$ implies both $M \in \Prod$ and $N \in \Prod$.
\end{lemma}

\begin{proof}
Consider first an invariant of head reduction on open $\{S, x_0\}$-terms: writing the left spine of $U$ as $U \equiv h\,U_1 \dots U_r$ ($h \in \{S, x_0\}$, $r \ge 0$), define $\Phi_{x_0}(U)$ to hold if either $h = x_0$ or there exists $i \in \{1, \dots, r\}$ whose leftmost leaf is $x_0$. Under a root head step $S A_1 A_2 A_3 \to_h (A_1 A_3)(A_2 A_3)$, if the leftmost leaf of $A_1$ is $x_0$, then $(A_1 A_3)(A_2 A_3)$ has head leaf $x_0$; if the leftmost leaf of $A_2$ is $x_0$, then the final spine argument $A_2 A_3$ has leftmost leaf $x_0$; and if the leftmost leaf of $A_3$ is $x_0$, then $A_3$ itself appears on the left spine of $(A_1 A_3)(A_2 A_3)$ with leftmost leaf $x_0$. Inductive left steps $L R \to_h L' R$ likewise preserve $\Phi_{x_0}$.

Now let $F, W \in \TS$ be closed with $\sem{F W}_{\HN} \neq \infty_{\HN}$. By Theorem~\ref{thm:question-10}(i) and Lemma~\ref{lem:head-conservation}, $F W \in \mathrm{SHN}$. Because substituting $\sigma_W(x_0) = W$ maps every open head step $U \to_h U'$ to a closed head step $\sigma_W(U) \to_h \sigma_W(U')$, accessibility pulls back from $\sigma_W(F x_0) \equiv F W \in \mathrm{SHN}$ to $F x_0 \in \mathrm{SHN}$ without referencing $\mu_{\HN}$. Thus open head reduction from $F x_0$ terminates at a head-irreducible $\{S, x_0\}$-term $U^*$ satisfying $\Phi_{x_0}(U^*)$, which must have: (a)~head leaf $x_0$; (b)~$U^* \equiv S A_1$ where $A_1$ has head leaf $x_0$; or (c)~$U^* \equiv S A_1 A_2$ where $A_1$ or $A_2$ has head leaf $x_0$. Hence left propagation at depth~$d$ implies right propagation at depth~$d$: if $\sem{F W}_{\HN} = \infty_{\HN}$, then $\bot_0(F W)$ gives $\bot_{d+1}(F W)$; otherwise $F W \redSs \sigma_W(U^*)$, and any subterm $K$ with head leaf $x_0$ substitutes to $W\,\sigma_W(K_1)\cdots\sigma_W(K_m)$, satisfying $\bot_d$ by $m$-fold left propagation at depth~$d$ and yielding $\bot_{d+1}(F W)$ by Lemma~\ref{lem:conv-invariance}.

We now prove both laws in~\eqref{eq:both-sink-laws} simultaneously by induction on $d \in \mathbb{N}$. For $d = 0$, if $\bot_0(Y)$, Theorem~\ref{thm:question-10}(i) gives $\sem{Y}_{\HN} = \infty_{\HN}$, so $\sem{Y Z}_{\HN} = \infty_{\HN} \odot_{\HN} \sem{Z}_{\HN} = \infty_{\HN}$ by~\eqref{eq:hn-asymmetry}, giving $\bot_0(Y Z)$ and hence right propagation at $d = 0$. For $d + 1$, assume~\eqref{eq:both-sink-laws} at depth~$d$ and suppose $\bot_{d+1}(Y)$. If $Y \notin \HN$, then $Y Z \notin \HN$ gives $\bot_{d+1}(Y Z)$. If $Y \redSs S A_1$ with $\bot_d(A_1)$, then $Y Z \redSs S A_1 Z \in \HNF$ gives $\bot_{d+1}(Y Z)$. If $Y \redSs S A_1 A_2$, then $Y Z \redSs (A_1 Z)(A_2 Z)$: when $\bot_d(A_1)$, two applications of depth-$d$ left propagation give $\bot_d((A_1 Z)(A_2 Z))$ and $\bot_{d+1}(Y Z)$; when $\bot_d(A_2)$, depth-$d$ left propagation gives $\bot_d(A_2 Z)$ and depth-$d$ right propagation gives $\bot_{d+1}((A_1 Z)(A_2 Z))$, whence $\bot_{d+1}(Y Z)$ by Lemma~\ref{lem:conv-invariance}. Right propagation at $d + 1$ then follows from left propagation at $d + 1$.
\end{proof}

\begin{lemma}[Zachos right-application $\bot$-absorption]\label{lem:zachos-bot-absorption}
Let $G_k$ ($k \in \mathbb{N}$) be Zachos's normal-form tower~\eqref{eq:zachos-tower}. For all $m, n \in \mathbb{N}$ and every closed $S$-term $U \in \TS$, there exists $d \in \mathbb{N}$ such that $\bot_d\bigl((G_m G_n) U\bigr)$; in particular, $(G_m G_n) U \notin \Prod$.
\end{lemma}

\begin{proof}
We first establish an auxiliary claim by induction on $n \in \mathbb{N}$: for any closed $Z, W \in \TS$, if $\sem{(G_0 Z) W}_{\HN} = \infty_{\HN}$, then there exists $d \in \mathbb{N}$ with $\bot_d((G_n Z) W)$. For $n = 0$, Theorem~\ref{thm:question-10}(i) directly gives $\bot_0((G_0 Z) W)$. For $n + 1$, expanding $G_{n+1} \equiv S S G_n$ gives the two-step reduction
\[
(G_{n+1} Z) W \;\equiv\; (S S G_n Z) W \;\redS\; \bigl(S Z (G_n Z)\bigr) W \;\redS\; (Z W)\bigl((G_n Z) W\bigr).
\]
By the induction hypothesis, $\bot_d((G_n Z) W)$ holds for some $d$; applying Lemma~\ref{lem:bot-propagation} to $(Z W)((G_n Z) W)$ yields $\bot_{d+1}\bigl((Z W)((G_n Z) W)\bigr)$, and Lemma~\ref{lem:conv-invariance} gives $\bot_{d+1}((G_{n+1} Z) W)$.

Now consider $(G_m G_n) U$ for $m = k + 1 \ge 1$. Two contractions yield $(G_{k+1} G_n) U \redS^2 (G_n U)\bigl((G_k G_n) U\bigr)$. Set $W \equiv (G_k G_n) U$. In $\mathcal{A}_{\HN}$, one computes $\sem{G_0}_{\HN} = q_{16}$ and $\sem{G_{r+1}}_{\HN} = q_{13}$ for all $r \ge 0$, whence $\sem{G_k G_n}_{\HN} \in \{q_{23}, q_{24}\}$. Inspection of rows $q_{23}$ and $q_{24}$ of $\mathcal{A}_{\HN}$ shows that for every $h = \sem{U}_{\HN} \in Q_{\HN}$, both $q_{23} \odot_{\HN} h$ and $q_{24} \odot_{\HN} h$ belong to $\{q_{23}, \infty_{\HN}\}$, so $\sem{W}_{\HN} \in \{q_{23}, \infty_{\HN}\}$. Moreover, for every $h \in Q_{\HN}$, the state $s = q_{16} \odot_{\HN} h = \sem{G_0 U}_{\HN}$ satisfies $s \odot_{\HN} q_{23} = \infty_{\HN}$ and $s \odot_{\HN} \infty_{\HN} = \infty_{\HN}$, giving $\sem{(G_0 U) W}_{\HN} = \infty_{\HN}$. By the auxiliary claim above, $(G_n U) W$ satisfies $\bot_d$ for some $d$, so $(G_{k+1} G_n) U$ satisfies $\bot_d$. Finally, for $m = 0$, we have $(G_0 G_n) U \equiv (S T T G_n) U \redS (G_{n+1} G_{n+1}) U$, reducing immediately to $m = n + 1 \ge 1$.
\end{proof}

We now define the $39$-state productivity $S$-algebra $\mathcal{A}_P = (Q_P, q_0, \odot_P)$ with carrier $Q_P = \{q_0, q_1, \dots, q_{36}, \infty_P, q_{38}\}$ (identifying the unique rejecting sink as $\infty_P = q_{37}$ and the accepting set as $F_P = Q_P \setminus \{\infty_P\}$, $|F_P| = 38$). Direct verification over all $39^3 = 59{,}319$ triples confirms~\eqref{eq:s-algebra-law}, and $\infty_P$ is a two-sided sink ($\infty_P \odot_P y = x \odot_P \infty_P = \infty_P$), matching Lemma~\ref{lem:bot-propagation} algebraically, while row $q_{29}$ satisfies $q_{29} \odot_P y = \infty_P$ for every $y \in Q_P$, matching Lemma~\ref{lem:zachos-bot-absorption} since $\sem{G_m G_n}_P = q_{29}$.

\begin{proposition}[Coinductive soundness of $\mathcal{A}_P$ and $34$-of-$38$ strong-normalization stratification]\label{prop:ap-soundness-stratification}
\begin{enumerate}[label=(\roman*),leftmargin=2.2em]
  \item \emph{HNF closure and soundness:} For all $x, y \in Q_P$, one has $q_0 \odot_P x \in F_P$ if and only if $x \in F_P$, and $(q_0 \odot_P x) \odot_P y \in F_P$ if and only if both $x \in F_P$ and $y \in F_P$. Every closed $S$-term $M$ with $\sem{M}_P \neq \infty_P$ satisfies $\sem{M}_{\HN} \neq \infty_{\HN}$ and $M \in \Prod$.
  \item \emph{Non-isomorphism with $\mathcal{A}_{\SN}$ and unique non-$\SN$ root-redex fiber:} Although $|Q_P| = |Q_{\SN}| = 39$, the reachable subalgebra of $\mathcal{A}_P \times \mathcal{A}_{\SN}$ has $44$ states. Thirty-four of the $38$ accepting states of $\mathcal{A}_P$---every $q \in F_P \setminus \{q_{18}, q_{24}, q_{29}, q_{33}\}$---pair exclusively with $F_{\SN}$ and therefore contain only strongly normalizing terms. Furthermore, every term in $(q_{18}, \infty_{\SN})$, $(q_{24}, \infty_{\SN})$, or $(q_{33}, \infty_{\SN})$ reduces to a spine-$1$ or spine-$2$ $\HNF$ without executing any root contraction, leaving $(q_{29}, \infty_{\SN})$ as the unique non-$\SN$ productive root-redex state pair in $\CLS$.
\end{enumerate}
\end{proposition}

\begin{proof}
For~(i), the $\HNF$ biconditionals hold by direct inspection of $\odot_P$. Reachable closure of $(q_0, q_0) \in Q_P \times Q_{\HN}$ yields $113$ pairs ($95$ with $\sem{M}_P \in F_P$ and $18$ with $\sem{M}_P = \infty_P$), and $(\infty_P, \infty_{\HN})$ is the sole pair whose second coordinate is $\infty_{\HN}$; hence $\sem{M}_P \neq \infty_P$ implies $\sem{M}_{\HN} \neq \infty_{\HN}$. Induction on $d \in \mathbb{N}$ shows that $\sem{M}_P \neq \infty_P$ implies $\neg\bot_d(M)$: at $d = 0$, $\sem{M}_{\HN} \neq \infty_{\HN}$ refutes $M \notin \HN$ by Theorem~\ref{thm:question-10}(i); at $d + 1$, if $M \redSs S A_1 \dots A_k$ ($1 \le k \le 2$), conversion invariance gives $\sem{S A_1 \dots A_k}_P = \sem{M}_P \in F_P$, so the $\HNF$ biconditionals give $\sem{A_i}_P \in F_P$ and $\neg\bot_d(A_i)$ by the depth-$d$ hypothesis.

For~(ii), reachable closure of $(q_0, q_0) \in Q_P \times Q_{\SN}$ produces $44$ pairs, of which only $(q_{18}, \infty_{\SN})$, $(q_{24}, \infty_{\SN})$, $(q_{29}, \infty_{\SN})$, $(q_{33}, \infty_{\SN})$, and $(\infty_P, \infty_{\SN})$ have second coordinate $\infty_{\SN}$; Theorem~\ref{thm:waldmann-sn} implies all other $34$ states of $F_P$ inhabit $\SN$. Inspecting all joint transitions landing in $(q_{18}, \infty_{\SN})$, $(q_{24}, \infty_{\SN})$, or $(q_{33}, \infty_{\SN})$ shows that in every such application $L R$, the left child $L$ satisfies $\sem{L}_{\SN} \neq \infty_{\SN}$ and normalizes to $S$ or $S U$, so $L R \redSs S R$ or $S U R \in \HNF$ without root contractions. Only $(q_{29}, \infty_{\SN})$ admits non-$\SN$ root-contracting terms, such as $Z_{\mathrm{Zac}} \equiv G_0 G_0$ and $Z_0 \equiv S(S T)(S T)T$.
\end{proof}

To prove completeness (that $\sem{M}_P = \infty_P$ implies $\bot_d(M)$ for some $d \in \mathbb{N}$), we replace normal forms $\mathcal{N}$ by a grammar of \emph{canonical productive values} $\mathcal{V}_P \subset \TS$ that internalizes the unique non-$\SN$ productive root-redex family $G_m G_n$ directly as a value constructor:
\begin{equation}\label{eq:pval-grammar}
\mathcal{V}_P \;::=\; S \;\mid\; S V \;\mid\; S V_1 V_2 \;\mid\; G_m G_n
\qquad (V, V_1, V_2 \in \mathcal{V}_P,\; m, n \in \mathbb{N}),
\end{equation}
with structural size $\mathrm{sz}(S) = \mathrm{sz}(G_m G_n) = 1$, $\mathrm{sz}(S V) = \mathrm{sz}(V) + 1$, and $\mathrm{sz}(S V_1 V_2) = \mathrm{sz}(V_1) + \mathrm{sz}(V_2) + 1$.

\begin{lemma}[Syntactic rigidity of $\mathcal{V}_P$, $35$-override $2$D collapse, and rank-$6$ normalization]\label{lem:pval-normalization}
Every $V \in \mathcal{V}_P$ satisfies $\sem{V}_P \in F_P$, and $\mathcal{V}_P$ obeys three syntactic rigidity laws:
\begin{equation}\label{eq:pval-rigidity}
\sem{V}_P = q_0 \Longrightarrow V \equiv S,
\qquad
\sem{V}_P = q_1 \Longrightarrow V \equiv T \equiv S S,
\qquad
\sem{V}_P \in \{q_{14}, q_{32}\} \Longrightarrow \exists k \in \mathbb{N},\; V \equiv G_k.
\end{equation}
Moreover, collapsing the $39^3 = 59{,}319$ triple-rank cube to a $39 \times 39$ matrix $\hat{\tau}_P((q_0 \odot_P x) \odot_P y,\, z) \in \{0, \dots, 6\}$ modified on $35$ explicit overrides ($34$ entries at $x = q_1$ assigned rank~$2$, plus the single triple $(q_6, q_0, q_1)$ assigned rank~$3$) yields a total rank-$6$ application operator $\mathrm{app}_{\mathcal{V}_P}(V_1, V_2) \in \mathcal{V}_P$ for all $V_1, V_2 \in \mathcal{V}_P$ with $\sem{V_1}_P \odot_P \sem{V_2}_P \neq \infty_P$ satisfying $V_1 V_2 \redSs \mathrm{app}_{\mathcal{V}_P}(V_1, V_2)$. Consequently, every closed $M \in \TS$ with $\sem{M}_P \neq \infty_P$ reduces constructively to some $V \in \mathcal{V}_P$.
\end{lemma}

\begin{proof}
Because $\sem{G_0}_P = q_{14}$, $\sem{G_{k+1}}_P = q_{32}$, and $\sem{G_m G_n}_P = q_{29} \in F_P$, structural induction on~\eqref{eq:pval-grammar} via Proposition~\ref{prop:ap-soundness-stratification}(i) confirms $\sem{V}_P \in F_P$ for all $V \in \mathcal{V}_P$. For~\eqref{eq:pval-rigidity}, induct on $V \in \mathcal{V}_P$: $q_0$ never lies in the image of $\odot_P$, forcing $V \equiv S$; the sole spine-$1$ or spine-$2$ transition into $q_1$ is $q_0 \odot_P q_0 = q_1$, forcing $V \equiv S S \equiv T$; no spine-$1$ state lies in $\{q_{14}, q_{32}\}$, while $(q_0 \odot_P x) \odot_P y = q_{14}$ forces $x = y = q_1$ ($V \equiv S T T \equiv G_0$), and $(q_0 \odot_P x) \odot_P y = q_{32}$ forces $x = q_0$ and $y \in \{q_{14}, q_{32}\}$ ($V \equiv S S G_k \equiv G_{k+1}$).

To define $\mathrm{app}_{\mathcal{V}_P}(V_1, V_2)$ when $\sem{V_1}_P \odot_P \sem{V_2}_P \neq \infty_P$, note that $V_1$ cannot be $G_m G_n$ since $q_{29} \odot_P z = \infty_P$. If $V_1 \equiv S$ or $V_1 \equiv S U$, then $V_1 V_2 \in \mathcal{V}_P$ directly. If $V_1 \equiv S X Y$ with $X, Y \in \mathcal{V}_P$, first test whether both $\sem{V_1}_P \in \{q_{14}, q_{32}\}$ and $\sem{V_2}_P \in \{q_{14}, q_{32}\}$: if so, \eqref{eq:pval-rigidity} identifies unique $m, n \in \mathbb{N}$ with $V_1 \equiv G_m$ and $V_2 \equiv G_n$, so $V_1 V_2 \equiv G_m G_n \in \mathcal{V}_P$ in zero steps. Otherwise, contract $(S X Y) V_2 \redS (X V_2)(Y V_2)$; across all non-Zachos accepted triples, the $35$-override collapsed rank $\tau_P(x, y, z) \in \{1, \dots, 6\}$ satisfies $\hat{\tau}_P(x, z) \le \tau_P(x, y, z)$ and $\hat{\tau}_P(y, z) \le \tau_P(x, y, z)$ (normalizing $X V_2 \redSs W_1$ and $Y V_2 \redSs W_2$ by inner induction on $\mathrm{sz}(V_1)$), and whenever $(W_1, W_2)$ is not itself a Zachos pair $(G_{m'}, G_{n'})$, satisfies strict outer rank decrease as in~\eqref{eq:sn-rank-1}--\eqref{eq:sn-rank-2}. Structural induction on $M \in \TS$ via the two-sided sink $\infty_P$ then normalizes every accepted $M$ into $\mathcal{V}_P$.
\end{proof}

\begin{theorem}[Regularity of B\"ohm-tree productivity and rank-$10$ resolution of Question~5 {\cite[p.~80]{BarendregtEndrullisKlopWaldmann2017}}]\label{thm:q5-main}
For every closed $S$-term $M \in \TS$,
\begin{equation}\label{eq:q5-iff}
M \in \Prod \;\Longleftrightarrow\; \sem{M}_P \neq \infty_P,
\end{equation}
so the set $\Prod$ of productive closed $S$-terms is a regular tree language decidable in $2|M| - 1$ deterministic table lookups. Constructively, for every closed $M \in \TS$, either $M \in \Prod$ or a unified $10$-rank induction over the $113$-state reachable subalgebra $\mathrm{PH} \subseteq Q_P \times Q_{\HN}$ extracts an explicit depth $d \in \mathbb{N}$ satisfying $\bot_d(M)$.
\end{theorem}

\begin{proof}
Soundness (that $\sem{M}_P \neq \infty_P$ implies $M \in \Prod$) was proved in Proposition~\ref{prop:ap-soundness-stratification}(i). For constructive completeness, structural induction on $M$ with $\sem{M}_P = \infty_P$ reduces via Lemma~\ref{lem:bot-propagation} to minimal critical applications $L R$ with $\sem{L}_P, \sem{R}_P \in F_P$ and $\sem{L}_P \odot_P \sem{R}_P = \infty_P$, which Lemma~\ref{lem:pval-normalization} normalizes to critical pairs $(F, Z) \in \mathcal{V}_P \times \mathcal{V}_P$ satisfying $\sem{F Z}_P = \infty_P$. If $F \equiv G_m G_n$, Lemma~\ref{lem:zachos-bot-absorption} yields $\bot_d((G_m G_n) Z)$ directly; since $F \equiv S$ and $F \equiv S U$ preserve $F_P$, only $F \equiv S X Y$ ($X, Y \in \mathcal{V}_P$) remains, contracting in one step to $(X Z)(Y Z)$.

Over the $113$-state reachable subalgebra $\mathrm{PH} \subseteq Q_P \times Q_{\HN}$ ($95$ accepting states $\mathrm{PH}_{\mathrm{acc}}$ and $18$ rejecting states), we unify accepted values $V \in \mathcal{V}_P$ (assigned value size $\mathrm{psz}(V) = \mathrm{sz}(V)$) and solvable rejected raw terms $R$ with $\sem{R}_P = \infty_P$ and $\sem{R}_{\HN} \neq \infty_{\HN}$ (assigned value size $\mathrm{psz}(R) = 0$) into a single argument sort. Recording for each $(\mathbf{x}, \mathbf{z}) \in \mathrm{PH}_{\mathrm{acc}}^2$ the finite set $\mathrm{out}_{2\mathrm{d}}(\mathbf{x}, \mathbf{z}) \subseteq \mathrm{PH}_{\mathrm{acc}}^2$ of spine-$2$ child state pairs $(\sem{A_1}_{\mathrm{PH}}, \sem{A_2}_{\mathrm{PH}})$ attainable when $\mathrm{app}_{\mathcal{V}_P}(X, V_2) \equiv S A_1 A_2$, and evaluating direct product ranks on $\sem{X Z}_{\mathrm{PH}} \odot_{\mathrm{PH}} \sem{Y Z}_{\mathrm{PH}}$ before expanding spine-$2$ joint critical transitions (shrinking the spine-$2$ joint critical table to $174$ entries across $98$ pairs $(\mathbf{x}, \mathbf{z})$), every critical pair $(\mathbf{f}, \mathbf{z}) \in \mathrm{PH}_{\mathrm{acc}} \times \mathrm{PH}$ with $\mathbf{f} \odot_{\mathrm{PH}} \mathbf{z} \notin \mathrm{PH}_{\mathrm{acc}}$ receives a rank $\rho_2(\mathbf{f}, \mathbf{z}) \in \{1, \dots, 10\}$ and every joint critical triple receives a rank $\rho_3(\mathbf{x}, \mathbf{z}, \mathbf{n}_2) \in \{1, \dots, 10\}$. Without parity splitting, every recursive branch either hits $\infty_{\HN}$ immediately, reduces to Lemma~\ref{lem:zachos-bot-absorption}, strictly decreases the outer rank $r \in \{1, \dots, 10\}$, or stays at rank~$r$ while strictly decreasing the unified inner measure ($2\,\mathrm{sz}(F)$ on pairs $(F, Z)$ and $2(\mathrm{sz}(X) + \mathrm{psz}(Y)) + 1$ on triples $(S X Y, Z)$), constructing an explicit finite $d$ with $\bot_d(F Z)$.
\end{proof}

\begin{proposition}[Myhill--Nerode minimality of the $113$-state joint algebra $\mathrm{PH}$]\label{prop:ph-113-minimal}
All $113$ states of $\mathrm{PH} \subseteq \mathcal{A}_P \times \mathcal{A}_{\HN}$ are inhabited by explicit closed $S$-terms $W_0, \dots, W_{112} \in \TS$, and there exists a fixed family of $17$ closed one-hole contexts $C_0[\square], \dots, C_{16}[\square]$ such that for any two closed $S$-terms $M_1, M_2 \in \TS$ with $\sem{M_1}_{\mathrm{PH}} \neq \sem{M_2}_{\mathrm{PH}}$, some $C_k[\square]$ ($0 \le k < 17$) separates $C_k[M_1]$ from $C_k[M_2]$ on $(\Prod, \HN)$:
\[
\neg\Bigl(\bigl(C_k[M_1] \in \Prod \iff C_k[M_2] \in \Prod\bigr) \;\wedge\; \bigl(C_k[M_1] \in \HN \iff C_k[M_2] \in \HN\bigr)\Bigr).
\]
Consequently, every deterministic bottom-up tree automaton or magma quotient $(Q, q_0, \odot)$ on $\TS$ through which $(\Prod, \HN)$ factors has at least $113$ states.
\end{proposition}

\begin{proof}
Each witness $W_i \in \TS$ ($0 \le i < 113$) is constructed by breadth-first generation from $S$, and each context $C_k[\square]$ ($0 \le k < 17$) is an iterated left and right application chain built from $\{W_i\}$. On $\mathrm{PH}$, Theorems~\ref{thm:question-10}(i) and~\ref{thm:q5-main} partition the $113$ states into three semantic classes: $95$ states in $\Prod$, $17$ states in $\HN \setminus \Prod$, and $1$ state $(\infty_P, \infty_{\HN})$ in $\overline{\HN}$. Direct evaluation over all $\binom{113}{2} = 6{,}328$ unordered state pairs $(i, j)$ confirms that at least one of the $17$ context transitions maps $i$ and $j$ into distinct semantic classes among these three. If $\sem{\cdot}_Q \colon \TS \to Q$ is any magma homomorphism deciding $(\Prod, \HN)$ and $\sem{W_i}_Q = \sem{W_j}_Q$ for $i \neq j$, congruence of $\sem{\cdot}_Q$ under every closed context $C_k[\square]$ gives $\sem{C_k[W_i]}_Q = \sem{C_k[W_j]}_Q$, contradicting semantic separation of $C_k[W_i]$ and $C_k[W_j]$. By the pigeonhole principle, $|Q| \ge 113$.
\end{proof}

Even at leaf size~$8$, completeness of $\mathcal{A}_P$ requires selective branch choice: chasing the leftmost unproductive child of each $\HNF$ can loop forever without visiting $\infty_{\HN}$.

\begin{proposition}[B\"ohm-game branch asymmetry at leaf size~$8$]\label{prop:bohm-game-asymmetry}
Consider the $8$-leaf closed $S$-term $L_{\mathrm{loop}} \equiv S T S T T \equiv S(S S)S(S S)(S S)$, which satisfies $\sem{L_{\mathrm{loop}}}_P = \infty_P$ and $\sem{L_{\mathrm{loop}}}_{\HN} = q_{24} \neq \infty_{\HN}$. Under the rightmost-unproductive B\"ohm strategy (head-normalizing and selecting the rightmost spine child having $P$-state $\infty_P$), $L_{\mathrm{loop}}$ reaches $\infty_{\HN}$ in $4$ steps, witnessing $\bot_4(L_{\mathrm{loop}})$. Under the leftmost-unproductive B\"ohm strategy, by contrast, $L_{\mathrm{loop}}$ never reaches $\infty_{\HN}$, cycling forever through solvable terms $K_0 K_n \mapsto^{n+2} K_0 K_{n+1}$ ($n \ge 1$), where $K_0 \equiv S S (S T)T$ and $K_{m+1} \equiv T K_m$ satisfy $|K_m| = 2m + 7 \to \infty$.
\end{proposition}

\begin{proof}
Set $G_0 \equiv S T T$, $K_0 \equiv S S (S T) T \equiv T(S T)T$, and $K_{m+1} \equiv T K_m \in \mathcal{N}$ ($\sem{K_m}_P \in F_P$). Five deterministic head steps reduce $L_{\mathrm{loop}}$ to a spine-$2$ head normal form:
\[
S T S T T
\;\longrightarrow_h\; (T T)(S T)T
\;\longrightarrow_h\; S(S T)\bigl(T(S T)\bigr)T
\;\longrightarrow_h\; G_0\,K_0
\;\longrightarrow_h\; K_1\,K_1
\;\longrightarrow_h\; S\,K_1\,(K_0 K_1) \;\in\; \HNF.
\]
Because $\sem{K_1}_P \in F_P$ and $\sem{K_0 K_1}_P = \infty_P$, the first step of both the leftmost- and rightmost-unproductive strategies enters $K_0 K_1$. For every $n \ge 1$, three head contractions yield
\[
K_0 K_n \;\equiv\; (S S (S T) T)\,K_n
\;\longrightarrow_h\; (S T G_0)\,K_n
\;\longrightarrow_h\; (T K_n)(G_0 K_n)
\;\longrightarrow_h\; S\,(G_0 K_n)\,\bigl(K_n(G_0 K_n)\bigr) \;\in\; \HNF,
\]
in which both spine children $G_0 K_n$ and $K_n(G_0 K_n)$ evaluate to $\infty_P$ in $\mathcal{A}_P$. At $n = 1$, the rightmost-unproductive strategy selects $K_1(G_0 K_1)$, which in two further B\"ohm transitions reaches $(G_0 K_1)\bigl(G_0(G_0 K_1)\bigr)$, a term of $\HN$ state $\infty_{\HN}$, certifying $\bot_4(L_{\mathrm{loop}})$. Conversely, for every $n \ge 1$, the leftmost-unproductive strategy selects $G_0 K_n \equiv S T T K_n \to_h^2 S\,K_{n+1}\,(K_n K_{n+1}) \in \HNF$, whose sole $\infty_P$-child is $K_n K_{n+1}$; $n$ subsequent head-normalizations $K_{m+1} K_{n+1} \equiv S S K_m K_{n+1} \to_h S\,K_{n+1}\,(K_m K_{n+1})$ ($m = n - 1, \dots, 0$) each have productive left child $K_{n+1} \in \mathcal{N}$ and return to $K_0 K_{n+1}$ in $n + 2$ B\"ohm steps without ever visiting $\infty_{\HN}$.
\end{proof}

\subsection{Automaton-Free Book Proofs, Complete B\"ohm Faithfulness Dichotomy, and Productive Convertibility}\label{subsec:q7-q31-q30}

Although finite $S$-algebras provide linear-time decision procedures across all of $\TS$, Questions~7 and~31 of Barendregt et al.~\cite[pp.~81, 93]{BarendregtEndrullisKlopWaldmann2017} admit short, self-contained proofs that require no finite-automaton multiplication tables at all. Following a suggestion of J\"org Endrullis to simplify these arguments as far as possible, we first present a three-constructor syntactic invariant $\mathcal{H}$ (condensing Waldmann's $Q Q Q$ generator~\cite{Waldmann2000}) that yields pen-and-paper proofs of both Question~7 (including the atomic $S \mapsto S S$ leaf-splitting law) and the unsolvable half of Question~31. We then combine canonical productive values $\mathcal{V}_P$ from Subsection~\ref{subsec:q5-productivity} with this unsolvable separation to establish a complete B\"ohm-tree faithfulness dichotomy on $\CLS$ and a total convertibility decider across $\Prod \times \TS$, advancing Scott's Question~30~\cite{Scott1975,Canal1978,Statman1986,Statman1989}.

\begin{definition}[Three-constructor invariant $\mathcal{H}$ and triple-application class $\mathcal{U}$]\label{def:heavy-invariant}
Let $\mathcal{H} \subset \TS$ be the smallest set of closed $S$-terms generated inductively by
\[
\mathcal{H} \;::=\; A \;\mid\; S X \;\mid\; X Y
\qquad (A \equiv S S S,\; X, Y \in \mathcal{H}),
\]
and define the triple-application class $\mathcal{U} := \{ P Q R : P, Q, R \in \mathcal{H} \} \subset \TS$.
\end{definition}

\begin{lemma}[Inversion, reduction closure, and divergence of $\mathcal{U}$]\label{lem:heavy-closure}
\begin{enumerate}[label=(\roman*),leftmargin=2.2em]
  \item \emph{Inversion:} One has $S \notin \mathcal{H}$ and $\ls(M) \ge 1$ for every $M \in \mathcal{H}$; if $S X \in \mathcal{H}$, then $X \in \mathcal{H}$; and if $S X Y \in \mathcal{H}$, then either $X \equiv Y \equiv S$ or $X, Y \in \mathcal{H}$.
  \item \emph{Forward $\redS$-closure:} Both $\mathcal{H}$ and $\mathcal{U}$ are closed forward under arbitrary single-step reduction $\redS$.
  \item \emph{Divergence:} Every $M \in \mathcal{U}$ satisfies $\ls(M) \ge 3$, $M \notin \HN$, and $M \notin \SN$.
\end{enumerate}
\end{lemma}

\begin{proof}
For~(i), $A \equiv (S S) S$, $S X$, and $X Y$ are all applications, so $S \notin \mathcal{H}$ and $\ls(M) \ge 1$. A term $S X$ cannot arise as $A$ (since $S \not\equiv S S$) or as $U V$ with $U, V \in \mathcal{H}$ (since $S \notin \mathcal{H}$), forcing $X \in \mathcal{H}$. Similarly, $S X Y \equiv (S X) Y$ cannot arise from the second constructor ($S X \not\equiv S$); if it is $A$, then $X \equiv Y \equiv S$, and if it arises as $(S X) Y$ with $S X, Y \in \mathcal{H}$, then $X \in \mathcal{H}$.

For~(ii), we first prove that $M \in \mathcal{H}$ and $M \redS M'$ imply $M' \in \mathcal{H}$ by induction on $M \in \mathcal{H}$. The term $A \in \mathcal{N}$ has no steps; steps inside $S X \redS S X'$ or inside $U$ or $V$ of $U V$ preserve $\mathcal{H}$ by the induction hypothesis; and if $U V \equiv (S X Y) V \to_r (X V)(Y V)$ with $S X Y, V \in \mathcal{H}$, then~(i) gives either $X \equiv Y \equiv S$ (so $(S V)(S V) \in \mathcal{H}$) or $X, Y \in \mathcal{H}$ (so $(X V)(Y V) \in \mathcal{H}$). Next, let $M \equiv P Q R \in \mathcal{U}$ with $P, Q, R \in \mathcal{H}$ and $M \redS M'$. Steps inside $P$, $Q$, or $R$ preserve $\mathcal{U}$ by $\mathcal{H}$-closure; if $P Q \equiv (S X Y) Q \to_r (X Q)(Y Q)$, then~(i) places $X Q, Y Q \in \mathcal{H}$, so $(X Q)(Y Q) R \in \mathcal{U}$; and if $(P Q) R \equiv (S X Y) R \to_r (X R)(Y R)$, then $P \equiv S X$ and $Q \equiv Y$, so~(i) gives $X \in \mathcal{H}$ and $M' \equiv X R (Q R) \in \mathcal{U}$.

For~(iii), any $M \equiv P Q R \in \mathcal{U}$ has $\ls(M) = \ls(P) + 2 \ge 3$ by~(i). By~(ii), every reduct $M'$ of $M$ remains in $\mathcal{U}$ and satisfies $\ls(M') \ge 3 > 2$, whence $M \notin \HN$. Since every closed term with $\ls \ge 3$ possesses a head redex $S X Y Z$, iterating head contraction inside $\mathcal{U}$ produces an infinite reduction sequence, proving $M \notin \SN$.
\end{proof}

In Question~7~\cite[p.~81]{BarendregtEndrullisKlopWaldmann2017}, Barendregt et al.\ color $\SN$ green and $\TS \setminus \SN$ red, asking whether successive single-leaf refinements starting from $S$ can alternate green and red infinitely often. Write $M \prec N$ if $N$ replaces a single $S$-leaf of $M$ by a composite closed $S$-term, and $M \prec_1 N$ for an \emph{atomic leaf split} replacing one $S$-leaf by $T \equiv S S$ (so $|N| = |M| + 1$). Table~\ref{tab:atomic-bridge} lists the $10$-term atomic bridge $P_0 \prec_1 \cdots \prec_1 P_9$.

\begin{table}[htbp]
\centering
\caption{The $10$-term atomic leaf-splitting bridge $P_0 \prec_1 \cdots \prec_1 P_9$ for Question~7 ($T \equiv S S$, $A \equiv S S S$).}
\label{tab:atomic-bridge}
\small
\begin{tabular}{@{}rll@{}}
\toprule
$r$ & $P_r$ & Automaton-free normalization status \\
\midrule
$0, 1, 2, 3, 4$ & $S,\; T,\; A,\; S T S,\; S A S$ & Closed normal forms in $\mathcal{N}$ \\
$5$ & $S(A S)S$ & $A S \redS T T \in \mathcal{N}$ ($2$-term reduction graph) \\
$6$ & $S(A T)S$ & $A T \redS (S T)(S T) \in \mathcal{N}$ ($2$-term reduction graph) \\
$7$ & $S(A A)S$ & $A A \redS (S A)(S A) \in \mathcal{N}$ ($2$-term reduction graph) \\
$8$ & $S(S T S A)S$ & $S T S A \redS^3 (S(S A))\bigl((S(S A))(S(S A))\bigr) \in \mathcal{N}$ ($4$-term chain) \\
$9$ & $S(S A S A)S \equiv S\,R(S)\,S$ & Non-normalizing via $R(S) \redS^2 ((S A)(S A))(S A) \in \mathcal{U}$ \\
\bottomrule
\end{tabular}
\end{table}

\begin{theorem}[Automaton-free refinement alternation, atomic period-$10$ law, and rescue classification, Question~7 {\cite[p.~81]{BarendregtEndrullisKlopWaldmann2017}}]\label{thm:question-7}
Define the gadget $R(X) \equiv X A S A$ and the right-nested spine $C^0(M) \equiv M$, $C^{k+1}(M) \equiv S\,R(T)\,C^k(M)$.
\begin{enumerate}[label=(\roman*),leftmargin=2.2em]
  \item \emph{Automaton-free $\SN / \overline{\SN}$ refinement chain:} Without invoking any finite $S$-algebra, $R(S) \equiv S A S A \notin \SN$ whereas $R(T) \equiv T A S A \in \SN$. Setting
  \begin{equation}\label{eq:q7-chain}
  F_{2k} \;\equiv\; C^k(S),
  \qquad
  F_{2k+1} \;\equiv\; C^k\bigl(S\,R(S)\,S\bigr)
  \qquad (k \in \mathbb{N}),
  \end{equation}
  one has $F_0 \equiv S$, $F_m \prec F_{m+1}$ for all $m \in \mathbb{N}$, $F_{2k} \in \SN$, and $F_{2k+1} \notin \SN$ for all $k \in \mathbb{N}$.
  \item \emph{Exact color pattern under atomic leaf splits $S \mapsto S S$:} With $P_0, \dots, P_9$ as in Table~\ref{tab:atomic-bridge}, set $T_{10k + r} \equiv C^k(P_r)$ for $k \in \mathbb{N}$ and $0 \le r \le 9$. Then $T_0 \equiv S$, $T_n \prec_1 T_{n+1}$ and $|T_n| = n + 1$ for every $n \in \mathbb{N}$, and
  \[
  T_n \in \SN \;\Longleftrightarrow\; n \not\equiv 9 \pmod{10}.
  \]
  \item \emph{Universal one-leaf rescue classification:} Across all $940$ one-hole context functions $f_C \colon Q_{\SN} \to Q_{\SN}$ on $\mathcal{A}_{\SN}$, if $C[S] \notin \SN$ and $C[P] \in \SN$ for $P \in \TS$, then $\sem{P}_{\SN} \in \{q_1, q_8\}$, forcing $P \equiv S S$ or $P \equiv S(S S S)$ syntactically; both plugs are realized ($C_1[\square] \equiv \square A S A$ and $C_2[\square] \equiv \square (S S) S S S S$).
\end{enumerate}
\end{theorem}

\begin{proof}
For~(i), two steps reduce $R(S) \equiv S A S A \redS (A A)(S A) \equiv ((S S S) A)(S A) \redS ((S A)(S A))(S A)$. Since $A \in \mathcal{H}$ implies $S A \in \mathcal{H}$, the term $((S A)(S A))(S A)$ belongs to $\mathcal{U}$, so $R(S) \notin \SN$ by Lemma~\ref{lem:heavy-closure}(iii). Conversely, since $T, A, S A, T A, T T \in \mathcal{N}$, the complete multiway reduction graph of $R(T) \equiv ((T A) S) A$ has only six terms: $g_0 \equiv ((T A) S) A \redS g_1 \equiv (T(A S))A$, branching to $g_2 \equiv (S A)((A S) A)$ and $g_3 \equiv (T(T T)) A$, both stepping to $g_4 \equiv (S A)((T T) A) \redS g_5 \equiv (S A)((S A)(T A)) \in \mathcal{N}$; hence every path terminates in at most $4$ steps and $R(T) \in \SN$. Because $S U V \in \SN$ holds if and only if both $U \in \SN$ and $V \in \SN$ by Lemma~\ref{lem:wn-eq-sn}(i), and because $C^{k+1}(S) \equiv C^k(S\,R(T)\,S)$, induction on $k$ gives $F_{2k} \in \SN$, $F_{2k+1} \notin \SN$, and $F_m \prec F_{m+1}$.

For~(ii), each step $P_r \prec_1 P_{r+1}$ ($0 \le r \le 8$) and $P_9 \equiv S(S A S A)S \prec_1 S(T A S A)S \equiv S\,R(T)\,S$ splits one $S$-leaf into $S S$ ($r = 4 \to 5$ splits the head of $A$ into $A S$; $r = 7 \to 8$ splits the middle $S$ of the left $A$ in $A A$ into $S T S$; $r = 8 \to 9$ splits the first leaf of $T$ in $S T S A$ into $S A S A$; and $r = 9 \to 10$ splits the head $S$ of $S A S A$ into $R(T)$, matching $C^{k+1}(S) \equiv C^k(S\,R(T)\,S)$). For any $V \in \mathcal{N}$, the application $A V \equiv S S S V$ has the unique step $A V \redS (S V)(S V) \in \mathcal{N}$, proving $P_5, P_6, P_7 \in \SN$, while $S T S A$ has the unique unbranched $3$-step reduction $S T S A \redS (T A)(S A) \redS (S(S A))(A(S A)) \redS (S(S A))((S(S A))(S(S A))) \in \mathcal{N}$, proving $P_8 \in \SN$. Thus $T_{10k + r} \in \SN$ for $0 \le r \le 8$ and $T_{10k + 9} \notin \SN$.

For~(iii), closing $\mathrm{id}_{Q_{\SN}}$ under left and right multiplication on $Q_{\SN}$ yields $940$ context maps, and collecting $\{ x \in Q_{\SN} : \exists f,\, f(q_0) = \infty_{\SN} \textrm{ and } f(x) \neq \infty_{\SN} \}$ gives $\{q_1, q_8\}$. The unique normal forms in $q_1$ and $q_8$ are $S S$ and $S(S S S)$; since the sole reducible closed term of size at most~$4$ is $S S S S \redS S S (S S)$ in $q_9 \notin \{q_1, q_8\}$, Proposition~\ref{prop:cr-conservation}(ii) forces $P \in \{S S, S(S S S)\}$. On $C_2[\square] \equiv \square(S S)S S S S$, we have $\sem{C_2[S]}_{\SN} = \infty_{\SN}$ whereas $\sem{C_2[S(S S S)]}_{\SN} = q_{30} \neq \infty_{\SN}$ (normalizing in $1{,}109$ leftmost-outermost steps).
\end{proof}

We now turn to Question~31~\cite[p.~93]{BarendregtEndrullisKlopWaldmann2017}, which asks whether closed $S$-terms with identical B\"ohm trees (or identical \emph{clocked} B\"ohm trees, recording at each node the exact number of deterministic head steps to $\HNF$) must be $\convS$-convertible, highlighting $A A A$ versus $A A A S$. Together, Theorems~\ref{thm:question-31} and~\ref{thm:productive-bohm-faithful} establish a sharp dichotomy: whenever $\bot$ occurs at any finite depth, even clocked B\"ohm trees fail infinitely to separate $\convS$-classes (provable without any finite algebra); conversely, on all productive ($\bot$-free) terms, ordinary unclocked B\"ohm-tree equality is completely faithful across $\Prod \times \TS$.

For any term $M$ and $k \in \mathbb{N}$, define $M S^0 \equiv M$ and $M S^{k+1} \equiv (M S^k) S$. Call $M$ \emph{right-compound} if $M \equiv U (V_1 V_2)$, and let $\tau(M)$ count outer trailing applications to $S$ ($\tau(U S) = 1 + \tau(U)$, and $\tau(M) = 0$ otherwise).

\begin{theorem}[Frozen-tail cancellation, clocked-B\"ohm unfaithfulness, and $\mathcal{A}_{36}$, Question~31 Part~I {\cite[p.~93]{BarendregtEndrullisKlopWaldmann2017}}]\label{thm:question-31}
\begin{enumerate}[label=(\roman*),leftmargin=2.2em]
  \item \emph{Automaton-free frozen-tail cancellation:} For any two right-compound unsolvable terms $M, N \in \TSV \setminus \HN$ and all $j, k \in \mathbb{N}$,
  \[
  M S^j \convS N S^k \;\Longleftrightarrow\; \bigl(j = k \;\textrm{ and }\; M \convS N\bigr).
  \]
  In particular, since $A A A \in \mathcal{U}$ satisfies $A A A \notin \HN$ by Lemma~\ref{lem:heavy-closure}(iii) and has right child $A \equiv (S S) S$, one has $A A A S^j \convS A A A S^k$ if and only if $j = k$, without invoking any finite $S$-algebra, and hence $A A A \neq_S A A A S$.
  \item \emph{Clocked-B\"ohm unfaithfulness at arbitrary $\bot$-depth $D$:} Define $W_0(U) \equiv U$ and $W_{D+1}(U) \equiv S\,W_D(U)$. For every depth $D \in \mathbb{N}$ and all $j \neq k$, the closed terms $W_D(A A A S^j)$ and $W_D(A A A S^k)$ are head-normalizing through depth $D - 1$, share identical B\"ohm and clocked B\"ohm trees at every depth $d \in \mathbb{N}$ (with a single $\bot$ at depth $D$), yet satisfy $W_D(A A A S^j) \neq_S W_D(A A A S^k)$.
  \item \emph{Homomorphic $36$-state separation:} Splitting $\infty_{\HN}$ into $\{q_{34}, q_{35}\}$ on $Q_{36} = \{q_0, \dots, q_{35}\}$ via
  \begin{equation}\label{eq:a36-def}
  x \odot_{36} y \;=\;
  \begin{cases}
  x \odot_{\HN} \tilde{y} & \textrm{if } x \in \{q_0, \dots, q_{33}\} \textrm{ and } x \odot_{\HN} \tilde{y} \neq \infty_{\HN}, \\
  q_{34} & \textrm{if } x \in \{q_0, \dots, q_{33}\} \textrm{ and } x \odot_{\HN} \tilde{y} = \infty_{\HN}, \\
  q_{35} & \textrm{if } x \in \{q_{34}, q_{35}\} \textrm{ and } y = q_0, \\
  q_{34} & \textrm{if } x \in \{q_{34}, q_{35}\} \textrm{ and } y \in \{q_1, \dots, q_{35}\},
  \end{cases}
  \end{equation}
  where $\tilde{y} = y$ for $y \in \{q_0, \dots, q_{33}\}$ and $\tilde{y} = \infty_{\HN}$ for $y \in \{q_{34}, q_{35}\}$, defines a $36$-state finite $S$-algebra $\mathcal{A}_{36}$ with $\sem{A A A}_{36} = q_{34} \neq q_{35} = \sem{A A A S}_{36}$ (while no $S$-algebra of size at most~$5$ separates $A A A$ from $A A A S$).
\end{enumerate}
\end{theorem}

\begin{proof}
For~(i), observe first that right-compoundness is preserved under $\redS$: a root step $(S X Y) Z \to_r (X Z)(Y Z)$ creates compound right child $Y Z$, a left step leaves the right child unchanged, and a right step reduces an application to an application. Thus every reduct $M'$ of $M$ has $\tau(M') = 0$ and $\tau(M' S^k) = k$. Next, we show by induction on $k$ that if $M \notin \HN$ and $M S^k \redSs W$, then $W \equiv M' S^k$ for some $M'$ with $M \redSs M'$. For $k + 1$, along any single step $(M_1 S^k) S \redS V$ (with $M \redSs M_1$), the right child $S$ is irreducible, and a root step would require $M_1 S^k \equiv S X Y$, which for $k \in \{0, 1, 2\}$ gives $\ls(M_1) \le 2$ (contradicting $M \notin \HN$) and for $k \ge 3$ equates an application $M_1 S^{k-2}$ with $S$; hence the step occurs inside $M_1 S^k$, and the outer hypothesis yields $V \equiv M_2 S^{k+1}$. Now if $M S^j \convS N S^k$, Church--Rosser gives a common reduct $W$ with $M' S^j \equiv W \equiv N' S^k$ where $M \redSs M'$ and $N \redSs N'$. Applying $\tau$ gives $j = \tau(M' S^j) = \tau(N' S^k) = k$, and syntactic cancellation of $S^j$ gives $M' \equiv N'$, whence $M \convS N$. The converse is immediate by congruence.

For~(ii), every $A A A S^k$ has $\sem{A A A S^k}_{\HN} = \infty_{\HN}$, so at depth $D = 0$ all $A A A S^k$ have clocked B\"ohm tree $\bot$. Inducting on $D$, each $W_{D+1}(U) \equiv S\,W_D(U)$ is already in $\HNF$ with head-step clock~$0$ and sole child $W_D(U)$, preserving equality of all clocked B\"ohm approximants; meanwhile, Church--Rosser inversion on $S A_1 \convS S A_2$ implies $A_1 \convS A_2$, lifting $A A A S^j \neq_S A A A S^k$ to $W_D(A A A S^j) \neq_S W_D(A A A S^k)$.

For~(iii), let $\pi \colon Q_{36} \to Q_{\HN}$ merge $\{q_{34}, q_{35}\}$ into $\infty_{\HN}$. Whenever $\pi(((q_0 \odot_{36} x) \odot_{36} y) \odot_{36} z) \neq \infty_{\HN}$, \eqref{eq:s-algebra-law} holds because $\pi$ is a homomorphism. When $\pi$ maps both sides to $\infty_{\HN}$, $(q_0 \odot_{36} x) \odot_{36} y \in \{q_0, \dots, q_{33}\}$ forces the left-hand side of~\eqref{eq:s-algebra-law} to be $q_{34}$, while $y \odot_{36} z \neq q_0$ forces $(x \odot_{36} z) \odot_{36} (y \odot_{36} z) = q_{34}$ on the right. Direct evaluation gives $\sem{A A A}_{36} = q_{34} \neq q_{35} = q_{34} \odot_{36} q_0 = \sem{A A A S}_{36}$, and constraint propagation rules out separating $S$-algebras of size at most~$5$.
\end{proof}

\begin{theorem}[Productive B\"ohm-tree faithfulness and clocked-B\"ohm fine structure, Question~31 Part~II]\label{thm:productive-bohm-faithful}
Let $\BT_d(M)$ ($d \in \mathbb{N}$) denote the computable depth-$d$ B\"ohm-tree approximant of $M \in \TS$ (returning truncation at $d = 0$, $\bot$ if $\sem{M}_{\HN} = \infty_{\HN}$, and $S\,\BT_{d-1}(A_1)\cdots\BT_{d-1}(A_r)$ if $M \to_h^* S A_1 \dots A_r \in \HNF$), and let $\mathrm{CBT}_d(M)$ additionally label each non-$\bot$ node by its exact head-step count to $\HNF$.
\begin{enumerate}[label=(\roman*),leftmargin=2.2em]
  \item \emph{One-sided B\"ohm faithfulness across $\Prod \times \TS$:} For every productive closed $S$-term $M \in \Prod$ and \emph{every} closed $S$-term $N \in \TS$,
  \[
  M \convS N \;\Longleftrightarrow\; \forall d \in \mathbb{N},\; \BT_d(M) = \BT_d(N).
  \]
  \item \emph{Clocked-B\"ohm non-invariance versus non-$\SN$ collision:} On $\Prod \setminus \SN$, the convertible terms $Z_{\mathrm{Zac}} \equiv G_0 G_0$ and $G_1 G_1$ satisfy $\mathrm{CBT}_1(G_0 G_0) \neq \mathrm{CBT}_1(G_1 G_1)$ (root head clocks $2$ versus $1$). Conversely, setting $E \equiv (S(S T)(S T))T$, the $12$-leaf term $M_{12} \equiv S T (S(S T)(S T)) T$ and the $13$-leaf term $M_{13} \equiv A S E$ lie in $\Prod \setminus \SN$, satisfy $M_{12} \not\to_S^* M_{13}$ and $M_{13} \not\to_S^* M_{12}$, yet share $\mathrm{CBT}_d(M_{12}) = \mathrm{CBT}_d(M_{13})$ for all $d \in \mathbb{N}$ and satisfy $M_{12} \convS M_{13}$.
\end{enumerate}
\end{theorem}

\begin{proof}
For~(i), forward soundness (that $M \convS N$ implies $\BT_d(M) = \BT_d(N)$) holds by induction on $d$: conversion preserves $\sem{\cdot}_{\HN}$, and if $M \to_h^* H_1 \in \HNF$ and $N \to_h^* H_2 \in \HNF$ with $H_1 \convS H_2$, Church--Rosser aligns the spine arities of $H_1$ and $H_2$ and gives componentwise $\convS$ between their spine arguments.

Conversely, suppose $M \in \Prod$ and $\BT_d(M) = \BT_d(N)$ for all $d \in \mathbb{N}$. A closed term lies in $\Prod$ if and only if $\BT_d$ contains no $\bot$ for every $d$; hence $N \in \Prod$ as well. By Theorem~\ref{thm:q5-main} and Lemma~\ref{lem:pval-normalization}, $M \redSs U \in \mathcal{V}_P$ and $N \redSs V \in \mathcal{V}_P$, and soundness gives $\BT_d(U) = \BT_d(V)$ for all $d$. First, observe \emph{finite-depth normal-form rigidity}: for any closed normal form $W \in \mathcal{N}$ and any $X \in \TS$, the equality $\BT_{|W|}(X) = \BT_{|W|}(W)$ implies $X \redSs W$ (by induction on $|W|$ across $S$, $S W_1$, and $S W_1 W_2$ using downward monotonicity of $\BT_d$ equality). Moreover, for $A_0 \in \mathcal{V}_P$ and $W \in \mathcal{N}$, one has $A_0 \convS W$ if and only if $A_0 \equiv W$: if $A_0$ contains no $G_m G_n$ constructor, then $A_0 \in \mathcal{N}$ and confluence forces $A_0 \equiv W$; if $A_0$ contains some $G_m G_n$, two-sided sink absorption in $\mathcal{A}_{\SN}$ gives $\sem{A_0}_{\SN} = \infty_{\SN} \neq \sem{W}_{\SN}$, ruling out $A_0 \convS W$.

Now normalize root $G_0 G_n$ values via the size-preserving head step $\nu(G_0 G_n) = G_{n+1} G_{n+1}$ ($\nu(U) = U$ otherwise, keeping $\mathrm{sz}(\nu(U)) = \mathrm{sz}(U)$ and $U \redSs \nu(U)$). We construct a common reduct $K \in \TS$ with $U \redSs K$ and $V \redSs K$ by induction on $\mathrm{sz}(U) + \mathrm{sz}(V)$ whenever $\BT_d(U) = \BT_d(V)$ holds for all $d \in \mathbb{N}$:
\begin{enumerate}[label=(\alph*),leftmargin=2.2em]
  \item If $\nu(U)$ and $\nu(V)$ are both $S$, both $S X_i$, or both $S A_i B_i$ (mismatching root arities are excluded at $d = 1$), $\BT_{d+1}$ projects equality to the corresponding children of smaller structural size, and the induction hypothesis joins them componentwise.
  \item In the cross-case $\nu(U) \equiv S A_1 B_1$ and $\nu(V) \equiv G_{k+1} G_n$ (or symmetrically), recall $G_{k+1} G_n \to_h S G_n (G_k G_n)$ where $G_n \in \mathcal{N}$. Comparing $\BT_{d+1}$ gives $\BT_{|G_n|}(A_1) = \BT_{|G_n|}(G_n)$ and $\BT_d(B_1) = \BT_d(G_k G_n)$ for all $d \in \mathbb{N}$. By normal-form rigidity, $A_1 \redSs G_n$ and hence $A_1 \equiv G_n$. Since $\mathrm{sz}(B_1) + \mathrm{sz}(G_k G_n) = \mathrm{sz}(B_1) + 1 < (\mathrm{sz}(A_1) + \mathrm{sz}(B_1) + 1) + 1 = \mathrm{sz}(U) + \mathrm{sz}(V)$, the induction hypothesis constructs $K_B$ with $B_1 \redSs K_B$ and $G_k G_n \redSs K_B$, yielding the common reduct $S G_n K_B$.
  \item When $\nu(U) \equiv G_{k_1+1} G_{n_1}$ and $\nu(V) \equiv G_{k_2+1} G_{n_2}$, comparing left children at depth $|G_{n_2}| + 1$ gives $G_{n_1} \redSs G_{n_2} \in \mathcal{N}$, so $G_{n_1} \equiv G_{n_2}$ and $n_1 = n_2 = n$ (since $|G_n| = 2n + 5$). Induction on $\min(k_1, k_2)$ then forces $k_1 = k_2$: if $k_1 = 0 < k_2$, unfolding two levels exposes $G_0 G_n \to_h G_{n+1} G_{n+1} \to_h S G_{n+1} (G_n G_{n+1})$ against $G_{k_2} G_n \to_h S G_n (G_{k_2-1} G_n)$, forcing $G_{n+1} \equiv G_n$, impossible. Thus $k_1 = k_2$ and $\nu(U) \equiv \nu(V)$ is already the common reduct.
\end{enumerate}
Hence $U \convS V$ and $M \convS N$.

For~(ii), $G_0 G_0 \to_h G_1 G_1 \to_h S G_1 (G_0 G_1) \in \HNF$ gives root clock~$2$ for $G_0 G_0$ versus~$1$ for $G_1 G_1$. Meanwhile, $M_{12} \equiv S T (S(S T)(S T)) T$ and $M_{13} \equiv (S S S) S E$ both step in a single head contraction to the common reduct $Q_{\mathrm{mid}} \equiv (T T) E$, so every node of $\mathrm{CBT}_d(M_{12})$ and $\mathrm{CBT}_d(M_{13})$ has identical clock $1 + \mathrm{clk}(Q_{\mathrm{mid}})$ at the root and identical subtrees below. Yet $|M_{13}| = 13 > 12 = |M_{12}|$ forbids $M_{13} \redSs M_{12}$, while the unique step from $M_{12}$ is $M_{12} \redS Q_{\mathrm{mid}}$ ($13$ leaves, $\not\equiv M_{13}$), and every redex inside $Q_{\mathrm{mid}}$ has a compound third argument ($T$ or $S T$), strictly increasing leaf count to at least $14$ and forbidding $M_{12} \redSs M_{13}$.
\end{proof}

\begin{proposition}[Finiteness of normalizing classes, isolated normal forms, and $\Prod \times \TS$ decidability, Question~30]\label{prop:q30-decidability}
Let $\mathrm{gk} \colon \TSV \to \TSV$ denote Takahashi's complete parallel development~\cite{Takahashi1995} (the Gross--Knuth operator $\Delta$ of Proposition~\ref{prop:cr-conservation}(i)).
\begin{enumerate}[label=(\roman*),leftmargin=2.2em]
  \item \emph{Cofinality of $\mathrm{gk}$ and exact leaf-growth dichotomy:} For all $M, N \in \TSV$, one has $M \convS N$ if and only if there exists $m \in \mathbb{N}$ such that $N \redSs \mathrm{gk}^m(M)$, and $M \redSs N$ is constructively decidable. On every root contraction, $|(X Z)(Y Z)| = |S X Y Z| + (|Z| - 1)$; hence on closed terms $M \redS M'$ preserves leaf count ($|M| = |M'|$) if and only if the contracted redex has atomic third argument $S$, and otherwise $|M'| \ge |M| + 1$.
  \item \emph{Strict finiteness of closed normalizing $\convS$-classes and isolated normal forms:} For every $V \in \mathcal{N}$, $[V]_{\convS} \cap \TS = \{ N \in \TS : N \redSs V \}$ with $|N| \le |V|$, so for every closed $M \in \SN$, its closed conversion class $[M]_{\convS} \cap \TS$ is strictly finite and constructively enumerable without duplicates. A normal form $V \in \NF$ is isolated ($[V]_{\convS} = \{V\}$) if and only if $V$ contains no subterm of the form $(X Z)(Y Z)$ with identical right subterms $Z$, witnessed by the infinite right-comb family $R_n$ ($n \ge 1$, $|R_n| = n$) and the spine-$2$ family $S S R_{k+3}$ ($k \in \mathbb{N}$, $|S S R_{k+3}| = k + 5$).
  \item \emph{RNF decomposition, total $\Prod \times \TS$ decidability, and failure of right-cancellation:} For all $A_1 B_1, A_2 B_2 \in \RNF$, one has $A_1 B_1 \convS A_2 B_2$ if and only if both $A_1 \convS A_2$ and $B_1 \convS B_2$. Across all pairs $(M, N) \in \Prod \times \TS$ (strictly extending $\SN \times \TS$ since $\SN \subsetneq \Prod$), convertibility $M \convS N$ is constructively decidable, and whenever $M \convS N$ holds, an explicit common reduct $W \in \TS$ satisfying $M \redSs W$ and $N \redSs W$ is constructible; by contrast, right-cancellation fails even on $\mathcal{N}$, witnessed by $N_1 \equiv S(S A S)S \neq_S S A (S A S) \equiv N_2$ with $N_1 S \convS N_2 S$, and universally by Waldmann's pair $U_1 \equiv S(S T S)A \neq_S S(T A)(S A S) \equiv U_2$ satisfying $U_1 Z \convS U_2 Z$ for every $Z \in \TSV$.
\end{enumerate}
\end{proposition}

\begin{proof}
For~(i), iterating Takahashi's triangle property (that $M \Rightarrow N$ implies $N \Rightarrow \mathrm{gk}(M)$) yields equivalence between $M \convS N$ and the existence of $m \in \mathbb{N}$ with $N \redSs \mathrm{gk}^m(M)$, while decidability of $\redSs$ holds by Proposition~\ref{prop:acyclicity-order}(i). Expanding $|(X Z)(Y Z)| = |X| + |Y| + 2|Z| = (1 + |X| + |Y| + |Z|) + (|Z| - 1)$ shows that leaf count is preserved if and only if $|Z| = 1$, which on closed terms means $Z \equiv S$.

For~(ii), if $V \in \mathcal{N}$ and $N \convS V$, Church--Rosser yields a common reduct $W$ with $N \redSs W$ and $V \redSs W$; irreducibility of $V$ gives $W \equiv V$, so $N \redSs V$ and $|N| \le |V|$ by Proposition~\ref{prop:cr-conservation}(ii). Filtering the finite duplicate-free list of closed $S$-terms of leaf count at most $|\mathrm{norm}(M)|$ by decidable reachability $\redSs \mathrm{norm}(M)$ enumerates $[M]_{\convS} \cap \TS$. Moreover, a term $V$ admits a one-step predecessor $U \redS V$ if and only if $V$ contains $(X Z)(Y Z)$ (contracted from $S X Y Z$ at that subterm); for $V \in \NF$, $[V]_{\convS} = \{N : N \redSs V\}$ equals $\{V\}$ if and only if $V$ has no one-step predecessor. In $R_{n+1} \equiv S R_n$, the left child of every application is $S$ (not an application), so no $(X Z)(Y Z)$ occurs; in $S S R_{k+3}$, the children $(S S)$ and $R_{k+3} \equiv S R_{k+2}$ have distinct right subterms $S \not\equiv R_{k+2}$ ($|R_{k+2}| = k + 2 \ge 2$), so $S S R_{k+3}$ is isolated for every $k \in \mathbb{N}$.

For~(iii), RNF decomposition and the counterexamples $(N_1, N_2)$ and $(U_1, U_2)$ follow by Church--Rosser (since root contractions cannot occur from $\RNF$, while $\mathrm{gk}^6(N_1 S) \equiv \mathrm{gk}^6(N_2 S)$ and $U_1 Z \redS^5 (P_Z(P_Z P_Z))(Z P_Z(P_Z P_Z)) \leftarrow_S^6 U_2 Z$ for $P_Z \equiv S Z$). Finally, given $M \in \Prod$ and arbitrary $N \in \TS$, test $\sem{N}_P \neq \infty_P$: if $\sem{N}_P = \infty_P$, conversion invariance of $\mathcal{A}_P$ refutes $M \convS N$; if $\sem{N}_P \neq \infty_P$, normalize $M \redSs U \in \mathcal{V}_P$ and $N \redSs V \in \mathcal{V}_P$ via Lemma~\ref{lem:pval-normalization}, and run the structural recursion of Theorem~\ref{thm:productive-bohm-faithful}(i) bounded by $\mathrm{sz}(U) + \mathrm{sz}(V)$, which decides $U \convS V$ (hence $M \convS N$) and returns an explicit common reduct $W$ whenever $M \convS N$.
\end{proof}

Whether convertibility $\convS$ remains decidable on general pairs of non-productive closed pure-$S$ terms $(\TS \setminus \Prod)^2$ (Scott's Question~30 in full generality) remains open.

\subsection{Functional Order \texorpdfstring{$\{0, 1, 2, 3\}$}{\{0, 1, 2, 3\}} Hierarchy and Minimal \texorpdfstring{$S$}{S}-Algebras}\label{subsec:order-hierarchy}

In Section~5.5 and Appendix~A.3 of Barendregt, Endrullis, Klop, and Waldmann~\cite[pp.~89--91, 99--101]{BarendregtEndrullisKlopWaldmann2017}, pure-$S$ combinators are stratified by their \emph{functional order} (in the sense of Curry and Feys~\cite{CurryFeys1958}), which measures how many fresh arguments a term can actively consume before all subsequent arguments play a purely dummy role. By Proposition~\ref{prop:acyclicity-order}(iii) (due to Visser and van~Oostrom), every term in $\CLS$ has order at most~$3$. Here we show that extending $\mathcal{A}_{\HN}$ to open terms by mapping variables to $\infty_{\HN}$ characterizes exact functional order across all open and closed $S$-terms, proves that closed order-$3$ terms are uniquely right combs, and determines the exact minimal $S$-algebras for every order class.

\begin{definition}[Open solvability evaluation and Curry--Feys functional order]\label{def:open-hn-order}
\begin{enumerate}[label=(\roman*),leftmargin=2.2em]
  \item Define the open evaluation homomorphism $\sem{\cdot}_{\HN}^{\mathrm{open}} \colon \TSV \to Q_{\HN}$ by $\sem{S}_{\HN}^{\mathrm{open}} = q_0$, $\sem{x}_{\HN}^{\mathrm{open}} = \infty_{\HN}$ for every $x \in \mathcal{V}$, and $\sem{M N}_{\HN}^{\mathrm{open}} = \sem{M}_{\HN}^{\mathrm{open}} \odot_{\HN} \sem{N}_{\HN}^{\mathrm{open}}$ (agreeing with $\sem{M}_{\HN}$ on $\TS$). Call $U \in \TSV$ an \emph{$S$-headed head normal form} if $U \equiv S$, $U \equiv S A_1$, or $U \equiv S A_1 A_2$.
  \item Following~\cite[Definition~37, p.~99]{BarendregtEndrullisKlopWaldmann2017}, a term $F \in \TSV$ has \emph{dummy variable tails} if for every nonempty list of variables $y_1, \dots, y_r$ ($r \ge 1$) not occurring in $F$ and every reduction $F y_1 \cdots y_r \redSs N$, there exists $F'$ with $F \redSs F'$, no $y_i$ occurring in $F'$, and $N \equiv F' y_1 \cdots y_r$. A term $M \in \TSV$ has \emph{order at most $k$} ($\ord(M) \le k$) if $M x_1 \cdots x_k$ has dummy variable tails for $k$ distinct fresh variables $x_1, \dots, x_k$ not occurring in $M$, and has \emph{exact order $k$} ($\ord(M) = k$) if $\ord(M) \le k$ and $\ord(M) \not\le j$ for all $j < k$.
\end{enumerate}
\end{definition}

\begin{lemma}[Universal open solvability characterization]\label{lem:open-solvability}
For every open or closed $S$-term $M \in \TSV$, the term $M$ reduces under $\redSs$ to an $S$-headed head normal form if and only if $\sem{M}_{\HN}^{\mathrm{open}} \neq \infty_{\HN}$. Consequently, $F \in \TSV$ has dummy variable tails if and only if $\sem{F}_{\HN}^{\mathrm{open}} = \infty_{\HN}$.
\end{lemma}

\begin{proof}
Because $\mathcal{A}_{\HN}$ satisfies~\eqref{eq:s-algebra-law}, $\sem{\cdot}_{\HN}^{\mathrm{open}}$ is invariant under $\convS$, and $q_0$, $q_0 \odot_{\HN} a$, and $(q_0 \odot_{\HN} a) \odot_{\HN} b$ never equal $\infty_{\HN}$; thus reducing to an $S$-headed $\HNF$ implies $\sem{M}_{\HN}^{\mathrm{open}} \neq \infty_{\HN}$. Conversely, suppose $\sem{M}_{\HN}^{\mathrm{open}} \neq \infty_{\HN}$, and fix the $8$-leaf closed unsolvable term $W_{34} \equiv S S S (S A) S \in \TS$ satisfying $\sem{W_{34}}_{\HN} = \infty_{\HN}$. Let $\sigma_{34} \colon \TSV \to \TS$ substitute $W_{34}$ for every variable. Structural induction gives $\sem{\sigma_{34}(P)}_{\HN} = \sem{P}_{\HN}^{\mathrm{open}}$ for all $P \in \TSV$, and every open head step $P \to_h P'$ induces a closed head step $\sigma_{34}(P) \to_h \sigma_{34}(P')$, decreasing $\mu_{\HN}(\sigma_{34}(P))$ by at least~$1$ via Lemma~\ref{lem:hn-potential}. Hence open head reduction from $M$ terminates within $\mu_{\HN}(\sigma_{34}(M))$ steps at a head-irreducible term $U^*$ with $\sem{U^*}_{\HN}^{\mathrm{open}} = \sem{M}_{\HN}^{\mathrm{open}} \neq \infty_{\HN}$. Any head-irreducible open term either has a variable leftmost leaf (forcing $\sem{U^*}_{\HN}^{\mathrm{open}} = \infty_{\HN}$ by left absorption $\infty_{\HN} \odot_{\HN} y = \infty_{\HN}$, a contradiction) or is an $S$-headed $\HNF$.

For the second equivalence, if $\sem{F}_{\HN}^{\mathrm{open}} \neq \infty_{\HN}$, then $F \redSs U$ where $U \equiv S$, $S A_1$, or $S A_1 A_2$. Appending $r = 3$, $2$, or $1$ fresh variables $y_1, \dots, y_r$ yields $F y_1 \cdots y_r \redSs (P y_r)(Q y_r)$ whose right child $Q y_r$ has leaf count $|Q| + 1 \ge 2$, which cannot equal $F' y_1 \cdots y_r$ (whose right child $y_r$ has leaf count~$1$). Conversely, if $\sem{F}_{\HN}^{\mathrm{open}} = \infty_{\HN}$, then neither $F$ nor any $F y_1 \cdots y_j$ ($\sem{F y_1 \cdots y_j}_{\HN}^{\mathrm{open}} = \infty_{\HN}$) can reduce to any $S P Q$, so no root contraction ever consumes any $y_i$, and induction on $r$ yields $N \equiv F' y_1 \cdots y_r$ with $F \redSs F'$ and $y_i$ absent from $F'$.
\end{proof}

\begin{theorem}[Order $\{0, 1, 2, 3\}$ classification, right-comb uniqueness, and minimal $S$-algebras]\label{thm:order-hierarchy}
Define the right-translation map $\tau_{\HN}(q) := q \odot_{\HN} \infty_{\HN}$ on $Q_{\HN}$.
\begin{enumerate}[label=(\roman*),leftmargin=2.2em]
  \item \emph{Nilpotency and $35 = 1 + 11 + 17 + 6$ state partition:} One has $\tau_{\HN}^3(q) = \infty_{\HN}$ for every $q \in Q_{\HN}$, partitioning $Q_{\HN}$ by nilpotency index $\mathrm{ord}_{\HN}(q) := \min\{k \in \mathbb{N} : \tau_{\HN}^k(q) = \infty_{\HN}\} \in \{0, 1, 2, 3\}$ into $1$ state of order~$0$ ($\{\infty_{\HN}\}$), $11$ states of order~$1$ ($\{q_{22}, \dots, q_{26}, q_{28}, \dots, q_{33}\}$), $17$ states of order~$2$ ($\{q_6, \dots, q_{21}, q_{27}\}$), and $6$ states of order~$3$ ($\{q_0, \dots, q_5\}$). For every open or closed $M \in \TSV$, one has $\ord(M) = \mathrm{ord}_{\HN}(\sem{M}_{\HN}^{\mathrm{open}}) \in \{0, 1, 2, 3\}$, which is $\convS$-invariant and agrees on closed $M \in \TS$ with single-variable consumption~\cite[Definition~21]{BarendregtEndrullisKlopWaldmann2017}:
  \begin{align*}
  \ord(M) = 0 &\;\Longleftrightarrow\; M \notin \HN, \\
  \ord(M) = 1 &\;\Longleftrightarrow\; M \in \HN,\; M \not\equiv S,\;\textrm{and}\; \neg\exists A_1, A_2,\; M x_0 \redSs S A_1 A_2, \\
  \ord(M) = 2 &\;\Longleftrightarrow\; M \not\equiv S,\; \bigl(\exists A_1, A_2,\; M x_0 \redSs S A_1 A_2\bigr),\;\textrm{and}\; \neg\exists A_1, A_2,\; M x_0 x_1 \redSs S A_1 A_2, \\
  \ord(M) = 3 &\;\Longleftrightarrow\; \exists A_1, A_2,\; M x_0 x_1 \redSs S A_1 A_2.
  \end{align*}
  \item \emph{Uniqueness of closed order-$3$ terms and sharp small-size thresholds:} A closed $S$-term $M \in \TS$ satisfies $\ord(M) = 3$ if and only if $M \equiv R_{|M|}$ is the unique right comb of leaf size $|M|$ (Definition~\ref{def:syntax}). Every closed $S$-term of leaf size at most~$2$ has order~$3$; at size~$3$, the term $A \equiv S S S$ is the unique closed term of order~$2$; every closed term of size at most~$4$ has order at least~$2$; at size~$5$, the normal form $S(S S S)S$ is the unique closed term of order~$1$; and order~$0$ first appears at size~$8$ ($15$ closed terms).
  \item \emph{Exact minimal $S$-algebras of sizes $35$, $12$, and $3$:} On $\TS$, any deterministic bottom-up tree automaton deciding $\{\ord = 0\}$, $\{\ord = 1\}$, $\{\ord \ge 1\}$, or $4$-valued $\ord$ has at least $35$ states, matching $\mathcal{A}_{\HN}$; $\{\ord = 2\}$ and $\{\ord \ge 2\}$ are both decided by a $12$-state quotient $S$-algebra $\mathcal{A}_{\mathrm{ord},2}$ of $\mathcal{A}_{\HN}$, and any deterministic bottom-up tree automaton deciding either set has at least $12$ states; and $\{\ord = 3\}$ is decided by a $3$-state quotient $S$-algebra $\mathcal{A}_{\mathrm{ord},3}$, with matching universal lower bound~$3$.
\end{enumerate}
\end{theorem}

\begin{proof}
For~(i), direct evaluation on $Q_{\HN}$ verifies $\tau_{\HN}^3(q) = \infty_{\HN}$ and the $1 + 11 + 17 + 6 = 35$ partition. For any $M \in \TSV$ and fresh variables $x_1, \dots, x_k$, we have $\sem{M x_1 \cdots x_k}_{\HN}^{\mathrm{open}} = \tau_{\HN}^k(\sem{M}_{\HN}^{\mathrm{open}})$, so Lemma~\ref{lem:open-solvability} shows that $\ord(M) \le k$ holds if and only if $\tau_{\HN}^k(\sem{M}_{\HN}^{\mathrm{open}}) = \infty_{\HN}$, which holds if and only if $\mathrm{ord}_{\HN}(\sem{M}_{\HN}^{\mathrm{open}}) \le k$, proving $\ord(M) = \mathrm{ord}_{\HN}(\sem{M}_{\HN}^{\mathrm{open}})$. On closed $M \in \TS$ with $M \not\equiv S$ (so $M$ is an application), Lemma~\ref{lem:visser-fibonacci}(i) excludes $M x_0 \redSs S$ and $M x_0 \redSs S W$ (and likewise for $M x_0 x_1$); hence Lemma~\ref{lem:open-solvability} equates the existence of $A_1, A_2$ with $M x_0 \redSs S A_1 A_2$ to $\tau_{\HN}(\sem{M}_{\HN}) \neq \infty_{\HN}$, and equates the existence of $A_1, A_2$ with $M x_0 x_1 \redSs S A_1 A_2$ to $\tau_{\HN}^2(\sem{M}_{\HN}) \neq \infty_{\HN}$, yielding the four single-variable clauses.

For~(ii), inspecting all $35^2$ pairs in $\mathcal{A}_{\HN}$ confirms that $\mathrm{ord}_{\HN}(a \odot_{\HN} b) = 3$ holds if and only if $a = q_0$ and $\mathrm{ord}_{\HN}(b) = 3$. Since $\sem{U}_{\HN} = q_0$ holds for $U \in \TS$ if and only if $U \equiv S$, structural induction on $M \in \TS$ gives $\ord(M) = 3$ if and only if $M \equiv R_{|M|}$. Enumeration of $\TS$ for $|M| \le 5$ ($1, 1, 2, 5, 14$ Catalan trees) confirms that $|M| \le 2$ gives $\{S, S S\}$ (order~$3$), $|M| = 3$ gives $S(S S)$ (order~$3$) and $S S S$ (order~$2$), $|M| = 4$ contains only orders~$2$ and~$3$, and among the $14$ closed terms of size~$5$, the term $S(S S S)S$ ($\sem{S(S S S)S}_{\HN} = q_{26}$) is the sole term of order~$1$.

For~(iii), define the $12$-state quotient $\pi_2 \colon Q_{\HN} \to Q_{\mathrm{ord},2} = \{o_0, \dots, o_{11}\}$ merging $\{q_4, q_5\} \mapsto o_{11}$, $\{q_6, q_{11}, q_{13}\} \mapsto o_2$, $\{q_7, q_{12}\} \mapsto o_4$, $\{q_8, q_9, q_{10}, q_{15}, q_{16}, q_{17}, q_{19}, q_{20}, q_{21}, q_{27}\} \mapsto o_5$, $\{q_{26}, q_{31}\} \mapsto o_9$, and $\{q_{22}, q_{23}, q_{24}, q_{25}, q_{28}, q_{29}, q_{30}, q_{32}, q_{33}, \infty_{\HN}\} \mapsto o_6$ (keeping $q_0 \mapsto o_0$, $q_1 \mapsto o_1$, $q_2 \mapsto o_3$, $q_3 \mapsto o_8$, $q_{14} \mapsto o_7$, $q_{18} \mapsto o_{10}$), and define the $3$-state quotient $\pi_3 \colon Q_{\HN} \to \{r_0, r_{\ge 1}, r_\bot\}$ mapping $q_0 \mapsto r_0$, $\{q_1, \dots, q_5\} \mapsto r_{\ge 1}$, and $\{q_6, \dots, \infty_{\HN}\} \mapsto r_\bot$ (with $r_0 \odot r_0 = r_0 \odot r_{\ge 1} = r_{\ge 1}$ and $x \odot y = r_\bot$ otherwise). Finite evaluation over $Q_{\HN}^2$ confirms that $\pi_2$ and $\pi_3$ are $S$-algebra homomorphisms deciding $\{\ord = 2\} = \{o_2, o_4, o_5, o_7, o_{10}\}$, $\{\ord \ge 2\} = Q_{\mathrm{ord},2} \setminus \{o_6, o_9\}$, and $\{\ord = 3\} = \{r_0, r_{\ge 1}\}$. Conversely, explicit closed witness terms inhabit all $35$ states of $Q_{\HN}$, all $12$ states of $Q_{\mathrm{ord},2}$, and all $3$ states of $Q_{\mathrm{ord},3}$ ($S, S S, S S S$); finite context families ($18$ closed contexts for $\{\ord = 0\}$, $13$ for $\{\ord = 1\}$, $6$ for $\{\ord = 2\}$, $8$ for $\{\ord \ge 2\}$, and $\{\square, \square S\}$ for $\{\ord = 3\}$) separate every distinct pair of states, forcing $|Q| \ge 35$, $|Q| \ge 12$, and $|Q| \ge 3$ on any recognizing deterministic bottom-up tree automaton by context congruence and the pigeonhole principle.
\end{proof}

\begin{remark}[Minimal order-$2$ and order-$1$ witnesses in Barendregt et al.\ {\cite{BarendregtEndrullisKlopWaldmann2017}}]\label{rem:order-starlings-p90}
In the informal overview of Section~5.5 in~\cite[p.~90]{BarendregtEndrullisKlopWaldmann2017}, $S S S$ is momentarily cited as having order~$1$, whereas in Appendix~A.3, Example~41(ii)~\cite[p.~102]{BarendregtEndrullisKlopWaldmann2017} of the same work, the two-step reduction $S S S x y \redS (S x)(S x) y \redS (x y)(S x y)$ reaches a variable-headed term of order~$0$ and assigns $S S S$ order~$2$. Theorem~\ref{thm:order-hierarchy}(ii) reconciles the two passages: $A \equiv S S S$ ($\sem{A}_{\HN} = q_6$) is the unique minimal closed $S$-term of order~$2$ ($3$ leaves), whereas no closed $S$-term of at most~$4$ leaves has order~$1$, the unique minimal closed order-$1$ combinator being the $5$-leaf normal form $S(S S S)S$ ($\sem{S(S S S)S}_{\HN} = q_{26}$, for which $S(S S S)S \in \HNF$ while $S(S S S)S x$ is already in normal form and cannot reduce to any $S A_1 A_2$).
\end{remark}

% ============================================================================
\section{Translation into the \texorpdfstring{$\lambda\mathrm{I}$}{lambda-I}-Calculus: Order Preservation, Normalization Divergences, Hereditary Saturation, and Equational Non-Axiomatizability}\label{sec:lambda-translation}
% ============================================================================

The classical correspondence between combinatory logic and the untyped $\lambda$-calculus embeds combinators homomorphically into $\lambda$-terms by mapping binary application to application, variables to themselves, and each primitive combinator to its defining $\lambda$-abstraction \cite{CurryFeys1958,Barendregt1984,HindleySeldin2008}. On the pure-$S$ subsystem $\CLS$, this canonical translation $(-)_\lambda$ assigns to $S$ the closed three-argument abstraction
\begin{equation}\label{eq:lambda-s-def}
  (S)_\lambda \;\equiv\; \lambda x y z.\, (x z)(y z).
\end{equation}
Because each of the three bound variables $x, y, z$ occurs free in the body $(x z)(y z)$, the translation $(M)_\lambda$ of every open or closed pure-$S$ term $M \in \TSV$ belongs to Alonzo Church's $\lambda\mathrm{I}$-calculus \cite{Church1941}, in which $\beta$-contraction never erases an argument and weak normalization coincides with strong normalization. Even within this non-erasing fragment, however, combinatory reduction and $\lambda\beta$-reduction diverge in subtle ways: whereas any two-argument pure-$S$ term $S P Q$ with $P, Q \in \NF$ is already a normal form of $\CLS$, its translation $(S P Q)_\lambda$ reduces in two weak-head $\beta$-steps to $\lambda z.\, ((P)_\lambda z)((Q)_\lambda z)$, exposing the application $(P z)(Q z)$ to full $\beta$-reduction underneath the $\lambda z$-binder. In Chapter~5 of \cite{BarendregtEndrullisKlopWaldmann2017}, Barendregt, Endrullis, Klop, and Waldmann examined this discrepancy in detail, posing Question~32 on equational axiomatizability and Question~34 on the $\lambda\beta$-normalizing closed pure-$S$ terms alongside Conjectures~40, 42, and~43 \cite[pp.~94--96, 98--102]{BarendregtEndrullisKlopWaldmann2017}. In Subsection~\ref{subsec:lambdaI-conservation-conj40} we establish constructive weak-head standardization and conservation on $\lambda\mathrm{I}$ and resolve Conjecture~40 in the affirmative across full combinatory logic $\mathrm{CL}(\{S, K, I\})$ and $\CLS$. In Subsection~\ref{subsec:conj42-43-refutation} we determine the minimal closed pure-$S$ combinators whose $\lambda$-translations diverge, refute Conjectures~42 and~43 simultaneously at $14$ leaves ($15$ leaves in normal form), and prove an exact pure-$\CLS$ hereditary saturation characterization of $\lambda\beta$-normalization. The remaining subsections then analyze the Myhill--Nerode congruence geometry of Question~34 and establish the equational non-axiomatizability of $\lambda\beta$-equality over $\CLS$.

\subsection{Constructive \texorpdfstring{$\lambda\mathrm{I}$}{lambda-I} Standardization, Conservation, and Affirmative Resolution of Conjecture 40}\label{subsec:lambdaI-conservation-conj40}

Because Conjecture~40 of Barendregt, Endrullis, Klop, and Waldmann \cite[p.~102]{BarendregtEndrullisKlopWaldmann2017} concerns functional order across general combinatory logic as well as its specialization to $\CLS$, we formulate the combinatory translation and the two notions of functional order simultaneously over the full basis $\{S, K, I\}$ and over the pure-$S$ basis.

\begin{definition}[Untyped $\lambda$-terms, the $\lambda\mathrm{I}$-fragment, combinatory translations, and functional order]\label{def:lambda-translation-order}
Let $\mathcal{V} = \{x_0, x_1, x_2, \dots\}$ be the countably infinite set of variables fixed in Definition~\ref{def:syntax}.
\begin{enumerate}[label=(\roman*),leftmargin=2.2em]
  \item \emph{Untyped $\lambda$-terms and $\lambda\mathrm{I}$-calculus:} Let $\Lambda$ denote the set of untyped $\lambda$-terms over $\mathcal{V}$ modulo $\alpha$-equivalence, generated by $P, Q ::= x \mid (P\,Q) \mid (\lambda x.\, P)$ ($x \in \mathcal{V}$), with free-variable set $\mathrm{FV}(P) \subset \mathcal{V}$ and capture-avoiding substitution $P[x := Q]$ (represented concretely in our constructive formalization via locally nameless de Bruijn indices for bound variables and $\mathcal{V}$ for free variables). \emph{One-step $\beta$-reduction} $\to_\beta$ is the full contextual closure (under left application, right application, and $\lambda$-abstraction) of $(\lambda x.\, P)\,Q \to_\beta P[x := Q]$, with transitive closure $\to_\beta^+$, reflexive-transitive closure $\to_\beta^*$, and $\beta$-convertibility $=_\beta$. \emph{Weak-head $\beta$-reduction} $\to_{\mathrm{wh}}$ restricts contraction to the top of the left application spine outside $\lambda$-abstractions: $(\lambda x.\, P)\,Q \to_{\mathrm{wh}} P[x := Q]$, and $P Q \to_{\mathrm{wh}} P' Q$ whenever $P \to_{\mathrm{wh}} P'$. A term $P \in \Lambda$ is in \emph{$\beta$-normal form} ($P \in \NF_\beta$) if there is no $P'$ with $P \to_\beta P'$, \emph{weakly $\beta$-normalizing} ($P \in \WN_\beta$) if $P \to_\beta^* N$ for some $N \in \NF_\beta$, and \emph{strongly $\beta$-normalizing} ($P \in \SN_\beta$) if $P$ is constructively well-founded (accessible) under $\to_\beta$. Following Church \cite{Church1941}, the \emph{$\lambda\mathrm{I}$-fragment} $\Lambda_{\mathrm{I}} \subset \Lambda$ is generated inductively by $x \in \Lambda_{\mathrm{I}}$ ($x \in \mathcal{V}$), $P Q \in \Lambda_{\mathrm{I}}$ whenever $P, Q \in \Lambda_{\mathrm{I}}$, and $\lambda x.\, P \in \Lambda_{\mathrm{I}}$ whenever $P \in \Lambda_{\mathrm{I}}$ and $x \in \mathrm{FV}(P)$.
  \item \emph{Full combinatory logic and canonical translation:} Let $\mathrm{Ter}(\mathrm{CL}(\{S, K, I\}))$ denote the set of open combinatory terms generated by $\{S, K, I\} \cup \mathcal{V}$ under binary application, equipped with one-step combinatory reduction $\to_{\mathrm{CL}}$ generated contextually by
  \begin{equation}\label{eq:cl-ski-rules}
    S X Y Z \;\to_{\mathrm{CL}}\; (X Z)(Y Z), \qquad K X Y \;\to_{\mathrm{CL}}\; X, \qquad I X \;\to_{\mathrm{CL}}\; X,
  \end{equation}
  and reflexive-transitive closure $\to_{\mathrm{CL}}^*$. Note that $\TSV \subset \mathrm{Ter}(\mathrm{CL}(\{S, K, I\}))$ embeds conservatively: for all $M, N \in \TSV$, one has $M \to_{\mathrm{CL}}^* N$ if and only if $M \redSs N$. Define the sets of \emph{partial combinator applications} in $\mathrm{CL}(\{S, K, I\})$ and $\CLS$ by
  \begin{equation}\label{eq:partial-comb-sets}
  \begin{aligned}
    \mathcal{P}_{\mathrm{CL}} &\;:=\; \{I, K, S\} \;\cup\; \{K A, S A : A \in \mathrm{Ter}(\mathrm{CL}(\{S, K, I\}))\} \\
    &\qquad\qquad\qquad \;\cup\; \{S A B : A, B \in \mathrm{Ter}(\mathrm{CL}(\{S, K, I\}))\}, \\
    \mathcal{P}_S &\;:=\; \mathcal{P}_{\mathrm{CL}} \cap \TSV \;=\; \{S\} \;\cup\; \{S A : A \in \TSV\} \;\cup\; \{S A B : A, B \in \TSV\}.
  \end{aligned}
  \end{equation}
  The canonical homomorphic translation $(-)_\lambda \colon \mathrm{Ter}(\mathrm{CL}(\{S, K, I\})) \to \Lambda$ is defined recursively by $(x)_\lambda \equiv x$ for $x \in \mathcal{V}$, $(M N)_\lambda \equiv (M)_\lambda\,(N)_\lambda$, $(S)_\lambda \equiv \lambda x y z.\, (x z)(y z)$, $(K)_\lambda \equiv \lambda x y.\, x$, and $(I)_\lambda \equiv \lambda x.\, x$. Observe that $\mathrm{FV}((M)_\lambda) = \mathrm{FV}(M)$ for all $M \in \mathrm{Ter}(\mathrm{CL}(\{S, K, I\}))$, and $(M)_\lambda \in \Lambda_{\mathrm{I}}$ for every $M \in \TSV$.
  \item \emph{Functional order in $\Lambda$ and $\mathrm{CL}$ \cite[Definitions~35 and~37, pp.~98--99]{BarendregtEndrullisKlopWaldmann2017}:} For $P \in \Lambda$, the \emph{$\lambda$-order} $\ord_\lambda(P) \in \mathbb{N} \cup \{\infty\}$ is defined inductively on $n \in \mathbb{N}$ by setting $\ord_\lambda(P) = 0$ if there is no $Q \in \Lambda$ with $P \to_\beta^* \lambda x.\, Q$; $\ord_\lambda(P) = n + 1$ if $P \to_\beta^* \lambda x.\, Q$ for some $Q \in \Lambda$ with $\ord_\lambda(Q) = n$; and $\ord_\lambda(P) = \infty$ if $\ord_\lambda(P) \ne n$ for every $n \in \mathbb{N}$. Following Curry and Feys \cite{CurryFeys1958} and \cite[Definition~37]{BarendregtEndrullisKlopWaldmann2017}, a combinatory term $M \in \mathrm{Ter}(\mathrm{CL}(\{S, K, I\}))$ \emph{never consumes fresh arguments} if for every $r \ge 1$, every tuple of distinct variables $y_1, \dots, y_r \in \mathcal{V} \setminus \mathrm{FV}(M)$, and every reduction $M y_1 \cdots y_r \to_{\mathrm{CL}}^* N$, there exists $M' \in \mathrm{Ter}(\mathrm{CL}(\{S, K, I\}))$ such that $N \equiv M' y_1 \cdots y_r$ and $M \to_{\mathrm{CL}}^* M'$. The \emph{combinatory order} $\ord_{\mathrm{CL}}(M) \in \mathbb{N} \cup \{\infty\}$ is then defined inductively on $n \in \mathbb{N}$ by setting $\ord_{\mathrm{CL}}(M) = 0$ if $M$ never consumes fresh arguments; $\ord_{\mathrm{CL}}(M) = n + 1$ if $M$ does not have order $0$ and $\ord_{\mathrm{CL}}(M z) = n$ for a fresh variable $z \in \mathcal{V} \setminus \mathrm{FV}(M)$; and $\ord_{\mathrm{CL}}(M) = \infty$ if $\ord_{\mathrm{CL}}(M) \ne n$ for every $n \in \mathbb{N}$.
\end{enumerate}
\end{definition}

Every combinatory root contraction in \eqref{eq:cl-ski-rules} is simulated in $\Lambda$ by weak-head $\beta$-steps: $((S)_\lambda X Y) Z \to_{\mathrm{wh}}^3 (X Z)(Y Z)$, $((K)_\lambda X) Y \to_{\mathrm{wh}}^2 X$, and $(I)_\lambda X \to_{\mathrm{wh}} X$. Lifting across application contexts shows that $M \to_{\mathrm{CL}} N$ implies $(M)_\lambda \to_\beta^+ (N)_\lambda$, and consequently $M \to_{\mathrm{CL}}^* N$ implies $(M)_\lambda \to_\beta^* (N)_\lambda$. Likewise, Takahashi's parallel $\beta$-reduction $P \Rightarrow_\beta Q$ and complete development map $\Delta_\beta$ \cite{Takahashi1995} establish the Church--Rosser property for $\to_\beta^*$ constructively: co-initial reductions $P \to_\beta^* Q_1$ and $P \to_\beta^* Q_2$ always admit a common reduct $R$ with $Q_1 \to_\beta^* R$ and $Q_2 \to_\beta^* R$. In particular, $\beta$-normal forms are unique, and $\ord_\lambda(P)$ is invariant under $=_\beta$ and uniquely determined whenever finite: if $P \to_\beta^* \lambda x.\, Q_1$ and $P \to_\beta^* \lambda x.\, Q_2$, confluence yields a common reduct $\lambda x.\, R$ with $Q_1 \to_\beta^* R$ and $Q_2 \to_\beta^* R$, whence $Q_1 =_\beta Q_2$.

To pull $\lambda$-abstractions and normalization properties back across $(-)_\lambda$ constructively without invoking classical choice or K\"onig's lemma, we next prove weak-head standardization for $\Lambda$ together with Church's conservation theorem on $\Lambda_{\mathrm{I}}$.

\begin{lemma}[Weak-head standardization and Church's conservation theorem on $\Lambda_{\mathrm{I}}$]\label{lem:lambdaI-conservation}
Let $P, Q \in \Lambda$ be untyped $\lambda$-terms.
\begin{enumerate}[label=(\roman*),leftmargin=2.2em]
  \item \emph{Weak-head standardization:} If $P \to_\beta^* \lambda x.\, Q$, there exists $B \in \Lambda$ such that $P \to_{\mathrm{wh}}^* \lambda x.\, B$ and $B \to_\beta^* Q$. Likewise, if $P \to_\beta^* y R_1 \cdots R_m$ for a variable $y$ (bound or free) and $m \ge 0$, there exist $R_1^0, \dots, R_m^0 \in \Lambda$ such that $P \to_{\mathrm{wh}}^* y R_1^0 \cdots R_m^0$ and $R_i^0 \to_\beta^* R_i$ for each $1 \le i \le m$.
  \item \emph{Spine head-expansion and $\Lambda_{\mathrm{I}}$ pullback:} For all $B, A, C_1, \dots, C_m \in \Lambda$ ($m \ge 0$), if $B \in \SN_\beta$, $A \in \SN_\beta$, $C_1, \dots, C_m \in \SN_\beta$, and $(B[x := A])\,C_1 \cdots C_m \in \SN_\beta$, then $((\lambda x.\, B)\,A)\,C_1 \cdots C_m \in \SN_\beta$. Furthermore, $\Lambda_{\mathrm{I}}$ is closed under $\to_\beta$, and whenever $P \in \Lambda_{\mathrm{I}}$ satisfies $P \to_{\mathrm{wh}}^* Q$ with $Q \in \SN_\beta$, one has $P \in \SN_\beta$.
  \item \emph{Church's conservation theorem and fresh-variable application \cite{Church1941,Barendregt1984}:} For every $P \in \Lambda_{\mathrm{I}}$, one has $P \in \WN_\beta$ if and only if $P \in \SN_\beta$. Moreover, for every $P \in \Lambda$ and every variable $z \in \mathcal{V}$, if $P \in \WN_\beta$, then $P z \in \WN_\beta$; hence on $\Lambda_{\mathrm{I}}$, if $P \in \SN_\beta$, then $P z \in \SN_\beta$.
\end{enumerate}
\end{lemma}

\begin{proof}
For~(i), recall that parallel $\beta$-reduction $M \Rightarrow_\beta N$ is generated inductively by $y \Rightarrow_\beta y$ on variables, $A B \Rightarrow_\beta A' B'$ when $A \Rightarrow_\beta A'$ and $B \Rightarrow_\beta B'$, $\lambda x.\, B \Rightarrow_\beta \lambda x.\, B'$ when $B \Rightarrow_\beta B'$, and $(\lambda x.\, U)\,W \Rightarrow_\beta U'[x := W']$ when $U \Rightarrow_\beta U'$ and $W \Rightarrow_\beta W'$, with $\Rightarrow_\beta^*$ equal to $\to_\beta^*$. Observe first that $\to_{\mathrm{wh}}^*$ is preserved under outer substitution: $A \to_{\mathrm{wh}}^* B$ implies $A[z := S] \to_{\mathrm{wh}}^* B[z := S]$. Structural induction on $M \Rightarrow_\beta N$ establishes three head-shape factorizations:
\begin{enumerate}[label=(\alph*),leftmargin=2.0em]
  \item If $N \equiv y$ is a variable, then $M \to_{\mathrm{wh}}^* y$. In the sole non-trivial clause $(\lambda x.\, U)\,W \Rightarrow_\beta U'[x := W'] \equiv y$, either $U' \equiv y$ with $y \ne x$ (so $U \to_{\mathrm{wh}}^* y$ by the induction hypothesis and $(\lambda x.\, U)\,W \to_{\mathrm{wh}} U[x := W] \to_{\mathrm{wh}}^* y$) or $U' \equiv x$ and $W' \equiv y$ (so $U \to_{\mathrm{wh}}^* x$ and $W \to_{\mathrm{wh}}^* y$, whence $(\lambda x.\, U)\,W \to_{\mathrm{wh}} U[x := W] \to_{\mathrm{wh}}^* W \to_{\mathrm{wh}}^* y$).
  \item If $N \equiv \lambda x.\, Q'$, there exists $B_0$ with $M \to_{\mathrm{wh}}^* \lambda x.\, B_0$ and $B_0 \Rightarrow_\beta Q'$. Under $(\lambda z.\, U)\,W \Rightarrow_\beta U'[z := W'] \equiv \lambda x.\, Q'$, either $U' \equiv \lambda x.\, U_1'$ with $U_1'[z := W'] \equiv Q'$ (giving $U \to_{\mathrm{wh}}^* \lambda x.\, U_0$ with $U_0 \Rightarrow_\beta U_1'$, so $(\lambda z.\, U)\,W \to_{\mathrm{wh}}^* \lambda x.\, U_0[z := W]$ and $U_0[z := W] \Rightarrow_\beta Q'$ by substitutivity of $\Rightarrow_\beta$) or $U' \equiv z$ and $W' \equiv \lambda x.\, Q'$ (giving $U \to_{\mathrm{wh}}^* z$ and $W \to_{\mathrm{wh}}^* \lambda x.\, B_0$ with $B_0 \Rightarrow_\beta Q'$, so $(\lambda z.\, U)\,W \to_{\mathrm{wh}} U[z := W] \to_{\mathrm{wh}}^* W \to_{\mathrm{wh}}^* \lambda x.\, B_0$).
  \item If $N \equiv A' B'$, an identical split on $U'$ yields $A_0, B_0$ with $M \to_{\mathrm{wh}}^* A_0 B_0$, $A_0 \Rightarrow_\beta A'$, and $B_0 \Rightarrow_\beta B'$.
\end{enumerate}
Using~(b) and~(c), induction on $M_1 \to_{\mathrm{wh}} M_2$ shows that whenever $M \Rightarrow_\beta M_1 \to_{\mathrm{wh}} M_2$, there exists $M_0$ with $M \to_{\mathrm{wh}}^* M_0 \Rightarrow_\beta M_2$: at a root step $M_1 \equiv (\lambda x.\, P')\,Q' \to_{\mathrm{wh}} P'[x := Q']$, applying~(c) then~(b) to $M \Rightarrow_\beta (\lambda x.\, P')\,Q'$ produces $M \to_{\mathrm{wh}}^* A_0 B_0 \to_{\mathrm{wh}}^* (\lambda x.\, P_0)\,B_0 \to_{\mathrm{wh}} P_0[x := B_0] \Rightarrow_\beta P'[x := Q']$; the left-application step follows immediately from~(c). Iterating across $M_1 \to_{\mathrm{wh}}^* M_2$ postpones $\Rightarrow_\beta$ past $\to_{\mathrm{wh}}^*$, and induction along $P \Rightarrow_\beta^* N$ yields both the abstraction standardization claim (from~(b)) and the binary-application standardization claim (from~(c)), which extends by induction on $m$ from~(a) to any variable-headed spine $y R_1 \cdots R_m$.

For~(ii), let $S_0 \equiv (B[x := A])\,C_1 \cdots C_m$. We prove $((\lambda x.\, B)\,A)\,C_1 \cdots C_m \in \SN_\beta$ by nested well-founded induction on $B \in \SN_\beta$, $A \in \SN_\beta$, componentwise reduction on $(C_1, \dots, C_m) \in \SN_\beta^m$, and $S_0 \in \SN_\beta$. Because $(\lambda x.\, B)\,A$ is not an abstraction, any single step $((\lambda x.\, B)\,A)\,C_1 \cdots C_m \to_\beta Y$ arises in one of four ways: contracting $(\lambda x.\, B)\,A$ at the head yields $Y \equiv S_0 \in \SN_\beta$ directly; stepping $B \to_\beta B'$ yields $S_0 \to_\beta (B'[x := A])\,C_1 \cdots C_m$, which is covered by the induction hypothesis on $B'$; stepping $A \to_\beta A'$ yields $S_0 \to_\beta^* (B[x := A'])\,C_1 \cdots C_m$, which lies in $\SN_\beta$ as a reduct of $S_0$ and is covered by the induction hypothesis on $A'$; and stepping $C_i \to_\beta C_i'$ steps $S_0$ in one step, covered by the tuple hypothesis on $(C_1, \dots, C_m)$. Next, induction on $P \to_\beta P'$ shows simultaneously that $P \in \Lambda_{\mathrm{I}}$ implies $P' \in \Lambda_{\mathrm{I}}$ and $\mathrm{FV}(P) \subseteq \mathrm{FV}(P')$: at $(\lambda x.\, B)\,A \to_\beta B[x := A]$, the hypothesis $\lambda x.\, B \in \Lambda_{\mathrm{I}}$ guarantees $x \in \mathrm{FV}(B)$, so $A$ occurs at least once in $B[x := A]$, preserving both $\Lambda_{\mathrm{I}}$-membership and all free variables of $A$ and $B \setminus \{x\}$. Now suppose $P \in \Lambda_{\mathrm{I}}$ and $P \to_{\mathrm{wh}} Q \in \SN_\beta$. Writing $P \equiv ((\lambda x.\, B)\,A)\,C_1 \cdots C_m$ and $Q \equiv (B[x := A])\,C_1 \cdots C_m$, subterm projection from $Q \in \SN_\beta$ gives $B[x := A] \in \SN_\beta$ and $C_1, \dots, C_m \in \SN_\beta$. Since every step $B \to_\beta B'$ induces $B[x := A] \to_\beta B'[x := A]$, we obtain $B \in \SN_\beta$; and because $\lambda x.\, B \in \Lambda_{\mathrm{I}}$ forces $x \in \mathrm{FV}(B)$, every step $A \to_\beta A'$ induces at least one step $B[x := A] \to_\beta^+ B[x := A']$, forcing $A \in \SN_\beta$ as well. The head-expansion property then yields $P \in \SN_\beta$, and induction along $P \to_{\mathrm{wh}}^* Q$ completes~(ii).

For~(iii), define the syntactic size $\|y\|_\lambda = 1$ on variables, $\|\lambda x.\, B\|_\lambda = 1 + \|B\|_\lambda$, and $\|A B\|_\lambda = 1 + \|A\|_\lambda + \|B\|_\lambda$. Given $P \in \Lambda_{\mathrm{I}}$ with $P \to_\beta^* N \in \NF_\beta$, we prove $P \in \SN_\beta$ by strong induction on $\|N\|_\lambda$. Because $N \in \NF_\beta$ contains no head $\beta$-redex, spine decomposition places $N$ into one of two shapes. If $N \equiv \lambda x.\, N_0$, then $N_0 \in \NF_\beta$ and $\|N_0\|_\lambda < \|N\|_\lambda$; part~(i) gives $P \to_{\mathrm{wh}}^* \lambda x.\, B_0$ with $B_0 \to_\beta^* N_0$, part~(ii) gives $\lambda x.\, B_0 \in \Lambda_{\mathrm{I}}$ and hence $B_0 \in \Lambda_{\mathrm{I}}$, the induction hypothesis gives $B_0 \in \SN_\beta$ (so $\lambda x.\, B_0 \in \SN_\beta$), and weak-head pullback from~(ii) yields $P \in \SN_\beta$. Otherwise $N \equiv y N_1 \cdots N_m$ ($m \ge 0$) for a bound or free variable $y$, where each $N_i \in \NF_\beta$ satisfies $\|N_i\|_\lambda < \|N\|_\lambda$; part~(i) gives $P \to_{\mathrm{wh}}^* y R_1^0 \cdots R_m^0 \in \Lambda_{\mathrm{I}}$ with $R_i^0 \to_\beta^* N_i$, the induction hypothesis gives $R_i^0 \in \SN_\beta$ for every $i$, whence $y R_1^0 \cdots R_m^0 \in \SN_\beta$ (as $y$ is irreducible and not an abstraction), and weak-head pullback again yields $P \in \SN_\beta$. Conversely, $P \in \SN_\beta$ implies $P \in \WN_\beta$ constructively on all of $\Lambda$ by iterating the total leftmost-outermost step detector until it returns normal. Finally, if $P \to_\beta^* N \in \NF_\beta$, then $P z \to_\beta^* N z$: when $N \equiv \lambda x.\, B$ ($B \in \NF_\beta$), one further step gives $N z \to_\beta B[x := z] \in \NF_\beta$ because substituting a free variable $z$ for $x$ never converts a non-abstraction into a $\lambda$-abstraction; when $N$ is not an abstraction, $N z \in \NF_\beta$ directly.
\end{proof}

We now establish Conjecture~40 of Barendregt, Endrullis, Klop, and Waldmann \cite[p.~102]{BarendregtEndrullisKlopWaldmann2017} in full generality by characterizing order~$0$ through $\mathcal{P}_{\mathrm{CL}}$ and bridging partial combinator reducts with $\lambda$-abstractions via weak-head standardization.

\begin{theorem}[Affirmative resolution of Conjecture~40]\label{thm:conjecture-40}
The canonical translation $(-)_\lambda \colon \mathrm{Ter}(\mathrm{CL}(\{S, K, I\})) \to \Lambda$ preserves functional order exactly:
\begin{enumerate}[label=(\roman*),leftmargin=2.2em]
  \item \emph{Partial-combinator and abstraction bridge:} For every $M \in \mathrm{Ter}(\mathrm{CL}(\{S, K, I\}))$, there exists $H \in \mathcal{P}_{\mathrm{CL}}$ with $M \to_{\mathrm{CL}}^* H$ if and only if there exists $Q \in \Lambda$ with $(M)_\lambda \to_\beta^* \lambda x.\, Q$. Moreover, $M$ never consumes fresh arguments ($\ord_{\mathrm{CL}}(M) = 0$) if and only if no reduct of $M$ belongs to $\mathcal{P}_{\mathrm{CL}}$.
  \item \emph{Exact order preservation on $\mathrm{CL}(\{S, K, I\})$ \cite[Conjecture~40, p.~102]{BarendregtEndrullisKlopWaldmann2017}:} For every combinatory term $M \in \mathrm{Ter}(\mathrm{CL}(\{S, K, I\}))$ and every $n \in \mathbb{N}$,
  \[
    \ord_{\mathrm{CL}}(M) = n \;\iff\; \ord_\lambda\bigl((M)_\lambda\bigr) = n,
  \]
  and consequently $\ord_{\mathrm{CL}}(M) = \infty$ if and only if $\ord_\lambda((M)_\lambda) = \infty$.
  \item \emph{Specialization to $\CLS$ and order-$3$ ceiling:} For every open or closed pure-$S$ term $M \in \TSV$ and every $n \in \mathbb{N}$, one has $\ord_{\mathrm{CL}}(M) = n$ if and only if $\ord_\lambda((M)_\lambda) = n$; furthermore, $\ord_{\mathrm{CL}}(M z_1 z_2 z_3) = \ord_\lambda((M z_1 z_2 z_3)_\lambda) = 0$ for all $z_1, z_2, z_3 \in \mathcal{V}$, so $\ord_{\mathrm{CL}}(M) = \ord_\lambda((M)_\lambda) \in \{0, 1, 2, 3\}$.
\end{enumerate}
\end{theorem}

\begin{proof}
For~(i), we first prove that $M$ never consumes fresh arguments if and only if no $H \in \mathcal{P}_{\mathrm{CL}}$ satisfies $M \to_{\mathrm{CL}}^* H$. Suppose $M \to_{\mathrm{CL}}^* H \in \mathcal{P}_{\mathrm{CL}}$, and choose distinct fresh variables $y_1, y_2, y_3 \in \mathcal{V} \setminus \mathrm{FV}(M)$; since $\to_{\mathrm{CL}}^*$ creates no new free variables, $y_1, y_2, y_3 \notin \mathrm{FV}(H)$. Direct inspection of the six cases of \eqref{eq:partial-comb-sets} refutes unconsumed-argument persistence in each case: $I y_1 \to_{\mathrm{CL}} y_1$ and $K y_1 y_2 \to_{\mathrm{CL}} y_1$ cannot have the form $M' y_1$ or $M' y_1 y_2$; $(K A) y_1 \to_{\mathrm{CL}} A$ satisfies $y_1 \notin \mathrm{FV}(A)$, whereas $y_1 \in \mathrm{FV}(M' y_1)$; and $S y_1 y_2 y_3 \to_{\mathrm{CL}} (y_1 y_3)(y_2 y_3)$, $(S A) y_1 y_2 \to_{\mathrm{CL}} (A y_2)(y_1 y_2)$, and $(S A B) y_1 \to_{\mathrm{CL}} (A y_1)(B y_1)$ all terminate in a compound right-hand factor $(-\,y_r) \not\equiv y_r$. Conversely, assume no reduct of $M$ belongs to $\mathcal{P}_{\mathrm{CL}}$. By induction on $r \ge 0$, we show that for any variables $y_1, \dots, y_r \in \mathcal{V}$, every reduction $M y_1 \cdots y_r \to_{\mathrm{CL}}^* N$ has the form $N \equiv M' y_1 \cdots y_r$ with $M \to_{\mathrm{CL}}^* M'$. The base $r = 0$ is immediate. For $r = r_0 + 1$, put $F \equiv M y_1 \cdots y_{r_0}$. By the induction hypothesis on $r_0$, the term $F$ cannot reduce to $I$, to any $K P$, or to any $S P Q$: for example, if $F \to_{\mathrm{CL}}^* S P Q$, then $S P Q \equiv M' y_1 \cdots y_{r_0}$ with $M \to_{\mathrm{CL}}^* M'$, which forces $M' \equiv S P Q \in \mathcal{P}_{\mathrm{CL}}$ ($r_0 = 0$), $M' \equiv S P \in \mathcal{P}_{\mathrm{CL}}$ ($r_0 = 1$), or $M' \equiv S \in \mathcal{P}_{\mathrm{CL}}$ ($r_0 = 2$), while $r_0 \ge 3$ is impossible since $\ls(S P Q) = 2$; the cases $F \to_{\mathrm{CL}}^* I$ and $F \to_{\mathrm{CL}}^* K P$ are excluded analogously. Hence no root contraction can occur along any reduction of $F y_{r_0+1}$, so every reduct has the form $F' y_{r_0+1}$ with $F \to_{\mathrm{CL}}^* F'$, and applying the induction hypothesis to $F \to_{\mathrm{CL}}^* F'$ completes the induction.

Next we prove that $M \to_{\mathrm{CL}}^* H$ for some $H \in \mathcal{P}_{\mathrm{CL}}$ if and only if $(M)_\lambda \to_\beta^* \lambda x.\, Q$ for some $Q \in \Lambda$. In the forward direction, $M \to_{\mathrm{CL}}^* H$ gives $(M)_\lambda \to_\beta^* (H)_\lambda$, and for every $H \in \mathcal{P}_{\mathrm{CL}}$ the closed-bound-variable term $(H)_\lambda$ reaches an abstraction in at most two weak-head steps: $(I)_\lambda, (K)_\lambda, (S)_\lambda$ are already abstractions, $(K A)_\lambda \to_{\mathrm{wh}} \lambda y.\, (A)_\lambda$, $(S A)_\lambda \to_{\mathrm{wh}} \lambda y z.\, ((A)_\lambda z)(y z)$, and $(S A B)_\lambda \to_{\mathrm{wh}}^2 \lambda z.\, ((A)_\lambda z)((B)_\lambda z)$. Conversely, suppose $(M)_\lambda \to_\beta^* \lambda x.\, Q$. By Lemma~\ref{lem:lambdaI-conservation}(i), $(M)_\lambda \to_{\mathrm{wh}}^k \lambda x.\, B_0$ for some step count $k \ge 0$. Notice that $\to_{\mathrm{wh}}$ is deterministic and stops on abstractions. We proceed by strong induction on $k$, decomposing $M$ along its left combinatory spine into five mutually exclusive cases:
\begin{enumerate}[label=(\arabic*),leftmargin=2.0em]
  \item If $M \equiv v C_1 \cdots C_m$ ($v \in \mathcal{V}$), then $(M)_\lambda \equiv v (C_1)_\lambda \cdots (C_m)_\lambda$ is neither an abstraction nor weak-head reducible, which is impossible.
  \item If $M \in \mathcal{P}_{\mathrm{CL}}$, then $M \to_{\mathrm{CL}}^0 M \in \mathcal{P}_{\mathrm{CL}}$ immediately.
  \item If $M \equiv (I X) C_1 \cdots C_m$, then $(M)_\lambda$ is an application (so $k \ge 1$) whose unique first weak-head step yields $(X C_1 \cdots C_m)_\lambda \to_{\mathrm{wh}}^{k-1} \lambda x.\, B_0$; the induction hypothesis at $k - 1 < k$ produces $X C_1 \cdots C_m \to_{\mathrm{CL}}^* H \in \mathcal{P}_{\mathrm{CL}}$, prepended by $M \to_{\mathrm{CL}} X C_1 \cdots C_m$.
  \item If $M \equiv (K X Y) C_1 \cdots C_m$, then $(M)_\lambda$ executes two forced weak-head steps through the non-abstraction $((\lambda y.\, (X)_\lambda)\,(Y)_\lambda)\,(C_1)_\lambda \cdots (C_m)_\lambda$ (so $k \ge 2$) to $(X C_1 \cdots C_m)_\lambda$ in $k - 2 < k$ remaining steps, lifting $M \to_{\mathrm{CL}} X C_1 \cdots C_m$.
  \item If $M \equiv (S X Y Z) C_1 \cdots C_m$, then $(M)_\lambda$ executes three forced weak-head steps through two intermediate non-abstractions (so $k \ge 3$) to $(((X Z)(Y Z)) C_1 \cdots C_m)_\lambda$ in $k - 3 < k$ remaining steps, lifting $M \to_{\mathrm{CL}} ((X Z)(Y Z)) C_1 \cdots C_m$.
\end{enumerate}

For~(ii), observe first the fresh-variable step-down property of $\ord_\lambda$ (\cite[Example~36(iv), p.~99]{BarendregtEndrullisKlopWaldmann2017}): whenever $P \to_\beta^* \lambda x.\, Q$ and $z \in \mathcal{V} \setminus \mathrm{FV}(P)$, one has $\ord_\lambda(P) = n + 1$ if and only if $\ord_\lambda(P z) = n$. Indeed, $z \notin \mathrm{FV}(P)$ implies $z \notin \mathrm{FV}(Q)$, and the syntactic variable-closing operation $\kappa_z^x$ replacing free $z$ by bound $x$ commutes with $\to_\beta^*$ and satisfies $\kappa_z^x(Q[x := z]) \equiv Q$. Induction on $n$ shows that $\ord_\lambda(Q) = n$ holds if and only if $\ord_\lambda(Q[x := z]) = n$: at $n = 0$, $Q \to_\beta^* \lambda y.\, R$ gives $Q[x := z] \to_\beta^* \lambda y.\, R[x := z]$, while $Q[x := z] \to_\beta^* \lambda y.\, R'$ applies $\kappa_z^x$ to yield $Q \to_\beta^* \lambda y.\, \kappa_z^x(R')$; at $n + 1$, applying $\kappa_z^x$ together with Church--Rosser confluence under $\lambda y$ transfers the inductive step. Because $P z \to_\beta^* (\lambda x.\, Q)\,z \to_\beta Q[x := z]$ gives $P z =_\beta Q[x := z]$ and $\ord_\lambda$ is $=_\beta$-invariant, $\ord_\lambda(P z) = n$ is equivalent to $\ord_\lambda(Q[x := z]) = n$, hence to $\ord_\lambda(Q) = n$, and therefore to $\ord_\lambda(P) = n + 1$.

We now prove $\ord_{\mathrm{CL}}(M) = n \iff \ord_\lambda((M)_\lambda) = n$ by induction on $n \in \mathbb{N}$ across all $M \in \mathrm{Ter}(\mathrm{CL}(\{S, K, I\}))$. For $n = 0$, part~(i) equates $\ord_{\mathrm{CL}}(M) = 0$ with absence of any $H \in \mathcal{P}_{\mathrm{CL}}$ satisfying $M \to_{\mathrm{CL}}^* H$, and hence with absence of any abstraction reduct of $(M)_\lambda$, which is $\ord_\lambda((M)_\lambda) = 0$. For $n + 1$, pick any fresh $z \in \mathcal{V} \setminus \mathrm{FV}(M) = \mathcal{V} \setminus \mathrm{FV}((M)_\lambda)$. By part~(i), $M$ reduces to an element of $\mathcal{P}_{\mathrm{CL}}$ if and only if $(M)_\lambda$ reduces to an abstraction; under this shared condition, the step-down law for $\ord_\lambda$ and the induction hypothesis on $M z$ yield
\[
  \ord_{\mathrm{CL}}(M) = n + 1 \;\iff\; \ord_{\mathrm{CL}}(M z) = n \;\iff\; \ord_\lambda\bigl((M z)_\lambda\bigr) = n \;\iff\; \ord_\lambda\bigl((M)_\lambda\bigr) = n + 1.
\]
Negating over all $n \in \mathbb{N}$ shows that $\ord_{\mathrm{CL}}(M) = \infty$ holds if and only if $\ord_\lambda((M)_\lambda) = \infty$.

For~(iii), conservative embedding of $\TSV$ into $\mathrm{Ter}(\mathrm{CL}(\{S, K, I\}))$ identifies $\to_{\mathrm{CL}}^*$ with $\redSs$ and $\mathcal{P}_{\mathrm{CL}} \cap \TSV$ with $\mathcal{P}_S$. By the Visser--van Oostrom invariant $\mathcal{W}_3$ of Proposition~\ref{prop:acyclicity-order}(iii), for all $M \in \TSV$ and $z_1, z_2, z_3 \in \mathcal{V}$, the term $M z_1 z_2 z_3$ belongs to $\mathcal{W}_3$, which is closed under $\redSs$ and disjoint from $\mathcal{P}_S$ (as $S \notin \mathcal{W}_1$). Part~(i) therefore gives $\ord_{\mathrm{CL}}(M z_1 z_2 z_3) = \ord_\lambda((M z_1 z_2 z_3)_\lambda) = 0$, bounding the chain of fresh-variable applications at depth at most~$3$ and ensuring $\ord_{\mathrm{CL}}(M) = \ord_\lambda((M)_\lambda) \in \{0, 1, 2, 3\}$.
\end{proof}

\subsection{Minimal \texorpdfstring{$\lambda\beta$}{lambda-beta}-Diverging Combinators, Simultaneous Refutation of Conjectures 42 and 43, and Hereditary Saturation}\label{subsec:conj42-43-refutation}

While Theorem~\ref{thm:conjecture-40} shows that $(-)_\lambda$ preserves functional order on $\CLS$, normalization behaves quite differently. Because $M \redS N$ implies $(M)_\lambda \to_\beta^+ (N)_\lambda$, strong normalization pulls back from $\Lambda_{\mathrm{I}}$ to $\CLS$: $(M)_\lambda \in \SN_\beta$ always implies $M \in \SN$. In the opposite direction, Barendregt, Endrullis, Klop, and Waldmann observed in Example~33 of \cite[p.~95]{BarendregtEndrullisKlopWaldmann2017} (following Waldmann \cite{Waldmann2000}) that the $13$-leaf closed normal form $W_{\mathrm{Wal}} \equiv S V V$, built from $T \equiv S S$, $Z \equiv S T T$, and $V \equiv S Z$, satisfies $(W_{\mathrm{Wal}})_\lambda \notin \WN_\beta$, and posed Question~34 asking for a characterization of $\{M \in \TS : (M)_\lambda \in \WN_\beta\}$. Our next result determines the exact minimum leaf sizes at which closed normalizing pure-$S$ combinators ($9$ leaves) and closed pure-$S$ normal forms ($10$ leaves) diverge under $\lambda\beta$-reduction.

\begin{proposition}[Minimal closed $\lambda\beta$-diverging pure-$S$ combinators]\label{prop:minimal-lambda-divergers}
Let $T \equiv S S$, $A \equiv S S S$, $Z \equiv S T T$, and $V \equiv S Z$, and define the closed pure-$S$ terms
\[
  W_9 \;\equiv\; S\bigl(S(A A)\bigr)S, \qquad N_{10} \;\equiv\; S\bigl(S(S A A)\bigr)S, \qquad W_9^{\mathrm{nf}} \;\equiv\; S\bigl(S(S A (S A))\bigr)S, \qquad W_{\mathrm{Wal}} \;\equiv\; S V V.
\]
Then the following properties hold:
\begin{enumerate}[label=(\roman*),leftmargin=2.2em]
  \item \emph{Sizes and $\CLS$-normalization:} One has $|W_9| = 9$, $|N_{10}| = 10$, $|W_9^{\mathrm{nf}}| = 11$, and $|W_{\mathrm{Wal}}| = 13$. Each of $N_{10}$, $W_9^{\mathrm{nf}}$, and $W_{\mathrm{Wal}}$ lies in $\TS \cap \NF$, whereas $W_9 \in \TS \cap \SN$ reduces in a single deterministic step $W_9 \redS W_9^{\mathrm{nf}}$ to its normal form $W_9^{\mathrm{nf}}$.
  \item \emph{Divergence in $\lambda\beta$:} For every variable $z \in \mathcal{V}$ and every $U \in \{W_9, N_{10}, W_9^{\mathrm{nf}}, W_{\mathrm{Wal}}\}$, the open application $U z$ satisfies $U z \notin \SN$ in $\CLS$, and consequently $(U)_\lambda \notin \SN_\beta$ and $(U)_\lambda \notin \WN_\beta$.
  \item \emph{Global minimality:} For every closed pure-$S$ term $M \in \TS$ with $|M| \le 8$, one has $(M)_\lambda \in \SN_\beta$ if and only if $M \in \SN$; and for every closed pure-$S$ normal form $N \in \TS \cap \NF$ with $|N| \le 9$, one has $(N)_\lambda \in \SN_\beta$.
\end{enumerate}
\end{proposition}

\begin{proof}
For~(i), counting $S$-leaves gives $|A| = 3$, $|A A| = 6$, $|W_9| = 1 + (1 + 6) + 1 = 9$, $|N_{10}| = 1 + (1 + 7) + 1 = 10$, $|W_9^{\mathrm{nf}}| = 1 + (1 + (4 + 4)) + 1 = 11$, and $|W_{\mathrm{Wal}}| = 1 + 6 + 6 = 13$. Every binary application node in $N_{10}$, $W_9^{\mathrm{nf}}$, and $W_{\mathrm{Wal}}$ has left-spine length $1$ or $2$, placing all three in $\NF$. In $W_9 \equiv S(S(A A))S$, the sole redex is $A A \equiv S S S A \redS (S A)(S A) \equiv S A (S A) \in \NF$, so $W_9 \redS W_9^{\mathrm{nf}} \in \NF$ and $W_9 \in \SN$.

For~(ii), observe first a general subterm-spiral criterion in $\CLS$: if a sequence $(F_n)_{n \ge 0}$ in $\TSV$ satisfies $F_n \redSp C_n[F_{n+1}]$ for some one-hole context $C_n[\,]$ at every $n \ge 0$, then $F_0 \notin \SN$, because accessibility of $X \in \SN$ under $\redSp$ lifts across subterm contexts to well-foundedness of the relation $X \succ Y$ defined by the existence of $C[\,]$ with $X \redSp C[Y]$. Moreover, whenever $U \in \TS$ and $U z \notin \SN$, we must have $(U)_\lambda \notin \WN_\beta$: otherwise Lemma~\ref{lem:lambdaI-conservation}(iii) would give $(U z)_\lambda \equiv (U)_\lambda z \in \WN_\beta = \SN_\beta$ (as $(U z)_\lambda \in \Lambda_{\mathrm{I}}$), which pulls back via forward simulation to $U z \in \SN$.

To exhibit the explicit $\beta$-loop beneath $\lambda z$ for $N_{10}$, note that two weak-head $\beta$-steps yield $(N_{10})_\lambda \equiv (S)_\lambda\,(S(S A A))_\lambda\,(S)_\lambda \to_{\mathrm{wh}}^2 \lambda z.\, (N_{10}\,z)_\lambda$. Inside $\lambda z$, three pure-$S$ contractions (nine $\beta$-steps) give
\[
  N_{10}\,z \equiv S\bigl(S(S A A)\bigr)S z \;\redS\; S(S A A)z(S z) \;\redS\; \bigl(S A A (S z)\bigr)\bigl(z(S z)\bigr) \;\redS\; \Bigl(\bigl(A(S z)\bigr)\bigl(A(S z)\bigr)\Bigr)\bigl(z(S z)\bigr).
\]
Define $G_0(z) \equiv A(S z)$, $G_{n+1}(z) \equiv S(S z)\,G_n(z)$, and $\Phi_n(z) \equiv G_0(z)\,G_n(z)$ for $n \ge 0$. For every $n \ge 0$, three pure-$S$ contractions yield
\begin{equation}\label{eq:n10-beta-spiral}
\begin{aligned}
  \Phi_n(z) \;\equiv\; \bigl(S S S (S z)\bigr)\,G_n(z)
  &\;\redS\; \bigl(S(S z)(S(S z))\bigr)\,G_n(z) \\
  &\;\redS\; \bigl(S z\,G_n(z)\bigr)\,\bigl(S(S z)\,G_n(z)\bigr) \;\equiv\; \bigl(S z\,G_n(z)\bigr)\,G_{n+1}(z) \\
  &\;\redS\; \bigl(z\,G_{n+1}(z)\bigr)\,\bigl(G_n(z)\,G_{n+1}(z)\bigr).
\end{aligned}
\end{equation}
By induction on $n \ge 0$, every $Y \in \TSV$ satisfies $G_n(z)\,Y \redS^n D_n[G_0(z)\,Y]$ with $D_0[\,] \equiv \square$ and $D_{n+1}[\,] \equiv (S z Y)\,D_n[\,]$ (since $G_{n+1}(z)\,Y \equiv S(S z)\,G_n(z)\,Y \redS (S z Y)(G_n(z)\,Y)$). Taking $Y \equiv G_{n+1}(z)$ in \eqref{eq:n10-beta-spiral} gives $\Phi_n(z) \redS^{n+3} (z\,G_{n+1}(z))\,D_n[\Phi_{n+1}(z)]$, so $(N_{10})_\lambda$ enters the non-terminating $\beta$-reduction
\[
  (N_{10})_\lambda \;\to_\beta^2\; \lambda z.\, (N_{10}\,z)_\lambda \;\to_\beta^9\; \lambda z.\, \bigl((\Phi_0(z))_\lambda\,(z\,(S z)_\lambda)\bigr) \;\to_\beta^9\; \lambda z.\, \Bigl(\bigl((z\,(G_1(z))_\lambda)\,(\Phi_1(z))_\lambda\bigr)\,\bigl(z\,(S z)_\lambda\bigr)\Bigr) \;\to_\beta^{12}\; \cdots
\]
underneath $\lambda z$, proving $N_{10}\,z \notin \SN$ and $(N_{10})_\lambda \notin \SN_\beta$.

Similarly, for $W_9$ and $W_9^{\mathrm{nf}}$, we have $W_9\,z \redS W_9^{\mathrm{nf}}\,z$ and $(W_9)_\lambda \to_\beta^3 (W_9^{\mathrm{nf}})_\lambda \to_{\mathrm{wh}}^2 \lambda z.\, (W_9^{\mathrm{nf}}\,z)_\lambda$. Put $P_z \equiv S A (S z)$, $E_0(z) \equiv S(S z)P_z$, $E_{n+1}(z) \equiv (S z)\,E_n(z)$, and $\Psi_n(z) \equiv P_z\,E_n(z)$. Six pure-$S$ contractions reduce $W_9^{\mathrm{nf}}\,z$ to $\bigl((z\,E_0(z))\,\Psi_0(z)\bigr)(z(S z))$:
\begin{equation}\label{eq:w9-beta-spiral}
\begin{aligned}
  W_9^{\mathrm{nf}}\,z
  &\;\redS^2\; \bigl(S A (S A)(S z)\bigr)\bigl(z(S z)\bigr) \;\redS\; \bigl(A(S z)\,P_z\bigr)\bigl(z(S z)\bigr) \\
  &\;\redS^2\; \bigl(S z P_z\,E_0(z)\bigr)\bigl(z(S z)\bigr) \;\redS\; \Bigl(\bigl(z\,E_0(z)\bigr)\,\Psi_0(z)\Bigr)\bigl(z(S z)\bigr).
\end{aligned}
\end{equation}
For each $n \ge 0$, three contractions give $\Psi_n(z) \redS (A\,E_n(z))\,E_{n+1}(z) \redS^2 (E_n(z)\,E_{n+1}(z))(S\,E_n(z)\,E_{n+1}(z))$, while induction on $n$ shows $E_n(z)\,Y \redS^{n+1} \widetilde{D}_n[P_z\,Y]$ (from $E_0(z)\,Y \redS (S z Y)(P_z Y)$ and $E_{n+1}(z)\,Y \redS (z Y)(E_n(z)\,Y)$), yielding $\Psi_n(z) \redS^{n+4} \widetilde{C}_n[\Psi_{n+1}(z)]$ and hence $W_9^{\mathrm{nf}}\,z, W_9\,z \notin \SN$ and $(W_9)_\lambda, (W_9^{\mathrm{nf}})_\lambda \notin \SN_\beta$. For $W_{\mathrm{Wal}} \equiv S V V$, setting $H_0(z) \equiv V z$ and $H_{n+1}(z) \equiv T\,H_n(z)$ gives $W_{\mathrm{Wal}}\,z \redS H_0(z)\,H_0(z)$ and $H_0(z)\,H_n(z) \redS^2 (H_{n+1}(z)\,H_{n+1}(z))(z\,H_n(z))$, while $H_m(z)\,Y \redS^m \widehat{D}_m[H_0(z)\,Y]$ yields $H_0(z)\,H_n(z) \redS^{n+3} \widehat{C}_n[H_0(z)\,H_{n+1}(z)]$, confirming $W_{\mathrm{Wal}}\,z \notin \SN$ and $(W_{\mathrm{Wal}})_\lambda \notin \SN_\beta$.

For~(iii), forward simulation already shows that $(M)_\lambda \in \SN_\beta$ implies $M \in \SN$ for all $M \in \TS$. Across the $626$ closed pure-$S$ terms with $|M| \le 8$ ($583$ of which belong to $\SN$) and the $323$ closed pure-$S$ normal forms of leaf size~$9$ ($539$ cumulative closed normal forms with $|N| \le 9$), executing deterministic leftmost-outermost $\beta$-reduction reaches a $\beta$-normal form on $(M)_\lambda$ for every $M \in \TS \cap \SN$ with $|M| \le 8$ and on $(N)_\lambda$ for every $N \in \TS \cap \NF$ with $|N| \le 9$ (for instance, on the $9$-leaf normal form $S(S(S A A))$, applying two fresh variables gives $S(S(S A A))\,x_0\,x_1 \redS^2 (S A A (x_0 x_1))(x_1(x_0 x_1))$, where the variable-headed term $x_0 x_1$ replaces $S z$ in \eqref{eq:n10-beta-spiral} and halts after four further steps at $((x_0 x_1)\,R)(S(x_0 x_1)\,R)$ with $R \equiv S(x_0 x_1)(S(x_0 x_1))$). Lemma~\ref{lem:lambdaI-conservation}(iii) upgrades $\WN_\beta$ to $\SN_\beta$ on all of these terms.
\end{proof}

Because every $M \in \TSV$ has functional order at most~$3$ (Proposition~\ref{prop:acyclicity-order}(iii) and Theorem~\ref{thm:conjecture-40}(iii)), applying $M$ to three distinct fresh variables $x_0, x_1, x_2 \notin \mathrm{FV}(M)$ exhausts all outer partial $S$-combinators along the top spine of $M$. On this basis, Barendregt, Endrullis, Klop, and Waldmann formulated two normalization conjectures on page~102 of \cite{BarendregtEndrullisKlopWaldmann2017}: Conjecture~42 asserts that for all $M \in \TSV$ and fresh $x_0, x_1, x_2 \notin \mathrm{FV}(M)$, the term $M x_0 x_1 x_2$ normalizes in $\CLS$ if and only if $(M)_\lambda$ normalizes in $\lambda\beta$; and Conjecture~43 asserts that for every finite-order combinator $M$ and fresh $x_1, \dots, x_{\ord_{\mathrm{CL}}(M)} \notin \mathrm{FV}(M)$, the term $M x_1 \cdots x_{\ord_{\mathrm{CL}}(M)}$ normalizes if and only if $(M)_\lambda$ normalizes in $\lambda\beta$ (specializing at $\ord_{\mathrm{CL}}(M) = 1$ to single-variable saturation $M x_0 \in \WN \implies (M)_\lambda \in \WN_\beta$).

Both conjectures fail in the forward direction because outer variable saturation inspects only the root spine: once $M x_0$ reduces to a variable-headed term $x_0\,U_1\,U_2$, any further outer variables $x_1, x_2$ accumulate passively on the right without probing whether an internal argument $U_i$ is itself a partial $S$-application $S W x_0 \in \NF$ whose $\lambda$-translation diverges beneath its own internal $\lambda$-binder. To capture nested binder saturation purely within $\CLS$, without $\lambda$-abstractions or substitution, we define \emph{hereditary fresh-variable saturation} $\mathrm{DeepNorm}(M)$ on $M \in \TSV$ inductively as the least predicate closed under the four inference rules
\begin{equation}\label{eq:deepnorm-rules}
\begin{gathered}
  \frac{M \redS M' \qquad \mathrm{DeepNorm}(M')}{\mathrm{DeepNorm}(M)} \;\textrm{(step)},
  \qquad\qquad
  \frac{x \in \mathcal{V}}{\mathrm{DeepNorm}(x)} \;\textrm{(var)}, \\[0.8ex]
  \frac{\mathrm{hd}(A) \in \mathcal{V} \qquad \mathrm{DeepNorm}(A) \qquad \mathrm{DeepNorm}(B)}{\mathrm{DeepNorm}(A B)} \;\textrm{(app)},
  \qquad\qquad
  \frac{M \in \mathcal{P}_S \qquad z \in \mathcal{V} \setminus \mathrm{FV}(M) \qquad \mathrm{DeepNorm}(M z)}{\mathrm{DeepNorm}(M)} \;\textrm{(sat)},
\end{gathered}
\end{equation}
where $\mathrm{hd}(a) = a$ on atoms $a \in \{S\} \cup \mathcal{V}$ and $\mathrm{hd}(A B) = \mathrm{hd}(A)$ denotes the leftmost leaf. Thus $\mathrm{DeepNorm}(M)$ holds if and only if $M$ reduces in $\CLS$ either to a variable-headed application $x\,A_1 \cdots A_r$ ($r \ge 0$) whose arguments $A_i$ satisfy $\mathrm{DeepNorm}(A_i)$, or to a partial $S$-application $H \in \mathcal{P}_S$ satisfying $\mathrm{DeepNorm}(H z)$ for a fresh variable $z \notin \mathrm{FV}(H)$.

\begin{theorem}[Simultaneous refutation of Conjectures~42 and~43 and hereditary saturation repair]\label{thm:conjectures-42-43-refuted}
Let $A \equiv S S S$, and for any $U \in \TS$ set $\mathcal{C}(U) \equiv S A (S U)$.
\begin{enumerate}[label=(\roman*),leftmargin=2.2em]
  \item \emph{Simultaneous refutation of Conjectures~42 and~43 \cite[p.~102]{BarendregtEndrullisKlopWaldmann2017}:} Instantiating $\mathcal{C}(U)$ with $W_9, N_{10}, W_9^{\mathrm{nf}}, W_{\mathrm{Wal}}$ from Proposition~\ref{prop:minimal-lambda-divergers} yields the closed pure-$S$ combinators
  \[
    M_{14} \;\equiv\; S A (S W_9), \qquad M_{15} \;\equiv\; S A (S N_{10}), \qquad M_{16} \;\equiv\; S A (S W_9^{\mathrm{nf}}), \qquad M_{18} \;\equiv\; S A (S W_{\mathrm{Wal}}),
  \]
  of leaf sizes $|M_{14}| = 14$, $|M_{15}| = 15$, $|M_{16}| = 16$, and $|M_{18}| = 18$, with $M_{15}, M_{16}, M_{18} \in \TS \cap \NF$ and $M_{14} \redS M_{16}$. Every $M \in \{M_{14}, M_{15}, M_{16}, M_{18}\}$ satisfies
  \[
    \ord_{\mathrm{CL}}(M) \;=\; \ord_\lambda\bigl((M)_\lambda\bigr) \;=\; 1, \qquad M x_0 \;\in\; \SN, \qquad M x_0 x_1 x_2 \;\in\; \SN
  \]
  for all $x_0, x_1, x_2 \in \mathcal{V}$, yet $(M)_\lambda \notin \SN_\beta$ and $(M)_\lambda \notin \WN_\beta$, refuting both Conjecture~42 and Conjecture~43 simultaneously. Moreover, $14$ leaves and $15$ leaves are the global minimum leaf sizes of closed pure-$S$ counterexamples in $\TS$ and $\TS \cap \NF$, respectively.
  \item \emph{Validity of the reverse direction:} For every $M \in \TSV$, every integer $k \ge 0$, and all variables $x_1, \dots, x_k \in \mathcal{V}$, if $(M)_\lambda \in \SN_\beta$ (equivalently $(M)_\lambda \in \WN_\beta$), then $M x_1 \cdots x_k \in \SN$ (equivalently $M x_1 \cdots x_k \in \WN$).
  \item \emph{Exact hereditary saturation theorem:} For every open or closed pure-$S$ term $M \in \TSV$,
  \[
    \mathrm{DeepNorm}(M) \;\iff\; (M)_\lambda \in \SN_\beta \;\iff\; (M)_\lambda \in \WN_\beta.
  \]
\end{enumerate}
\end{theorem}

\begin{proof}
For~(i), let $U \in \TS$ be any closed term satisfying $U \redSs U^{\mathrm{nf}} \in \NF$ and $U y \notin \SN$ for $y \in \mathcal{V}$, and put $M \equiv S A (S U)$. Clearly $|M| = 1 + 3 + (1 + |U|) = 5 + |U|$ and $M \in \NF$ whenever $U \in \NF$. For any variable $x_0 \in \mathcal{V}$, three pure-$S$ contractions yield
\begin{equation}\label{eq:mkce-mx0}
  M x_0 \;\equiv\; S A (S U)\,x_0 \;\redS\; (A x_0)(S U x_0) \;\redS\; \bigl(S x_0 (S x_0)\bigr)(S U x_0) \;\redS\; \bigl(x_0\,(S U x_0)\bigr)\,\bigl(S x_0\,(S U x_0)\bigr).
\end{equation}
Reducing $U \redSs U^{\mathrm{nf}}$ inside both occurrences of $S U x_0$ gives $M x_0 \redSs N_1(x_0) \equiv (x_0\,(S U^{\mathrm{nf}} x_0))(S x_0\,(S U^{\mathrm{nf}} x_0))$. Since $U^{\mathrm{nf}} \in \NF$, both $S U^{\mathrm{nf}} x_0$ and $S x_0 (S U^{\mathrm{nf}} x_0)$ have spine length~$2$ and lie in $\NF$, while $\mathrm{hd}(N_1(x_0)) = x_0 \ne S$ precludes any root redex; hence $N_1(x_0) \in \NF$ and $M x_0 \in \WN = \SN$ by Proposition~\ref{prop:cr-conservation}(iii). Appending $x_1, x_2 \in \mathcal{V}$ likewise gives $M x_0 x_1 x_2 \redSs N_1(x_0)\,x_1\,x_2 \in \NF$, so $M x_0 x_1 x_2 \in \WN = \SN$. Next, $M \equiv S A (S U) \in \mathcal{P}_S$ is closed, so to verify $\ord_{\mathrm{CL}}(M) = 1$ (and hence $\ord_\lambda((M)_\lambda) = 1$ by Theorem~\ref{thm:conjecture-40}(iii)) it suffices by Theorem~\ref{thm:conjecture-40}(i) to check that $M x_0$ cannot reduce to any $H \in \mathcal{P}_S$. Indeed, if $M x_0 \redSs H \in \mathcal{P}_S$, Church--Rosser confluence of $\redSs$ (Proposition~\ref{prop:cr-conservation}(i)) with $M x_0 \redSs N_1(x_0) \in \NF$ forces $H \redSs N_1(x_0)$; yet every $H \in \mathcal{P}_S$ satisfies $\ls(H) \le 2$ and $\mathrm{hd}(H) = S$, both of which are invariant under $\redSs$, contradicting $\mathrm{hd}(N_1(x_0)) = x_0$. Now suppose toward a contradiction that $(M)_\lambda \in \WN_\beta$. Applying Lemma~\ref{lem:lambdaI-conservation}(iii) at $x_0$ gives $(M x_0)_\lambda \in \WN_\beta = \SN_\beta$; forward simulation along \eqref{eq:mkce-mx0} places $((x_0\,(S U x_0)_\lambda)((S x_0)_\lambda\,(S U x_0)_\lambda)) \in \SN_\beta$ and hence its subterm $(S U x_0)_\lambda \in \SN_\beta = \WN_\beta$. A second application of Lemma~\ref{lem:lambdaI-conservation}(iii) at $x_1$ gives $(S U x_0 x_1)_\lambda \in \WN_\beta = \SN_\beta$, and simulating $S U x_0 x_1 \redS (U x_1)(x_0 x_1)$ places $(U x_1)_\lambda \in \SN_\beta$, forcing $U x_1 \in \SN$ and contradicting Proposition~\ref{prop:minimal-lambda-divergers}(ii). Note that replacing $A \equiv S S S$ by $T \equiv S S$ at $13$ leaves fails because $S T (S W_9)\,x_0\,x_1 \redS^4 (W_9 x_1)(x_0 x_1)(\dots) \notin \SN$ exposes $S W_9 x_0$ directly to $x_1$; checking all closed $M \in \TS \cap \SN$ with $|M| \le 13$ and all $N \in \TS \cap \NF$ with $|N| \le 14$ confirms global minimality of $M_{14}$ and $M_{15}$.

For~(ii), since $(M)_\lambda \in \Lambda_{\mathrm{I}}$ and $(M x_1 \cdots x_k)_\lambda \equiv (M)_\lambda x_1 \cdots x_k \in \Lambda_{\mathrm{I}}$, iterating Lemma~\ref{lem:lambdaI-conservation}(iii) $k$ times shows that $(M)_\lambda \in \SN_\beta$ (equivalently $(M)_\lambda \in \WN_\beta$) implies $(M x_1 \cdots x_k)_\lambda \in \WN_\beta = \SN_\beta$, which pulls back via forward simulation to $M x_1 \cdots x_k \in \SN = \WN$.

For~(iii), Lemma~\ref{lem:lambdaI-conservation}(iii) already equates $(M)_\lambda \in \SN_\beta$ with $(M)_\lambda \in \WN_\beta$ since $(M)_\lambda \in \Lambda_{\mathrm{I}}$. To prove that $\mathrm{DeepNorm}(M)$ implies $(M)_\lambda \in \WN_\beta$, proceed by induction on \eqref{eq:deepnorm-rules}. Rule~(step) follows from $(M)_\lambda \to_\beta^+ (M')_\lambda$, and rule~(var) from $(x)_\lambda \equiv x \in \NF_\beta$. In rule~(app), $\mathrm{hd}(A) = x \in \mathcal{V}$ implies $(A)_\lambda \equiv x C_1 \cdots C_m$, so Lemma~\ref{lem:lambdaI-conservation}(i) prevents $(A)_\lambda$ from reducing to any abstraction; thus if $(A)_\lambda \to_\beta^* N_A \in \NF_\beta$ and $(B)_\lambda \to_\beta^* N_B \in \NF_\beta$, then $N_A$ is not an abstraction and $(A B)_\lambda \to_\beta^* N_A N_B \in \NF_\beta$. In rule~(sat), $M \in \mathcal{P}_S$ gives $(M)_\lambda \to_\beta^* \lambda x.\, B$ by Theorem~\ref{thm:conjecture-40}(i) with $z \notin \mathrm{FV}(B)$; if $(M z)_\lambda \to_\beta^* N_z \in \NF_\beta$, confluence with $(M z)_\lambda \to_\beta^* (\lambda x.\, B)\,z \to_\beta B[x := z]$ forces $B[x := z] \to_\beta^* N_z$. Applying the syntactic variable-closing map $\kappa_z^x$ from the proof of Theorem~\ref{thm:conjecture-40}(ii) yields $B \equiv \kappa_z^x(B[x := z]) \to_\beta^* \kappa_z^x(N_z) \in \NF_\beta$ (as replacing free $z$ by bound $x$ never creates a $\lambda$-abstraction), whence $(M)_\lambda \to_\beta^* \lambda x.\, \kappa_z^x(N_z) \in \NF_\beta$.

Conversely, suppose $(M)_\lambda \to_\beta^* N \in \NF_\beta$. We prove $\mathrm{DeepNorm}(M)$ by strong induction on $\|N\|_\lambda$, noting that $\|B[x := z]\|_\lambda = \|B\|_\lambda$. By part~(ii) at $k = 0$, we have $M \in \SN = \WN$, so $M \redSs M^{\mathrm{nf}} \in \NF$; by rule~(step) and confluence $(M^{\mathrm{nf}})_\lambda \to_\beta^* N$, we may assume $M \equiv M^{\mathrm{nf}} \in \NF$. Every $M^{\mathrm{nf}} \in \TSV \cap \NF$ falls into three cases: if $M^{\mathrm{nf}} \equiv x \in \mathcal{V}$, rule~(var) applies; if $M^{\mathrm{nf}} \equiv A B$ with $\mathrm{hd}(A) = x \in \mathcal{V}$ and $A, B \in \NF$, then $(A)_\lambda$ never reduces to an abstraction, so $(A B)_\lambda \to_\beta^* N$ factors as $N \equiv N_A N_B$ with $(A)_\lambda \to_\beta^* N_A \in \NF_\beta$, $(B)_\lambda \to_\beta^* N_B \in \NF_\beta$, and $\|N_A\|_\lambda, \|N_B\|_\lambda < \|N\|_\lambda$, allowing rule~(app) to combine the induction hypotheses on $A$ and $B$; finally, if $\mathrm{hd}(M^{\mathrm{nf}}) = S$, normality forces $\ls(M^{\mathrm{nf}}) \le 2$ and hence $M^{\mathrm{nf}} \in \mathcal{P}_S$, so $(M^{\mathrm{nf}})_\lambda \to_\beta^* \lambda x.\, B$ and confluence forces $N \equiv \lambda x.\, N_0$ ($N_0 \in \NF_\beta$). Choosing a fresh variable $z \in \mathcal{V} \setminus \mathrm{FV}(M^{\mathrm{nf}})$ gives $(M^{\mathrm{nf}} z)_\lambda \to_\beta^* (\lambda x.\, N_0)\,z \to_\beta N_0[x := z] \in \NF_\beta$ with $\|N_0[x := z]\|_\lambda = \|N_0\|_\lambda < \|N\|_\lambda$, so the induction hypothesis gives $\mathrm{DeepNorm}(M^{\mathrm{nf}} z)$ and rule~(sat) yields $\mathrm{DeepNorm}(M^{\mathrm{nf}})$.
\end{proof}

\subsection{Open Question 34: Single-Constant Combinatory Saturation, the Unique \texorpdfstring{$\lambda\mathrm{I}$}{lambda-I} Sink, and Congruence Incomparability}\label{subsec:question-34}

Theorem~\ref{thm:conjectures-42-43-refuted}(iii) characterizes $\lambda\beta$-normalization of $(M)_\lambda$ for any $M \in \TSV$ internally within $\CLS$ via the hereditary fresh-variable saturation predicate $\mathrm{DeepNorm}(M)$ of \eqref{eq:deepnorm-rules}. Restricting to closed pure-$S$ terms $\TS$, Waldmann \cite{Waldmann2000} established that the intrinsic strong normalization language $\SN \subset \TS$ is recognized by the $39$-state $S$-algebra $\mathcal{A}_{\SN}$, whereas the $9$-leaf combinator $W_9 \equiv S(S(A A))S$ and the $10$-leaf normal form $N_{10} \equiv S(S(S A A))S$ of Proposition~\ref{prop:minimal-lambda-divergers} witness the strict inclusion $\{M \in \TS : (M)_\lambda \in \SN_\beta\} \subsetneq \SN$. Because $(M)_\lambda \in \Lambda_{\mathrm{I}}$ for every $M \in \TSV$, Lemma~\ref{lem:lambdaI-conservation}(iii) equates weak and strong $\beta$-normalization on $(M)_\lambda$. Concerning the tree-automaton complexity of this language, Barendregt, Endrullis, Klop, and Waldmann \cite[Question~34, p.~96]{BarendregtEndrullisKlopWaldmann2017} introduced
\begin{equation}\label{eq:l34-def}
  L_{34} \;:=\; \bigl\{M \in \TS : (M)_\lambda \in \SN_\beta\bigr\} \;=\; \bigl\{M \in \TS : (M)_\lambda \in \WN_\beta\bigr\},
\end{equation}
and asked whether membership in $L_{34}$ is decidable, noting empirical evidence that $L_{34}$ is not a regular tree language over $\{S, \cdot\}$.

By the Myhill--Nerode theorem for binary trees, $L_{34}$ is regular if and only if its syntactic congruence $\equiv_{34}$ has finite index $|\TS / {\equiv_{34}}| < \infty$, where for $M, N \in \TS$ we write
\begin{equation}\label{eq:equiv34-def}
  M \equiv_{34} N \quad:\Longleftrightarrow\quad \forall\,\text{closed pure-$S$ one-hole contexts } C[\square],\;\; \bigl(C[M] \in L_{34} \iff C[N] \in L_{34}\bigr).
\end{equation}
Analyzing $L_{34}$ directly presents two apparent obstacles: untyped $\lambda\beta$-reduction involves bound-variable capture avoidance and substitution under binders, while the combinatory characterization $\mathrm{DeepNorm}(M)$ in \eqref{eq:deepnorm-rules} introduces a fresh variable $z \notin \mathrm{FV}(U)$ at every saturation step $U \mapsto U z$, seemingly requiring an unbounded variable alphabet along deep syntactic branches.

Our next theorem removes this fresh-variable requirement completely. Let $\mathrm{CL}(S, x_0)$ denote the left-linear combinatory rewriting system obtained by adjoining to $\{S\}$ a single inert constant $x_0$, with term set $\TS(\{x_0\})$ and one-step reduction $\redS$ generated contextually by $S X Y Z \redS (X Z)(Y Z)$ ($x_0$ participates in no root rule). Define the variable-erasing applicative homomorphism $\pi_0 \colon \TSV \to \TS(\{x_0\})$ by
\[
  \pi_0(S) \;\equiv\; S, \qquad \pi_0(v) \;\equiv\; x_0 \quad (v \in \mathcal{V}), \qquad \pi_0(A B) \;\equiv\; \pi_0(A)\,\pi_0(B).
\]
On $\TS(\{x_0\})$, define \emph{single-constant hereditary saturation} $\mathrm{DeepNorm}_1(U)$ inductively as the least predicate closed under four monotone clauses: (i) if $U \redS U'$ and $\mathrm{DeepNorm}_1(U')$, then $\mathrm{DeepNorm}_1(U)$; (ii) $\mathrm{DeepNorm}_1(x_0)$ holds; (iii) if $\mathrm{hd}(A) = x_0$ and both $\mathrm{DeepNorm}_1(A)$ and $\mathrm{DeepNorm}_1(B)$ hold, then $\mathrm{DeepNorm}_1(A B)$ holds; and (iv) if $U$ is a partial $S$-application ($U \equiv S$, $S P$, or $S P Q$) and $\mathrm{DeepNorm}_1(U x_0)$ holds, then $\mathrm{DeepNorm}_1(U)$ holds. Notice that clause~(iv) feeds the same fixed atom $x_0$ at every saturation depth without any freshness side condition. Throughout what follows we retain the abbreviations $T \equiv S S$, $A \equiv S S S$, $V \equiv S(S T T)$, $W_9 \equiv S(S(A A))S$, $N_{10} \equiv S(S(S A A))S$, and $W_{\mathrm{Wal}} \equiv S V V$.

\begin{theorem}[Single-constant reduction, unique $\lambda\mathrm{I}$ sink, and congruence incomparability for Question~34]\label{thm:question-34-structure}
The language $L_{34} \subset \TS$ and its syntactic congruence $\equiv_{34}$ satisfy the following properties:
\begin{enumerate}[label=(\roman*),leftmargin=2.2em]
  \item \emph{Single-constant combinatory reduction:} For every open or closed pure-$S$ term $M \in \TSV$, one has $\mathrm{DeepNorm}(M)$ if and only if $\mathrm{DeepNorm}_1(\pi_0(M))$. In particular, for every closed term $M \in \TS$,
  \[
    M \in L_{34} \;\iff\; (M)_\lambda \in \SN_\beta \;\iff\; (M)_\lambda \in \WN_\beta \;\iff\; \mathrm{DeepNorm}_1(M).
  \]
  \item \emph{Unique two-sided absorbing sink $\infty_{34}$:} For every closed pure-$S$ one-hole context $C[\square]$ and every $M \in \TS$ with $(M)_\lambda \notin \SN_\beta$, one has $(C[M])_\lambda \notin \SN_\beta$. Consequently, $\infty_{34} := \TS \setminus L_{34}$ forms a single $\equiv_{34}$-congruence class and acts as the unique two-sided absorbing zero of the syntactic quotient magma $\TS / {\equiv_{34}}$: for all $M, N \in \infty_{34}$ and $P \in \TS$, $M \equiv_{34} N$ and $M P, P M \in \infty_{34}$.
  \item \emph{Mutual incomparability with intrinsic $S$-algebra congruences and four-way splitting of $q_6$:} The $9$-leaf combinators $A A A \notin \SN$ and $W_9 \in \SN$ satisfy $A A A \equiv_{34} W_9$ while $A A A \not\equiv_{\SN} W_9$, so ${\equiv_{34}} \not\subseteq {\equiv_{\SN}}$. Conversely, inside the single accepting state $q_6$ of $\mathcal{A}_{\SN}$ (and $p_{18}$ of $\mathcal{A}_P$), the four order-$2$ closed pure-$S$ normal forms
  \begin{equation}\label{eq:q34-u0123}
    u_0 \;\equiv\; S(S(S A)), \qquad u_1 \;\equiv\; V \;\equiv\; S(S T T), \qquad u_2 \;\equiv\; S V, \qquad u_3 \;\equiv\; S(S(S A A))
  \end{equation}
  all lie in $L_{34}$ (with $u_0$ and $u_2$ sharing the identical joint state $(q_6, h_9, p_{18})$ in $\mathcal{A}_{\SN} \times \mathcal{A}_{\HN} \times \mathcal{A}_P$), yet are pairwise $\equiv_{34}$-inequivalent under the three closed contexts $C_1[\square] \equiv \square \cdot S$, $C_2[\square] \equiv \square \cdot V$, and $C_3[\square] \equiv (S V) \cdot \square$ (and inequivalent to $\infty_{34}$ in the empty context $\square$); hence ${\equiv_{\SN \times \HN \times \Prod}} \not\subseteq {\equiv_{34}}$.
\end{enumerate}
\end{theorem}

\begin{proof}
For~(i), we first verify that $\pi_0$ both preserves and reflects one-step reduction $\redS$. Because $\pi_0$ is an applicative homomorphism fixing $S$, $\pi_0(S X Y Z) \equiv S\,\pi_0(X)\,\pi_0(Y)\,\pi_0(Z) \redS (\pi_0(X)\,\pi_0(Z))(\pi_0(Y)\,\pi_0(Z)) \equiv \pi_0((X Z)(Y Z))$, so structural induction gives $M \redS M'$ implies $\pi_0(M) \redS \pi_0(M')$. Conversely, suppose $\pi_0(M) \redS U'$. By definition of $\pi_0$, $\pi_0(E) \equiv S$ forces $E \equiv S$, $\pi_0(E) \equiv x_0$ forces $E \in \mathcal{V}$, and $\pi_0(E) \equiv P_1 P_2$ forces $E \equiv E_1 E_2$ with $\pi_0(E_i) \equiv P_i$. If $\pi_0(M) \redS U'$ is a root contraction $\pi_0(M) \equiv S X' Y' Z'$ with $U' \equiv (X' Z')(Y' Z')$, inverting $\pi_0$ three times along the left spine forces $M \equiv S X Y Z$ with $\pi_0(X) \equiv X'$, $\pi_0(Y) \equiv Y'$, and $\pi_0(Z) \equiv Z'$; because the pure-$S$ rule $S X Y Z \redS (X Z)(Y Z)$ is left-linear (no non-linear equality test between subterms appears on the left-hand side), $M$ contracts directly to $M' \equiv (X Z)(Y Z)$ with $\pi_0(M') \equiv U'$, even when distinct variables inside $X, Y, Z$ are identified to $x_0$ by $\pi_0$. Context steps lift directly by the induction hypothesis. Similarly, $\mathrm{hd}(M) \in \mathcal{V}$ holds if and only if $\mathrm{hd}(\pi_0(M)) = x_0$, and $M \in \mathcal{P}_S$ holds if and only if $\pi_0(M)$ is a partial $S$-application in $\TS(\{x_0\})$.

An immediate induction on \eqref{eq:deepnorm-rules} now maps $\mathrm{DeepNorm}(M)$ to $\mathrm{DeepNorm}_1(\pi_0(M))$, since $\pi_0(M z) \equiv \pi_0(M)\,x_0$ in rule~(sat). Conversely, induction on a derivation of $\mathrm{DeepNorm}_1(U)$ proves that every $M \in \TSV$ with $\pi_0(M) \equiv U$ satisfies $\mathrm{DeepNorm}(M)$: clause~(i) uses backward step reflection; clause~(ii) forces $M \in \mathcal{V}$ and applies rule~(var); clause~(iii) splits $M \equiv A B$ with $\mathrm{hd}(A) \in \mathcal{V}$ and applies rule~(app); and in clause~(iv), where $U \equiv \pi_0(M)$ is a partial $S$-application with $\mathrm{DeepNorm}_1(U x_0)$, we have $M \in \mathcal{P}_S$, and choosing a fresh variable $z \in \mathcal{V} \setminus \mathrm{FV}(M)$ (for instance, of index strictly exceeding the maximum variable index occurring in the finite tree $M$) gives $\pi_0(M z) \equiv \pi_0(M)\,\pi_0(z) \equiv U x_0$. The induction hypothesis on $U x_0$ then yields $\mathrm{DeepNorm}(M z)$, and rule~(sat) gives $\mathrm{DeepNorm}(M)$. Since $\pi_0(M) \equiv M$ for $M \in \TS$, Theorem~\ref{thm:conjectures-42-43-refuted}(iii) completes~(i).

For~(ii), in the untyped $\lambda$-calculus any binary application $P Q \in \SN_\beta$ satisfies $P \in \SN_\beta$ and $Q \in \SN_\beta$, by well-founded induction on reachable reducts of $P Q$. Because closed pure-$S$ one-hole contexts are built solely from $\square$ and left/right binary applications without outer $\lambda K$-abstractions, induction on $C[\square]$ shows that $(C[M])_\lambda \in \SN_\beta$ implies $(M)_\lambda \in \SN_\beta$. Contrapositively, $M \in \infty_{34}$ implies $C[M] \in \infty_{34}$ for every closed context $C[\square]$. Specializing to $\square \cdot P$ and $P \cdot \square$ shows that $\infty_{34}$ is a two-sided ideal, and for all $M, N \in \infty_{34}$ both $C[M] \in L_{34}$ and $C[N] \in L_{34}$ evaluate to false in every context $C[\square]$, so $M \equiv_{34} N$.

For~(iii), $A A A \notin \SN$ implies $(A A A)_\lambda \notin \SN_\beta$ by forward simulation, whereas Proposition~\ref{prop:minimal-lambda-divergers}(ii) gives $W_9 \in \SN$ and $(W_9)_\lambda \notin \SN_\beta$. Thus $A A A, W_9 \in \infty_{34}$, so $A A A \equiv_{34} W_9$ by~(ii), whereas $A A A \not\equiv_{\SN} W_9$ in the empty context $\square$. Next, each term in \eqref{eq:q34-u0123} has spine length at most~$2$ at every subterm and hence belongs to $\TS \cap \NF$. Writing $(q_i, h_j, p_k)$ for non-absorbing states of $\mathcal{A}_{\SN} \times \mathcal{A}_{\HN} \times \mathcal{A}_P$, bottom-up evaluation from $\sem{S} = (q_0, h_0, p_0)$, $\sem{T} = (q_1, h_1, p_1)$, and $\sem{A} = (q_7, h_6, p_3)$ gives $\sem{u_0} = (q_6, h_9, p_{18})$, $\sem{u_1} = (q_6, h_8, p_{18})$, $\sem{u_2} = (q_6, h_9, p_{18})$, and $\sem{u_3} = (q_6, h_{10}, p_{18})$; in particular, all four map to $q_6$ in $\mathcal{A}_{\SN}$ and $p_{18}$ in $\mathcal{A}_P$, and $u_0, u_2$ share $(q_6, h_9, p_{18})$. Applying part~(i), single-constant saturation reduces $u_0, u_1, u_2, u_3$ to pure $x_0$-headed normal trees, placing all four in $L_{34}$. Under $(C_1, C_2, C_3)$, saturation reduces $C_1[u_0], C_2[u_0], C_3[u_0]$ completely to inert $x_0$-trees, so $(C_1[u_0], C_2[u_0], C_3[u_0]) \in L_{34}^3$. For $u_1 \equiv V$, one has $V S, V V \in L_{34}$, whereas $C_3[u_1] \equiv S V V \equiv W_{\mathrm{Wal}} \notin L_{34}$ by Proposition~\ref{prop:minimal-lambda-divergers}(ii), separating $u_1$ from $u_0$. For $u_2 \equiv S V$, one has $S V S \in L_{34}$, whereas $C_2[u_2] \equiv S V V \equiv W_{\mathrm{Wal}} \notin L_{34}$, separating $u_2$ from both $u_0$ and $u_1$ (and demonstrating ${\equiv_{\SN \times \HN \times \Prod}} \not\subseteq {\equiv_{34}}$ via $u_0$ and $u_2$). For $u_3 \equiv S(S(S A A))$, already $C_1[u_3] \equiv S(S(S A A))S \equiv N_{10} \notin L_{34}$ by Proposition~\ref{prop:minimal-lambda-divergers}(i), separating $u_3$ from $u_0, u_1, u_2$. The resulting membership rows $(1, 1, 1)$, $(1, 1, 0)$, $(1, 0, *)$, and $(0, *, *)$ form an upper-triangular separation matrix, while the empty context $\square$ separates all four from $N_{10} \in \infty_{34}$.
\end{proof}

\begin{remark}[Exact single-constant frontier of Open Question~34]\label{rem:question-34-frontier}
Whether the cardinal $|\TS / {\equiv_{34}}|$ is finite (making $L_{34}$ a regular tree language and hence decidable in linear time, answering Question~34 affirmatively) or countably infinite (making $L_{34}$ non-regular, with general decidability still open) remains open. Theorem~\ref{thm:question-34-structure} isolates the exact combinatory content of this question: part~(i) eliminates $\lambda$-binders, $\alpha$-renaming, and fresh-variable bookkeeping altogether, reducing Question~34 to whether single-constant saturation $\mathrm{DeepNorm}_1$ in the two-atom left-linear system $\mathrm{CL}(S, x_0)$ induces finitely many contextual equivalence classes on $\TS$. Parts~(ii) and~(iii) explain why existing finite pure-$S$ algebras do not transfer directly: on the rejecting side, non-erasing $\lambda\mathrm{I}$ contexts collapse $\SN \setminus L_{34}$ and $\TS \setminus \SN$ into the single sink $\infty_{34}$, while on the accepting side, repeated probing by $x_0$ penetrates arbitrarily nested partial $S$-combinators and activates frozen latent subterms (such as $S A A$ inside $u_3$), shattering even single joint fibers such as $(q_6, h_9, p_{18})$ of $\mathcal{A}_{\SN} \times \mathcal{A}_{\HN} \times \mathcal{A}_P$.
\end{remark}

\subsection{Open Question 32: Small-Term \texorpdfstring{$\lambda\beta$}{lambda-beta}-Rigidity and Equational Non-Finite Axiomatizability}\label{subsec:question-32}

We now turn from tree-automaton complexity to equational completeness, comparing external $\lambda\beta$-convertibility $(M)_\lambda =_\beta (N)_\lambda$ on $\TSV$ with intrinsic combinatory convertibility $M \convS N$. Because every pure-$S$ contraction $S X Y Z \redS (X Z)(Y Z)$ is simulated by $\beta$-reduction, induction on conversion steps yields the one-way implication
\begin{equation}\label{eq:convS-implies-convBeta}
  M \convS N \;\implies\; (M)_\lambda =_\beta (N)_\lambda \qquad \bigl(M, N \in \TSV\bigr).
\end{equation}
Over the full combinatory basis $\{S, K\}$, intrinsic equality $=_{S,K}$ is likewise strictly weaker than $=_\beta$, yet Curry \cite{Curry1930,CurryFeys1958} and Barendregt \cite[Theorem~7.3.10]{Barendregt1984} proved that adjoining finitely many closed combinatory equations (Curry's five axioms) suffices to axiomatize $(M)_\lambda =_\beta (N)_\lambda$ across all open and closed $\mathrm{CL}(S, K)$-terms.

In $\CLS$, where the erasing combinator $K$ is absent, constant bracket abstractions $\lambda^* x.\, M$ ($x \notin \mathrm{FV}(M)$) cannot be synthesized internally. Waldmann \cite{Waldmann2000} observed that $\convS$ is strictly weaker than $=_\beta$ even on closed pure-$S$ normal forms, exhibiting a pair of closed normal forms of leaf sizes $8$ and $11$. Asking whether a finite Curry-style basis nevertheless exists for pure $S$, Barendregt, Endrullis, Klop, and Waldmann \cite[Question~32, p.~94]{BarendregtEndrullisKlopWaldmann2017} posed the problem of whether there is a finite set $\mathcal{E}$ of equations between pure-$S$ terms such that, for all $M, N \in \TSV$,
\begin{equation}\label{eq:q32-question}
  \CLS + \mathcal{E} \vdash M = N \;\iff\; (M)_\lambda =_\beta (N)_\lambda.
\end{equation}
Here, given any finite set of schematic equations $\mathcal{E} = \{P_1 = Q_1, \dots, P_r = Q_r\} \subset \TSV \times \TSV$ (whose members $P_i, Q_i$ may contain arbitrary free variables), equational derivability $\CLS + \mathcal{E} \vdash M = N$ is the least equivalence relation on $\TSV$ containing one-step reduction $A \redS B$ together with every substitution instance $P_i^\sigma = Q_i^\sigma$ ($\sigma \colon \mathcal{V} \to \TSV$) of $(P_i = Q_i) \in \mathcal{E}$, and closed under left and right application congruence ($A_1 = A_2$ implies $A_1 C = A_2 C$ and $C A_1 = C A_2$). We call $\mathcal{E}$ \emph{$\lambda\beta$-sound} if $(P)_\lambda =_\beta (Q)_\lambda$ for every $(P = Q) \in \mathcal{E}$; since $=_\beta$ is closed under substitution and application contexts, $\lambda\beta$-soundness of $\mathcal{E}$ guarantees that $\CLS + \mathcal{E} \vdash M = N$ implies $(M)_\lambda =_\beta (N)_\lambda$.

Two syntactic notions govern local rigidity under $\convS$ and $=_\beta$. First, a term $M \in \TSV$ is \emph{globally $\lambda\beta$-rigid} if $(M)_\lambda =_\beta (N)_\lambda$ holds if and only if $M \convS N$ for every $N \in \TSV$. Second, a subterm of the form $(X Z)(Y Z)$, whose left and right factors are applications sharing the identical right argument $Z$, is called a \emph{co-redex}; a term $W \in \TSV$ is an \emph{isolated normal form} if $W \in \NF$ and $W$ contains no co-redex subterm (cf.\ Subsection~\ref{subsec:goe-isolated}). Because every contraction $C[S X Y Z] \redS C[(X Z)(Y Z)]$ creates a co-redex $(X Z)(Y Z)$, an isolated normal form $W$ has neither forward reducts nor backward predecessors, and Church--Rosser confluence (Proposition~\ref{prop:cr-conservation}(i)) collapses its $\convS$-class to the singleton $[W]_{\convS} = \{W\}$.

\begin{proposition}[Universal $S$-free rigidity, small-term global rigidity, and uniqueness of the minimal $(8, 11)$ pair]\label{prop:q32-rigidity-waldmann-pair}
The following rigidity and separation laws hold on $\TSV$:
\begin{enumerate}[label=(\roman*),leftmargin=2.2em]
  \item \emph{Universal backward-cone rigidity of $S$-free trees:} Let $F \in \TSV$ be an $S$-free term (built solely from variables in $\mathcal{V}$ and binary application). For every open or closed pure-$S$ term $M \in \TSV$,
  \begin{equation}\label{eq:sfree-rigidity}
    (F)_\lambda =_\beta (M)_\lambda \;\iff\; M \redSs F.
  \end{equation}
  In particular, every variable $x \in \mathcal{V}$ satisfies $(x)_\lambda =_\beta (M)_\lambda$ if and only if $M \equiv x$, and every co-redex-free $S$-free term $F$ satisfies $(F)_\lambda =_\beta (M)_\lambda$ if and only if $M \equiv F$.
  \item \emph{Small-term global rigidity and the minimal $(8, 11)$ pairs:} The combinators $S$, $T \equiv S S$, $A \equiv S S S$, and every closed strongly normalizing term $M \in \TS \cap \SN$ with $|M| \le 7$ (comprising all $65$ closed terms with $|M| \le 6$ and all $89$ closed normal forms with $|M| \le 7$) are globally $\lambda\beta$-rigid in $\TSV$; across all $197$ closed pure-$S$ terms with $|M|, |N| \le 7$, $(M)_\lambda =_\beta (N)_\lambda$ holds if and only if $M \convS N$. At sizes $(8, 11)$, for any $x, y \in \TSV$, the two-variable schematic terms
  \begin{equation}\label{eq:open-8-11-pair}
    L(x, y) \;\equiv\; S\,(S\,T\,y)\,(S\,x\,y), \qquad R(x, y) \;\equiv\; S\,\bigl(S\,y\,(S\,x\,y)\bigr)\,\bigl(S\,(T\,x)\,y\bigr)
  \end{equation}
  satisfy $(L(x, y))_\lambda =_\beta (R(x, y))_\lambda$; when $x, y \in \mathcal{V}$, $L(x, y)$ ($8$ leaves) is an isolated normal form and $R(x, y)$ ($11$ leaves) lies in $\NF$, so $L(x, y) \not\convS R(x, y)$. Specializing to $x \equiv y \equiv S$ recovers Waldmann's $(8, 11)$ closed normal-form pair \cite{Waldmann2000}
  \begin{equation}\label{eq:waldmann-8-11-pair}
    U_1 \;\equiv\; S\,(S\,T\,S)\,A \quad (|U_1| = 8), \qquad U_2 \;\equiv\; S\,(T\,A)\,(S\,A\,S) \quad (|U_2| = 11),
  \end{equation}
  satisfying $(U_1)_\lambda =_\beta (U_2)_\lambda$, $[U_1]_{\convS} = \{U_1\}$, and $U_1 \not\convS U_2$; moreover, $\{U_1, U_2\}$ is the unique unordered pair of distinct closed pure-$S$ normal forms $M, N \in \TS \cap \NF$ with $|M|, |N| \le 11$ (among all $3{,}562$ such normal forms) satisfying $(M)_\lambda =_\beta (N)_\lambda$.
\end{enumerate}
\end{proposition}

\begin{proof}
For~(i), $M \redSs F$ implies $(M)_\lambda \to_\beta^* (F)_\lambda$, giving $(F)_\lambda =_\beta (M)_\lambda$. For the converse, we first establish quantitative invariants of $\beta$-reduction on $\lambda$-translations of pure-$S$ terms. For $E \in \Lambda$, define the variable-leaf count $\mathrm{vl}(E)$ ($\mathrm{vl}(v) = 1$ on bound or free variable occurrences, $\mathrm{vl}(\lambda z.\, P) = \mathrm{vl}(P)$, $\mathrm{vl}(P Q) = \mathrm{vl}(P) + \mathrm{vl}(Q)$), the binder count $\#_\lambda(E)$ ($\#_\lambda(v) = 0$, $\#_\lambda(\lambda z.\, P) = 1 + \#_\lambda(P)$, $\#_\lambda(P Q) = \#_\lambda(P) + \#_\lambda(Q)$), and the free-occurrence count $\#_v(E)$ for $v \in \mathcal{V}$. Let $\mathcal{I}_S \subset \Lambda_{\mathrm{I}}$ be the class of $\lambda$-terms $E$ such that at every abstraction subterm $\lambda z.\, P$ of $E$, the bound variable $z$ occurs $c_z(P) \ge 1$ times in $P$, $\#_\lambda(P) + 2 \le \mathrm{vl}(P)$, and either $P$ is itself an abstraction or $c_z(P) \ge 2$. Structural induction shows $\mathrm{vl}(E) - \#_\lambda(E) \ge 1$ for all $E \in \mathcal{I}_S$; since $(S)_\lambda \equiv \lambda x y z.\, (x z)(y z) \in \mathcal{I}_S$ (its three nested bodies have $(\#_\lambda, \mathrm{vl}) = (0, 4), (1, 4), (2, 4)$ and innermost multiplicity $c_z = 2$), every $M \in \TSV$ satisfies $(M)_\lambda \in \mathcal{I}_S$ with $|M| + \#_\lambda((M)_\lambda) = \mathrm{vl}((M)_\lambda)$ and $\#_v((M)_\lambda) = \#_v(M)$.

For any root $\beta$-step $(\lambda z.\, P)\,Q \to_\beta P[z := Q]$ inside $\mathcal{I}_S$, writing $c := c_z(P) = d + 1$ ($d \ge 0$) gives $\mathrm{vl}(P[z := Q]) + c = \mathrm{vl}(P) + c\,\mathrm{vl}(Q)$, $\#_\lambda(P[z := Q]) = \#_\lambda(P) + c\,\#_\lambda(Q)$, and $\#_v(P[z := Q]) = \#_v(P) + c\,\#_v(Q) \ge \#_v(P) + \#_v(Q)$. Subtracting the binder count from the variable-leaf count yields
\[
  \bigl(\mathrm{vl}(P[z := Q]) - \#_\lambda(P[z := Q])\bigr) - \bigl(\mathrm{vl}((\lambda z.\, P)\,Q) - \#_\lambda((\lambda z.\, P)\,Q)\bigr) \;=\; d\,\bigl(\mathrm{vl}(Q) - \#_\lambda(Q) - 1\bigr) \;\ge\; 0,
\]
and $\mathcal{I}_S$ is preserved under substitution and contextual $\beta$-steps. Induction along $(M)_\lambda \to_\beta^* B$ therefore yields, for every $M \in \TSV$,
\begin{equation}\label{eq:lambdaI-size-bound}
  |M| + \#_\lambda(B) \;\le\; \mathrm{vl}(B), \qquad \#_v(M) \;\le\; \#_v(B) \quad (v \in \mathcal{V}).
\end{equation}
From \eqref{eq:lambdaI-size-bound} we obtain constructive $\CLS$ normal-form extraction: whenever $(M)_\lambda \to_\beta^* B$ with $B \in \NF_\beta$, there exists $N \in \NF$ such that $M \redSs N$ and $(N)_\lambda \to_\beta^* B$. Indeed, if $M \notin \NF$, picking any step $M \redS M'$ gives $(M)_\lambda \to_\beta^+ (M')_\lambda$, and confluence of $\to_\beta^*$ with $B \in \NF_\beta$ forces $(M')_\lambda \to_\beta^* B$, whence $|M'| \le \mathrm{vl}(B)$ by \eqref{eq:lambdaI-size-bound}. By Proposition~\ref{prop:cr-conservation}(ii) and Theorem~\ref{thm:leaf-preserving}(i), each step $M \redS M'$ either strictly decreases $\mathrm{vl}(B) - |M| \in \mathbb{N}$ (when $|M| < |M'|$) or preserves $|M'| = |M|$ while strictly decreasing Bergstra's weight $\|M\|_{\mathrm{B}}$ by at least~$6$, so well-founded induction on $(\mathrm{vl}(B) - |M|, \|M\|_{\mathrm{B}}) \in \mathbb{N}^2$ reaches $N \in \NF$ with $(N)_\lambda \to_\beta^* B$.

Now let $F$ be $S$-free and suppose $(F)_\lambda =_\beta (M)_\lambda$. Because $(F)_\lambda \equiv F$ contains no $\lambda$-binders, $(F)_\lambda \in \NF_\beta$, so confluence gives $(M)_\lambda \to_\beta^* (F)_\lambda$ and normal-form extraction gives $M \redSs N \in \NF$ with $(N)_\lambda \to_\beta^* (F)_\lambda$. If $\mathrm{hd}(N) = S$, then $N \in \mathcal{P}_S$ reduces in $0, 1$, or $2$ weak-head $\beta$-steps to an abstraction, which is preserved under $\to_\beta^*$ and contradicts $(F)_\lambda$ not being an abstraction. Hence $\mathrm{hd}(N) \in \mathcal{V}$, so $(N)_\lambda$ has a neutral free-variable head that never creates a root $\beta$-redex: structural induction on $F$ immediately forces $N \equiv F$, establishing $M \redSs F$.

For~(ii), computing the $\beta$-normal forms $B_{M_0}$ of $M_0 \in \{S, T, A\}$ gives $(\#_\lambda, \mathrm{vl}) = (3, 4)$ for $S$, $(3, 5)$ for $T$ ($B_T \equiv \lambda y z w.\, (z w)((y z) w)$), and $(3, 6)$ for $A$ ($B_A \equiv \lambda z w.\, (z w)(\lambda w'.\, (z w')(w w'))$). Whenever $(M_0)_\lambda =_\beta (N)_\lambda$ for $N \in \TSV$, normal-form extraction and \eqref{eq:lambdaI-size-bound} yield $N \redSs N_0 \in \NF$ with $\#_v(N_0) \le \#_v(B_{M_0}) = 0$ (so $N_0 \in \TS$ is closed) and $|N_0| \le \mathrm{vl}(B_{M_0}) - \#_\lambda(B_{M_0}) \in \{1, 2, 3\}$; comparing the distinct $\beta$-normal forms of closed normal forms up to size~$3$ ($S, T, A$, and $S T$ whose $\beta$-normal form has three top-level $\lambda$-abstractions unlike $B_A$) forces $N_0 \equiv M_0$. Likewise, because the minimum leaf size in $\TS \cap \SN \setminus L_{34}$ is $|W_9| = 9 > 7$ (Proposition~\ref{prop:minimal-lambda-divergers}(iii)), every $M \in \TS \cap \SN$ with $|M| \le 7$ has a closed $\beta$-normal form $B_M$ satisfying $\mathrm{vl}(B_M) - \#_\lambda(B_M) \le 7$; any $N \in \TSV$ with $(M)_\lambda =_\beta (N)_\lambda$ reduces in $\CLS$ to a closed normal form of size at most~$7$, and the $89$ closed pure-$S$ normal forms of size at most~$7$ have pairwise distinct $\beta$-normal forms, forcing $N \convS M$. The only two non-$\SN$ closed terms of size~$7$, $S T S S S S$ and $A T S S$, satisfy $S T S S S S \redS A T S S$ and have non-normalizing $\lambda$-translations, so $=_\beta$ and $\convS$ coincide across all $197$ closed terms of size at most~$7$.

For \eqref{eq:open-8-11-pair}, pick a fresh variable $z \in \mathcal{V}$ not occurring in $x, y$. Five contractions in $\CLS$ give
\begin{align*}
  L(x, y)\,z &\;\redS\; (S T y z)(S x y z) \;\redS^2\; \bigl((T z)(y z)\bigr)\bigl((x z)(y z)\bigr) \\
  &\;\redS\; \Bigl(S\,(y z)\,\bigl(z\,(y z)\bigr)\Bigr)\bigl((x z)(y z)\bigr) \;\redS\; \Bigl((y z)\,\bigl((x z)(y z)\bigr)\Bigr)\Bigl(\bigl(z\,(y z)\bigr)\,\bigl((x z)(y z)\bigr)\Bigr),
\end{align*}
while six contractions give
\begin{align*}
  R(x, y)\,z &\;\redS\; \bigl(S y (S x y) z\bigr)\bigl(S(T x) y z\bigr) \;\redS^4\; \Bigl((y z)\,\bigl((x z)(y z)\bigr)\Bigr)\Bigl(\bigl((S z)(x z)\bigr)\,(y z)\Bigr) \\
  &\;\redS\; \Bigl((y z)\,\bigl((x z)(y z)\bigr)\Bigr)\Bigl(\bigl(z\,(y z)\bigr)\,\bigl((x z)(y z)\bigr)\Bigr),
\end{align*}
so $L(x, y)\,z \convS R(x, y)\,z$. Abstracting the fresh variable $z$ via $(S P Q)_\lambda \to_\beta^2 \lambda z.\, ((P)_\lambda z)((Q)_\lambda z)$ yields $(L(x, y))_\lambda =_\beta (R(x, y))_\lambda$. Direct subterm inspection shows that $L(x, y)$ (for $x, y \in \mathcal{V}$) and $U_1 \equiv L(S, S)$ have spine length at most~$2$ everywhere and unequal right arguments ($4$ leaves versus $1$ leaf) at their sole two-application node, making both isolated normal forms. Finally, among the $3{,}562$ closed normal forms of size at most~$11$, $25$ lie outside $L_{34}$ ($4$ of size~$10$ and $21$ of size~$11$, pairwise $\beta$-inconvertible via distinct finite B\"ohm-tree prefixes and $\beta$-inconvertible to every term in $L_{34}$), and pairwise comparison of the $\beta$-normal forms of the remaining $3{,}537$ closed normal forms yields $\mathrm{nf}_\beta((U_1)_\lambda) \equiv \mathrm{nf}_\beta((U_2)_\lambda)$ as the unique collision.
\end{proof}

Proposition~\ref{prop:q32-rigidity-waldmann-pair} shows that adjoining $U_1 = U_2$ (or even $L(x, y) = R(x, y)$) repairs all discrepancies up to $11$ leaves, yet tower-delayed distributivity generates infinitely many independent equational gaps at unbounded depths. For any $U \in \TSV$ and $k \in \mathbb{N}$, define the right-nested unary $S$-tower recursively by $S^0(U) \equiv U$ and $S^{k+1}(U) \equiv S(S^k(U))$ ($|S^k(U)| = k + |U|$), and for each $m \in \mathbb{N}$ and $Y \in \TSV$ set
\begin{equation}\label{eq:q32-mm-nm}
\begin{aligned}
  M_m(Y) &\;\equiv\; S\,\Bigl(S\,\bigl(S^{m+1}(T)\bigr)\,Y\Bigr)\,\Bigl(S\,\bigl(S^m(T)\bigr)\,Y\Bigr), \\
  N_m(Y) &\;\equiv\; S\,\Bigl(S\,\bigl(S^m(T T)\bigr)\,Y\Bigr)\,\Bigl(S\,\bigl(S^{m+1}(A)\bigr)\,Y\Bigr).
\end{aligned}
\end{equation}
When $Y \equiv y \in \mathcal{V}$ (or $Y \equiv S$), $|M_m(Y)| = 2m + 10$ whereas $|N_m(Y)| = 2m + 13$. Combining $S$-free probe rigidity, exact finite backward reduction cones, and deep $T$-core variable abstraction yields the negative resolution of Open Question~32.

\begin{theorem}[Equational non-finite axiomatizability of $\lambda\beta$-equality and negative resolution of Question~32]\label{thm:question-32}
Let $\mathcal{E} = \{P_1 = Q_1, \dots, P_r = Q_r\} \subset \TSV \times \TSV$ ($r \ge 0$) be any finite $\lambda\beta$-sound set of schematic pure-$S$ equations (allowing arbitrary free variables), and set $m := \max(\{0\} \cup \{|P_i|, |Q_i| : 1 \le i \le r\})$. For every variable $y \in \mathcal{V}$, the pair $(M_m(y), N_m(y))$ of \eqref{eq:q32-mm-nm} satisfies:
\begin{enumerate}[label=(\roman*),leftmargin=2.2em]
  \item The term $M_m(y)$ is an isolated normal form of leaf size $2m + 10$, $N_m(y) \in \NF$ has leaf size $2m + 13$, and $(M_m(y))_\lambda =_\beta (N_m(y))_\lambda$.
  \item The $\CLS + \mathcal{E}$ equational class of $M_m(y)$ collapses to the singleton $\{R \in \TSV : \CLS + \mathcal{E} \vdash M_m(y) = R\} = \{M_m(y)\}$, whence $\CLS + \mathcal{E} \nvdash M_m(y) = N_m(y)$ (and if $\mathcal{E} \subset \TS \times \TS$ consists of closed equations, $\CLS + \mathcal{E} \nvdash M_m(S) = N_m(S)$ despite $(M_m(S))_\lambda =_\beta (N_m(S))_\lambda$).
\end{enumerate}
Consequently, there exists no finite set $\mathcal{E}$ of pure-$S$ equations satisfying \eqref{eq:q32-question} for all $M, N \in \TSV$, resolving Open Question~32 of Barendregt, Endrullis, Klop, and Waldmann \cite[p.~94]{BarendregtEndrullisKlopWaldmann2017} in the negative.
\end{theorem}

\begin{proof}
We divide the argument into four steps.

\emph{Step~1: Fibonacci probe expansion and $(M_m(Y))_\lambda =_\beta (N_m(Y))_\lambda$.} For $A_0, B_0 \in \TSV$, define Fibonacci applicative words by $F_0(A_0, B_0) \equiv A_0$, $F_1(A_0, B_0) \equiv B_0$, and $F_{k+2}(A_0, B_0) \equiv F_k(A_0, B_0)\,F_{k+1}(A_0, B_0)$, so $F_k(B_0, A_0 B_0) \equiv F_{k+1}(A_0, B_0)$ by induction on $k$ (cf.\ Lemma~\ref{lem:visser-fibonacci}). Induction on $k$ from $S(S^{k'}(U)) A_0 B_0 \redS S^{k'}(U)\,B_0\,(A_0 B_0)$ gives $S^k(U)\,A_0\,B_0 \redS^k U\,F_k(A_0, B_0)\,F_{k+1}(A_0, B_0)$, and one outer contraction gives $S(S^k(U))\,Y\,X \redS^{k+1} U\,F_k(X, Y X)\,F_{k+1}(X, Y X)$. Fix a fresh variable $a \in \mathcal{V} \setminus \mathrm{FV}(Y)$, put $d_j \equiv F_j(a, Y a)$ and $E \equiv S\,d_{m+1}\,d_{m+2}$, and write $M_m(Y) \equiv S L_A R_A$ and $N_m(Y) \equiv S L_B R_B$. Applying the single-arm identity gives $L_A\,a \redS^{m+2} T\,d_{m+1}\,d_{m+2} \redS S\,d_{m+2}\,d_{m+3}$ and $R_A\,a \redS^{m+1} T\,d_m\,d_{m+1} \redS E$, whence $(L_A a)(R_A a) \redSs (S\,d_{m+2}\,d_{m+3})\,E \redS (d_{m+2} E)(d_{m+3} E)$. Similarly, $L_B\,a \redS^{m+1} (T T)\,d_m\,d_{m+1} \redS^2 (d_m d_{m+1})(T\,d_m\,d_{m+1}) \redS d_{m+2}\,E$ and $R_B\,a \redS^{m+2} A\,d_{m+1}\,d_{m+2} \redS^2 (d_{m+1} d_{m+2})(S\,d_{m+1}\,d_{m+2}) \equiv d_{m+3}\,E$, so $(L_B a)(R_B a) \redSs (d_{m+2} E)(d_{m+3} E)$. Abstracting the fresh variable $a$ as in Proposition~\ref{prop:q32-rigidity-waldmann-pair}(ii) yields $(M_m(Y))_\lambda =_\beta (N_m(Y))_\lambda$.

\emph{Step~2: Isolated-subterm equational rigidity criterion.} For $y \in \mathcal{V}$, every subterm of $M_m(y)$ and $N_m(y)$ has left-spine length at most~$2$, and at the root of $M_m(y) \equiv (S L_A)\,R_A$ the right arguments of $S L_A$ and $R_A$ have leaf sizes $m + 5 \ne 1 = |y|$, while all other compound subterms have left child $S$ or right child $y$; thus $M_m(y)$ is an isolated normal form. In general, for any isolated normal form $W \in \TSV$ (writing $U \preceq W$ for subterms), if every $(P = Q) \in \mathcal{E}$ and substitution $\sigma$ satisfy $(P^\sigma \preceq W \implies Q^\sigma \equiv P^\sigma)$ and $(Q^\sigma \preceq W \implies P^\sigma \equiv Q^\sigma)$, structural induction on $\CLS + \mathcal{E} \vdash A = B$ proves $(A \preceq W \implies B \equiv A) \land (B \preceq W \implies A \equiv B)$ (reduction steps are vacuous on isolated subterms, and congruence steps preserve subterm containment). By symmetry of $(P = Q)$, it suffices to show that whenever $(P)_\lambda =_\beta (Q)_\lambda$, $|P| \le m$, and $P^\sigma \preceq M_m(y)$, one has $Q^\sigma \equiv P^\sigma$.

\emph{Step~3: Finite backward-cone rigidity of proper subterms and the variable-core template.} Every subterm $U \preceq M_m(y)$ falls into six classes: (1) $S$; (2) $y$; (3) $S^j(T)$ ($0 \le j \le m + 2$); (4) $P_k \equiv S(S^k(T))y$ ($k \in \{m, m + 1\}$); (5) $S P_{m+1}$; or (6) $M_m(y)$ itself. Classes~(1)--(2) are globally rigid by Proposition~\ref{prop:q32-rigidity-waldmann-pair}. For classes~(3)--(5) and the two-variable core template
\[
  A_{\mathrm{var}}(m + 1, u, y) \;\equiv\; S\,\Bigl(S\,\bigl(S^{m+1}(u)\bigr)\,y\Bigr)\,\Bigl(S\,\bigl(S^m(u)\bigr)\,y\Bigr) \qquad (u, y \in \mathcal{V},\; u \ne y),
\]
we prove exact syntactic rigidity ($(U)_\lambda =_\beta (R)_\lambda$ implies $R \equiv U$) via $S$-free probes. Note that for distinct variables $a, b$ (and $y, a$), the $S$-free words $F_k(a, b)$ ($k \ge 1$) and $d_k \equiv F_k(a, y a)$ ($k \ge 0$, with $k + 1 \le |d_k| < |d_{k+1}|$ and $|d_{k+1}| \ge 2$) have strictly increasing leaf counts and hence contain no co-redexes.
\begin{itemize}[leftmargin=1.8em]
  \item For $S^k(T)$: if $(S^k(T))_\lambda =_\beta (R)_\lambda$, choose distinct $a, b, c \in \mathcal{V}$ and set $F_j \equiv F_j(a, b)$. Then $S^k(T)\,a\,b\,c \redSs (F_{k+1} c)(F_{k+2} c) \equiv: \Theta_1$, which is $S$-free, so $R\,a\,b\,c \redSs \Theta_1$ by Proposition~\ref{prop:q32-rigidity-waldmann-pair}(i). Tracing backward one-step predecessors from $\Theta_1$ (using co-redex-freeness of $F_j$ and disjointness of closed right subterms of $S^j(T)$ from variable-bearing $F_r$) closes the finite backward cone at $\mathcal{C}_1 = \{\Theta_1,\, S\,F_{k+1}\,F_{k+2}\,c\} \cup \{S^j(T)\,F_r\,F_{r+1}\,c : j + r = k\}$. Matching $R\,a\,b\,c \in \mathcal{C}_1$ against $|b| = 1$ forces $r = 0$ and $j = k$, whence $R \equiv S^k(T)$. (For closed $\mathcal{E}$, all subterms of $M_m(S)$ of size $\le m$ have the form $S$ or $S^j(T)$, already completing the closed-equation case.)
  \item For $P_k \equiv S(S^k(T))y$ and $S P_{m+1}$: picking distinct $a, b \in \mathcal{V} \setminus \{y\}$ gives $P_k\,a\,b \redSs (d_{k+1} b)(d_{k+2} b) \equiv: \Theta_2$, whose backward cone $\{\Theta_2,\, S\,d_{k+1}\,d_{k+2}\,b\} \cup \{S^j(T)\,d_r\,d_{r+1}\,b : j + r = k\} \cup \{P_k\,a\,b\}$ contains $P_k\,a\,b$ as the unique term whose second spine argument from the right has leaf size $1 = |a|$ (since $|d_{r+1}| \ge 2$), forcing $R \equiv P_k$. Likewise, for distinct $w, a \in \mathcal{V} \setminus \{y\}$ with $B \equiv w a$, $(S P_{m+1})\,w\,a \redSs (d_{m+2} B)(d_{m+3} B) \equiv: \Theta_3$ has backward cone $\{\Theta_3,\, S\,d_{m+2}\,d_{m+3}\,B\} \cup \{S^j(T)\,d_r\,d_{r+1}\,B : j + r = m + 1\} \cup \{(P_{m+1} a)\,B,\, (S P_{m+1})\,w\,a\}$, where only $(S P_{m+1})\,w\,a$ has outer right argument of leaf size $1 = |a| < 2 = |B|$, forcing $R \equiv S P_{m+1}$.
  \item For $A_{\mathrm{var}}(m + 1, u, y)$ ($u \ne y$): picking fresh $a \notin \{u, y\}$ gives $A_{\mathrm{var}}(m + 1, u, y)\,a \redS^{2m+4} ((u\,d_{m+1})\,d_{m+2})((u\,d_m)\,d_{m+1}) \equiv: \Theta_4$, so $(A_{\mathrm{var}}(m + 1, u, y))_\lambda =_\beta (R)_\lambda$ implies $R\,a \redSs \Theta_4$. Since $u$ occurs in $S^j(u)$ and never in $d_r$, each arm $(u\,d_k)\,d_{k+1}$ has backward cone $\mathcal{D}_k = \{S(S^k(u))\,y\,a\} \cup \{(S^j(u)\,d_r)\,d_{r+1} : j + r = k\}$ (all of leaf size $\ge 4$). Root-uncontracting $X_L \cdot X_R \in \mathcal{D}_{m+1} \cdot \mathcal{D}_m$ requires identical right arguments, yielding only $A_{\mathrm{var}}(m + 1, u, y)\,a$ (at shared argument $a$) and $\{S(S^{m+1-r}(u)\,d_r)(S^{m-r}(u)\,d_r)\,d_{r+1} : 0 \le r \le m\}$ (at shared argument $d_{r+1}$, by injectivity of $r \mapsto |d_{r+1}|$), none of which contains a co-redex. Among this closed backward cone $\mathcal{C}_4$, only $A_{\mathrm{var}}(m + 1, u, y)\,a$ has outer right argument of size $1 = |a|$, forcing $R \equiv A_{\mathrm{var}}(m + 1, u, y)$.
\end{itemize}

\emph{Step~4: Deep $T$-core variable abstraction at the root.} Finally, suppose $(P)_\lambda =_\beta (Q)_\lambda$, $|P| \le m$, and $P^\sigma \equiv M_m(y)$ at the root. Choose two distinct fresh variables $x_0, x_2 \in \mathcal{V} \setminus \{y\}$ (so $x_0 \notin \mathrm{FV}(M_m(y))$ and $x_2 \notin \{x_0, y\}$). Define the $T$-core replacement map $\rho_{x_0} \colon \TSV \to \TSV$ recursively by $\rho_{x_0}(S) \equiv S$, $\rho_{x_0}(v) \equiv v$, $\rho_{x_0}(S S) \equiv x_0$, and $\rho_{x_0}(A B) \equiv \rho_{x_0}(A)\,\rho_{x_0}(B)$ whenever $A B \not\equiv S S$. The only subterms of $M_m(y)$ literally equal to $S S$ are the two $T$-cores at the bottoms of the two $S$-towers, which sit at binary application depths $m + 4$ and $m + 3$ from the root of $M_m(y)$, so $\rho_{x_0}(M_m(y)) \equiv A_{\mathrm{var}}(m + 1, x_0, y)$. Because $|P| \le m < m + 3$, every internal application node of the pattern $P$ lies at depth at most $|P| - 1 < m$, so no application node of $P$ can assemble $S S$ in $P^\sigma \equiv M_m(y)$; both $T$-cores lie strictly inside substitution images $\sigma(v)$, giving $\rho_{x_0}(P^\sigma) \equiv P^{\tau_0}$ where $\tau_0(v) \equiv \rho_{x_0}(\sigma(v))$.

Define the modified substitution $\tau \colon \mathcal{V} \to \TSV$ by $\tau(v) \equiv \rho_{x_0}(\sigma(v))$ for $v \in \mathrm{FV}(P)$ and $\tau(v) \equiv x_2$ for $v \notin \mathrm{FV}(P)$. Then $P^\tau \equiv P^{\tau_0} \equiv \rho_{x_0}(M_m(y)) \equiv A_{\mathrm{var}}(m + 1, x_0, y)$, and substitution invariance of $=_\beta$ gives $(A_{\mathrm{var}}(m + 1, x_0, y))_\lambda \equiv (P^\tau)_\lambda =_\beta (Q^\tau)_\lambda$. By Step~3, syntactic rigidity of $A_{\mathrm{var}}(m + 1, x_0, y)$ forces $Q^\tau \equiv A_{\mathrm{var}}(m + 1, x_0, y)$. Because $x_2$ does not occur in $A_{\mathrm{var}}(m + 1, x_0, y)$, this immediately implies $\mathrm{FV}(Q) \subseteq \mathrm{FV}(P)$. Letting $\psi \colon \mathcal{V} \to \TSV$ be the restoration substitution mapping $x_0 \mapsto T \equiv S S$ and fixing all other variables, every $v \in \mathrm{FV}(Q) \subseteq \mathrm{FV}(P)$ satisfies $x_0 \notin \mathrm{FV}(\sigma(v))$ (as $\sigma(v) \preceq M_m(y)$) and hence $(\rho_{x_0}(\sigma(v)))^\psi \equiv \sigma(v)$. Applying $\psi$ to $Q^\tau \equiv A_{\mathrm{var}}(m + 1, x_0, y)$ yields
\[
  Q^\sigma \;\equiv\; (Q^\tau)^\psi \;\equiv\; \bigl(A_{\mathrm{var}}(m + 1, x_0, y)\bigr)^\psi \;\equiv\; M_m(y) \;\equiv\; P^\sigma.
\]
By Step~2, $[M_m(y)]_{\CLS + \mathcal{E}} = \{M_m(y)\}$, and since $|N_m(y)| = 2m + 13 \ne 2m + 10 = |M_m(y)|$, we conclude $\CLS + \mathcal{E} \nvdash M_m(y) = N_m(y)$. Finally, any $\mathcal{E}$ satisfying \eqref{eq:q32-question} is necessarily $\lambda\beta$-sound (taking identity substitutions), contradicting $(M_m(y))_\lambda =_\beta (N_m(y))_\lambda$.
\end{proof}

\section{Fixed-Order Dynamics, Super-Pushdown Traces, and Clone-Theoretic Cellular Barriers}
\label{sec:fixed-cellular}

Fixed deterministic reduction strategies---leftmost-outermost ($\to_{\mathrm{lo}}$), rightmost-outermost ($\to_{\mathrm{ro}}$), leftmost-innermost ($\to_{\mathrm{li}}$), rightmost-innermost ($\to_{\mathrm{ri}}$), shallowest-leftmost breadth-first ($\to_{\mathrm{bfs}}$), and Gross--Knuth parallel reduction ($\to_{\mathrm{gk}}$)---govern step-by-step evaluation in $\mathrm{CL}(S)$, while clone-theoretic monotonicity separates the semantic $S$-algebras $\mathcal{A}_{\mathrm{HN}}$, $\mathcal{A}_{\mathrm{SN}}$, and $\mathcal{A}_P$ from normal-form cellular computation.

\subsection{Pure-Line Strategy Coincidence, Quadratic Automaton Traces, and Innermost \texorpdfstring{$A A A$}{AAA}}
\label{subsec:pure-line-zachos}

Let $\#\mathrm{Redex}(M) \in \mathbb{N}$ denote the number of redex occurrences $S X Y Z$ in $M$, and let $\mathcal{N}_2 = \{ S X Y : X, Y \in \mathcal{N} \} \subset \mathcal{N}$ denote the set of Spine-$2$ closed normal forms. Following~\cite{BarendregtEndrullisKlopWaldmann2017}, a term $M$ is \emph{pure-line} along a trajectory if every visited reduct satisfies $\#\mathrm{Redex}(M_t) = 1$.

\begin{proposition}[Pure-line six-strategy coincidence, quadratic traces, and innermost $A A A$]
\label{prop:pure-line-and-aaa}
\begin{enumerate}[label=(\roman*)]
\item \emph{(Six-strategy coincidence on single-redex terms)} For every $M \in \mathcal{T}_S(\mathcal{V})$ satisfying $\#\mathrm{Redex}(M) = 1$, all six deterministic evaluators $\to_{\mathrm{lo}}, \to_{\mathrm{ro}}, \to_{\mathrm{li}}, \to_{\mathrm{ri}}, \to_{\mathrm{bfs}}, \to_{\mathrm{gk}}$ select the unique redex of $M$ and yield the identical one-step reduct $M'$.
\item \emph{(Super-pushdown quadratic checkpoint language on $Z_{\mathrm{Zac}}$)} Along the reduction sequence $Z_{\mathrm{Zac}} = T_0 \to T_1 \to T_2 \to \cdots$ starting from Zachos's $10$-leaf combinator $Z_{\mathrm{Zac}} \equiv G_0 G_0 \equiv (S T T)(S T T)$, every reduct $T_t$ satisfies $\#\mathrm{Redex}(T_t) = 1$ (so all six strategies generate the identical infinite sequence $(T_t)_{t \ge 0}$). Moreover, there is an explicit $8$-state deterministic bottom-up tree automaton $\mathcal{B}_Z = (\{b_0, \dots, b_7\}, b_0, \odot_{\mathcal{B}_Z}, \{b_6\})$ such that for all $t \in \mathbb{N}$,
\begin{equation}
\label{eq:zachos-milestones}
\llbracket T_t \rrbracket_{\mathcal{B}_Z} = b_6
\;\Longleftrightarrow\;
t \in \Bigl\{ \frac{n(n+3)}{2} : n \in \mathbb{N} \Bigr\}
\;=\; \{0, 2, 5, 9, 14, 20, 27, \dots\}
\end{equation}
(shifted by $+1$ to the triangular numbers $\binom{n+2}{2} = \{1, 3, 6, 10, 15, \dots\}$ when starting from the $9$-leaf precursor $Z_0 \equiv S(S T)(S T)T$, while on $D_{\mathrm{seed}} \equiv Z_0 S$ an explicit $68$-state automaton $\mathcal{B}_{\mathrm{quad}}$ recognizes $\{2n^2 + 3n + 1 : n \in \mathbb{N}\}$ under $\to_{\mathrm{lo}}$). By Parikh's theorem~\cite{Parikh1966}, the unary trace language $\{ 1^t : \llbracket T_t \rrbracket_{\mathcal{B}_Z} = b_6 \}$ is not context-free and cannot be generated by any pushdown or one-counter machine.
\item \emph{(Closed-form exponential size growth of $A A A$ under $\to_{\mathrm{li}}$, versus $\to_{\mathrm{ro}}$ trapping)} Under rightmost-outermost reduction $\to_{\mathrm{ro}}$, any non-root-redex application $L R$ with $R \in \mathcal{T}_S \setminus \mathrm{SN}$ traps forever on the right: $(L R)_t \equiv L\,R_t$ for all $t \ge 0$. Under leftmost-innermost reduction $\to_{\mathrm{li}}$, every non-normalizing closed term reduces once its minimal leftmost-innermost $\infty_{\mathrm{SN}}$-subterm freezes its outer context to iterating the normal-form map $F_{\mathrm{li}} : \mathcal{N}_2 \times \mathcal{N} \to \mathcal{N}_2 \times \mathcal{N}$ on critical pairs $(L, Z)$ ($\llbracket L Z \rrbracket_{\mathrm{SN}} = \infty_{\mathrm{SN}}$). For $A A A \equiv (S S S)(S S S)(S S S)$, setting $M_1 \equiv S A A$, $H_0 \equiv S A M_1$, and $H_{n+1} \equiv S H_n (S H_n) \in \mathcal{N}_2$ ($n \ge 0$), a $4$-macro-step transient from $(S A (S A), A)$ reaches $(H_1, H_1)$, after which every cycle $n \ge 1$ executes the $(n+2)$-macro-step loop $(H_n, H_n) \mapsto^n (H_0, H_n) \mapsto (M_1, H_n) \mapsto (H_{n+1}, H_{n+1})$ with exact exponential leaf count $|H_n| = 13 \cdot 2^n - 2$.
\end{enumerate}
\end{proposition}

\begin{proof}
For~(i), structural induction on $M$ shows that $\#\mathrm{Redex}(M) = 0$ if and only if $M \in \mathrm{NF}$, on which every strategy is idle ($\mathrm{gk}(M) = M$). When $\#\mathrm{Redex}(M) = 1$, either $M \equiv S X Y Z$ is itself a root redex with $X, Y, Z \in \mathrm{NF}$---where $\to_{\mathrm{lo}}, \to_{\mathrm{ro}}, \to_{\mathrm{bfs}}$ contract the root directly, $\to_{\mathrm{li}}, \to_{\mathrm{ri}}$ find $S X Y, Z \in \mathrm{NF}$ and contract the root, and $\mathrm{gk}(S X Y Z) = (\mathrm{gk}(X)\,\mathrm{gk}(Z))(\mathrm{gk}(Y)\,\mathrm{gk}(Z)) = (X Z)(Y Z)$---or $M \equiv L R$ is not a root redex and $(\#\mathrm{Redex}(L), \#\mathrm{Redex}(R)) \in \{(1, 0), (0, 1)\}$, where the idle child contains no redexes and the induction hypothesis equates all six evaluators on the active child.

For~(ii), recall $G_0 \equiv S T T$ and $G_{k+1} \equiv T G_k$, which contain no redexes ($\#\mathrm{Redex}(G_k) = 0$), whereas $\#\mathrm{Redex}(G_k G_m) = 1$ for all $k, m \in \mathbb{N}$ (at the root $S T T G_m$ when $k = 0$, and at $S S G_{k-1} G_m$ when $k \ge 1$). For any sequence of indices $\vec{m} = (m_1, \dots, m_r)$, define the normal-form wrapper $W_{\vec{m}}(C) \equiv S G_{m_1} (S G_{m_2} (\cdots (S G_{m_r} C)\cdots))$; since $\#\mathrm{Redex}(S G_m) = 0$ and $\mathrm{ls}(S G_m C) = 2$, we have $\#\mathrm{Redex}(W_{\vec{m}}(C)) = \#\mathrm{Redex}(C)$. Starting at $W_{\vec{m}}(G_0 G_n)$, one contraction yields $W_{\vec{m}}(G_{n+1} G_{n+1})$, and $n + 1$ subsequent contractions $G_{k+1} G_{n+1} \equiv S S G_k G_{n+1} \to_S S G_{n+1} (G_k G_{n+1})$ ($k = n, \dots, 0$) yield $W_{\vec{m} \cdot (n+1)^{n+1}}(G_0 G_{n+1})$ in $(n + 2)$ steps total. Summing $\sum_{j=0}^{n-1} (j + 2) = \frac{1}{2}n(n+3)$ from $T_0 \equiv G_0 G_0$ ($n = 0$, empty $\vec{m}$) shows $\#\mathrm{Redex}(T_t) = 1$ for all $t \ge 0$, with $T_{n(n+3)/2} \equiv W_{\vec{\mu}_n}(G_0 G_n)$ and intermediate steps $T_{n(n+3)/2 + j} \equiv W_{\vec{\mu}'}(G_{n+2-j} G_{n+1})$ ($1 \le j \le n + 1$). Define $\odot_{\mathcal{B}_Z}$ on $\{b_0, \dots, b_7\}$ by $b_0 \odot b_0 = b_1$ ($T$), $b_0 \odot b_1 = b_2$ ($S T$), $b_2 \odot b_1 = b_3$ ($G_0$), $b_1 \odot b_3 = b_1 \odot b_4 = b_4$ ($G_{k+1}$ for $k \ge 0$), $b_0 \odot b_3 = b_0 \odot b_4 = b_5$ ($S G_m$), $b_3 \odot b_3 = b_3 \odot b_4 = b_6$ ($G_0 G_n$), $b_5 \odot b_6 = b_6$ (preserving $b_6$ across $W_{\vec{m}}$), and $b_7$ elsewhere (so $G_{k+1} G_m \mapsto b_4 \odot b_{3,4} = b_7$ and $b_5 \odot b_7 = b_7$). Then $\llbracket T_t \rrbracket_{\mathcal{B}_Z} = b_6$ holds if and only if the core is $G_0 G_n$, proving~\eqref{eq:zachos-milestones}. Because consecutive gaps $\frac{1}{2}(n+1)(n+4) - \frac{1}{2}n(n+3) = n + 2 \to \infty$ are unbounded while the set is infinite, $\{ \frac{1}{2}n(n+3) : n \in \mathbb{N} \}$ is not a finite union of arithmetic progressions (not semilinear); by Parikh's theorem~\cite{Parikh1966}, $\{ 1^t : \llbracket T_t \rrbracket_{\mathcal{B}_Z} = b_6 \}$ is not context-free.

For~(iii), if $L R$ is not a root redex and $R \in \mathcal{T}_S \setminus \mathrm{SN}$, Theorem~\ref{thm:waldmann-sn} ensures $R$ is not in normal form, so $\to_{\mathrm{ro}}$ steps inside $R \to_{\mathrm{ro}} R'$ while preserving non-root-redex shape and $\llbracket R' \rrbracket_{\mathrm{SN}} = \infty_{\mathrm{SN}}$ inductively. Under $\to_{\mathrm{li}}$, $A A A \equiv ((S S S) A) A \to_{\mathrm{li}} (S A (S A)) A$ reaches the critical pair $(L_0, Z_0) = (S A (S A), A)$ with $\llbracket L_0 \rrbracket_{\mathrm{SN}} = q_{29}$, $\llbracket A \rrbracket_{\mathrm{SN}} = q_7$, and $q_{29} \odot_{\mathrm{SN}} q_7 = \infty_{\mathrm{SN}}$. Since $A U \equiv S S S U \to_{\mathrm{li}} S U (S U) \in \mathcal{N}_2$ for every $U \in \mathcal{N}$, induction on $k \ge 0$ gives $\llbracket H_0 \rrbracket_{\mathrm{SN}} = q_{29}$ and $\llbracket H_{k+1} \rrbracket_{\mathrm{SN}} = q_{30}$ (from $\llbracket M_1 \rrbracket_{\mathrm{SN}} = q_{34}$, $\llbracket H_0 \rrbracket_{\mathrm{SN}} = (q_0 \odot_{\mathrm{SN}} q_7) \odot_{\mathrm{SN}} q_{34} = q_{29}$, and $q_0 \odot_{\mathrm{SN}} q_{29} = q_0 \odot_{\mathrm{SN}} q_{30} = q_6$ with $q_6 \odot_{\mathrm{SN}} q_6 = q_{30}$), while $\{q_{29}, q_{30}\} \odot_{\mathrm{SN}} \{q_{29}, q_{30}\} = \{\infty_{\mathrm{SN}}\}$ and $q_{34} \odot_{\mathrm{SN}} \{q_{29}, q_{30}\} = \{\infty_{\mathrm{SN}}\}$. After $(S A (S A), A) \mapsto (S A (S A), M_1) \mapsto (S M_1 (S M_1), H_0) \mapsto (M_1, H_0) \mapsto (H_1, H_1)$, each $(H_k, H_n)$ with $k, n \ge 1$ has left contraction child $H_{k-1} H_n$ of state $\llbracket H_{k-1} \rrbracket_{\mathrm{SN}} \odot_{\mathrm{SN}} q_{30} = \infty_{\mathrm{SN}}$ (Case~1, peeling $H_k \mapsto H_{k-1}$ in $n$ steps to $(H_0, H_n)$); at $(H_0, H_n) = (S A M_1, H_n)$, the left contraction child satisfies $A H_n \to_{\mathrm{li}} H_{n+1} \in \mathcal{N} \subseteq \mathrm{SN}$ while $M_1 H_n$ has state $q_{34} \odot_{\mathrm{SN}} q_{30} = \infty_{\mathrm{SN}}$ (Case~2, advancing to $(M_1, H_n) = (S A A, H_n)$); and at $(S A A, H_n)$, both $A H_n$ children normalize in one step to $H_{n+1}$, yielding $(H_{n+1}, H_{n+1})$ (Case~3). Finally, $|H_0| = 1 + 3 + (1 + 3 + 3) = 11$ and $|H_{n+1}| = 2|H_n| + 2$, solving to $|H_n| = 13 \cdot 2^n - 2$.
\end{proof}

\subsection{The Sharp Five-Step Innermost Case-3 Ceiling}
\label{subsec:case3-ceiling}

Along any innermost reduction ($\to_{\mathrm{li}}$ or $\to_{\mathrm{ri}}$) of a non-normalizing closed $S$-term, each macro-step on the active critical pair $(L, Z) \in \mathcal{N}_2 \times \mathcal{N}$ ($L \equiv S X Y$, $\llbracket L Z \rrbracket_{\mathrm{SN}} = \infty_{\mathrm{SN}}$) falls into Case~1 ($\llbracket X Z \rrbracket_{\mathrm{SN}} = \infty_{\mathrm{SN}}$, descending to $(X, Z)$), Case~2 ($\llbracket X Z \rrbracket_{\mathrm{SN}} \neq \infty_{\mathrm{SN}}$ and $\llbracket Y Z \rrbracket_{\mathrm{SN}} = \infty_{\mathrm{SN}}$, descending to $(Y, Z)$ under $\to_{\mathrm{li}}$), or \emph{Case~3}, in which both $X Z$ and $Y Z$ are strongly normalizing---normalizing to $L' = \mathrm{nf}(X Z) \in \mathcal{N}_2$ and $Z' = \mathrm{nf}(Y Z) \in \mathcal{N}$---while $\llbracket L' Z' \rrbracket_{\mathrm{SN}} = \infty_{\mathrm{SN}}$. By global top-termination (Theorem~\ref{thm:waldmann-top-term}), Case~3 cannot occur infinitely many times consecutively, because consecutive Case~3 steps contract repeatedly at the root of the fixed minimal non-$\mathrm{SN}$ locus without descending into a proper subterm. Although the abstract Case-3 transition graph over state pairs $(u, v) \in F_{\mathrm{SN}}^2$ contains directed cycles---such as $(q_{25}, q_{25}) \to (q_{25}, q_{30}) \to (q_{30}, q_{30}) \to (q_{30}, q_6) \to (q_{29}, q_{25}) \to (q_{25}, q_{25})$---refining $L$ by a two-level normal-form descriptor yields a universal uniform bound of $5$ consecutive Case-3 steps across all of $\mathrm{CL}(S)$, dropping to $3$ asymptotically.

\begin{theorem}[Sharp consecutive Case-3 ceiling under $\to_{\mathrm{li}}$ and $\to_{\mathrm{ri}}$]
\label{thm:case3-ceiling}
Under both leftmost-innermost ($\to_{\mathrm{li}}$) and rightmost-innermost ($\to_{\mathrm{ri}}$) evaluation:
\begin{enumerate}[label=(\roman*)]
\item \emph{(Universal five-step ceiling)} No closed $S$-term admits $6$ consecutive Case-3 macro-steps: there exists no chain $(L_0, Z_0) \mapsto_{\mathrm{C3}} (L_1, Z_1) \mapsto_{\mathrm{C3}} \dots \mapsto_{\mathrm{C3}} (L_6, Z_6)$.
\item \emph{(Sharpness at five steps)} Five consecutive Case-3 macro-steps are achieved by the $9$-leaf witness $S(S(S A)) S S$ ($A \equiv S S S$), starting from $L_0 \equiv S(S(S A)) S$ and $Z_0 \equiv S$.
\item \emph{(Monotonicity of right-spine depth and asymptotic ceiling $3$)} Every Case-3 step $(L, Z) \mapsto_{\mathrm{C3}} (L', Z')$ satisfies $\mathrm{rs}(Z') \ge \mathrm{rs}(Z) + 1$ and $\mathrm{rs}(L') \ge \mathrm{rs}(Z) + 1$. Once $\mathrm{rs}(Z) \ge 2$ (which holds after at most two Case-3 steps along any trajectory), at most $3$ consecutive Case-3 macro-steps can occur.
\end{enumerate}
\end{theorem}

\begin{proof}
Because $\mathrm{CL}(S)$ is confluent and $X Z, Y Z \in \mathrm{SN}$ in Case~3 have unique normal forms $L' = \mathrm{nf}(X Z)$ and $Z' = \mathrm{nf}(Y Z)$ independent of whether $X Z$ and $Y Z$ are normalized via $\to_{\mathrm{li}}$ or $\to_{\mathrm{ri}}$, the Case-3 relation $(L, Z) \mapsto_{\mathrm{C3}} (L', Z')$ is identical for $\to_{\mathrm{li}}$ and $\to_{\mathrm{ri}}$.

For~(i), to eliminate the spurious state-graph cycles in $F_{\mathrm{SN}}^2$, define a $5$-valued child tag $\kappa : Q_{\mathrm{SN}} \to \{0, 1, 2, 3, 4\}$ grouping normal-form states by structural spine capability, and assign to each $V \in \mathcal{N}$ its finite \emph{normal-form descriptor}
\[
\delta(V) \;=\; \bigl(\mathrm{ls}(V),\; \llbracket V \rrbracket_{\mathrm{SN}},\; \kappa_1(V),\; \kappa_2(V)\bigr)
\;\in\; \{0, 1, 2\} \times F_{\mathrm{SN}} \times \{0, \dots, 4\}^2,
\]
where $(\kappa_1, \kappa_2) = (0, 0)$ for $V \equiv S$, $(\kappa(\llbracket U \rrbracket_{\mathrm{SN}}), 0)$ for $V \equiv S U$, and $(\kappa(\llbracket X \rrbracket_{\mathrm{SN}}), \kappa(\llbracket Y \rrbracket_{\mathrm{SN}}))$ for $V \equiv S X Y$. Only $96$ descriptors are inhabited across all $V \in \mathcal{N}$. Using abstract interpretation of $\mathrm{app}_{\mathcal{N}}(X, Z)$ over descriptors (defining a finite superset $\hat{\alpha}(\delta(X), z)$ of the descriptors reachable from $\delta(X)$ and $z = \llbracket Z \rrbracket_{\mathrm{SN}}$), we construct an integer rank $\rho_{\mathrm{c3}}(\delta(L), \llbracket Z \rrbracket_{\mathrm{SN}}) \in \{1, 2, 3, 4, 5\}$ verified across all $(x, y, z) \in F_{\mathrm{SN}}^3$ satisfying $((q_0 \odot x) \odot y) \odot z = \infty_{\mathrm{SN}}$, $x \odot z \neq \infty_{\mathrm{SN}}$, and $y \odot z \neq \infty_{\mathrm{SN}}$ to obey: (a)~$1 \le \rho_{\mathrm{c3}}(\delta(S X Y), z) \le 5$; and (b)~for every reachable Spine-$2$ descriptor $d' = (2, x', \kappa'_1, \kappa'_2) \in \hat{\alpha}(\delta(X), z)$ with $x' \odot_{\mathrm{SN}} (y \odot_{\mathrm{SN}} z) = \infty_{\mathrm{SN}}$, we have $\rho_{\mathrm{c3}}(d', y \odot_{\mathrm{SN}} z) < \rho_{\mathrm{c3}}(\delta(S X Y), z)$. Six consecutive Case-3 steps $(L_0, Z_0) \mapsto_{\mathrm{C3}} \dots \mapsto_{\mathrm{C3}} (L_6, Z_6)$ would therefore yield $5 \ge \rho_{\mathrm{c3}}(\delta(L_0), \llbracket Z_0 \rrbracket_{\mathrm{SN}}) > \dots > \rho_{\mathrm{c3}}(\delta(L_5), \llbracket Z_5 \rrbracket_{\mathrm{SN}}) \ge 1$, requiring $6$ strictly decreasing integers in $\{1, 2, 3, 4, 5\}$, a contradiction.

For~(ii), letting $B \equiv S T \equiv S(S S)$, direct normalization verifies the $5$-step Case-3 chain:
\begin{align*}
(L_0, Z_0) &\;=\; \bigl(S(S(S A)) S,\; S\bigr)
\;\mapsto_{\mathrm{C3}}\; \bigl(S(S A) S,\; T\bigr)
\;\mapsto_{\mathrm{C3}}\; \bigl(S A T,\; B\bigr) \\
&\;\mapsto_{\mathrm{C3}}\; \bigl(S B (S B),\; S S B\bigr)
\;\mapsto_{\mathrm{C3}}\; \bigl(S T (S S B),\; S B (S S B)\bigr) \\
&\;\mapsto_{\mathrm{C3}}\; \bigl(S S Z_4,\; S Z_4 (S T Z_4)\bigr)
\qquad \bigl(Z_4 \equiv S B (S S B)\bigr),
\end{align*}
where every pair $(L_i, Z_i)$ ($0 \le i \le 5$) satisfies $L_i \in \mathcal{N}_2$, $Z_i \in \mathcal{N}$, and $\llbracket L_i Z_i \rrbracket_{\mathrm{SN}} = \infty_{\mathrm{SN}}$.

For~(iii), since $\mathrm{rs}(M) \le \mathrm{rs}(M')$ holds for every multiway step $M \to_S M'$ (non-root steps preserve or increase $\mathrm{rs}$, and root steps add $+1$), normalizing $X Z \to_S^* L'$ and $Y Z \to_S^* Z'$ gives $\mathrm{rs}(L') \ge \mathrm{rs}(X Z) = 1 + \mathrm{rs}(Z)$ and $\mathrm{rs}(Z') \ge \mathrm{rs}(Y Z) = 1 + \mathrm{rs}(Z)$. Restricting $z = \llbracket Z \rrbracket_{\mathrm{SN}}$ to states inhabited by normal forms with $\mathrm{rs}(Z) \ge 2$ (which excludes $q_0 = \llbracket S \rrbracket$ and all states $\llbracket U S \rrbracket$ of $\mathrm{rs} = 1$) bounds $\rho_{\mathrm{c3}}(\delta(L), z) \le 3$, limiting any run of Case-3 steps from that point onward to at most~$3$.
\end{proof}

\subsection{Post's Lattice Monotonicity versus Minimal Normal-Form Rule 110}
\label{subsec:posts-lattice-rule110}

We next examine whether a fixed open $S$-polynomial $P(x_1, \dots, x_k)$ built from $S$ and variables $x_1, \dots, x_k$ can evaluate Boolean logic or evolve a universal one-dimensional cellular automaton such as Cook's universal, $\mathrm{P}$-complete elementary cellular automaton Rule~110~\cite{Cook2004Rule110,NearyWoods2006Rule110PComplete,Wolfram2002NKS} over two distinguished closed equivalence classes of $\mathrm{CL}(S)$. Using Post's functional completeness theorem and lattice of Boolean clones~\cite{Post1941}, we establish a sharp algebraic dichotomy: inside the semantic $S$-algebras $\mathcal{A}_{\mathrm{HN}}$, $\mathcal{A}_{\mathrm{SN}}$, and $\mathcal{A}_P$, compatible partial orders confine all induced Boolean operations to Post's monotone clone $M$, forbidding negation and Rule~110; yet over a dedicated $4$-state conversion-invariant $S$-algebra $\mathcal{D}_4$ (Table~\ref{tab:d4-algebra}), non-monotonicity emerges cleanly and yields a globally minimal $11$-leaf normal-form Rule~110 combinator.

\begin{table}[htbp]
\centering
\caption{Multiplication table of the $4$-state functionally complete $S$-algebra $\mathcal{D}_4 = (\{d_0, d_1, d_2, d_3\}, d_0, \odot_{\mathcal{D}_4})$ and minimal normal-form ground representatives.}
\label{tab:d4-algebra}
\small
\begin{tabular}{@{}cccccll@{}}
\toprule
$\odot_{\mathcal{D}_4}$ & $d_0$ & $d_1$ & $d_2$ & $d_3$ & Minimal Representative & Boolean / Algebraic Role \\
\midrule
$d_0$ & $d_1$ & $d_1$ & $d_1$ & $d_2$ & $S$ ($1$ leaf, $\mathcal{N}$) & Generator $d_0 = \llbracket S \rrbracket_{\mathcal{D}_4}$ \\
$d_1$ & $d_3$ & $d_3$ & $d_3$ & $d_3$ & $S S$ ($2$ leaves, Spine-$1$ $\mathcal{N}$) & Encoded Boolean $\mathsf{false}$ ($0$) \\
$d_2$ & $d_1$ & $d_1$ & $d_3$ & $d_3$ & $S(S S S)$ ($4$ leaves, Spine-$1$ $\mathcal{N}$) & Encoded Boolean $\mathsf{true}$ ($1$) \\
$d_3$ & $d_3$ & $d_3$ & $d_3$ & $d_3$ & $S S S$ ($3$ leaves, Spine-$2$ $\mathcal{N}$) & Auxiliary absorbing mediator \\
\bottomrule
\end{tabular}
\end{table}

\begin{theorem}[Post's lattice monotonicity barrier and minimal $11$-leaf normal-form Rule~110]
\label{thm:clones-rule110}
\begin{enumerate}[label=(\roman*)]
\item \emph{(Monotone-clone barrier in $\mathcal{A}_{\mathrm{HN}}, \mathcal{A}_{\mathrm{SN}}, \mathcal{A}_P$)} Each of $\mathcal{A}_{\mathrm{HN}}$, $\mathcal{A}_{\mathrm{SN}}$, and $\mathcal{A}_P$ admits non-trivial compatible preorders $\sqsubseteq$ satisfying $x_1 \odot y_1 \sqsubseteq x_2 \odot y_2$ whenever $x_1 \sqsubseteq x_2$ and $y_1 \sqsubseteq y_2$ ($6$ non-trivial compatible preorders on $\mathcal{A}_{\mathrm{HN}}$, and $5$ each on $\mathcal{A}_{\mathrm{SN}}$ and $\mathcal{A}_P$). Across every pair of distinct states $(e_0, e_1)$ in $\mathcal{A}_{\mathrm{HN}}$, $\mathcal{A}_{\mathrm{SN}}$, or $\mathcal{A}_P$, every open $S$-polynomial $P(x_1, \dots, x_k)$ mapping $\{e_0, e_1\}^k$ into $\{e_0, e_1\}$ computes a monotone Boolean function belonging to Post's clone $M$~\cite{Post1941}; in particular, none of $\mathrm{NOT}$, $\mathrm{NAND}$, $\mathrm{XOR}$, or Rule~110 is polynomial-definable over any $2$-state encoding in $\mathcal{A}_{\mathrm{HN}}$, $\mathcal{A}_{\mathrm{SN}}$, or $\mathcal{A}_P$.
\item \emph{(Functional completeness of $\mathcal{D}_4$ and minimal $11$-leaf normal-form Rule~110)} The $4$-state structure $\mathcal{D}_4$ of Table~\ref{tab:d4-algebra} is a conversion-invariant $S$-algebra satisfying~\eqref{eq:s-algebra-law} on all $64$ triples, breaking monotonicity via $d_0 \odot_{\mathcal{D}_4} d_3 = d_2 \neq d_1 = d_0 \odot_{\mathcal{D}_4} d_1$. Encoding $\mathsf{false}$ by $F_{\flat} \equiv S S$ ($\llbracket F_{\flat} \rrbracket_{\mathcal{D}_4} = d_1$) and $\mathsf{true}$ by $T_{\sharp} \equiv S(S S S)$ ($\llbracket T_{\sharp} \rrbracket_{\mathcal{D}_4} = d_2$), the $4$-leaf polynomial $\mathrm{NAND}(x, y) \equiv S(x(y S))$ realizes Boolean $\mathrm{NAND}$ on $\{d_1, d_2\}$, and the $11$-leaf polynomial
\begin{equation}
\label{eq:rule110-poly}
R_{110}(x, y, z) \;\equiv\; S\Bigl(S(y\,x)\Bigl(S(y\,z)\bigl(S(z\,y)\,S\bigr)\Bigr)\Bigr)
\end{equation}
evaluates elementary cellular automaton Rule~110 ($b'_i = (b_y + b_z + b_y b_z + b_x b_y b_z) \bmod 2$, truth vector $01101110_2$ on $111 \dots 000$) over $\{d_1, d_2\}$. Furthermore:
\begin{enumerate}
\item[(a)] Whenever $x, y, z \in \mathcal{T}_S$ are Spine-$1$ normal forms ($S W$ with $W \in \mathcal{N}$, including $F_{\flat}$ and $T_{\sharp}$), $R_{110}(x, y, z)$ is already a Spine-$1$ normal form ($0$ reduction steps), so synchronous iteration of $R_{110}$ evolves arbitrary Rule~110 configurations entirely within Spine-$1$ normal forms.
\item[(b)] No open $\{S, x, y, z\}$-polynomial with fewer than $11$ leaves realizes Rule~110 on $\{d_1, d_2\}$ in $\mathcal{D}_4$, and no curried closed combinator $G \in \mathcal{T}_S$ can satisfy $\llbracket G x y z \rrbracket_{\mathcal{D}_4} = \llbracket R_{110}(x, y, z) \rrbracket_{\mathcal{D}_4}$ (nor reduce $G x y z \to_S^* S W$).
\end{enumerate}
\end{enumerate}
\end{theorem}

\begin{proof}
For~(i), if $\sqsubseteq$ is a compatible preorder on $(Q, q_0, \odot)$, structural induction on any open polynomial $P(x_1, \dots, x_k)$ built from constants in $Q$ and variables shows that $\vec{u} \sqsubseteq \vec{v}$ componentwise implies $P(\vec{u}) \sqsubseteq P(\vec{v})$. Whenever $e_0 \sqsubset e_1$ (meaning $e_0 \sqsubseteq e_1$ and $e_1 \not\sqsubseteq e_0$), any induced Boolean operation $f : \{0, 1\}^k \to \{0, 1\}$ with $P(e_{b_1}, \dots, e_{b_k}) = e_{f(b_1, \dots, b_k)}$ must satisfy $f(\vec{b}) \le f(\vec{b}')$ whenever $\vec{b} \le \vec{b}'$, placing $f$ inside Post's monotone clone $M \subsetneq \mathrm{Pol}(\{0, 1\})$~\cite{Post1941}. On $\mathcal{A}_{\mathrm{HN}}$, for example, ordering $q \sqsubseteq \infty_{\mathrm{HN}}$ for all $q$ (and refining along the $6$ compatible preorders of $\mathcal{A}_{\mathrm{HN}}$ and $5$ of $\mathcal{A}_{\mathrm{SN}}, \mathcal{A}_P$), together with complete finite closure of unary and binary operations on every incomparable pair $(e_0, e_1)$, confirms that every unary polynomial in $\mathcal{A}_{\mathrm{HN}}, \mathcal{A}_{\mathrm{SN}}, \mathcal{A}_P$ is aperiodic ($f(e_0) = e_1$ and $f(e_1) = e_0$ never hold simultaneously), ruling out $\mathrm{NOT}$ ($\neg x$), $\mathrm{NAND}$ ($\mathrm{NAND}(x, x) = \neg x$), $\mathrm{XOR}$ ($\mathrm{XOR}(x, 1) = \neg x$), and Rule~110 ($R_{110}(1, x, 1) = \neg x$).

For~(ii), we verify~\eqref{eq:s-algebra-law} on $\mathcal{D}_4$ in four lines: for all $x, y, z \in Q_{\mathcal{D}_4}$, $d_0 \odot x \in \{d_1, d_2\}$, whence $(d_0 \odot x) \odot y \in \{d_1, d_3\}$, and since rows $d_1$ and $d_3$ are identically $d_3$, the left-hand side $((d_0 \odot x) \odot y) \odot z$ equals $d_3$. On the right-hand side $(x \odot z) \odot (y \odot z)$, note $d_0 \notin \mathrm{im}(\odot)$, so $x \odot z \in \{d_1, d_2, d_3\}$; if $x \odot z \in \{d_1, d_3\}$ the product is $d_3$, while the sole entry giving $x \odot z = d_2$ is $(x, z) = (d_0, d_3)$, for which $y \odot d_3 \in \{d_2, d_3\}$ and $d_2 \odot \{d_2, d_3\} = \{d_3\}$. Thus both sides of~\eqref{eq:s-algebra-law} equal $d_3$ identically.

Evaluating $\mathrm{NAND}(x, y) \equiv S(x(y S))$ on $x, y \in \{d_1, d_2\}$: if $y = d_1$, then $y \odot d_0 = d_3$, $x \odot d_3 = d_3$, and $d_0 \odot d_3 = d_2$; if $y = d_2$, then $y \odot d_0 = d_1$, so $d_1 \odot d_1 = d_3 \mapsto d_2$ whereas $d_2 \odot d_1 = d_1 \mapsto d_0 \odot d_1 = d_1$, matching $\mathrm{NAND}$. Similarly, for $x, y, z \in \{d_1, d_2\}$, observe that $u \odot v = d_1$ if $(u, v) = (d_2, d_1)$ and $d_3$ otherwise on $\{d_1, d_2\}^2$. Direct substitution into~\eqref{eq:rule110-poly} over the eight triples $(x, y, z) \in \{d_1, d_2\}^3$ ordered from $(d_1, d_1, d_1)$ ($000$) to $(d_2, d_2, d_2)$ ($111$) yields $(d_1, d_2, d_2, d_2, d_1, d_2, d_2, d_1)$, matching $0, 1, 1, 1, 0, 1, 1, 0$ (Rule~110). For~(a), if $x \equiv S W_x$, $y \equiv S W_y$, $z \equiv S W_z$ with $W_x, W_y, W_z \in \mathcal{N}$, then $y x \equiv S W_y (S W_x) \in \mathcal{N}_2 \subset \mathcal{N}$ (and likewise $y z, z y \in \mathcal{N}$), whence $S(z y)S \in \mathcal{N}$, $S(y z)(S(z y)S) \in \mathcal{N}$, $S(y x)(\dots) \in \mathcal{N}$, and $R_{110}(x, y, z) \equiv S(\dots)$ is a Spine-$1$ normal form. For~(b), dynamic programming over all $4^8 = 65{,}536$ truth vectors in $Q_{\mathcal{D}_4}^8$ reachable from $\{S, x, y, z\}$ by leaf size $s = 1, \dots, 11$ confirms that $(d_1, d_2, d_2, d_2, d_1, d_2, d_2, d_1)$ is absent for all $s \le 10$ and first appears at $s = 11$; finally, for any $g \in Q_{\mathcal{D}_4}$, $((g \odot x) \odot y) \odot z = d_3$ identically as above, while Lemma~\ref{lem:visser-fibonacci}(i) forbids $(G x y) z \to_S^* S W$.
\end{proof}

% ============================================================================
% SECTION 9
% ============================================================================

\section{Wolfram's Benchmark Evolutions, All-Step Jump-Ahead, Garden-of-Eden Normal Forms, and Exact Catalan Phase Transitions}
\label{sec:asymptotics}

In his centennial combinator monograph~\cite{Wolfram2021CombinatorsBook} and challenge announcement~\cite{Wolfram2021Challenge}, Stephen Wolfram surveyed the empirical dynamics of small closed $S$-terms across the Catalan tree hierarchy, noting the minimal $7$-leaf threshold for non-termination, presenting twelve non-terminating benchmark expressions $W_1, \dots, W_{12}$ of leaf sizes $8$ through~$11$ whose leftmost-outermost leaf-count differences exhibit varied transient and staircase patterns, and asking whether bounded-time ``jump-ahead'' procedures can predict halting or future structural states without step-by-step simulation. Beyond individual trajectories, Wolfram posed quantitative questions concerning the prevalence of Garden-of-Eden expressions (terms of in-degree zero in the multiway rewrite graph), strategy-dependent subterm appearance, and the limiting proportions of dynamical regimes among large random $S$-terms.

In this section we resolve each of these questions constructively. Subsection~\ref{subsec:wolfram-14-benchmarks} classifies Dubou\'e's $7$-leaf terms and Wolfram's twelve benchmarks algebraically and dynamically, isolating $W_2$ from Zachos's productive combinator $Z_{\mathrm{Zac}}$. Subsection~\ref{subsec:jump-ahead} establishes linear-time multi-hole context halting, explicit all-step leftmost-outermost state and cubic leaf-count formulas for $Z_{\mathrm{Zac}}$, $W_2$ (together with a constant $+14$ leaf-shift law linking $W_2$ to an $\convS$-inconvertible twin), and the $10$-leaf universal seed $D_{\mathrm{seed}}$, followed by a three-rule multiway grammar on $A A A$ that yields sharp subterm-appearance separations between $\redSs$ and $\to_{\mathrm{lo}}$. Subsection~\ref{subsec:goe-isolated} characterizes Garden-of-Eden normal forms, equates them on $\mathcal{N} := \NF \cap \TS$ with isolated normal forms ($[P]_{\convS} = \{P\}$), proves a constructive tower law and linear-time decidability, and derives the exact Fibonacci--Mahler functional equation governing their growth constant $\gamma_{\mathrm{Iso}} \approx 2.907622 < 3$. Subsection~\ref{subsec:catalan-minimality} establishes sharp small-term minimality and confluence thresholds across all $23{,}714$ closed $S$-terms up to $11$ leaves. Finally, Subsection~\ref{subsec:analytic-combinatorics} combines analytic combinatorics with constructive arithmetic certificates over $\mathbb{N}$ to prove universal positive-density bounds and exact dyadic-ternary rational Catalan densities $(c_0, c_1, c_2, c_3)$ for the functional order hierarchy $\{0, 1, 2, 3\}$, completing a six-tier exponential growth spectrum for $\CLS$.

\subsection{Algebraic and Dynamical Classification of Wolfram's Benchmark Expressions}
\label{subsec:wolfram-14-benchmarks}

Alongside Dubou\'e's two minimal $7$-leaf non-normalizing $S$-terms $X_1 \equiv S(S S)S S S S$ and $X_2 \equiv S S S (S S) S S$ (communicated by Dubou\'e and Baron in 1975; see~\cite{BarendregtBergstraKlopVolken1976}, \cite[Exercise~7.4.5]{Barendregt1984}, and~\cite[line~131]{Wolfram2021Challenge}), let $W_1, \dots, W_{12}$ denote the twelve non-normalizing expressions displayed in lines~147--153 of~\cite{Wolfram2021Challenge}. Throughout this section we retain the standard abbreviations $T \equiv S S$, $A \equiv S S S$, $G_0 \equiv S T T$, $G_{k+1} \equiv T G_k$, $Z_0 \equiv S(S T)(S T)T$, $Z_{\mathrm{Zac}} \equiv G_0 G_0$, and $D_{\mathrm{seed}} \equiv Z_0 S$. Under leftmost-outermost reduction $\to_{\mathrm{lo}}$, whenever a closed $S$-term reaches a head normal form $S A_1$ or $S A_1 A_2$, the outer $S$-application freezes permanently as a \emph{head-normal shell}: for $S A_1$, reduction descends into $A_1$; for $S A_1 A_2$, reduction descends into $A_1$ if $A_1 \notin \SN$, and otherwise normalizes $A_1 \to_{\mathrm{lo}}^* \mathrm{nf}(A_1) \in \mathcal{N}$ before descending permanently into $A_2$. Table~\ref{tab:wolfram-benchmarks} records the exact algebraic coordinates and $\to_{\mathrm{lo}}$ shell-stabilization behavior of all fourteen terms.

\begin{table}[htbp]
\centering
\caption{Complete algebraic and $\to_{\mathrm{lo}}$ dynamical classification of Dubou\'e's $7$-leaf terms $X_1, X_2$ and Wolfram's twelve challenge benchmarks $W_1, \dots, W_{12}$~\cite{Wolfram2021Challenge}. State indices refer to $(\sem{M}_{\SN}, \sem{M}_{\HN}, \sem{M}_P)$ with sinks $(q_{38}, q_{34}, q_{37}) = (\infty_{\SN}, \infty_{\HN}, \infty_P)$.}
\label{tab:wolfram-benchmarks}
\footnotesize
\setlength{\tabcolsep}{4.2pt}
\begin{tabular}{@{}lllcccl@{}}
\toprule
Term & Combinatory Form & Wolfram Announcement Syntax & $|M|$ & $(\SN, \HN, P)$ & Shells / Lock & $\to_{\mathrm{lo}}$ Dynamical Regime \\
\midrule
$X_2$ & $S S S (S S) S S$ & \texttt{s[s][s][s[s]][s][s]} & $7$ & $(q_{38}, q_{23}, q_{37})$ & $1$ / step~$8$ & Stabilizes $1$ shell, then $\infty_{\HN}$ \\
$X_1$ & $S(S S)S S S S$ & \texttt{s[s[s]][s][s][s][s]} & $7$ & $(q_{38}, q_{23}, q_{37})$ & $1$ / step~$9$ & Steps in $1$ step to $X_2$; locks at step~$9$ \\
$W_1$ & $S(S(S(S S)S))S S$ & \texttt{s[s[s[s[s]][s]]][s][s]} & $8$ & $(q_{38}, q_{23}, q_{37})$ & $4$ / step~$18$ & Stabilizes $4$ shells, then $\infty_{\HN}$ \\
$W_2$ & $S S S (S(S S))(S S)$ & \texttt{s[s][s][s[s[s]]][s[s]]} & $8$ & $(q_{38}, q_{24}, q_{37})$ & $\infty$ / never & Perpetual $\HN\setminus\SN$; depth-$5$ $\bot$ \\
$W_3$ & $S S S (S(S(S S S)))S$ & \texttt{s[s][s][s[s[s[s][s]]]][s]} & $9$ & $(q_{38}, q_{34}, q_{37})$ & $0$ / step~$0$ & Unsolvable ($\BT = \bot$) \\
$W_4$ & $S(S(S S)(S(S S)S))S S$ & \texttt{s[s[s[s]][s[s[s]][s]]][s][s]} & $10$ & $(q_{38}, q_{34}, q_{37})$ & $0$ / step~$0$ & Unsolvable ($\BT = \bot$) \\
$W_5$ & $S S S S S S (S(S S))S$ & \texttt{s[s][s][s][s][s][s[s[s]]][s]} & $10$ & $(q_{38}, q_{23}, q_{37})$ & $1$ / step~$21$ & Stabilizes $1$ shell, then $\infty_{\HN}$ \\
$W_6$ & $S S S (S(S(S S S)S))S$ & \texttt{s[s][s][s[s[s[s][s]][s]]][s]} & $10$ & $(q_{38}, q_{34}, q_{37})$ & $0$ / step~$0$ & Unsolvable ($\BT = \bot$) \\
$W_7$ & $S(S S S (S S))(S S S)S$ & \texttt{s[s[s][s][s[s]]][s[s][s]][s]} & $10$ & $(q_{38}, q_{34}, q_{37})$ & $0$ / step~$0$ & Unsolvable ($\BT = \bot$) \\
$W_8$ & $S S S (S(S S))(S S)(S S)$ & \texttt{s[s][s][s[s[s]]][s[s]][s[s]]} & $10$ & $(q_{38}, q_{23}, q_{37})$ & $1$ / step~$8$ & Stabilizes $1$ shell, then $\infty_{\HN}$ \\
$W_9$ & $S S S S (S S S) S S S S$ & \texttt{s[s][s][s][s[s][s]][s][s][s][s]} & $11$ & $(q_{38}, q_{34}, q_{37})$ & $0$ / step~$0$ & Unsolvable ($\BT = \bot$) \\
$W_{10}$ & $S S S (S(S(S(S S)S)S))S$ & \texttt{s[s][s][s[s[s[s[s]][s]][s]]][s]} & $11$ & $(q_{38}, q_{34}, q_{37})$ & $0$ / step~$0$ & Unsolvable ($\BT = \bot$) \\
$W_{11}$ & $S S S (S S (S(S(S S S))))S$ & \texttt{s[s][s][s[s][s[s[s[s][s]]]]][s]} & $11$ & $(q_{38}, q_{25}, q_{37})$ & $1$ / step~$4$ & Stabilizes $1$ shell, then $\infty_{\HN}$ \\
$W_{12}$ & $S(S(S S)S S)(S(S S)S)S$ & \texttt{s[s[s[s]][s][s]][s[s[s]][s]][s]} & $11$ & $(q_{38}, q_{34}, q_{37})$ & $0$ / step~$0$ & Unsolvable ($\BT = \bot$) \\
\bottomrule
\end{tabular}
\end{table}

\begin{theorem}[Classification of Wolfram's benchmarks and separation of $W_2$ from $Z_{\mathrm{Zac}}$]
\label{thm:wolfram-14-dichotomy}
Every term $M$ among the fourteen benchmark $S$-terms $X_1, X_2, W_1, \dots, W_{12}$ of Table~\ref{tab:wolfram-benchmarks} satisfies $\sem{M}_{\SN} = \infty_{\SN}$ and $\sem{M}_P = \infty_P$, and the fourteen terms partition into three disjoint $\to_{\mathrm{lo}}$ regimes:
\begin{enumerate}[label=(\roman*),leftmargin=2.2em]
\item \emph{(Immediate unsolvability)} Seven benchmarks---$W_3, W_4, W_6, W_7, W_9, W_{10}, W_{12}$---satisfy $\sem{W_i}_{\HN} = \infty_{\HN}$ ($\ord(W_i) = 0$), executing pure head reduction indefinitely with $\BT(W_i) = \bot$.
\item \emph{(Finite shell stabilization into $\infty_{\HN}$)} Six benchmarks---$X_1, X_2, W_1, W_5, W_8, W_{11}$---lie in $\HN \setminus \Prod$ ($\ord(M) = 1$) and emit $4$ head-normal shells for $W_1$ and $1$ shell for each of the other five before their active $\to_{\mathrm{lo}}$ focus locks permanently into $\infty_{\HN}$ at steps $9, 8, 18, 21, 8, 4$, respectively.
\item \emph{(Perpetual solvable focus with hidden $\bot$)} Alone among the fourteen benchmarks, $W_2 \equiv S S S (S T) T$ ($8$ leaves, $\ord(W_2) = 1$) maintains an active $\to_{\mathrm{lo}}$ focus in $\HN \setminus \SN$ at every step $t \in \mathbb{N}$: defining $E_0 \equiv S(S T)T$ and $E_{m+1} \equiv T E_m$, after $\frac{1}{2}(n+1)(n+6) + 1$ steps ($n \ge 0$), $W_2$ reaches active focus $E_0 E_{n+1}$ ($\HN$ state $q_{23} \neq \infty_{\HN}$) beneath $\frac{1}{2}(n+1)(n+2)$ head-normal shells. Nevertheless, $\BT(W_2)$ contains an unsolvable subterm ($\infty_{\HN}$) at depth~$5$, whence $W_2 \notin \Prod$, $W_2 \neq_S Z_{\mathrm{Zac}}$, and $W_2 \neq_S Z_0$.
\end{enumerate}
\end{theorem}

\begin{proof}
Regimes~(i) and~(ii) follow from Theorem~\ref{thm:question-10}(i) and Theorem~\ref{thm:order-hierarchy}(i) by evaluating $\sem{\cdot}_{\HN}$ on the initial terms and tracing $4$ to $21$ deterministic $\to_{\mathrm{lo}}$ steps until the unique active subterm beneath the stabilized $\HNF$ context reaches $\HN$ state $\infty_{\HN}$, after which pure head reduction continues inside that focus forever because $\infty_{\HN}$ admits no $\HNF$.

For~(iii), set $E_0 \equiv S(S T)T \in \mathcal{N}$ ($6$ leaves) and $E_{m+1} \equiv T E_m \equiv S S E_m \in \mathcal{N}$ ($m \ge 0$, $|E_m| = 2m + 6$). Tracing the first four $\to_{\mathrm{lo}}$ steps of $W_2 \equiv ((S S S)(S T)) T$ yields
\begin{align*}
W_2 &\;\equiv\; ((S S S)(S T)) T
\;\longrightarrow_{\mathrm{lo}}\; \bigl(S(S T)(S(S T))\bigr) T
\;\longrightarrow_{\mathrm{lo}}\; (S T T)\bigl(S(S T) T\bigr) \;\equiv\; (S T T)\,E_0 \\
&\;\longrightarrow_{\mathrm{lo}}\; (T E_0)(T E_0) \;\equiv\; (S S E_0)\,E_1
\;\longrightarrow_{\mathrm{lo}}\; S\,E_1\,(E_0 E_1).
\end{align*}
Because $E_1 \in \mathcal{N}$ is in normal form, step~$4$ stabilizes the Spine-$2$ right-hole shell $S\,E_1\,\square$ (shell depth~$1$ at $n = 0$, matching $\frac{1}{2}(0+1)(0+6) + 1 = 4$) around the active focus $E_0 E_1$. From any focus $E_0 E_{k+1}$ ($k \ge 0$), three $\to_{\mathrm{lo}}$ steps give
\[
E_0 E_{k+1} \;\equiv\; S(S T)T\,E_{k+1}
\;\longrightarrow_{\mathrm{lo}}\; (S T E_{k+1})\,E_{k+2}
\;\longrightarrow_{\mathrm{lo}}\; E_{k+3}\,(E_{k+1} E_{k+2})
\;\longrightarrow_{\mathrm{lo}}\; S\,\bigl(E_{k+1} E_{k+2}\bigr)\,\bigl(E_{k+2}(E_{k+1} E_{k+2})\bigr).
\]
Here the outer shell $S\,\square\,[E_{k+2}(E_{k+1} E_{k+2})]$ permanently freezes its right sibling $R_{\mathrm{fr}}(k) \equiv E_{k+2}(E_{k+1} E_{k+2})$ and enters the left child $E_{k+1} E_{k+2}$, which then executes $k + 1$ single-step countdowns $E_{m+1} E_{k+2} \equiv S S E_m E_{k+2} \to_{\mathrm{lo}} S E_{k+2} (E_m E_{k+2})$ ($m = k, \dots, 0$, each $E_{k+2} \in \mathcal{N}$) down to $E_0 E_{k+2}$. Summing $4 + \sum_{j=0}^{n-1} (j + 4) = \frac{1}{2}(n+1)(n+6) + 1$ steps and $1 + \sum_{j=0}^{n-1} (j + 2) = \frac{1}{2}(n+1)(n+2)$ shells establishes the checkpoint trajectory, with every intermediate focus in $F_{\HN}$. Yet at $k = 0$, the first orphaned right sibling $R_{\mathrm{fr}}(0) \equiv E_2(E_1 E_2)$ sits at shell depth~$2$ inside $S\,E_1\,(S\,\square\,R_{\mathrm{fr}}(0))$, and unfolding three further B\"ohm levels inside $R_{\mathrm{fr}}(0)$ reaches an unsolvable subterm ($\infty_{\HN}$) at B\"ohm depth~$5$. Thus $\sem{W_2}_P = \infty_P \neq q_{29} = \sem{Z_{\mathrm{Zac}}}_P = \sem{Z_0}_P$, proving $W_2 \neq_S Z_{\mathrm{Zac}}$ and $W_2 \neq_S Z_0$.
\end{proof}

\subsection{Compositional Context Halting, All-Step Jump-Ahead Closed Forms, and Multiway Subterm Separations}
\label{subsec:jump-ahead}

In line~179 of~\cite{Wolfram2021Challenge}, Wolfram asks whether ``jump-ahead'' procedures---algorithms guaranteed to halt in bounded time---can decide termination or compute distant future states of $S$-terms without simulating intermediate reductions. Theorems~\ref{thm:waldmann-sn} and~\ref{thm:question-10} furnish linear-time halting deciders for $\to_{\mathrm{lo}}$ and $\to_h$, and homomorphic compositionality extends these deciders to multi-hole contexts whose explicit substitution would grow exponentially. Beyond halting decisions, we show in Theorem~\ref{thm:all-step-jump-ahead} that for the principal non-terminating combinators $Z_{\mathrm{Zac}}$, $W_2$, and $D_{\mathrm{seed}}$, both the exact $t$-th leftmost-outermost reduct $\mathrm{lo}^t(M)$ and its exact cubic leaf count $|\mathrm{lo}^t(M)|$ admit closed-form descriptions at \emph{every} integer step $t \in \mathbb{N}$.

\begin{proposition}[Linear-time multi-hole context halting and checkpoint jump-ahead]
\label{prop:jump-ahead}
\begin{enumerate}[label=(\roman*),leftmargin=2.2em]
\item \emph{(Compositional $k$-hole context halting)} Let $C[\square_1, \dots, \square_k]$ be an arbitrary $k$-hole $S$-context (allowing each hole $\square_i$ to appear arbitrarily many times at arbitrary depths) of syntactic size $|C|$, and let $M_1, \dots, M_k \in \TS$ be closed arguments. Whether $C[M_1, \dots, M_k]$ halts under $\to_{\mathrm{lo}}$ (respectively $\to_h$) is decidable in $O\bigl(|C| + \sum_{i=1}^k |M_i|\bigr)$ time without constructing $C[M_1, \dots, M_k]$.
\item \emph{(Two-parameter counter checkpoints on $D_{\mathrm{seed}}$)} On the $10$-leaf seed $D_{\mathrm{seed}} \equiv S(S(S S))(S(S S))(S S)S$, setting $C_0 \equiv S T T \equiv G_0$ and $C_{m+1} \equiv T C_m \equiv G_{m+1}$, for every $n \ge 1$ and $0 \le r \le n$, step $(2n^2 - n + 1) + 4r$ under $\to_{\mathrm{lo}}$ places $D_{\mathrm{seed}}$ at active counter focus $C_{n-r} C_n S$ beneath $\frac{1}{2}n(n-1) + r$ left-active head-normal shells (reaching $C_0 C_n S$ at step $2n^2 + 3n + 1$ when $r = n$).
\item \emph{(Exact quarter-square normal-form step count on left spines)} Define the left-associated spine $P_0 \equiv S$, $P_{m+1} \equiv P_m S$ ($m + 1$ leaves, arity $m \ge 0$). Then $P_m$ reaches its normal form under $\to_{\mathrm{lo}}$ in exact step count
\begin{equation}
\label{eq:sprefix-quarter-square}
s(m) \;=\; \Bigl\lfloor \frac{m-1}{2} \Bigr\rfloor \cdot \Bigl\lfloor \frac{m}{2} \Bigr\rfloor
\qquad (m \in \mathbb{N}),
\end{equation}
whereas the right-associated comb $R_m$ ($R_0 \equiv S$, $R_{m+1} \equiv S R_m$) is already in normal form ($0$ steps).
\end{enumerate}
\end{proposition}

\begin{proof}
For~(i), evaluate $q_i = \sem{M_i}_{\SN} \in Q_{\SN}$ (or $\sem{M_i}_{\HN} \in Q_{\HN}$) in $2|M_i| - 1$ table lookups for $1 \le i \le k$, then evaluate $C$ bottom-up in $|C|$ lookups by assigning state $q_i$ to every occurrence of $\square_i$ and $q_0$ to $S$. By homomorphic compositionality of $\sem{\cdot}_{\mathcal{A}}$, the root state equals $\sem{C[M_1, \dots, M_k]}_{\mathcal{A}}$, deciding $\SN$ (hence $\to_{\mathrm{lo}}$ halting by standardization and Proposition~\ref{prop:cr-conservation}) and $\HN$ ($\to_h$ halting) in $O(|C| + \sum |M_i|)$ time.

For~(ii), two initial contractions reduce $D_{\mathrm{seed}} \equiv S(S T)(S T)T S \to_{\mathrm{lo}} (S T T)(S T T) S \equiv C_0 C_0 S \to_{\mathrm{lo}} (T C_0)(T C_0) S \equiv C_1 C_1 S$, reaching $(1, 1)$ at step $2 = 2(1)^2 - 1 + 1$ with $0$ shells. From any counter state $C_{m+1} C_p S \equiv (S S C_m) C_p S$ ($p \ge 1$), four $\to_{\mathrm{lo}}$ contractions via $((S C_p)(C_m C_p)) S \to_{\mathrm{lo}} (C_p S)(C_m C_p S) \equiv (S S (C_{p-1} S))(C_m C_p S)$ produce $S\,(C_m C_p S)\,\bigl((C_{p-1} S)(C_m C_p S)\bigr)$, adding one left-active Spine-$2$ shell around $C_m C_p S$; at $m = 0$, a single contraction resets $C_0 C_p S \equiv S T T C_p S \to_{\mathrm{lo}} (T C_p)(T C_p) S \equiv C_{p+1} C_{p+1} S$ with $0$ additional shells. Induction over $n \ge 1$ sums $2 + \sum_{j=1}^{n-1} (4j + 1) = 2n^2 - n + 1$ steps and $\sum_{j=1}^{n-1} j = \frac{1}{2}n(n-1)$ shells to reach $C_n C_n S$, after which $r$ further countdown transitions ($0 \le r \le n$) add $4r$ steps and $r$ shells to reach $C_{n-r} C_n S$.

For~(iii), $P_0 \equiv S$, $P_1 \equiv S S$, and $P_2 \equiv S S S$ lie in $\mathcal{N}$, matching $s(0) = s(1) = s(2) = 0$. For $m = d + 3$ ($d \ge 0$), $P_{d+3} \equiv S S S S S^d \to_{\mathrm{lo}} (S S)(S S) S^d$, which peels in $d$ further $\to_{\mathrm{lo}}$ steps via $(S S) U S \to_{\mathrm{lo}} S S (U S)$ to $(S S)\,P_{d+1} \equiv T\,P_{d+1}$, giving $P_{d+3} \to_{\mathrm{lo}}^{d+1} T\,P_{d+1}$ with spine arity at least $3$ at all $j < d + 1$ steps. Because $T\,U \equiv S S U$ is a Spine-$2$ right-hole shell $S\,S\,\square$ with $S \in \mathcal{N}$, subsequent $\to_{\mathrm{lo}}$ steps act strictly inside $P_{d+1}$, yielding the recurrence $s(d + 3) = (d + 1) + s(d + 1)$. Since $\lfloor \frac{d+2}{2} \rfloor \lfloor \frac{d+3}{2} \rfloor - \lfloor \frac{d}{2} \rfloor \lfloor \frac{d+1}{2} \rfloor = (\lfloor \frac{d}{2} \rfloor + 1)(\lfloor \frac{d+1}{2} \rfloor + 1) - \lfloor \frac{d}{2} \rfloor \lfloor \frac{d+1}{2} \rfloor = \lfloor \frac{d}{2} \rfloor + \lfloor \frac{d+1}{2} \rfloor + 1 = d + 1$, induction on $m$ establishes~\eqref{eq:sprefix-quarter-square}.
\end{proof}

To extend checkpoint descriptions to \emph{all} steps $t \in \mathbb{N}$ without existential context quantifiers, we define the accumulated head-normal shell contexts explicitly.

\begin{definition}[Explicit leftmost-outermost shell contexts for $Z_{\mathrm{Zac}}$, $W_2$, and $D_{\mathrm{seed}}$]
\label{def:lo-explicit-contexts}
For any closed normal form $L \in \mathcal{N}$, write $(S L)^0(U) \equiv U$ and $(S L)^{k+1}(U) \equiv S\,L\,\bigl((S L)^k(U)\bigr)$ ($k \in \mathbb{N}$). Define three families of one-hole head-normal contexts inductively:
\begin{enumerate}[label=(\roman*),leftmargin=2.2em]
\item \emph{(Zachos context tower)} Set $\mathcal{C}_Z(0)[\square] \equiv \square$ and, for $n \in \mathbb{N}$,
\[
\mathcal{C}_Z(n+1)[\square] \;\equiv\; \mathcal{C}_Z(n)\Bigl[(S G_{n+1})^{n+1}(\square)\Bigr]
\qquad (|G_k| = 2k + 5).
\]
\item \emph{(Benchmark $W_2$ context tower)} With $E_0 \equiv S(S T)T$ and $E_{k+1} \equiv T E_k$ ($|E_k| = 2k + 6$), set $\mathcal{C}_{W_2}(0)[\square] \equiv S\,E_1\,\square$ and, for $n \in \mathbb{N}$,
\[
\mathcal{C}_{W_2}(n+1)[\square] \;\equiv\; \mathcal{C}_{W_2}(n)\Bigl[S\,\bigl((S E_{n+2})^{n+1}(\square)\bigr)\,\bigl(E_{n+2}(E_{n+1} E_{n+2})\bigr)\Bigr].
\]
\item \emph{(Universal seed context tower)} With $C_k \equiv G_k$ and $D_{k,p} \equiv C_k C_p S$ ($|D_{k,p}| = 2k + 2p + 11$), for each $p \in \mathbb{N}$ and $0 \le r \le p + 1$ define the partial cycle wrapper $\mathcal{F}_{p,0}[\square] \equiv \square$ and $\mathcal{F}_{p, r+1}[\square] \equiv \mathcal{F}_{p,r}\bigl[S\,\square\,\bigl((C_p S)\,D_{p-r, p+1}\bigr)\bigr]$, and set $\mathcal{C}_{\mathrm{seed}}(0)[\square] \equiv \square$ and $\mathcal{C}_{\mathrm{seed}}(m+1)[\square] \equiv \mathcal{C}_{\mathrm{seed}}(m)\bigl[\mathcal{F}_{m, m+1}[\square]\bigr]$.
\end{enumerate}
For any context $\mathcal{C}[\square]$ built from Spine-$1$ frames $S\,\square$, Spine-$2$ right-hole frames $S\,L\,\square$ ($L \in \mathcal{N}$), and Spine-$2$ left-hole frames $S\,\square\,R$, let $\lambda(\mathcal{C}) := |\mathcal{C}[S]| - 1$ denote the frozen shell leaf contribution, so that $|\mathcal{C}[U]| = \lambda(\mathcal{C}) + |U|$ and $\mathrm{lo}^k(\mathcal{C}[U]) \equiv \mathcal{C}[\mathrm{lo}^k(U)]$ whenever $\mathrm{lo}^j(U)$ is reducible for all $j < k$.
\end{definition}

\begin{theorem}[All-step $\to_{\mathrm{lo}}$ closed forms, triangular staircase law, and $+14$ twin shift]
\label{thm:all-step-jump-ahead}
Let $\mathrm{lo}^t(M)$ denote the $t$-th iterate of deterministic leftmost-outermost reduction starting from $M \in \TS$.
\begin{enumerate}[label=(\roman*),leftmargin=2.2em]
\item \emph{(All-step closed form and triangular staircase law on $Z_{\mathrm{Zac}}$)} Every step $t \in \mathbb{N}$ decomposes uniquely as $t = \frac{1}{2}n(n+3) + r$ with $n \in \mathbb{N}$ and $0 \le r \le n + 1$. At offset $r = 0$, $\mathrm{lo}^t(Z_{\mathrm{Zac}}) \equiv \mathcal{C}_Z(n)[G_0 G_n]$, whereas at offset $r = k + 1$ ($0 \le k \le n$), $\mathrm{lo}^t(Z_{\mathrm{Zac}}) \equiv \mathcal{C}_Z(n)\bigl[(S G_{n+1})^k(G_{n+1-k}\,G_{n+1})\bigr]$. For every $t = \frac{1}{2}n(n+3) + r$ ($0 \le r \le n + 1$), the exact leaf count is
\begin{equation}
\label{eq:zachos-cubic-leaves}
\bigl|\mathrm{lo}^t(Z_{\mathrm{Zac}})\bigr|
\;=\; 10 \;+\; \frac{2n(n+2)(n+4)}{3} \;+\; \mathbf{1}_{r \ge 1}\,\bigl(r(2n+6) - 2\bigr).
\end{equation}
Moreover, re-indexing $t$ along triangular blocks $t = \frac{1}{2}m(m+1) + j$ ($m \in \mathbb{N}$, $0 \le j \le m$), the first difference obeys the exact staircase law
\begin{equation}
\label{eq:zachos-staircase}
\bigl|\mathrm{lo}^{t+1}(Z_{\mathrm{Zac}})\bigr| - \bigl|\mathrm{lo}^t(Z_{\mathrm{Zac}})\bigr|
\;=\; 2(m + 2)
\qquad (0 \le j \le m),
\end{equation}
generating the blockwise step increments $+4;\; +6, +6;\; +8, +8, +8;\; +10, +10, +10, +10;\; \dots$
\item \emph{(All-step closed form on $W_2$ and constant $+14$ twin shift law)} On $W_2 \equiv S S S (S T) T$, initial leaf counts at $t = 0, 1, 2, 3$ are $8, 10, 11, 16$. Every step $t \ge 4$ decomposes uniquely as $t = \frac{1}{2}(n+1)(n+6) + 1 + r$ with $n \in \mathbb{N}$ and $0 \le r \le n + 3$, where
\[
\mathrm{lo}^t(W_2) \;\equiv\;
\begin{cases}
\mathcal{C}_{W_2}(n)[E_0 E_{n+1}] & \textrm{if } r = 0, \\
\mathcal{C}_{W_2}(n)\bigl[(S T E_{n+1})\,E_{n+2}\bigr] & \textrm{if } r = 1, \\
\mathcal{C}_{W_2}(n)\bigl[E_{n+3}(E_{n+1} E_{n+2})\bigr] & \textrm{if } r = 2, \\
\mathcal{C}_{W_2}(n)\Bigl[S\,\bigl((S E_{n+2})^j(E_{n+1-j} E_{n+2})\bigr)\,\bigl(E_{n+2}(E_{n+1} E_{n+2})\bigr)\Bigr] & \textrm{if } r = j + 3 \; (0 \le j \le n),
\end{cases}
\]
with exact cubic leaf count
\begin{equation}
\label{eq:w2-cubic-leaves}
\bigl|\mathrm{lo}^t(W_2)\bigr|
\;=\; \frac{4n^3 + 51n^2 + 197n + 138}{6} \;+\; \delta_{W_2}(n, r),
\end{equation}
for $\delta_{W_2}(n, 0) = 0$, $\delta_{W_2}(n, 1) = 2n + 7$, $\delta_{W_2}(n, 2) = 4n + 16$, and $\delta_{W_2}(n, j + 3) = (8n + 33) + j(2n + 9)$. Furthermore, the $12$-leaf twin $W_2^{\mathrm{twin}} \equiv E_0 E_0$ shares identical $(\mathcal{A}_{\SN}, \mathcal{A}_P)$ evaluations $(\infty_{\SN}, \infty_P)$ and order-$1$ solvability ($\sem{W_2}_{\HN} = q_{24}$, $\sem{W_2^{\mathrm{twin}}}_{\HN} = q_{23}$) with $W_2$, satisfies $W_2 \neq_S W_2^{\mathrm{twin}}$, yet obeys for every $t \ge 3$ the exact constant-offset identity
\begin{equation}
\label{eq:w2-twin-shift-14}
\bigl|\mathrm{lo}^{t+1}(W_2)\bigr| + 14 \;=\; \bigl|\mathrm{lo}^t(W_2^{\mathrm{twin}})\bigr|.
\end{equation}
\item \emph{(All-step exact leaf formula on $D_{\mathrm{seed}}$)} On $D_{\mathrm{seed}}$, initial leaf counts at $t = 0, 1$ are $10, 11$. Every step $t \ge 2$ decomposes uniquely as $t = (2N^2 - N + 1) + d$ with $N \ge 1$ and $0 \le d \le 4N$; writing $r = \lfloor d / 4 \rfloor$, one has
\begin{equation}
\label{eq:seed-cubic-leaves}
\bigl|\mathrm{lo}^t(D_{\mathrm{seed}})\bigr|
\;=\; \frac{5N^3 + 15N^2 - 8N + 33}{3} \;+\; r(6N + 13 - r) \;+\; \mathbf{1}_{d \not\equiv 0 \pmod{4}}\,(2N + 4).
\end{equation}
\end{enumerate}
\end{theorem}

\begin{proof}
For~(i), because $\tau_Z(n+1) - \tau_Z(n) = \frac{1}{2}(n+1)(n+4) - \frac{1}{2}n(n+3) = n + 2$, the intervals $\{\tau_Z(n) + r : 0 \le r \le n + 1\}$ partition $\mathbb{N}$. From $\mathcal{C}_Z(n)[G_0 G_n]$, one contraction inside the active focus gives $G_0 G_n \equiv S T T G_n \to_{\mathrm{lo}} (T G_n)(T G_n) \equiv G_{n+1} G_{n+1}$, and $k$ subsequent contractions ($0 \le k \le n + 1$) peel Spine-$2$ right-hole shells $G_{m+1} G_{n+1} \equiv S S G_m G_{n+1} \to_{\mathrm{lo}} S G_{n+1} (G_m G_{n+1})$ since $G_{n+1} \in \mathcal{N}$; at $k = n + 1$, the accumulated shells merge into $\mathcal{C}_Z(n+1)[G_0 G_{n+1}]$. Induction on $n$ gives $3\,\lambda(\mathcal{C}_Z(n)) = \sum_{i=1}^n 3i(1 + |G_i|) = \sum_{i=1}^n 6i(i + 3) = n(n+1)(2n+10)$. Adding $|G_0 G_n| = 5 + (2n + 5) = 2n + 10$ gives $3\,\lambda(\mathcal{C}_Z(n)) + 3(2n + 10) = 2n(n+2)(n+4) + 30$, proving~\eqref{eq:zachos-cubic-leaves} at $r = 0$. For $r = k + 1 \ge 1$, the $k$ inner frames $(S G_{n+1})^k$ and focus $G_{n+1-k} G_{n+1}$ contribute $k(2n + 8) + (2(n + 1 - k) + 5) + (2n + 7) = (2n + 10) + (k + 1)(2n + 6) - 2$, yielding~\eqref{eq:zachos-cubic-leaves} for all $r$.

To deduce~\eqref{eq:zachos-staircase}, at $m = 0$ ($j = 0, t = 0$) we have $14 - 10 = 4 = 2(0 + 2)$. For $m = n + 1 \ge 1$, note that $\frac{1}{2}(n+1)(n+2) = \frac{1}{2}n(n+3) + 1 = \tau_Z(n) + 1$. When $0 \le j \le n$, step $t = \tau_Z(n) + (j + 1)$ advances to $t + 1 = \tau_Z(n) + (j + 2)$ within cycle $n$ ($r = j + 1 \to j + 2 \le n + 2$, valid since the formula at $r = n + 2$ agrees with cycle $n + 1$ at $r = 0$), increasing leaf count by $((j + 2)(2n + 6) - 2) - ((j + 1)(2n + 6) - 2) = 2n + 6 = 2(m + 2)$. At the block endpoint $j = n + 1$, $t = \tau_Z(n) + (n + 2) = \tau_Z(n + 1)$ transitions from $r = 0$ to $r = 1$ in cycle $n + 1$, increasing leaf count by $1 \cdot (2(n + 1) + 6) - 2 = 2n + 6 = 2(m + 2)$.

For~(ii), Theorem~\ref{thm:wolfram-14-dichotomy}(iii) yields the intra-cycle transitions of $W_2$ at offsets $r = 0, 1, 2$ and $r = j + 3$ ($0 \le j \le n + 1$), with $\tau_{W_2}(n+1) - \tau_{W_2}(n) = n + 4$. Each cycle $i \to i + 1$ appends one left-hole frame $S\,\square\,[E_{i+2}(E_{i+1} E_{i+2})]$ of extra leaf weight $1 + (2i + 10) + (2i + 8) + (2i + 10) = 6i + 29$ and $i + 1$ right-hole frames $S E_{i+2} \square$ of weight $(i + 1)(2i + 11)$, so induction on $n$ gives $6\,\lambda(\mathcal{C}_{W_2}(n)) = 4n^3 + 51n^2 + 185n + 54$. Adding $|E_0 E_{n+1}| = 6 + (2n + 8) = 2n + 14$ at $r = 0$ gives $(4n^3 + 51n^2 + 197n + 138)/6$, and inspecting the focus and local frame weights at $r = 1, 2, j + 3$ yields $\delta_{W_2}(n, r)$ and~\eqref{eq:w2-cubic-leaves}. Next, for $W_2^{\mathrm{twin}} \equiv E_0 E_0$, direct automaton evaluation gives $(\sem{W_2^{\mathrm{twin}}}_{\SN}, \sem{W_2^{\mathrm{twin}}}_{\HN}, \sem{W_2^{\mathrm{twin}}}_P) = (\infty_{\SN}, q_{23}, \infty_P)$. Three leftmost-outermost contractions reduce
\[
W_2^{\mathrm{twin}} \;\equiv\; S(S T)T\,E_0
\;\longrightarrow_{\mathrm{lo}}\; (S T E_0)\,E_1
\;\longrightarrow_{\mathrm{lo}}\; E_2\,(E_0 E_1)
\;\longrightarrow_{\mathrm{lo}}\; S\,(E_0 E_1)\,\bigl(E_1(E_0 E_1)\bigr),
\]
whereas $W_2 \to_{\mathrm{lo}}^4 S\,E_1\,(E_0 E_1)$. If $W_2 \convS W_2^{\mathrm{twin}}$, then the Spine-$2$ root normal forms $S\,E_1\,(E_0 E_1)$ and $S\,(E_0 E_1)\,(E_1(E_0 E_1))$ would be $\convS$-convertible, forcing $S E_1 \convS S(E_0 E_1)$ by Proposition~\ref{prop:q30-decidability}(iii); this is impossible because $S E_1 \in \mathcal{N} \subseteq \SN$ while $\sem{S(E_0 E_1)}_{\SN} = \infty_{\SN}$ ($E_0 E_1 \notin \SN$). Yet both $\mathrm{lo}^4(W_2) \equiv \mathcal{K}_1[E_0 E_1]$ ($\mathcal{K}_1[\square] \equiv S\,E_1\,\square$, $\lambda(\mathcal{K}_1) = 1 + 8 = 9$) and $\mathrm{lo}^3(W_2^{\mathrm{twin}}) \equiv \mathcal{K}_2[E_0 E_1]$ ($\mathcal{K}_2[\square] \equiv S\,\square\,[E_1(E_0 E_1)]$, $\lambda(\mathcal{K}_2) = 1 + 8 + 6 + 8 = 23$) wrap valid head-normal shells around the identical non-normalizing active focus $E_0 E_1$. For every $k \in \mathbb{N}$ ($t = 3 + k$), $\mathrm{lo}^{4+k}(W_2) \equiv \mathcal{K}_1[\mathrm{lo}^k(E_0 E_1)]$ and $\mathrm{lo}^{3+k}(W_2^{\mathrm{twin}}) \equiv \mathcal{K}_2[\mathrm{lo}^k(E_0 E_1)]$, so their leaf counts differ constantly by $\lambda(\mathcal{K}_2) - \lambda(\mathcal{K}_1) = 23 - 9 = 14$, proving~\eqref{eq:w2-twin-shift-14}.

For~(iii), in cycle $N = m + 1 \ge 1$ starting at step $2N^2 - N + 1$ with $\mathcal{C}_{\mathrm{seed}}(m)[D_{N,N}]$ (Proposition~\ref{prop:jump-ahead}(ii)), the $r$-th four-step countdown ($0 \le r \le m$) peels the frozen right sibling $(C_m S)\,D_{m-j, m+1}$ for $j = 0, \dots, r - 1$ of weight $1 + (2m + 6) + (2(m - j) + 2(m + 1) + 11) = 6m + 20 - 2j$, summing over $0 \le j < r$ to $r(6m + 21 - r) = r(6N + 15 - r)$. Summing full cycles $1 \le i \le m$ yields $3\,\lambda(\mathcal{C}_{\mathrm{seed}}(m)) = 5m^3 + 30m^2 + 25m$, and adding $|D_{N,N}| = 4N + 11$ gives $(5N^3 + 15N^2 - 8N + 33)/3$ at $d = 0$. Writing $d = 4r + u$ ($0 \le u \le 3$), decreasing $D_{N,N}$ to $D_{N-r, N}$ subtracts $2r$ from the focus while adding $r(6N + 15 - r)$ to the shell, for a net gain of $r(6N + 13 - r)$ at $u = 0$. When $u \in \{1, 2, 3\}$ ($r < N$, $k = N - 1 - r$), the first contraction $D_{k+1, N} \equiv (S S C_k) C_N S \to_{\mathrm{lo}} (S C_N)(C_k C_N) S$ duplicates $C_N$ ($|C_N| = 2N + 5$), adding $(2N + 5) - 1 = 2N + 4$ leaves, while the second and third contractions $((S C_N)(C_k C_N)) S \to_{\mathrm{lo}} (C_N S)(C_k C_N S) \to_{\mathrm{lo}} (S S (C_{N-1} S))(C_k C_N S)$ have atomic third argument $S$ and therefore preserve leaf count identically at $+(2N + 4)$, proving~\eqref{eq:seed-cubic-leaves}.
\end{proof}

Wolfram also asked whether subterms reachable somewhere in the multiway reduction graph $\redSs$ must eventually appear along deterministic leftmost-outermost evolution $\to_{\mathrm{lo}}$. Analyzing the multiway invariant of Barendregt's unsolvable term $A A A$ yields negative answers in both the normalizing and perpetually evolving regimes.

\begin{proposition}[Multiway grammar on $A A A$, universal $S(S S)$ exclusion, and subterm separations]
\label{prop:aaa-grammar-and-subterm-sep}
Let $\mathcal{G}_A \subseteq \TS$ be the inductive class generated by the three rules $A \in \mathcal{G}_A$ ($A \equiv S S S$), $U \in \mathcal{G}_A \implies S U \in \mathcal{G}_A$, and $U, V \in \mathcal{G}_A \implies U V \in \mathcal{G}_A$ (coinciding on closed terms with $\mathcal{H}$ of Definition~\ref{def:heavy-invariant}).
\begin{enumerate}[label=(\roman*),leftmargin=2.2em]
\item \emph{(Multiway closure, $+2$ leaf growth, and universal $S(S S)$ exclusion on $A A A$)} Every $U \in \mathcal{G}_A$ satisfies $|U| \ge 3$, and every one-step multiway contraction $U \redS U'$ from $U \in \mathcal{G}_A$ satisfies $U' \in \mathcal{G}_A$ and $|U'| \ge |U| + 2$. Consequently, along $\to_{\mathrm{lo}}$, $A A A \in \mathcal{G}_A$ executes pure head reduction with $|\mathrm{lo}^{k+1}(A A A)| \ge |\mathrm{lo}^k(A A A)| + 2$ and $|\mathrm{lo}^k(A A A)| \ge 9 + 2k$ for all $k \in \mathbb{N}$, and across the entire infinite multiway reduction set $\{ N \in \TS : A A A \redSs N \}$, the $3$-leaf normal form $S(S S)$ never occurs as a subterm.
\item \emph{(Halting and non-halting subterm-appearance separations)} On the $7$-leaf normalizing term $M_{\mathrm{sep}}^{\SN} \equiv T(S S S S)S$, the $6$-leaf normal form $P_{\mathrm{sep}}^{\SN} \equiv T(T T)$ occurs as a subterm in one multiway step ($M_{\mathrm{sep}}^{\SN} \redS T(T T)S$), yet never occurs in $\mathrm{lo}^k(M_{\mathrm{sep}}^{\SN})$ for any $k \in \mathbb{N}$. On the $14$-leaf non-normalizing term $M_{\mathrm{sep}}^{\infty} \equiv S(A A A)(A T) \in \HN \setminus \SN$, the $3$-leaf normal form $S(S S)$ occurs as a subterm in one multiway step ($M_{\mathrm{sep}}^{\infty} \redS S(A A A)\bigl(S(S S)(S(S S))\bigr)$), yet $\mathrm{lo}^k(M_{\mathrm{sep}}^{\infty})$ reduces forever without $S(S S)$ ever occurring as a subterm at any step $k \in \mathbb{N}$.
\end{enumerate}
\end{proposition}

\begin{proof}
For~(i), structural induction on $U \in \mathcal{G}_A$ gives $|U| \ge 3$ (hence $U \not\equiv S$ and $U \not\equiv T$). Inversion on $\mathcal{G}_A$ shows $S U \in \mathcal{G}_A \implies U \in \mathcal{G}_A$ (since $S \notin \mathcal{G}_A$ rules out the binary-application constructor) and $S X Y \in \mathcal{G}_A \implies (X \equiv S \land Y \equiv S) \vee (X \in \mathcal{G}_A \land Y \in \mathcal{G}_A)$. If $U \in \mathcal{G}_A$ contracts at the root as $U \equiv S X Y Z \redS (X Z)(Y Z) \equiv U'$, then $Z \in \mathcal{G}_A$ ($|Z| \ge 3$), and either $(X, Y) = (S, S)$ giving $(S Z)(S Z) \in \mathcal{G}_A$ or $X, Y \in \mathcal{G}_A$ giving $(X Z)(Y Z) \in \mathcal{G}_A$; in both cases $|U'| - |U| = |Z| - 1 \ge 2$. Contextual steps lift directly by induction because neither $S$ nor $S S$ admits an internal step in $A \equiv (S S) S$ or $S U$. Since $A A A \in \mathcal{G}_A$ and $\sem{A A A}_{\HN} = \infty_{\HN}$, every $\to_{\mathrm{lo}}$ step is a head step satisfying $|\mathrm{lo}^{k+1}(A A A)| \ge |\mathrm{lo}^k(A A A)| + 2$, giving $|\mathrm{lo}^k(A A A)| \ge 9 + 2k$. Finally, we check by induction on $U \in \mathcal{G}_A$ that $S(S S)$ is not a subterm of $U$: the subterms of $A \equiv (S S) S$ are $A, S S, S$; a node $S W$ with $W \in \mathcal{G}_A$ cannot equal $S(S S)$ at the root because $|W| \ge 3 > 2 = |S S|$; and a node $W_1 W_2$ with $W_1, W_2 \in \mathcal{G}_A$ cannot equal $S(S S)$ at the root because $|W_1| \ge 3 > 1 = |S|$. Since $\mathcal{G}_A$ is closed under $\redSs$, no multiway reduct $N$ of $A A A$ contains $S(S S)$.

For~(ii), contracting $S S S S \redS T T$ inside $M_{\mathrm{sep}}^{\SN} \equiv T(S S S S)S$ produces $T(T T)S$ in one step, whereas $\to_{\mathrm{lo}}$ reduces $T(S S S S)S \equiv S S (S S S S) S \to_{\mathrm{lo}} T((S S S S)S) \to_{\mathrm{lo}} T(T T S) \to_{\mathrm{lo}} T(S S(T S)) \in \mathcal{N}$, none of which contains $T(T T)$. Similarly, contracting $A T \equiv S S S (S S) \redS S(S S)(S(S S))$ on the right of $M_{\mathrm{sep}}^{\infty}$ produces $S(S S)$ in one step; under $\to_{\mathrm{lo}}$, however, the Spine-$2$ left-hole shell $S\,\square\,(A T)$ keeps the right sibling $A T$ permanently unreduced while stepping forever inside $A A A$, so $\mathrm{lo}^k(M_{\mathrm{sep}}^{\infty}) \equiv S\,\mathrm{lo}^k(A A A)\,(A T)$ for all $k \in \mathbb{N}$. Neither $A T \equiv ((S S) S)(S S)$ nor $\mathrm{lo}^k(A A A) \in \mathcal{G}_A$ contains $S(S S)$ by~(i), and $|\mathrm{lo}^k(A A A)| \ge 9$ precludes $S\,\mathrm{lo}^k(A A A) \equiv S(S S)$.
\end{proof}

\subsection{Garden-of-Eden and Isolated Normal Forms: Characterization, Decidability, and Mahler Growth}
\label{subsec:goe-isolated}

In cellular automata and combinatory rewriting~\cite{Wolfram2002NKS,Wolfram2021CombinatorsBook}, a configuration or term is called \emph{Garden-of-Eden} if it cannot be reached by one or more rewrite steps from any predecessor. On closed normal forms $\mathcal{N} = \NF \cap \TS$, Garden-of-Eden status connects directly to Dana Scott's convertibility problem (Question~30, Proposition~\ref{prop:q30-decidability}(ii)): a closed normal form has in-degree zero in $\CLS$ if and only if its $\convS$-equivalence class is a singleton.

\begin{definition}[Garden-of-Eden terms and isolated normal forms]
\label{def:goe-isolated}
A closed pure-$S$ term $P \in \TS$ is \emph{Garden-of-Eden} ($P \in \mathrm{GoE}$) if there is no closed term $M \in \TS$ such that $M \redSp P$. A term $P \in \TSV$ is an \emph{isolated normal form} ($P \in \mathrm{IsoNF}$) if $P \in \NF$ and there is no $M \in \TSV$ such that $M \redS P$. We write $\#\mathrm{IsoNF}(n) := \#\{ P \in \TS \cap \mathrm{IsoNF} : |P| = n \}$ and $\#\mathcal{N}(n) := \#\{ P \in \mathcal{N} : |P| = n \}$.
\end{definition}

\begin{theorem}[Characterization, tower law, Motzkin hierarchy, and Mahler asymptotics of Garden-of-Eden normal forms]
\label{thm:goe-characterization-and-mahler}
\begin{enumerate}[label=(\roman*),leftmargin=2.2em]
\item \emph{(Equivalence, backward decomposition, and syntactic decidability)} For every $P \in \TSV$, $[P]_{\convS} = \{P\}$ holds if and only if $P \in \mathrm{IsoNF}$; for closed normal forms $P \in \mathcal{N}$, $P \in \mathrm{GoE} \iff P \in \mathrm{IsoNF}$. Every closed term with $|P| \le 3$ belongs to $\mathrm{GoE}$, and at $|P| = 4$, $S S (S S)$ is the unique closed normal form outside $\mathrm{GoE}$. For closed normal forms $U, V \in \mathcal{N}$, one has $S U \in \mathrm{GoE} \iff U \in \mathrm{GoE}$, whereas $S U V \in \mathrm{GoE}$ holds if and only if $U \in \mathrm{GoE}$, $V \in \mathrm{GoE}$, and there exists no $Y_0 \in \mathcal{N}$ with $|Y_0| + |U| \le |V|$ and $Y_0 U \redSs V$---a condition that, given $V \in \mathrm{GoE}$, reduces to the purely syntactic check $\forall Y \in \TS,\, V \not\equiv Y U$. Consequently, membership in $\mathcal{N} \cap \mathrm{GoE}$ is constructively decidable in linear time on tree representations with subterm equality (or quadratic leaf-comparison time).
\item \emph{(Constructive tower law and isolation of order-$3$ combs)} Whenever $U, V \in \mathcal{N} \cap \mathrm{GoE}$ satisfy $|V| \le |U|$, the Spine-$2$ normal form $S U V$ lies in $\mathrm{GoE}$. Moreover, every closed order-$3$ right comb $R_{n-1}$ ($n \ge 1$) lies in $\mathrm{IsoNF}$.
\item \emph{(Shifted Motzkin law for $\mathcal{N}$, small-$n$ census, and constructive exponential decay)} For every $n \ge 1$, $\#\mathcal{N}(n) = M_{n-1}$ equals the shifted Motzkin number defined by $M_0 = 1$ and $\#\mathcal{N}(m+2) = \#\mathcal{N}(m+1) + \sum_{a+b=m+1} \#\mathcal{N}(a)\,\#\mathcal{N}(b)$ ($m \ge 0$, dominant growth rate $\gamma_{\NF} = 3$). Constructively for every $n \ge 1$,
\begin{equation}
\label{eq:iso-cnf-sn-hierarchy}
1 \;\le\; \#\mathrm{IsoNF}(n) \;\le\; \#\mathcal{N}(n) \;\le\; \#\SN(n),
\qquad
\frac{\#\mathrm{IsoNF}(n)}{C_{n-1}} \;\le\; \frac{\#\mathcal{N}(n)}{C_{n-1}} \;\le\; 51.39438 \cdot \Bigl(\frac{250}{259}\Bigr)^n.
\end{equation}
For leaf sizes $1 \le n \le 7$, the exact ratio $\#\mathrm{IsoNF}(n) / \#\mathcal{N}(n)$ takes the values $1/1,\; 1/1,\; 2/2,\; 3/4,\; 7/9,\; 15/21,\; 37/51$ (continuing $87/127,\; 218/323,\; 546/835,\; 1393/2188$ for $8 \le n \le 11$).
\item \emph{(Bilinear exclusion recurrence and Fibonacci--Mahler growth constant)} Setting $I_n := \#\mathrm{IsoNF}(n)$, the count $V(a, b)$ of valid Spine-$2$ pairs $(U, V) \in (\mathcal{N} \cap \mathrm{IsoNF})^2$ with $|U| = a$, $|V| = b$, and $S U V \in \mathrm{IsoNF}$ satisfies
\begin{equation}
\label{eq:iso-bilinear-recurrence}
V(a, b) \;=\; I_a I_b \;-\; \mathbf{1}_{b = a + 1}\,I_a \;-\; \mathbf{1}_{b \ge a + 2}\,V(b - a - 1, a),
\qquad
I_n \;=\; I_{n-1} + \sum_{a=1}^{n-2} V(a, n - 1 - a) \quad (n \ge 2).
\end{equation}
Defining Fibonacci numbers $F_1 = F_2 = 1$ and $F_{k+2} = F_{k+1} + F_k$, the ordinary generating function $F_{\mathrm{Iso}}(z) = \sum_{n=1}^{\infty} I_n z^n$ satisfies the exact Fibonacci--Mahler functional equation
\begin{equation}
\label{eq:iso-mahler}
F_{\mathrm{Iso}}(z) \;=\; z + z\,F_{\mathrm{Iso}}(z) + \sum_{k=1}^{\infty} (-1)^{k-1} z^{F_{k+2}-1}\Bigl(F_{\mathrm{Iso}}\bigl(z^{F_k}\bigr)\,F_{\mathrm{Iso}}\bigl(z^{F_{k+1}}\bigr) - z^{F_{k+1}} F_{\mathrm{Iso}}\bigl(z^{F_{k+2}}\bigr)\Bigr),
\end{equation}
whose dominant singularity is a square-root branch point at $r_{\mathrm{Iso}} \approx 0.3439236475 > \frac{1}{3}$, yielding the exponential growth constant $\gamma_{\mathrm{Iso}} = r_{\mathrm{Iso}}^{-1} \approx 2.9076221057 < 3$.
\end{enumerate}
\end{theorem}

\begin{proof}
For~(i), if $[P]_{\convS} = \{P\}$, acyclicity of $\redS$ (Proposition~\ref{prop:acyclicity-order}(i)) forbids both $P \redS Q$ and $M \redS P$, so $P \in \mathrm{IsoNF}$; conversely, if $P \in \mathrm{IsoNF}$ and $P \convS Q$, Church--Rosser gives a common reduct $D$ with $P \redSs D$ and $Q \redSs D$, forcing $D \equiv P$ by $P \in \NF$ and $Q \equiv P$ because any non-empty reduction $Q \redSp P$ ends in a final single step $M' \redS P$. Since any predecessor $M \redS P$ of a closed term $P \in \TS$ is itself closed ($S X Y Z \redS (X Z)(Y Z)$ preserves variable sets from contractum to redex), $\mathrm{GoE}$ and absence of single-step predecessors coincide on $\TS$, so $P \in \mathrm{GoE} \iff P \in \mathrm{IsoNF}$ for $P \in \mathcal{N}$. Because every redex $S X Y Z$ and contractum $(X Z)(Y Z)$ has leaf count at least~$4$, all closed terms of size $\le 3$ ($S, S S, S(S S), S S S$) lie in $\mathrm{GoE}$; at size~$4$, the four normal forms are $S(S(S S))$, $S(S S S)$, $S(S S)S$, and $S S (S S)$, among which only $S S (S S)$ has a predecessor ($S S S S \redS S S (S S)$).

For $S U \in \mathcal{N}$, no root contraction can produce $S U$ because any root contractum $(X Z)(Y Z)$ has compound left child $X Z$ ($|X Z| \ge 2 > 1 = |S|$), and $S$ is irreducible; hence $M \redSp S U \iff \exists U_0,\, M \equiv S U_0 \land U_0 \redSp U$, giving $S U \in \mathrm{GoE} \iff U \in \mathrm{GoE}$. For $S U V \in \mathcal{N}$ ($U, V \in \mathcal{N}$), backward induction along $M \redSs S U V$ shows that either no root step occurs ($M \equiv S U_0 V_0$ with $U_0 \redSs U$ and $V_0 \redSs V$) or a last root step $S S Y Z \redS (S Z)(Y Z)$ occurs ($M \redSs S S Y Z$ with $Z \redSs U$ and $Y Z \redSs V$). In the second case, $Y Z \in \WN = \SN$ implies $Y \in \SN$, so $Y \to_{\mathrm{lo}}^* Y_0 \in \mathcal{N}$, and confluence on $Y_0 U \leftarrow_S^* Y Z \redSs V \in \mathcal{N}$ gives $Y_0 U \redSs V$ with $|Y_0| + |U| \le |V|$ (conversely, any such $Y_0$ gives $S S Y_0 U \redS (S U)(Y_0 U) \redSs S U V$, so $S S Y_0 U \redSp S U V$). Moreover, once $V \in \mathrm{GoE}$, no closed term can reduce in one or more steps to $V$, so $Y_0 U \redSs V$ holds if and only if $Y_0 U \equiv V$, which is equivalent to $V$ having syntactic form $Y U$ (since any subterm $Y$ of $V \in \mathcal{N}$ automatically lies in $\mathcal{N}$).

For~(ii), if $|V| \le |U|$, then for every $Y \in \TS$ we have $|Y U| = |Y| + |U| \ge 1 + |U| > |V|$, ruling out $V \equiv Y U$ and giving $S U V \in \mathrm{GoE}$. Induction on $n$ using $S U \in \mathrm{GoE} \iff U \in \mathrm{GoE}$ immediately places $R_{n-1} \in \mathrm{IsoNF}$ for every $n \ge 1$.

For~(iii), every closed normal form of size $m + 2 \ge 2$ has spine length~$1$ or~$2$, decomposing disjointly into $S U$ ($U \in \mathcal{N}$ of size $m + 1$) and $(S A_1) A_2$ ($A_1, A_2 \in \mathcal{N}$ with $|A_1| + |A_2| = m + 1$), establishing the shifted Motzkin recurrence $\#\mathcal{N}(n) = M_{n-1}$ whose generating function $M_{\NF}(z) = z + z M_{\NF}(z) + z M_{\NF}(z)^2 = \frac{1 - z - \sqrt{(1+z)(1-3z)}}{2z}$ has dominant singularity $z = \frac{1}{3}$ ($\gamma_{\NF} = 3$). Because $R_{n-1} \in \mathrm{IsoNF} \subseteq \mathcal{N} \subseteq \SN$, we have $1 \le \#\mathrm{IsoNF}(n) \le \#\mathcal{N}(n) \le \#\SN(n)$, and~\eqref{eq:iso-cnf-sn-hierarchy} follows from the constructive $\SN$ decay bound proved in Theorem~\ref{thm:asymptotics}(ii) below; direct recursive evaluation verifies the exact counts through $n = 11$.

For~(iv), let $\mathcal{I}_m := \{ P \in \mathcal{N} \cap \mathrm{IsoNF} : |P| = m \}$ ($I_m = \#\mathcal{I}_m$). By~(i), $(U, V) \in \mathcal{I}_a \times \mathcal{I}_b$ fails $S U V \in \mathrm{IsoNF}$ if and only if $V \equiv Y U$ for some $Y$. When $b \le a$, no such $Y$ exists, so $V(a, b) = I_a I_b$. When $b = a + 1$, $|Y| = 1$ forces $Y \equiv S$ and $V \equiv S U \in \mathcal{I}_{a+1}$ for each $U \in \mathcal{I}_a$, subtracting $I_a$. When $b \ge a + 2$, $|Y| = b - a \ge 2$ and $V \equiv Y U \in \mathcal{N}$ force $Y \equiv S W$ with $|W| = b - a - 1 \ge 1$, so $V \equiv S W U$; by~(i), $S W U \in \mathcal{I}_b$ (given $U \in \mathcal{I}_a$) holds if and only if $W \in \mathcal{I}_{b-a-1}$ and $U \not\equiv Z W$, which bijects the excluded pairs $(U, V)$ onto the $V(b - a - 1, a)$ valid pairs $(W, U) \in \mathcal{I}_{b-a-1} \times \mathcal{I}_a$, proving~\eqref{eq:iso-bilinear-recurrence}. Setting $S_1(x, y) := \sum_{a, b \ge 1} V(a, b)\,x^a y^b$, multiplying~\eqref{eq:iso-bilinear-recurrence} by $x^a y^b$ and summing over $a, b \ge 1$ (with $w = b - a - 1 \ge 1$ in the third term) gives
\[
S_1(x, y) \;=\; F_{\mathrm{Iso}}(x)\,F_{\mathrm{Iso}}(y) \;-\; y\,F_{\mathrm{Iso}}(x y) \;-\; y\,S_1(y, x y).
\]
Iterating $(x_k, y_k) \mapsto (y_k, x_k y_k)$ from $(x_1, y_1) = (z, z) = (z^{F_1}, z^{F_2})$ produces $(x_k, y_k) = (z^{F_k}, z^{F_{k+1}})$ with $x_k y_k = z^{F_{k+2}}$ and prefactor $(-1)^{k-1} z \prod_{j=1}^{k-1} y_j = (-1)^{k-1} z^{1 + \sum_{j=2}^k F_j} = (-1)^{k-1} z^{F_{k+2}-1}$. Substituting $F_{\mathrm{Iso}}(z) = z + z F_{\mathrm{Iso}}(z) + z S_1(z, z)$ yields~\eqref{eq:iso-mahler}. Separating $k = 1$ writes~\eqref{eq:iso-mahler} as $z\,F_{\mathrm{Iso}}(z)^2 - B(z)\,F_{\mathrm{Iso}}(z) + C(z) = 0$, where $B(z) = 1 - z + z^2 F_{\mathrm{Iso}}(z^2)$ and $C(z) = z - z^2 F_{\mathrm{Iso}}(z^2) + z^4 F_{\mathrm{Iso}}(z^3) + \sum_{k=3}^{\infty} (-1)^{k-1} z^{F_{k+2}-1}(F_{\mathrm{Iso}}(z^{F_k}) F_{\mathrm{Iso}}(z^{F_{k+1}}) - z^{F_{k+1}} F_{\mathrm{Iso}}(z^{F_{k+2}}))$. Since $I_n \le M_{n-1} \le 3^n$ gives $r_{\mathrm{Iso}} \ge \frac{1}{3}$, for $|z| < 0.58 < \sqrt{1/3} \le \sqrt{r_{\mathrm{Iso}}}$ every higher Mahler argument satisfies $|z^{F_k}| \le |z|^2 < r_{\mathrm{Iso}}$, so $B(z)$ and $C(z)$ are analytic across a neighborhood of $r_{\mathrm{Iso}}$. Solving the discriminant vanishing condition $B(r_{\mathrm{Iso}})^2 - 4 r_{\mathrm{Iso}} C(r_{\mathrm{Iso}}) = 0$ locates the simple root $r_{\mathrm{Iso}} = 0.3439236474519548\dots > \frac{1}{3}$, yielding $\gamma_{\mathrm{Iso}} = r_{\mathrm{Iso}}^{-1} = 2.9076221056875632\dots < 3$.
\end{proof}

\subsection{Small-Term Catalan Census and Sharp Minimality Thresholds}
\label{subsec:catalan-minimality}

By Section~\ref{subsec:prelim-syntax}, the closed $S$-terms of leaf size $n \ge 1$ have cardinality $C_{n-1} = \frac{1}{n}\binom{2n-2}{n-1}$. Table~\ref{tab:catalan-census} reports the exact cardinalities of isolated normal forms ($\mathrm{IsoNF}$), closed normal forms ($\mathcal{N}$), strongly normalizing terms ($\SN$), productive terms ($\Prod$), and the four functional order classes $\mathrm{Ord}_k := \{ M \in \TS : \ord(M) = k \}$ ($k \in \{0, 1, 2, 3\}$, where $\mathrm{Ord}_0 = \overline{\HN}$ and $\mathrm{Ord}_1 \sqcup \mathrm{Ord}_2 \sqcup \mathrm{Ord}_3 = \HN$) for $1 \le n \le 11$ (a cumulative total of $23{,}714$ trees, extended by dynamic programming through $n = 150$).

\begin{table}[htbp]
\centering
\caption{Exact Catalan phase and functional-order census of closed $S$-terms for leaf sizes $1 \le n \le 11$ and sharp onset thresholds.}
\label{tab:catalan-census}
\footnotesize
\setlength{\tabcolsep}{3.6pt}
\begin{tabular}{@{}rrrrrrrrrrl@{}}
\toprule
$n$ & $C_{n-1}$ & $\#\mathrm{IsoNF}$ & $\#\mathcal{N}$ & $\#\SN$ & $\#\Prod$ & $\#\mathrm{Ord}_3$ & $\#\mathrm{Ord}_2$ & $\#\mathrm{Ord}_1$ & $\#\mathrm{Ord}_0$ & Sharp Structural Threshold \\
\midrule
$1$ & $1$ & $1$ & $1$ & $1$ & $1$ & $1$ & $0$ & $0$ & $0$ & Unique atom $S$ ($\ord = 3$) \\
$2$ & $1$ & $1$ & $1$ & $1$ & $1$ & $1$ & $0$ & $0$ & $0$ & Unique Spine-$1$ comb $S S$ ($\ord = 3$) \\
$3$ & $2$ & $2$ & $2$ & $2$ & $2$ & $1$ & $1$ & $0$ & $0$ & Minimal $\ord = 2$ ($S S S$); all in $\mathrm{IsoNF}$ \\
$4$ & $5$ & $3$ & $4$ & $5$ & $5$ & $1$ & $4$ & $0$ & $0$ & Minimal redex $S S S S$; non-$\mathrm{GoE}$ $S S(S S)$ \\
$5$ & $14$ & $7$ & $9$ & $14$ & $14$ & $1$ & $12$ & $1$ & $0$ & Unique minimal $\ord = 1$ ($S(S S S)S$) \\
$6$ & $42$ & $15$ & $21$ & $42$ & $42$ & $1$ & $32$ & $9$ & $0$ & Universal $\SN$ and deterministic $\redS$ ($n \le 6$) \\
$7$ & $132$ & $37$ & $51$ & $130$ & $130$ & $1$ & $83$ & $48$ & $0$ & Minimal $\overline{\SN}$ ($X_1, X_2$); $1$-step diamond break \\
$8$ & $429$ & $87$ & $127$ & $388$ & $388$ & $1$ & $240$ & $173$ & $15$ & Minimal $\mathrm{Ord}_0 = \overline{\HN}$ ($15$); $\Prod = \SN$ ($n \le 8$) \\
$9$ & $1{,}430$ & $218$ & $323$ & $1{,}154$ & $1{,}155$ & $1$ & $739$ & $554$ & $136$ & Unique minimal $\Prod \setminus \SN$ ($Z_0$); $15$ in $\RNF\setminus\HN$ \\
$10$ & $4{,}862$ & $546$ & $835$ & $3{,}381$ & $3{,}383$ & $1$ & $2{,}380$ & $1{,}740$ & $741$ & $S Z_0$ and Zachos's $Z_{\mathrm{Zac}} \equiv G_0 G_0$ appear \\
$11$ & $16{,}796$ & $1{,}393$ & $2{,}188$ & $9{,}967$ & $9{,}972$ & $1$ & $7{,}910$ & $5{,}685$ & $3{,}200$ & Normalizing fraction $\#\SN / C_{10} \approx 59.34\%$ \\
\bottomrule
\end{tabular}
\end{table}

\begin{theorem}[Sharp small-term minimality and confluence thresholds]
\label{thm:small-term-minimality}
Across all closed $S$-terms $M \in \TS$:
\begin{enumerate}[label=(\roman*),leftmargin=2.2em]
\item \emph{(Universal $\SN$ and determinism up to size $6$)} Every closed $S$-term with $|M| \le 6$ ($65$ terms in total) is strongly normalizing and satisfies $\#\mathrm{Redex}(M) \le 1$, so one-step multiway reduction $M \redS N$ is deterministic for $|M| \le 6$.
\item \emph{(Size-$7$ non-halting and strict diamond thresholds)} Among the $132$ closed $S$-terms with $|M| = 7$, all $132$ lie in $\HN$, $130$ lie in $\SN$, and exclusively $X_1 \equiv S(S S)S S S S$ and $X_2 \equiv S S S (S S) S S$ fail to normalize. Exactly $3$ terms of size~$7$ admit multiway branching ($\#\mathrm{Redex}(M) = 2$): $S S S (S S S S)$, $S S (S S S S) S$, and $S(S S S S) S S$. Among all $197$ terms with $|M| \le 7$, $S S S (S S S S)$ is the unique counterexample to the strict one-step $\redS$ diamond property (joining instead in one parallel step at $\mathrm{gk}(S S S (S S S S))$).
\item \emph{(Size-$8$ unsolvability threshold and $\Prod = \SN$)} Every closed $S$-term with $|M| \le 7$ is solvable; at $|M| = 8$ ($429$ terms), $15$ terms lie in $\overline{\HN}$ (all evaluating to $q_{34}$ in $\mathcal{A}_{36}$, such as $S A S A$, $S S S(S A)S$, and $S(S S)S S S S S$). Moreover, $M \in \Prod$ if and only if $M \in \SN$ for all $626$ closed $S$-terms with $|M| \le 8$.
\item \emph{(Unique size-$9$ productive non-$\SN$ combinator and minimal $\RNF \setminus \HN$)} Out of all $2{,}056$ closed $S$-terms with $|M| \le 9$, $Z_0 \equiv S(S(S S))(S(S S))(S S)$ ($9$ leaves, reducing in one step to $Z_{\mathrm{Zac}}$) is the \emph{unique} productive non-normalizing $S$-term. Size~$9$ is also the minimal leaf size of $\RNF \setminus \HN$, inhabited by the $15$ terms $U S$ ($q_{35}$ in $\mathcal{A}_{36}$) obtained by applying the $15$ size-$8$ unsolvable terms $U$ to $S$.
\end{enumerate}
\end{theorem}

\begin{proof}
Constructive tree enumeration generates for each $n \ge 1$ the duplicate-free list $\mathcal{L}_n$ of all $C_{n-1}$ closed $S$-trees via binary splits $n = n_1 + n_2$ ($n_1, n_2 \ge 1$). Applying the sound and complete algebraic characterizations of Theorems~\ref{thm:waldmann-sn}, \ref{thm:question-10}, \ref{thm:q5-main}, and~\ref{thm:order-hierarchy} across $\mathcal{L}_1, \dots, \mathcal{L}_{11}$ establishes all $\SN$, $\HN$, $\Prod$, $\mathrm{Ord}_k$, and $\RNF \setminus \HN$ counts and isolates $\{X_1, X_2\}$ at $n = 7$, the $15$ unsolvable witnesses at $n = 8$ (notice that by~\eqref{eq:footnote4-rnf}, any $M \equiv P Q \in \RNF \setminus \HN$ requires $P \in \overline{\HN}$, hence $|P| \ge 8$ and $|P Q| = |P| + |Q| \ge 8 + 1 = 9$, with equality if and only if $|P| = 8$ and $Q \equiv S$), and the unique inhabitant $Z_0$ of $\Prod \setminus \SN$ at $n \le 9$.

For the multiway branching and diamond claims in~(i)--(ii), every redex $S X Y Z$ has at least $4$ leaves; hence a tree with $\#\mathrm{Redex}(M) \ge 2$ either contains two disjoint redexes (requiring at least $4 + 4 = 8$ leaves) or nests one redex strictly inside an argument of another (requiring at least $3 + 4 = 7$ leaves, achieved exclusively by placing $S S S S$ in one of the three argument slots of $S \square_1 \square_2 \square_3$ with the remaining two slots equal to $S$). Thus $\#\mathrm{Redex}(M) \le 1$ for $|M| \le 6$. At $|M| = 7$, both $S(S S S S)S S$ and $S S (S S S S) S$ satisfy the strict one-step diamond property because $S S S S$ occurs linearly ($1$ copy) after root contraction, whereas $S S S (S S S S)$ duplicates its third argument:
\[
S S S (S S S S) \;\longrightarrow_S\; S(S S S S)\bigl(S(S S S S)\bigr) \;\equiv\; N_L,
\qquad
S S S (S S S S) \;\longrightarrow_S\; S S S \bigl(S S (S S)\bigr) \;\equiv\; N_R.
\]
Contracting the outer root of $N_R$ yields $S(S S (S S))\bigl(S(S S (S S))\bigr) \equiv \mathrm{gk}(S S S (S S S S))$, which requires \emph{two} sequential $\redS$ steps from $N_L$ (one in each duplicated copy of $S S S S$), refuting a $1$-step $\redS$ join while confirming a $1$-step parallel join.
\end{proof}

\subsection{Analytic Combinatorics, Certified Rational Order Densities, and Catalan Phase Transitions}
\label{subsec:analytic-combinatorics}

We conclude our analysis of unguided pure-$S$ rewriting by determining the asymptotic phase diagram of a uniformly random closed $S$-term of $n$ leaves as $n \to \infty$. In quantitative studies of untyped $\lambda$-calculus~\cite{BodiniGardyGittenberger2011,GrygielLescanne2013,BodiniGardyGittenbergerJacquot2013,DavidGrygielKozikRaffalliTheyssierZaionc2013}, David et al.~\cite{DavidGrygielKozikRaffalliTheyssierZaionc2013} proved that asymptotically almost all $\lambda$-terms are strongly normalizing, whereas in full combinatory logic $\mathrm{CL}(S, K)$ Bendkowski, Grygiel, Lescanne, and Zaionc~\cite{BendkowskiGrygielZaionc2015,BendkowskiGrygielZaionc2017,BendkowskiGrygielLescanneZaionc2017} showed via Woods's subgraph-containment theorem~\cite{Woods1998} that almost no $\{S, K\}$-combinators normalize. In $\CLS$, because $\SN$, $\Prod$, $\HN$, and exact functional order $\ord \in \{0, 1, 2, 3\}$ are all recognized by explicit finite deterministic $S$-algebras, their counting sequences obey positive quadratic generating-function systems~\cite{Drmota1997,FlajoletSedgewick2009} amenable to both exact rational singularity analysis and elementary all-$n$ constructive bounds over $\mathbb{N}$.

Forming the reachable product of $\mathcal{A}_{\SN} \times \mathcal{A}_P \times \mathcal{A}_{\HN}$ from $(q_0, q_0, q_0)$ yields $146$ reachable state triples: $98$ in $\SN$, $30$ in $\Prod \setminus \SN$, $17$ in $\HN \setminus \Prod$, and $1$ in $\overline{\HN}$, confirming the strict global inclusions $\mathrm{IsoNF} \cap \TS \subsetneq \mathcal{N} \subsetneq \SN \subsetneq \Prod \subsetneq \HN \subsetneq \TS$. Simultaneously, partitioning $Q_{\HN}$ by nilpotency index under $\tau_{\HN}(q) = q \odot_{\HN} \infty_{\HN}$ (Theorem~\ref{thm:order-hierarchy}(i)) decomposes $Q_{\HN} = Q_{\mathrm{ord}=0} \sqcup Q_{\mathrm{ord}=1} \sqcup Q_{\mathrm{ord}=2} \sqcup Q_{\mathrm{ord}=3}$ into $1 + 11 + 17 + 6 = 35$ states. We first establish exact all-$n$ combinatorial identities and constructive Catalan lower bounds for every order stratum.

\begin{proposition}[Catalan order partition, uniqueness of order $3$, and universal lower bounds]
\label{prop:order-catalan-partition-bounds}
Let $\#\mathrm{Ord}_k(n) := \#\{ M \in \TS : |M| = n,\; \ord(M) = k \}$ for $k \in \{0, 1, 2, 3\}$ and $n \ge 1$.
\begin{enumerate}[label=(\roman*),leftmargin=2.2em]
\item \emph{(Exact Catalan partition and universal uniqueness of order~$3$)} For every $n \ge 1$,
\[
\#\mathrm{Ord}_0(n) = \#\overline{\HN}(n), \qquad
\#\mathrm{Ord}_1(n) + \#\mathrm{Ord}_2(n) + \#\mathrm{Ord}_3(n) = \#\HN(n), \qquad
\sum_{k=0}^3 \#\mathrm{Ord}_k(n) = C_{n-1},
\]
and $\#\mathrm{Ord}_3(n) = 1$ holds identically for every $n \ge 1$ (inhabited solely by the right comb $R_{n-1}$).
\item \emph{(Universal all-$n$ positive-density lower bounds for orders $0, 1, 2$ and $\HN$)} Unconditionally for all $n$ above each sharp onset threshold:
\begin{equation}
\label{eq:order-universal-lowers}
C_{n-2} \;\le\; \#\mathrm{Ord}_2(n) + 1 \;\; (n \ge 2), \qquad
C_{n-5} \;\le\; \#\mathrm{Ord}_1(n) \;\; (n \ge 5), \qquad
C_{n-8} \;\le\; \#\mathrm{Ord}_0(n) \;\; (n \ge 8),
\end{equation}
with equality $1 = 1$ at $n = 2$ (for $\#\mathrm{Ord}_2(2) + 1$) and $n = 5$ (for $\#\mathrm{Ord}_1(5)$). Furthermore,
\begin{equation}
\label{eq:hn-elementary-bounds}
2\,C_{n-2} \;\le\; \#\HN(n) \quad (n \ge 3)
\qquad\textrm{and}\qquad
\#\HN(n) + C_{n-8} \;\le\; C_{n-1} \quad (n \ge 8),
\end{equation}
which implies $\#\HN(n) + C_{n-9} \le C_{n-1}$ and $\frac{1}{2} < \frac{\#\HN(n)}{C_{n-1}} \le 1 - \frac{C_{n-8}}{C_{n-1}}$ for every $n \ge 8$.
\end{enumerate}
\end{proposition}

\begin{proof}
Part~(i) follows directly from Theorem~\ref{thm:order-hierarchy}(i)--(ii): $\ord(M) = \mathrm{ord}_{\HN}(\sem{M}_{\HN})$ partitions the duplicate-free tree list $\mathcal{L}_n$ of length $C_{n-1}$ into the four disjoint fibers of $Q_{\mathrm{ord}=k}$ ($k = 0, 1, 2, 3$), with $Q_{\mathrm{ord}=0} = \{\infty_{\HN}\}$ and $\mathcal{L}_n \cap \mathrm{Ord}_3 = \{R_{n-1}\}$.

For~(ii), consider first $n \ge 2$: for every closed tree $R \in \mathcal{L}_{n-1}$ ($C_{n-2}$ trees), the Spine-$1$ normal form $S R$ has $\sem{S R}_{\HN} = q_0 \odot_{\HN} \sem{R}_{\HN} \in \{q_1, \dots, q_5, q_{10}, q_{14}, q_{18}, q_{27}\} \subseteq Q_{\mathrm{ord}=2} \sqcup Q_{\mathrm{ord}=3}$ across all $35$ values of $\sem{R}_{\HN} \in Q_{\HN}$. Injectivity of $R \mapsto S R$ yields $C_{n-2} \le \#\mathrm{Ord}_2(n) + \#\mathrm{Ord}_3(n) = \#\mathrm{Ord}_2(n) + 1$. Second, for $n \ge 5$, fix the $4$-leaf generator $L_1 \equiv S(S S S)$, which satisfies $\sem{L_1}_{\HN} = q_7$. Direct inspection across all $35$ states $r \in Q_{\HN}$ confirms $q_7 \odot_{\HN} r \in Q_{\mathrm{ord}=1}$ identically; hence $R \mapsto L_1 R$ injects $\mathcal{L}_{n-4}$ ($C_{n-5}$ trees) into $\mathrm{Ord}_1$ of size $n$. Third, for $n \ge 8$, fix the $7$-leaf generator $L_0 \equiv S S S (S(S S S))$, which satisfies $\sem{L_0}_{\HN} = q_{23}$. Inspection across all $35$ states $r \in Q_{\HN}$ confirms $q_{23} \odot_{\HN} r = \infty_{\HN}$ identically; hence $R \mapsto L_0 R$ injects $\mathcal{L}_{n-7}$ ($C_{n-8}$ trees) into $\mathrm{Ord}_0 = \overline{\HN}$ of size $n$, proving both $C_{n-8} \le \#\mathrm{Ord}_0(n)$ and $\#\HN(n) + C_{n-8} = C_{n-1} - \#\mathrm{Ord}_0(n) + C_{n-8} \le C_{n-1}$. Finally, for $n \ge 3$, every tree $U \in \mathcal{L}_{n-1}$ is an application $U \equiv A_1 A_2$, and the disjoint injections $U \mapsto S U \in \HNF_1$ and $(A_1, A_2) \mapsto S A_1 A_2 \in \HNF_2$ yield $\#\HN(n) \ge C_{n-2} + C_{n-2} = 2 C_{n-2}$.
\end{proof}

For any finite $S$-algebra $\mathcal{A} = (Q, q_0, \odot)$, let $L_q(z) = \sum_{n=1}^{\infty} a_q(n)\,z^n$ denote the ordinary generating function in which $z$ marks $S$-leaves and $a_q(n) = \#\{ M \in \TS : |M| = n,\; \sem{M}_{\mathcal{A}} = q \}$. Root decomposition translates into the positive quadratic polynomial system
\begin{equation}
\label{eq:ogf-system}
L_q(z) \;=\; \delta_{q, q_0}\,z \;+\; \sum_{\substack{(r, s) \in Q^2 \\ r \odot s = q}} L_r(z)\,L_s(z)
\qquad (q \in Q),
\end{equation}
whose total sum satisfies $\sum_{q \in Q} L_q(z) = T(z) = \sum_{n=1}^{\infty} C_{n-1} z^n = \frac{1 - \sqrt{1 - 4z}}{2}$, singular at the Catalan radius $z = \frac{1}{4}$ where $[z^n] T(z) = C_{n-1} \sim \frac{4^{n-1}}{\sqrt{\pi}\,n^{3/2}}$.

\begin{theorem}[Certified rational Catalan order densities, subcritical decay, and six-tier growth spectrum]
\label{thm:asymptotics}
Let $\#\SN(n)$, $\#\Prod(n)$, $\#\HN(n)$, and $\#\mathrm{Ord}_k(n)$ ($k \in \{0, 1, 2, 3\}$) denote the exact leaf-size-$n$ cardinalities in $\TS$.
\begin{enumerate}[label=(\roman*),leftmargin=2.2em]
\item \emph{(Certified rational stationary algebra of $\mathcal{A}_{\HN}$ and exact Catalan order densities)} At the Catalan singularity $z = \frac{1}{4}$ ($u = \sqrt{1 - 4z} = 0$), the $35$-variable system~\eqref{eq:ogf-system} for $\mathcal{A}_{\HN}$ admits an exact dyadic-ternary rational fixed point $a_q = L_q(\frac{1}{4}) \in 2^{-22} 3^{-3}\mathbb{Z}$ ($2^{22}\cdot 3^3 = 113{,}246{,}208$, with $\sum_{q \in Q_{\HN}} a_q = \frac{1}{2}$ and $a_{\infty_{\HN}} = \frac{3{,}366{,}067}{113{,}246{,}208}$), and the singular Puiseux expansion $L_q(z) = a_q - \frac{1}{2}\pi_q \sqrt{1 - 4z} + O(1 - 4z)$ is governed by the unique normalized rational eigenvector $\pi \in 2^{-18} 3^{-3}\mathbb{Z}^{35}$ ($2^{18}\cdot 3^3 = 7{,}077{,}888$) satisfying $\sum_{q \in Q_{\HN}} \pi_q = 1$ and
\[
\pi_q \;=\; \sum_{\substack{(i, j) \in Q_{\HN}^2 \\ i \odot_{\HN} j = q}} (a_i \pi_j + \pi_i a_j)
\qquad (q \in Q_{\HN}).
\]
Summing $\pi_q$ over $Q_{\mathrm{ord}=k}$ yields the exact limiting Catalan densities $c_k := \lim_{n \to \infty} \frac{\#\mathrm{Ord}_k(n)}{C_{n-1}}$:
\begin{equation}
\label{eq:exact-order-densities}
c_0 \;=\; \frac{2{,}648{,}215}{7{,}077{,}888}, \qquad
c_1 \;=\; \frac{1{,}864{,}553}{7{,}077{,}888}, \qquad
c_2 \;=\; \frac{835}{2{,}304} \;=\; \frac{2{,}565{,}120}{7{,}077{,}888}, \qquad
c_3 \;=\; 0,
\end{equation}
satisfying $c_0 + c_1 + c_2 + c_3 = 1$ ($c_0 \approx 0.3741532785$, $c_1 \approx 0.2634335271$, $c_2 \approx 0.3624131944$) and total solvable density
\begin{equation}
\label{eq:chn-exact}
c_{\HN} \;:=\; \lim_{n \to \infty} \frac{\#\HN(n)}{C_{n-1}}
\;=\; c_1 + c_2 + c_3
\;=\; \frac{4{,}429{,}673}{7{,}077{,}888}
\;=\; \frac{4{,}429{,}673}{2^{18} \cdot 3^3}
\;\approx\; 0.6258467215.
\end{equation}
\item \emph{(Algebraic singularities, constructive decay of $\SN$ and $\Prod$, and six-tier growth spectrum)} Because $\infty_{\SN}$ in $\mathcal{A}_{\SN}$ and $\infty_P$ in $\mathcal{A}_P$ are two-sided absorbing sinks coupled into top $4$-state strongly connected components ($\{q_6, q_{25}, q_{29}, q_{30}\}$ with pivot $q_{30}$ in $\mathcal{A}_{\SN}$; $\{q_{18}, q_{24}, q_{29}, q_{33}\}$ with pivot $q_{29}$ in $\mathcal{A}_P$), the generating functions $\sum_{q \in F_{\SN}} L_q(z)$ and $\sum_{q \in F_P} L_q(z)$ have dominant square-root singularities $r_{\SN} < \frac{1}{3}$ and $r_P < r_{\SN}$ strictly exceeding $\frac{1}{4}$, governed by irreducible degree-$35$ minimal polynomials $P_{\SN}(z), P_P(z) \in \mathbb{Z}[z]$:
\begin{align*}
P_{\SN}(z) &\;=\; 8 z^{35} + 28 z^{34} - 76 z^{33} - 384 z^{32} + \cdots - 5550 z^{10} + 3898 z^9 + \cdots - 73 z^3 - 3 z^2 + 7 z - 1, \\
P_P(z) &\;=\; 12 z^{35} + 48 z^{34} - 88 z^{33} - 584 z^{32} - \cdots - 5546 z^{10} + 3898 z^9 + \cdots - 73 z^3 - 3 z^2 + 7 z - 1,
\end{align*}
with exponential growth constants $\gamma_{\SN} = r_{\SN}^{-1} \approx 3.3699195015 < \gamma_P = r_P^{-1} \approx 3.3699837893 < 4$. Constructively over $\mathbb{N}$ (at rational test point $z_0 = \frac{37}{125} = 0.296 < r_P < r_{\SN}$, where $\frac{1}{4 z_0} = \frac{125}{148}$ and certified majorants yield base $\frac{250}{259} < 1$), for every $n \ge 1$:
\begin{equation}
\label{eq:sn-prod-exp-decay}
\frac{\#\SN(n)}{C_{n-1}} \;\le\; 51.39438 \cdot \Bigl(\frac{250}{259}\Bigr)^n,
\qquad
\frac{\#\Prod(n)}{C_{n-1}} \;\le\; 51.42860 \cdot \Bigl(\frac{250}{259}\Bigr)^n.
\end{equation}
Combining Parts~(i)--(ii) with Theorem~\ref{thm:goe-characterization-and-mahler}(iii)--(iv) and Proposition~\ref{prop:order-catalan-partition-bounds}(i) yields six strictly separated exponential growth rates $\gamma_{\mathcal{K}} := \lim_{n \to \infty} \bigl(\#\mathcal{K}(n)\bigr)^{1/n}$ across $\CLS$:
\begin{equation}
\label{eq:six-tier-growth}
1 \;=\; \gamma_{\mathrm{Ord}_3}
\;<\; \gamma_{\mathrm{Iso}} \approx 2.907622
\;<\; \gamma_{\NF} = 3
\;<\; \gamma_{\SN} \approx 3.369920
\;<\; \gamma_P \approx 3.369984
\;<\; \gamma_{\mathrm{Ord}_0} = \gamma_{\mathrm{Ord}_1} = \gamma_{\mathrm{Ord}_2} = 4.
\end{equation}
\end{enumerate}
\end{theorem}

\begin{proof}
For~(i), in $\mathcal{A}_{\HN}$, because $\infty_{\HN}$ is a left sink ($\infty_{\HN} \odot_{\HN} y = \infty_{\HN}$) rather than a two-sided sink ($q_0 \odot_{\HN} \infty_{\HN} = q_{10} \in F_{\HN}$), substituting $\sum_{s \in Q_{\HN}} L_s(z) = T(z) = \frac{1 - u}{2}$ ($u = \sqrt{1 - 4z}$, $z = \frac{1 - u^2}{4}$) eliminates $L_{\infty_{\HN}}(z)$ and expresses the $34$-variable system over $F_{\HN}$ in a lower-triangular block hierarchy whose diagonal equations are at most quadratic with rational square roots at $u = 0$: specifically, $a_{q_0} = \frac{1}{4}$, $a_{q_1} = \frac{1}{16}$, $\dots$, down to exact integer numerators $N_q = 113{,}246{,}208\,a_q \in \mathbb{N}$ satisfying $2\sum_{q \in Q_{\HN}} N_q = 113{,}246{,}208$ and $113{,}246{,}208\,N_q = 113{,}246{,}208 \cdot 28{,}311{,}552\,\delta_{q, q_0} + \sum_{i \odot_{\HN} j = q} N_i N_j$ across all $35$ states $q \in Q_{\HN}$. Differentiating~\eqref{eq:ogf-system} with respect to $u$ at $u = 0$ (where $\frac{dz}{du}\big|_{u=0} = 0$ and $\sum_{q \in Q_{\HN}} L_q(z) = \frac{1}{2} - \frac{1}{2}u$) shows that $\pi_q := -2\frac{d L_q}{du}\big|_{u=0}$ satisfies the linear fixed-point system $\pi = J(a)\,\pi$ with normalization $\sum_{q \in Q_{\HN}} \pi_q = 1$; because the sub-Jacobian of $J(a)$ restricted to $F_{\HN}$ has spectral radius strictly below~$1$ (a positive fraction of products leak into $\infty_{\HN}$), $(I - J(a)|_{F_{\HN}})$ is invertible and determines $\pi$ uniquely. Clearing denominators via $\Pi_q = 7{,}077{,}888\,\pi_q \in \mathbb{N}$ verifies $113{,}246{,}208\,\Pi_q = \sum_{i \odot_{\HN} j = q} (N_i \Pi_j + \Pi_i N_j)$ identically across all $35$ states (in particular, $\Pi_q = 0$ on all $6$ states of $Q_{\mathrm{ord}=3}$, whose generating sum $\sum_{q \in Q_{\mathrm{ord}=3}} L_q(z) = \frac{z}{1 - z}$ is analytic at $z = \frac{1}{4}$) and confirms the exact partition sums $\sum_{Q_{\mathrm{ord}=0}} \Pi_q = 2{,}648{,}215$, $\sum_{Q_{\mathrm{ord}=1}} \Pi_q = 1{,}864{,}553$, $\sum_{Q_{\mathrm{ord}=2}} \Pi_q = 2{,}565{,}120$ ($2{,}565{,}120 \cdot 2{,}304 = 835 \cdot 7{,}077{,}888$), and $\sum_{Q_{\mathrm{ord}=3}} \Pi_q = 0$, summing to $7{,}077{,}888$. Because each $L_q(z)$ is an algebraic power series with non-negative coefficients dominated coefficientwise by $T(z)$ (hence analytic in a $\Delta$-domain around $z = \frac{1}{4}$ with no other singularities on $|z| = \frac{1}{4}$), Flajolet and Sedgewick's Transfer Theorem~\cite[Theorem~VI.1 and Section~VII.6]{FlajoletSedgewick2009} yields $[z^n] \sum_{q \in Q_{\mathrm{ord}=k}} L_q(z) = c_k\,C_{n-1}\bigl(1 + O(n^{-1})\bigr)$ for $k \in \{0, 1, 2\}$, proving~\eqref{eq:exact-order-densities} and~\eqref{eq:chn-exact}.

For~(ii), eliminating the $34$ lower transient and acyclic states of $F_{\SN}$ (respectively $F_P$) in~\eqref{eq:ogf-system} reduces the top $4$-state strongly connected component to a single quadratic equation $A_2(z)\,Y^2 - B_1(z)\,Y + C_0(z) = 0$ ($\deg A_2 = 30, \deg B_1 = 34, \deg C_0 = 38$) for the pivot generating function $Y = L_{q_{30}}(z)$ (respectively $L_{q_{29}}(z)$), whose discriminant $B_1(z)^2 - 4 A_2(z) C_0(z)$ factors over $\mathbb{Z}[z]$ with minimal positive root $r_{\SN}$ of $P_{\SN}(z)$ (respectively $r_P$ of $P_P(z)$, sharing identical coefficients up to $z^9$ because $\Prod$ and $\SN$ agree for all $n \le 8$ by Theorem~\ref{thm:small-term-minimality}(iii)). By the Drmota--Lalley--Woods Theorem~\cite{Drmota1997,Woods1998}, \cite[Theorem~VII.6]{FlajoletSedgewick2009}, $[z^n] \sum_{q \in F_{\SN}} L_q(z) \sim \alpha_{\SN}\,\gamma_{\SN}^n n^{-3/2}$ and $[z^n] \sum_{q \in F_P} L_q(z) \sim \alpha_P\,\gamma_P^n n^{-3/2}$. To prove~\eqref{eq:sn-prod-exp-decay} constructively over $\mathbb{N}$ without complex analysis, choose positive rational upper-envelope vectors $\mathbf{v}_{\SN} \in \mathbb{Q}_{>0}^{38}$ and $\mathbf{v}_P \in \mathbb{Q}_{>0}^{38}$ satisfying quadratic sub-invariance $\delta_{q, q_0}\,z_0 + \sum_{r \odot s = q} v_r v_s \le v_q$ at $z_0 = \frac{37}{125}$. Defining truncated partial sums $S_{q, N}(z_0) = \sum_{n=1}^N a_q(n)\,z_0^n$, induction on $N \ge 1$ shows $S_{q, N+1}(z_0) \le \delta_{q, q_0}\,z_0 + \sum_{r \odot s = q} S_{r, N}(z_0)\,S_{s, N}(z_0) \le v_q$, whence $a_q(n) \le v_q\,z_0^{-n} = v_q (\frac{125}{37})^n$ for every $n \ge 1$. Pairing this upper bound with the inductive Catalan lower bound $C_{n-1} \ge c_0 (\frac{7}{2})^n$ (since $\frac{C_n}{C_{n-1}} = \frac{4n-2}{n+1} \ge \frac{7}{2}$ for all $n \ge 11$, with $1 \le n \le 11$ checked directly) and simplifying
\[
\frac{(125 / 37)^n}{(7 / 2)^n} \;=\; \Bigl(\frac{250}{259}\Bigr)^n
\]
yields the constructive integer inequalities $100{,}000 \cdot 259^n \cdot \#\SN(n) \le 5{,}139{,}438 \cdot 250^n \cdot C_{n-1}$ and $100{,}000 \cdot 259^n \cdot \#\Prod(n) \le 5{,}142{,}860 \cdot 250^n \cdot C_{n-1}$ for all $n \ge 1$, completing~\eqref{eq:sn-prod-exp-decay} and~\eqref{eq:six-tier-growth}.
\end{proof}

% ============================================================================
% SECTION 10
% ============================================================================

\section{Constructive Formalization in Lean 4 and Cross-Language Verification}
\label{sec:formalization}

Following the methodology established in large-scale mechanized mathematics and theoretical computer science---such as Gonthier's formalizations of the Four-Color and Odd Order Theorems~\cite{Gonthier2008FourColor,Gonthier2013OddOrder}, the Flyspeck verification of the Kepler conjecture by Hales et al.~\cite{Hales2017Kepler}, and Forster, Kunze, and Roth's mechanized computability theory~\cite{ForsterKunzeRoth2020}---every core theorem, compiler stage, contextual tree transducer, finite $S$-algebra, $\lambda\mathrm{I}$-calculus theorem, and constructive asymptotic inequality in Sections~\ref{sec:prelim}--\ref{sec:asymptotics} has been formalized and kernel-checked in Lean~4 (version~4.33.1; with finite small-term minimality and collision censuses cross-validated by the companion verification suite), unified in the master module \texttt{s\_challenge\_master\_headline.lean} across its transitive dependency graph of 202 self-contained formal modules. In this closing section we document the logical axiom footprint, kernel-reduction engineering, and cross-language verification architecture that connect the mathematical proofs of this monograph to the formal artifacts.

\subsection{Axiom Footprint and Constructive Elimination of Non-Constructive Principles}
\label{subsec:axiom-footprint}

In Lean~4's Calculus of Inductive Constructions, standard library reasoning may invoke three foundational axioms: propositional extensionality (\texttt{propext}), quotient soundness (\texttt{Quot.sound}, used internally by Lean's linear arithmetic decision procedure \texttt{omega} on \texttt{Nat} and \texttt{Int}), and Hilbert's choice principle (\texttt{Classical.choice}). Across the entire formalization---verified programmatically via \texttt{\#print axioms} in \texttt{TRUST\_ZERO\_SPEC\_CHECK.lean} and \texttt{adversarial\_referee\_audit/axiom\_deep\_audit.lean}---there are no occurrences of \texttt{sorry}, \texttt{admit}, user-declared \texttt{axiom}s, \texttt{partial}, \texttt{unsafe}, or \texttt{extern} declarations, and every exported theorem in \texttt{SChallengeMaster} depends solely on a subset of \texttt{[propext, Quot.sound]}, without relying on \texttt{Classical.choice}.

Achieving \texttt{Classical.choice}-free proofs across infinitary rewriting, B\"ohm-tree semantics, and $\lambda\mathrm{I}$-standardization required replacing four classical idioms common in paper rewriting theory by explicit constructive witnesses:
\begin{enumerate}[label=(\arabic*),leftmargin=2.2em]
\item \emph{Eliminating K\"onig's lemma and pigeonhole arguments in termination and non-spiralling:} Classical local-termination proofs~\cite{EndrullisDeVrijerWaldmann2010} infer termination from absence of infinite paths in finitely branching trees via K\"onig's lemma, and classical unbounded-growth proofs invoke infinite pigeonhole principles. We replaced these in Section~\ref{sec:prelim} by the exact $\Delta L = -2$ left-subtree drop on leaf-preserving steps, constructive lexicographic path-ordering (LPO) and weight-pair accessibility for Waldmann's $\mathcal{R}_m$ system and Klop's flat $\{S, K\}$-terms, and computable step thresholds for unbounded leaf and redex growth; in Section~\ref{subsec:sn-algebra} by the rank-$6$ canonical normalizer $\mathrm{app}_{\mathcal{N}}$; in Section~\ref{subsec:hn-q10} by automaton-free head conservation paired with the pure-multiplier potential $\mu_{\mathrm{HN}}$ ($B \equiv 0$, $\max A = 30$); and in Section~\ref{subsec:top-termination} by well-founded recursion (\texttt{Acc.rec}) on $\nu_{\mathrm{HN}}$ and right-spine strata.
\item \emph{Constructive $\bot$-extraction and productive B\"ohm joinability:} On closed $S$-terms, productivity $\mathrm{Prod}(M) \equiv \forall d,\, \neg\bot_d(M)$ is a $\Pi_1^0$ predicate whose classical negation $\neg\mathrm{Prod}(M) \implies \exists d,\, \bot_d(M)$ is Markov's principle. Over the Myhill--Nerode minimal $113$-state joint algebra $\mathrm{PH} \subseteq Q_P \times Q_{\mathrm{HN}}$, rank-$10$ critical-pair stratification constructs an explicit depth $d \in \mathbb{N}$ witnessing $\bot_d(M)$ whenever $\llbracket M \rrbracket_P = \infty_P$, yielding the constructive disjunction $\mathrm{Prod}(M) \vee \exists d,\, \bot_d(M)$, while root-normal-form spine peeling constructs explicit common reducts witnessing B\"ohm-tree faithfulness and decidability of $\convS$ across $\mathrm{Prod} \times \mathcal{T}_S$.
\item \emph{Constructive $\lambda\mathrm{I}$ standardization and $S$-free backward cones:} In Section~\ref{sec:lambda-translation}, Church's conservation theorem ($\mathrm{WN}_\beta \iff \mathrm{SN}_\beta$ on $\lambda\mathrm{I}$)~\cite{Church1941}, Conjecture~40, the hereditary saturation repair $\mathrm{DeepNorm}(M) \iff (M)_\lambda \in \mathrm{SN}_\beta$, and single-constant collapse $\mathrm{DeepNorm}_1$ for Question~34 are proved constructively via de Bruijn parallel $\beta$-reduction and weak-head standardization without classical choice, while Question~32 is resolved via constructive $S$-free anti-redex backward-cone inversion.
\item \emph{Step-indexed deterministic traces and contextual transducers:} In the universal Turing machine, program-free Cyclic Tag System, and pocket tree transducer simulations (Sections~\ref{sec:tm-sim}--\ref{sec:cts-sim}), existential halting witnesses are carried by explicit fuel-indexed evaluators (\texttt{uRun}, \texttt{InvokeTrace}, \texttt{pfIter}, and the two-pass $(25, 23, 136)$ transducer \texttt{pocketPFSelect}), avoiding non-constructive witness extraction from transitive closures.
\end{enumerate}

\subsection{Kernel Reduction Without \texorpdfstring{\texttt{native\_decide}}{native\_decide}}
\label{subsec:kernel-reduction}

Lean~4 provides a tactic \texttt{native\_decide} that compiles decidable propositions to native machine code, adding the compiler and runtime to the trusted computing base. To keep the trusted computing base confined strictly to Lean~4's small logical kernel, every finite algebraic, Myhill--Nerode, and small-term census check reduces directly inside the kernel's definitional equality checker (\texttt{rfl} / \texttt{decide}), making no appeal to \texttt{native\_decide}.

Checking $39^3 = 59{,}319$ $S$-law triples for $\mathcal{A}_{\mathrm{SN}}$ and $\mathcal{A}_P$, $42{,}875$ triples for $\mathcal{A}_{\mathrm{HN}}$, the $35$-override two-dimensional collapse and rank-$10$ critical-pair tables over the $113$-state algebra $\mathrm{PH}$, the $17$-context Myhill--Nerode separation of $\mathrm{PH}$, and the $(25, 23, 136)$ pocket transducer equivalence inside Lean's kernel without stack overflow or quadratic pattern-compilation overhead uses three standard reflection techniques~\cite{Gonthier2008FourColor}:
\begin{enumerate}[label=(\arabic*),leftmargin=2.2em]
\item \emph{Balanced binary tree quantifiers and row factoring:} Universal quantification over finite inductive state types ($3$ to $39$ constructors) and bounded natural intervals is expressed via balanced Boolean conjunction trees (\texttt{checkAll1}, \texttt{checkAll2}, \texttt{checkNat1}, \texttt{checkNat2}) accompanied by projection soundness lemmas, and three-variable checks $(x, y, z) \in Q^3$ are factored into $|Q|$ separate two-variable row theorems each proved by \texttt{rfl}.
\item \emph{Depth-$6$ and depth-$7$ binary Patricia lookup trees (\texttt{PTree64}, \texttt{PTree}):} Transition matrices, $113$-state witness tables, and rank tables over $\mathrm{PH}$ are stored as balanced binary trees of depth $6$ or $7$, reducing every state-pair or rank lookup in at most $7$ pointer steps rather than $113$ sequential linear comparisons.
\item \emph{Explicit constructor case-splitting on pattern matches:} Lean~4's equation compiler and \texttt{split} tactic occasionally synthesize \texttt{Classical.em} on overlapping multi-constructor wildcard branches; replacing wildcard \texttt{split} steps with explicit constructor \texttt{cases} decompositions eliminates \texttt{Classical.choice} throughout the pocket transducer and algebraic parsers.
\end{enumerate}
Similarly, for the root-restarted zipper controller \texttt{invoke} (\texttt{formal\_controller/finite\_tree\_walker.lean}), whose nested register type \texttt{Control} has cardinality $4{,}376{,}490{,}905$, finiteness is proved symbolically via product-list membership (\texttt{mem\_allControls}: $\forall c,\, c \in \mathrm{allControls}$) without materializing a four-billion-element list, while \texttt{pocketPFTransducer} achieves an explicit $(25, 23, 136)$ finite-state realization running in at most $2|M| - 1 + \mathrm{depth}(M)$ transitions.

\subsection{Syntax Bijections and Cross-Language Parity}
\label{subsec:cross-language}

Because the formal modules were developed in modular namespaces (\texttt{FormalAlgebra}, \texttt{FormalTopTermination}, \texttt{IntrinsicTheorems}, \texttt{Q5SNStratifiedGame}, \texttt{FormalQ5Final}, \texttt{SimplifiedQ31Q7}, \texttt{SimplifiedQ10Q5}, \texttt{FormalQ30Q31Deep}, \texttt{FormalOrderAndIsolated}, \texttt{FormalNoSpiral}, \texttt{Conjecture17Obstructions}, \texttt{FormalLambdaConjectures}, \texttt{FormalQ32}, \texttt{FormalWolframFrontier}, \texttt{ZeroNavigatorFixedOrder}, \texttt{ZeroNavigatorMultiway}, \texttt{ZeroNavigatorCellular}, \texttt{FormalWolfram12}, \texttt{FormalMinimality}, and \texttt{FormalAsymptotics}), the master file \texttt{s\_challenge\_master\_headline.lean} defines explicit recursive syntax translations (\texttt{toAlgebraTerm}, \texttt{ofAlgebraTerm}, \dots) and proves bidirectional round-trip identities together with closedness preservation and one-step rewrite equivalence ($\mathrm{Step}\;M\;N \iff \mathrm{Step}_{\text{mod}}(\tau(M), \tau(N))$) across all pure-$S$ syntax definitions. Thus no discrepancy in associativity, variable indexing, or redex semantics can arise between modules.

In parallel with the Lean~4 development, all compilers, tree-walking selectors, pocket tree transducers, halting automata, finite $S$-algebras (sizes $39, 35, 39, 36, 12, 3, 4$, and $113$), $\lambda\mathrm{I}$ divergence witnesses ($W_9, N_{10}, M_{14}, M_{15}$), Rule~110 normal-form iterators, all-step $\to_{\mathrm{lo}}$ jump-ahead formulas, Garden-of-Eden normal-form recognizers, and Catalan order-density constants are implemented in the standalone Wolfram Language package \texttt{wolfram\_submission/PureSUniversality.wl} (accompanied by \texttt{wolfram\_submission/VerifyPureSUniversality.wls}) and an independent Python verification suite (\texttt{verify\_all\_theorems.py} and \texttt{wolfram\_submission/cross\_validate\_wl.py}). Automated comparison confirms exact agreement across all shared transition matrices, pure-multiplier table $A$ ($\max A = 30$), compiled Turing machine leaf counts ($1{,}996{,}921$ and $7{,}477{,}337$ leaves), $912$-contraction execution of the $2$-state Busy Beaver, pocket transducer trajectories from $D_{\mathrm{seed}}$, and $130$-digit dynamic programming verification of the Catalan phase transitions and rational order densities through $n = 150$.

% ============================================================================
% APPENDIX A
% ============================================================================

\appendix

\section{Concordance of Mathematical Results and Formal Artifacts}
\label{app:concordance}

Table~\ref{tab:concordance} maps the principal theorems, propositions, and lemmas of Sections~\ref{sec:prelim}--\ref{sec:asymptotics} to their corresponding Lean~4 source files, kernel-verified declaration identifiers, exact foundational axiom footprints reported by \texttt{\#print axioms} (where $\emptyset$ denotes no axioms, $\mathsf{P}$ denotes \texttt{[propext]}, and $\mathsf{PQ}$ denotes \texttt{[propext, Quot.sound]}), and exported Wolfram Language symbols in \texttt{wolfram\_submission/PureSUniversality.wl}.

\begin{table}[htbp]
\centering
\caption{Concordance between monograph results (Sections~\ref{sec:prelim}--\ref{sec:asymptotics}), Lean~4 declarations, axiom footprints, and Wolfram Language symbols.}
\label{tab:concordance}
\scriptsize
\setlength{\tabcolsep}{3.0pt}
\renewcommand{\arraystretch}{1.12}
\begin{tabular}{@{}lllcl@{}}
\toprule
Monograph Result & Lean~4 Module Path & Lean~4 Declaration Identifier & Axioms & Wolfram Language Symbol \\
\midrule
Prop.~\ref{prop:cr-conservation}: CR \& $\WN = \SN$ & \texttt{formal\_algebra/s\_structural\_laws.lean} & \texttt{FormalAlgebra.church\_rosser}, \texttt{wn\_implies\_sn} & $\mathsf{P}$ & \texttt{PureSContractAt} \\
Thm.~\ref{thm:leaf-preserving}: $\Delta L = -2$ \& $\lfloor(n{-}2)^2/4\rfloor$ & \texttt{formal\_order\_and\_isolated/leaf\_preserving\_and\_isolated.lean} & \texttt{FormalOrderAndIsolated.leftSubtreeSum\_leafPresStep}, \texttt{leafPres\_sn} & $\mathsf{PQ}$ & \texttt{PureSLeftSubtreeSum} \\
Prop.~\ref{prop:acyclicity-order},~\ref{prop:remark16-klop18}: No-spiral \& Klop~18 & \texttt{formal\_no\_spiral/waldmann\_no\_spiral.lean}, \texttt{formal\_no\_spiral/bergstra\_and\_klop.lean} & \texttt{FormalNoSpiral.waldmann\_no\_spiral}, \texttt{klop\_theorem\_18\_flat\_sk\_sn} & $\mathsf{P}$ & --- \\
Prop.~\ref{prop:open-spiral-divergence-threshold}: $8$-leaf open $\overline{\SN}$ & \texttt{formal\_conjecture17\_obstructions/conjecture17\_obstructions.lean} & \texttt{Conjecture17Obstructions.open\_spiral\_not\_sn} & $\emptyset$ & --- \\
Thm.~\ref{thm:conj17-dp-weighted}: Self-right-tail Conj.~17 & \texttt{formal\_no\_spiral/conjecture17\_open\_spirals.lean} & \texttt{FormalNoSpiral.conjecture17\_no\_open\_spiral\_unless\_self\_right\_tail} & $\mathsf{P}$ & --- \\
Thm.~\ref{thm:conj17-unconditional-and-coupling}: Short-loop \& near-miss & \texttt{formal\_conjecture17\_obstructions/conjecture17\_obstructions.lean} & \texttt{Conjecture17Obstructions.no\_two\_step\_open\_spiral}, \texttt{waldmann\_open\_parametric\_near\_miss} & $\mathsf{PQ}$ & --- \\
Thm.~\ref{thm:conj17-minmaxrs-monotonicity}: $\mathrm{minRS}_x/\mathrm{minLS}_x$ \& $\mathcal{DP}$ & \texttt{formal\_conjecture17\_deep/minmaxrs\_deep\_obstructions.lean} & \texttt{FormalNoSpiral.step\_minMaxRS\_le}, \texttt{conjecture17\_iff\_no\_dpChain\_spiral} & $\mathsf{P}$ & --- \\
Prop.~\ref{prop:u311-sim}: $M(w) \leftrightarrow U_{3,11}$ sim. & \texttt{formal\_utm/neary\_u311\_sim.lean}, \texttt{formal\_utm/neary\_encoder.lean} & \texttt{source\_halts\_iff\_u\_stops}, \texttt{encodeSourceInstance\_poly\_bound} & $\mathsf{PQ}$ & \texttt{EncodeNearyU311} \\
Lem.~\ref{lem:microtick-bound}: Zipper walker bound & \texttt{formal\_controller/bare\_output\_and\_invoke\_trace.lean} & \texttt{exists\_microtick\_le\_node\_bound} & $\mathsf{PQ}$ & \texttt{RootRestartedSelectAddress} \\
Prop.~\ref{prop:halt-automaton}: $11$-state halt DTA $\mathcal{H}_{11}$ & \texttt{formal\_controller/halt\_tree\_automaton.lean} & \texttt{haltAutomaton\_accepts\_eq\_rootHalt} & $\mathsf{P}$ & \texttt{PureSHaltDetectorQ} \\
Thm.~\ref{thm:tm-universality}: Root-restarted UTM & \texttt{formal\_controller/constructive\_endpoint.lean}, \texttt{formal\_controller/bare\_output\_and\_invoke\_trace.lean} & \texttt{constructive\_fixed\_u\_to\_s\_halting}, \texttt{decodeBareSTermOutput\_correct} & $\mathsf{PQ}$ & \texttt{CompileTuringMachineToPureS} \\
Prop.~\ref{prop:seed-genesis}: Seed genesis $D_{\mathrm{seed}} \to_S^* t_0$ & \texttt{formal\_program\_free/program\_free\_cts.lean} & \texttt{FormalProgramFree.seed\_reduces\_to\_encodeProgramFreeCTS} & $\mathsf{PQ}$ & \texttt{EncodeProgramFreeCTS} \\
Thm.~\ref{thm:cts-universality}: Program-free CTS univ. & \texttt{formal\_program\_free/program\_free\_cts.lean} & \texttt{FormalProgramFree.program\_free\_cts\_halting\_iff\_automaton} & $\mathsf{PQ}$ & \texttt{ProgramFreePureSSelect} \\
Thm.~\ref{thm:pocket-transducer-universality}: $(25,23,136)$ $\mathcal{T}_{\mathrm{pocket}}$ & \texttt{formal\_wolfram\_frontier/pocket\_pf\_transducer.lean} & \texttt{FormalWolframFrontier.pocketPFSelect\_eq\_pfSelect}, \texttt{pocket\_transducer\_cts\_universality} & $\mathsf{PQ}$ & \texttt{ProgramFreePureSSelect} \\
Thm.~\ref{thm:leaf-split-univ}: Unary leaf-split univ. & \texttt{formal\_algebra/trivial\_rule\_separation.lean} & \texttt{FormalAlgebra.TrivialRuleSeparation.fixed\_universal\_trivial\_rule\_separation} & $\mathsf{PQ}$ & --- \\
Prop.~\ref{prop:multiway-shortcuts}: $7$/$5$-step shortcuts & \texttt{zero\_navigator\_multiway/multiway\_and\_q30\_theorems.lean} & \texttt{ZeroNavigatorMultiway.protected\_trie\_nonhalting\_shortcut\_7steps} & $\mathsf{P}$ & --- \\
Thm.~\ref{thm:no-church-bool}: No pure-$S$ Church gate & \texttt{zero\_navigator\_multiway/multiway\_and\_q30\_theorems.lean} & \texttt{ZeroNavigatorMultiway.no\_church\_boolean\_pair\_in\_pure\_s} & $\mathsf{P}$ & --- \\
Thm.~\ref{thm:waldmann-sn}: $39$-state $\mathcal{A}_{\SN}$ & \texttt{formal\_algebra/sn\_hn\_algebras.lean} & \texttt{FormalAlgebra.waldmann\_theorem\_9} & $\mathsf{P}$ & \texttt{EvalSNState} \\
Lem.~\ref{lem:head-conservation},~\ref{lem:hn-potential}: Head cons.\ \& $B\equiv 0$ & \texttt{formal\_q10\_q5\_simplified/q10\_thm9\_simplified.lean} & \texttt{SimplifiedQ10Q5.hasHNF\_iff\_strongHeadNormalizing} & $\mathsf{PQ}$ & --- \\
Prop.~\ref{prop:hna-minimality}: Pointwise minimal $A$ & \texttt{formal\_q10\_q5\_simplified/q10\_thm9\_simplified.lean} & \texttt{SimplifiedQ10Q5.hnA\_pointwise\_minimal}, \texttt{no\_masked\_rank1\_hnA} & $\mathsf{PQ}$ & --- \\
Thm.~\ref{thm:question-10}: Question~10 \& Footnote~4 & \texttt{formal\_q10\_q5\_simplified/q10\_thm9\_simplified.lean}, \texttt{formal\_algebra/sn\_hn\_algebras.lean} & \texttt{SimplifiedQ10Q5.barendregt\_open\_problem\_10\_simplified}, \texttt{FormalAlgebra.footnote4\_rnf\_iff} & $\mathsf{PQ}$ & \texttt{EvalHNState} \\
Thm.~\ref{thm:waldmann-top-term}: Top-termination & \texttt{formal\_top\_termination/top\_termination.lean} & \texttt{FormalTopTermination.waldmann\_theorem\_29\_top\_termination} & $\mathsf{PQ}$ & \texttt{IsRootNormalFormQ} \\
Lem.~\ref{lem:bot-propagation},~\ref{lem:zachos-bot-absorption}: $\bot$-laws & \texttt{formal\_q10\_q5\_simplified/q10\_thm9\_simplified.lean}, \texttt{formal\_q10\_q5\_simplified/q5\_simplified.lean} & \texttt{SimplifiedQ10Q5.hasBotWithin\_app\_left}, \texttt{hasBot\_zachos\_app} & $\mathsf{PQ}$ & --- \\
Prop.~\ref{prop:ap-soundness-stratification}: $34$-of-$38$ $\SN$ strat. & \texttt{q5\_sn\_stratified\_game/sn\_stratified\_game.lean} & \texttt{Q5SNStratifiedGame.unique\_nonSN\_productive\_root\_redex\_state} & $\mathsf{P}$ & \texttt{EvalProductivityState} \\
Thm.~\ref{thm:q5-main}: Question~5 (rank-$10$ $\mathcal{A}_P$) & \texttt{formal\_q10\_q5\_simplified/q5\_simplified.lean}, \texttt{formal\_q5\_final/q5\_completeness.lean} & \texttt{SimplifiedQ10Q5.barendregt\_open\_question\_5\_simplified} & $\mathsf{PQ}$ & \texttt{IsProductivePureSQ} \\
Prop.~\ref{prop:ph-113-minimal}: $113$-state $\mathrm{PH}$ minimal & \texttt{formal\_q10\_q5\_simplified/q5\_simplified.lean} & \texttt{SimplifiedQ10Q5.ph\_113\_nerode\_minimal}, \texttt{ph\_113\_quotient\_lower\_bound} & $\mathsf{PQ}$ & --- \\
Prop.~\ref{prop:bohm-game-asymmetry}: $8$-leaf $L_{\mathrm{loop}}$ game & \texttt{q5\_sn\_stratified\_game/sn\_stratified\_game.lean} & \texttt{Q5SNStratifiedGame.bohm\_game\_asymmetry\_size8} & $\mathsf{P}$ & --- \\
Thm.~\ref{thm:question-7}: Question~7 ($\mathcal{H}/\mathcal{U}$ \& mod-$10$) & \texttt{SIMPLIFIED\_Q31\_Q7\_BOOK\_PROOFS.lean}, \texttt{formal\_algebra/question7\_flipflop.lean} & \texttt{SimplifiedQ31Q7.question\_7\_atomic\_leaf\_splits}, \texttt{FormalAlgebra.one\_step\_rescue\_classification} & $\mathsf{PQ}$ & \texttt{Question7AtomicSplitTerm}, \texttt{Question7FlipFlopTerm} \\
Thm.~\ref{thm:question-31}: Question~31 ($\{A A A S^k\}$ \& $\mathcal{A}_{36}$) & \texttt{SIMPLIFIED\_Q31\_Q7\_BOOK\_PROOFS.lean}, \texttt{intrinsic\_research/intrinsic\_theorems.lean} & \texttt{SimplifiedQ31Q7.question\_31\_infinite\_family}, \texttt{IntrinsicTheorems.aaa\_not\_conv\_aaas} & $\mathsf{PQ}$ & \texttt{EvalQuestion31State} \\
Thm.~\ref{thm:productive-bohm-faithful}: $\Prod \times \TS$ B\"ohm faithful & \texttt{formal\_q30\_q31\_deep/q31\_productive\_bohm\_faithfulness.lean} & \texttt{FormalQ30Q31Deep.productive\_oneSided\_bohm\_iff\_convS} & $\mathsf{PQ}$ & --- \\
Prop.~\ref{prop:q30-decidability}: Question~30 ($\Prod \times \TS$ \& finiteness) & \texttt{formal\_q30\_q31\_deep/q31\_productive\_bohm\_faithfulness.lean}, \texttt{formal\_q30\_q31\_deep/q30\_finiteness\_and\_sn\_decidability.lean} & \texttt{FormalQ30Q31Deep.decideConvProd}, \texttt{closed\_normalizing\_equivClass\_finite} & $\mathsf{PQ}$ & --- \\
Thm.~\ref{thm:order-hierarchy}: Orders $\{0,1,2,3\}$ ($35,12,3$) & \texttt{formal\_order\_and\_isolated/order\_classification.lean} & \texttt{FormalOrderAndIsolated.visser\_van\_oostrom\_exact\_order\_exhaustion} & $\mathsf{PQ}$ & \texttt{PureSFunctionalOrder} \\
Lem.~\ref{lem:lambdaI-conservation}: $\lambda\mathrm{I}$ $\WN_\beta \iff \SN_\beta$ & \texttt{formal\_lambda\_conjectures/lambda\_standardization\_conservation.lean} & \texttt{FormalLambdaConjectures.lambdaI\_wn\_iff\_sn} & $\mathsf{PQ}$ & --- \\
Thm.~\ref{thm:conjecture-40}: Conjecture~40 ($\ord_{\mathrm{CL}} = \ord_\lambda$) & \texttt{formal\_lambda\_conjectures/conjecture\_40\_order.lean} & \texttt{FormalLambdaConjectures.conjecture\_40\_full\_cl}, \texttt{conjecture\_40\_pure\_s} & $\mathsf{PQ}$ & --- \\
Prop.~\ref{prop:minimal-lambda-divergers}: $W_9, N_{10}$ $\lambda\beta$-divergers & \texttt{formal\_lambda\_conjectures/conjectures\_42\_43\_and\_repair.lean} & \texttt{FormalLambdaConjectures.w9\_toLam\_no\_betaNF}, \texttt{n10\_toLam\_no\_betaNF} & $\mathsf{PQ}$ & \texttt{MinimalLambdaDivergingPureSTerms} \\
Thm.~\ref{thm:conjectures-42-43-refuted}: Conj.~42/43 refuted \& $\mathrm{DeepNorm}$ & \texttt{formal\_lambda\_conjectures/conjectures\_42\_43\_and\_repair.lean} & \texttt{FormalLambdaConjectures.conjecture\_42\_and\_43\_refuted\_14\_leaves}, \texttt{repaired\_conjecture\_42} & $\mathsf{PQ}$ & \texttt{Conjecture42And43Counterexamples} \\
Thm.~\ref{thm:question-34-structure}: Question~34 ($\mathrm{DeepNorm}_1$ \& $\infty_{34}$) & \texttt{formal\_lambda\_conjectures/question\_34\_hereditary\_and\_shattering.lean} & \texttt{FormalLambdaConjectures.question\_34\_iff\_single\_constant\_saturation} & $\mathsf{PQ}$ & --- \\
Prop.~\ref{prop:q32-rigidity-waldmann-pair}: $S$-free rigidity \& $(8,11)$ pair & \texttt{formal\_q32\_equational/lambda\_i\_rigidity.lean} & \texttt{FormalQ32.sFree\_rigidity}, \texttt{theoremA\_waldmann\_8\_11\_separation} & $\mathsf{PQ}$ & --- \\
Thm.~\ref{thm:question-32}: Question~32 refuted & \texttt{formal\_q32\_equational/question\_32\_negative\_resolution.lean} & \texttt{FormalQ32.no\_finite\_equational\_axiomatization\_of\_lambda\_beta} & $\mathsf{PQ}$ & --- \\
Prop.~\ref{prop:pure-line-and-aaa}: $6$ strategies \& $\mathcal{B}_Z$ & \texttt{zero\_navigator\_fixed\_order/zero\_nav\_fixed\_theorems.lean} & \texttt{ZeroNavigatorFixedOrder.zachos\_unified\_quadratic\_trace} & $\mathsf{PQ}$ & \texttt{EvalBZachosState} \\
Thm.~\ref{thm:case3-ceiling}: $5$-step Case-$3$ ceiling & \texttt{zero\_navigator\_fixed\_order/zero\_nav\_fixed\_theorems.lean} & \texttt{ZeroNavigatorFixedOrder.no\_six\_consecutive\_case3\_li} & $\mathsf{PQ}$ & --- \\
Thm.~\ref{thm:clones-rule110}: Clone barrier \& Rule~110 & \texttt{zero\_navigator\_cellular/cellular\_and\_clones.lean} & \texttt{ZeroNavigatorCellular.no\_rule110\_SN}, \texttt{dnand\_rule110\_correct} & $\mathsf{PQ}$ & \texttt{PureSRule110Term} \\
Thm.~\ref{thm:wolfram-14-dichotomy}: $14$ benchmarks & \texttt{formal\_wolfram12/wolfram\_12\_benchmarks.lean} & \texttt{FormalWolfram12.wolfram\_14\_dichotomy\_master} & $\mathsf{P}$ & \texttt{WolframChallenge12Benchmarks} \\
Prop.~\ref{prop:jump-ahead}: Context \& $s(m)$ jump-ahead & \texttt{formal\_wolfram12/wolfram\_12\_benchmarks.lean} & \texttt{FormalWolfram12.decideContextHalts\_exact}, \texttt{jump\_ahead\_sPrefix\_exact} & $\mathsf{P}$ & \texttt{PureSJumpAheadOutcome} \\
Thm.~\ref{thm:all-step-jump-ahead}: All-step $\to_{\mathrm{lo}}$ cubic laws & \texttt{formal\_wolfram\_frontier/lo\_jump\_ahead\_and\_goe.lean} & \texttt{FormalWolframFrontier.zachos\_all\_step\_closed\_form}, \texttt{wolfram2\_vs\_w2Twin\_shift\_14} & $\mathsf{PQ}$ & \texttt{ZachosExactLeafCount}, \texttt{Wolfram2ExactLeafCount} \\
Prop.~\ref{prop:aaa-grammar-and-subterm-sep}: $A A A$ grammar \& subterms & \texttt{formal\_wolfram\_frontier/lo\_jump\_ahead\_and\_goe.lean} & \texttt{FormalWolframFrontier.aaa\_lo\_and\_multiway\_master} & $\mathsf{PQ}$ & --- \\
Thm.~\ref{thm:goe-characterization-and-mahler}: Garden-of-Eden NFs & \texttt{formal\_wolfram\_frontier/lo\_jump\_ahead\_and\_goe.lean} & \texttt{FormalWolframFrontier.goe\_s2\_iff}, \texttt{decideGardenOfEdenNF\_iff} & $\mathsf{PQ}$ & \texttt{IsGardenOfEdenNormalFormQ} \\
Thm.~\ref{thm:small-term-minimality}: $n \le 9$ minimality & \texttt{formal\_census\_and\_minimality/minimality\_theorems.lean} & \texttt{FormalMinimality.theorem\_A\_universal\_sn\_le\_6}, \texttt{theorem\_F\_confluence\_and\_diamond\_thresholds} & $\mathsf{PQ}$ & \texttt{PureSCatalanPhaseCensus} \\
Prop.~\ref{prop:order-catalan-partition-bounds}, Thm.~\ref{thm:asymptotics}: Catalan laws & \texttt{formal\_order\_and\_isolated/higher\_order\_catalan\_laws.lean}, \texttt{formal\_asymptotics/asymptotic\_laws.lean} & \texttt{FormalOrderAndIsolated.exact\_rational\_order\_catalan\_densities\_certified}, \texttt{FormalAsymptotics.sn\_and\_acceptP\_exponential\_decay\_to\_zero} & $\mathsf{PQ}$ & \texttt{PureSAsymptoticConstants} \\
\bottomrule
\end{tabular}
\end{table}

\begin{thebibliography}{99}



\bibitem[Bar84]{Barendregt1984}
Henk~P. Barendregt,
\emph{The Lambda Calculus: Its Syntax and Semantics},
Revised edition, Studies in Logic and the Foundations of Mathematics, vol.~103,
North-Holland, Amsterdam, 1984.

\bibitem[BBKV76]{BarendregtBergstraKlopVolken1976}
Henk~P. Barendregt, Jan~A. Bergstra, Jan~Willem Klop, and Henri Volken,
\emph{Degrees, Reductions and Representability in the Lambda Calculus},
Preprint No.~22, Department of Mathematics, Utrecht University, 1976.

\bibitem[BDS13]{BarendregtDekkersStatman2013}
Henk~P. Barendregt, Wil Dekkers, and Richard Statman,
\emph{Lambda Calculus with Types},
Perspectives in Logic, Cambridge University Press, Cambridge, 2013,
\texttt{doi:10.1017/CBO9781139032636}.

\bibitem[BEKW17]{BarendregtEndrullisKlopWaldmann2017}
Henk~P. Barendregt, J{\"o}rg Endrullis, Jan~Willem Klop, and Johannes Waldmann,
``Dance of the Starlings,''
in Melvin Fitting and Brian Rayman (eds.), \emph{Raymond Smullyan on Self Reference},
Outstanding Contributions to Logic, vol.~14, Springer, Cham, 2017, pp.~67--111,
\texttt{doi:10.1007/978-3-319-68732-2\_5}.

\bibitem[BGLZ17]{BendkowskiGrygielLescanneZaionc2017}
Maciej Bendkowski, Katarzyna Grygiel, Pierre Lescanne, and Marek Zaionc,
``Combinatorics of $\lambda$-terms: a natural approach,''
\emph{Journal of Logic and Computation} \textbf{27} (2017), no.~8, 2611--2630,
\texttt{doi:10.1093/logcom/exx018}.

\bibitem[BGZ15]{BendkowskiGrygielZaionc2015}
Maciej Bendkowski, Katarzyna Grygiel, and Marek Zaionc,
``Asymptotic properties of combinatory logic,''
in \emph{Theory and Applications of Models of Computation (TAMC 2015)},
Lecture Notes in Computer Science, vol.~9076, Springer, Cham, 2015, pp.~62--72,
\texttt{doi:10.1007/978-3-319-17142-5\_7}.

\bibitem[BGZ17]{BendkowskiGrygielZaionc2017}
Maciej Bendkowski, Katarzyna Grygiel, and Marek Zaionc,
``On the likelihood of normalisation in combinatory logic,''
\emph{Journal of Logic and Computation} \textbf{27} (2017), no.~7, 2251--2269,
\texttt{doi:10.1093/logcom/exx005}.

\bibitem[Ber96]{Berarducci1996}
Alessandro Berarducci,
``Infinite $\lambda$-calculus and non-sensible models,''
in Aldo Ursini and Paolo Aglian{\`o} (eds.), \emph{Logic and Algebra},
Lecture Notes in Pure and Applied Mathematics, vol.~180, Marcel Dekker, New York, 1996, pp.~339--377,
\texttt{doi:10.1201/9780203748671-17}.

\bibitem[BB93]{BerarducciBohm1993}
Alessandro Berarducci and Corrado B{\"o}hm,
``A self-interpreter of lambda calculus having a normal form,''
in \emph{Computer Science Logic (CSL '92)},
Lecture Notes in Computer Science, vol.~702, Springer, Berlin, 1993, pp.~85--99,
\texttt{doi:10.1007/3-540-57182-5\_7}.

\bibitem[BGG11]{BodiniGardyGittenberger2011}
Olivier Bodini, Dani{\`e}le Gardy, and Bernhard Gittenberger,
``Lambda-terms of bounded unary height,''
in \emph{Proceedings of the Eighth Workshop on Analytic Algorithmics and Combinatorics (ANALCO 2011)},
SIAM, 2011, pp.~23--32, \texttt{doi:10.1137/1.9781611973013.3}.

\bibitem[BGGJ13]{BodiniGardyGittenbergerJacquot2013}
Olivier Bodini, Dani{\`e}le Gardy, Bernhard Gittenberger, and Alice Jacquot,
``Enumeration of generalized $BCI$ lambda-terms,''
\emph{The Electronic Journal of Combinatorics} \textbf{20} (2013), no.~4, Paper~P30,
\texttt{doi:10.37236/3051}.

\bibitem[BC08]{BojanczykColcombet2008}
Miko{\l}aj Boja{\'n}czyk and Thomas Colcombet,
``Tree-walking automata do not recognize all regular languages,''
\emph{SIAM Journal on Computing} \textbf{38} (2008), no.~2, 658--701,
\texttt{doi:10.1137/050645427}.

\bibitem[B{\"o}h68]{Bohm1968}
Corrado B{\"o}hm,
``Alcune propriet{\`a} delle forme $\beta$-$\eta$-normali nel $\lambda$-$K$-calcolo,''
Pubblicazioni dell'Istituto per le Applicazioni del Calcolo, no.~696, Rome, 1968, pp.~1--19.

\bibitem[Can78]{Canal1978}
Richard Canal,
``Complexit{\'e} de la r{\'e}duction en logique combinatoire,''
\emph{RAIRO -- Informatique th{\'e}orique} \textbf{12} (1978), no.~4, 339--367,
\texttt{doi:10.1051/ita/1978120403391}.

\bibitem[CRZ11]{CheilarisRamirezZachos2011}
Panagiotis Cheilaris, John Ramirez, and Stathis Zachos,
``Checking in linear time if an $S$-term normalizes,''
in \emph{8th Panhellenic Logic Symposium}, Ioannina, Greece, July 2011.

\bibitem[Chu41]{Church1941}
Alonzo Church, \emph{The Calculi of Lambda-Conversion}, Annals of Mathematics Studies,
no.~6, Princeton University Press, Princeton, NJ, 1941.

\bibitem[CM64]{CockeMinsky1964Tag}
John Cocke and Marvin~L. Minsky,
``Universality of tag systems with $P = 2$,''
\emph{Journal of the ACM} \textbf{11} (1964), no.~1, 15--20,
\texttt{doi:10.1145/321203.321206}.

\bibitem[Coo04]{Cook2004Rule110}
Matthew Cook,
``Universality in elementary cellular automata,''
\emph{Complex Systems} \textbf{15} (2004), no.~1, 1--40,
\texttt{doi:10.25088/ComplexSystems.15.1.1}.

\bibitem[Cur30]{Curry1930}
Haskell~B. Curry,
``Grundlagen der kombinatorischen Logik,''
\emph{American Journal of Mathematics} \textbf{52} (1930), no.~3--4, 509--536, 789--834,
\texttt{doi:10.2307/2370619}.

\bibitem[CF58]{CurryFeys1958}
Haskell~B. Curry and Robert Feys,
\emph{Combinatory Logic, Volume~I},
Studies in Logic and the Foundations of Mathematics, North-Holland, Amsterdam, 1958.

\bibitem[DGK+13]{DavidGrygielKozikRaffalliTheyssierZaionc2013}
Ren{\'e} David, Katarzyna Grygiel, Jakub Kozik, Christophe Raffalli, Guillaume Theyssier, and Marek Zaionc,
``Asymptotically almost all $\lambda$-terms are strongly normalizing,''
\emph{Logical Methods in Computer Science} \textbf{9} (2013), no.~1:2, 1--30,
\texttt{doi:10.2168/LMCS-9(1:2)2013}.

\bibitem[Dav56]{Davis1956NoteUTM}
Martin Davis,
``A note on universal Turing machines,''
in Claude~E. Shannon and John McCarthy (eds.), \emph{Automata Studies},
Annals of Mathematics Studies, vol.~34, Princeton University Press, Princeton, NJ, 1956, pp.~167--175.

\bibitem[Dav57]{Davis1957DefUTM}
Martin Davis,
``The definition of universal Turing machine,''
\emph{Proceedings of the American Mathematical Society} \textbf{8} (1957), no.~6, 1125--1126,
\texttt{doi:10.1090/S0002-9939-1957-0095781-8}.

\bibitem[Drm97]{Drmota1997}
Michael Drmota,
``Systems of functional equations,''
\emph{Random Structures \& Algorithms} \textbf{10} (1997), no.~1--2, 103--124,
\texttt{doi:10.1002/(SICI)1098-2418(199701/03)10:1/2<103::AID-RSA5>3.0.CO;2-Z}.

\bibitem[EdVW10]{EndrullisDeVrijerWaldmann2010}
J{\"o}rg Endrullis, Roel~C. de~Vrijer, and Johannes Waldmann,
``Local termination: theory and practice,''
\emph{Logical Methods in Computer Science} \textbf{6} (2010), no.~3:20, 1--37,
\texttt{doi:10.2168/LMCS-6(3:20)2010}.

\bibitem[EGH+10]{EndrullisGrabmayerHendriksIsiharaKlop2010}
J{\"o}rg Endrullis, Clemens Grabmayer, Dimitri Hendriks, Ariya Isihara, and Jan~Willem Klop,
``Productivity of stream definitions,''
\emph{Theoretical Computer Science} \textbf{411} (2010), no.~4--5, 765--782,
\texttt{doi:10.1016/j.tcs.2009.10.014}.

\bibitem[Fil26]{Filin2026}
Alexander Filin,
\emph{Pure S Is Computationally Universal Under a Fixed Root-Restarted Finite Controller},
Technical manuscript, Predictive Universe Project, September~14, 2026.

\bibitem[Fis65]{Fischer1965Formalisms}
Patrick~C. Fischer,
``On formalisms for Turing machines,''
\emph{Journal of the ACM} \textbf{12} (1965), no.~4, 570--580,
\texttt{doi:10.1145/321296.321308}.

\bibitem[FS09]{FlajoletSedgewick2009}
Philippe Flajolet and Robert Sedgewick,
\emph{Analytic Combinatorics},
Cambridge University Press, Cambridge, 2009,
\texttt{doi:10.1017/CBO9780511801655}.

\bibitem[FKR20]{ForsterKunzeRoth2020}
Yannick Forster, Fabian Kunze, and Marc Roth,
``The weak call-by-value $\lambda$-calculus is reasonable for both time and space,''
\emph{Proceedings of the ACM on Programming Languages} \textbf{4} (POPL, 2020), Article~27, 27:1--27:23,
\texttt{doi:10.1145/3371095}.

\bibitem[Gon08]{Gonthier2008FourColor}
Georges Gonthier,
``Formal proof---the Four-Color Theorem,''
\emph{Notices of the American Mathematical Society} \textbf{55} (2008), no.~11, 1382--1393.

\bibitem[GAA+13]{Gonthier2013OddOrder}
Georges Gonthier, Andrea Asperti, Jeremy Avigad, Yves Bertot, Cyril Cohen, Fran{\c{c}}ois Garillot, St{\'e}phane Le~Roux, Assia Mahboubi, Russell O'Connor, Sidi Ould~Biha, Ioana Pasca, Laurence Rideau, Alexey Solovyev, Enrico Tassi, and Laurent Th{\'e}ry,
``A machine-checked proof of the Odd Order Theorem,''
in \emph{Interactive Theorem Proving (ITP 2013)},
Lecture Notes in Computer Science, vol.~7998, Springer, Berlin, 2013, pp.~163--179,
\texttt{doi:10.1007/978-3-642-39634-2\_14}.

\bibitem[GL13]{GrygielLescanne2013}
Katarzyna Grygiel and Pierre Lescanne,
``Counting and generating lambda terms,''
\emph{Journal of Functional Programming} \textbf{23} (2013), no.~5, 594--628,
\texttt{doi:10.1017/S0956796813000157}.

\bibitem[HAB+17]{Hales2017Kepler}
Thomas Hales, Mark Adams, Gertrud Bauer, Tat~Dat Dang, John Harrison, Le~Truong Hoang, Cezary Kaliszyk, Victor Magron, Sean McLaughlin, Tat~Thang Nguyen, Quang~Truong Nguyen, Tobias Nipkow, Steven Obua, Joseph Pleso, Jason Rute, Alexey Solovyev, Ta~Thi~Hoai An, Tran~Nam Trung, Thi~Diep Trieu, Josef Urban, Ky~Vu, and Roland Zumkeller,
``A formal proof of the Kepler conjecture,''
\emph{Forum of Mathematics, Pi} \textbf{5} (2017), e2, 1--29,
\texttt{doi:10.1017/fmp.2017.1}.

\bibitem[Her68]{Herman1968Simulation}
Gabor~T. Herman,
``Simulation of one abstract computing machine by another,''
\emph{Communications of the ACM} \textbf{11} (1968), no.~12, 802, 813,
\texttt{doi:10.1145/364175.364181}.

\bibitem[HS08]{HindleySeldin2008}
J.~Roger Hindley and Jonathan~P. Seldin,
\emph{Lambda-Calculus and Combinators: An Introduction},
2nd edition, Cambridge University Press, Cambridge, 2008,
\texttt{doi:10.1017/CBO9780511809835}.

\bibitem[Hue97]{Huet1997}
G{\'e}rard Huet,
``The Zipper,''
\emph{Journal of Functional Programming} \textbf{7} (1997), no.~5, 549--554,
\texttt{doi:10.1017/S0956796897002864}.

\bibitem[KvOvR07]{KlopVanOostromVanRaamsdonk2007}
Jan~Willem Klop, Vincent van~Oostrom, and Femke van~Raamsdonk,
``Reduction strategies and acyclicity,''
in \emph{Rewriting, Computation and Proof},
Lecture Notes in Computer Science, vol.~4600, Springer, Berlin, 2007, pp.~89--112,
\texttt{doi:10.1007/978-3-540-73147-4\_5}.

\bibitem[KR02]{KudlekRogozhin2002UTM3x9}
Manfred Kudlek and Yurii Rogozhin,
``A universal Turing machine with $3$ states and $9$ symbols,''
in \emph{Developments in Language Theory (DLT 2001)},
Lecture Notes in Computer Science, vol.~2295, Springer, Berlin, 2002, pp.~311--318,
\texttt{doi:10.1007/3-540-46011-X\_25}.

\bibitem[KM23]{KutribMeyer2023}
Martin Kutrib and Uwe Meyer,
``Deterministic real-time tree-walking-storage automata,''
in \emph{Proceedings of the 13th International Workshop on Non-Classical Models of Automata and Applications (NCMA 2023)},
Electronic Proceedings in Theoretical Computer Science (EPTCS), vol.~388, 2023, pp.~48--62,
\texttt{doi:10.4204/EPTCS.388.6}.

\bibitem[KM26]{KutribMeyer2026}
Martin Kutrib and Uwe Meyer,
``Deterministic tree-walking-storage automata,''
\emph{Acta Informatica} \textbf{63} (2026), no.~1, Article~9,
\texttt{doi:10.1007/s00236-026-00522-5}.

\bibitem[Min61]{Minsky1961Tag}
Marvin~L. Minsky,
``Recursive unsolvability of Post's problem of `Tag' and other topics in theory of Turing machines,''
\emph{Annals of Mathematics} (2) \textbf{74} (1961), no.~3, 437--455,
\texttt{doi:10.2307/1970290}.

\bibitem[Min67]{Minsky1967Computation}
Marvin~L. Minsky,
\emph{Computation: Finite and Infinite Machines},
Prentice-Hall, Englewood Cliffs, NJ, 1967.

\bibitem[Nea08]{Neary2008Thesis}
Turlough Neary,
\emph{Small Universal Turing Machines},
Ph.D.\ thesis, Department of Computer Science, National University of Ireland, Maynooth, October 2008.

\bibitem[NW06a]{NearyWoods2006SmallFast}
Turlough Neary and Damien Woods,
``Small fast universal Turing machines,''
\emph{Theoretical Computer Science} \textbf{362} (2006), no.~1--3, 171--195,
\texttt{doi:10.1016/j.tcs.2006.06.002}.

\bibitem[NW06b]{NearyWoods2006Rule110PComplete}
Turlough Neary and Damien Woods,
``$\mathrm{P}$-completeness of cellular automaton Rule~110,''
in \emph{Automata, Languages and Programming (ICALP 2006), Part~I},
Lecture Notes in Computer Science, vol.~4051, Springer, Berlin, 2006, pp.~132--143,
\texttt{doi:10.1007/11786986\_13}.

\bibitem[Par66]{Parikh1966}
Rohit~J. Parikh,
``On context-free languages,''
\emph{Journal of the ACM} \textbf{13} (1966), no.~4, 570--581,
\texttt{doi:10.1145/321356.321364}.

\bibitem[Pos36]{Post1936FormulationI}
Emil~L. Post,
``Finite combinatory processes---Formulation~1,''
\emph{The Journal of Symbolic Logic} \textbf{1} (1936), no.~3, 103--105,
\texttt{doi:10.2307/2269031}.

\bibitem[Pos41]{Post1941}
Emil~L. Post,
\emph{The Two-Valued Iterative Systems of Mathematical Logic},
Annals of Mathematics Studies, no.~5, Princeton University Press, Princeton, NJ, 1941.

\bibitem[Pos43]{Post1943Reductions}
Emil~L. Post,
``Formal reductions of the general combinatorial decision problem,''
\emph{American Journal of Mathematics} \textbf{65} (1943), no.~2, 197--215,
\texttt{doi:10.2307/2371809}.

\bibitem[Pri79]{Priese1979ComplexityUTM}
Lutz Priese,
``Towards a precise characterization of the complexity of universal and nonuniversal Turing machines,''
\emph{SIAM Journal on Computing} \textbf{8} (1979), no.~4, 508--523,
\texttt{doi:10.1137/0208039}.

\bibitem[Rog96]{Rogozhin1996SmallUTM}
Yurii Rogozhin,
``Small universal Turing machines,''
\emph{Theoretical Computer Science} \textbf{168} (1996), no.~2, 215--240,
\texttt{doi:10.1016/S0304-3975(96)00077-1}.

\bibitem[Rog98]{Rogozhin1998UTM22x2}
Yurii Rogozhin,
``A universal Turing machine with $22$ states and $2$ symbols,''
\emph{Romanian Journal of Information Science and Technology} \textbf{1} (1998), no.~3, 259--265.

\bibitem[Sch24]{Schonfinkel1924}
Moses Sch{\"o}nfinkel,
``{\"U}ber die Bausteine der mathematischen Logik,''
\emph{Mathematische Annalen} \textbf{92} (1924), no.~3--4, 305--316,
\texttt{doi:10.1007/BF01448013}.

\bibitem[Sco75]{Scott1975}
Dana~S. Scott,
``Open problem,''
in Corrado B{\"o}hm (ed.), \emph{$\lambda$-Calculus and Computer Science Theory},
Lecture Notes in Computer Science, vol.~37, Springer, Berlin, 1975, p.~368.

\bibitem[Smi20]{Smith2020TM23}
Alex Smith,
``Universality of Wolfram's $2, 3$ Turing machine,''
\emph{Complex Systems} \textbf{29} (2020), no.~1, 1--44,
\texttt{doi:10.25088/ComplexSystems.29.1.1}.

\bibitem[Sta86]{Statman1986}
Richard Statman,
``On translating lambda terms into combinators; the basis problem,''
in \emph{Proceedings of the First Annual IEEE Symposium on Logic in Computer Science (LICS 1986)},
IEEE Computer Society Press, 1986, pp.~378--382.

\bibitem[Sta89]{Statman1989}
Richard Statman,
``The word problem for Smullyan's lark combinator is decidable,''
\emph{Journal of Symbolic Computation} \textbf{7} (1989), no.~2, 103--112,
\texttt{doi:10.1016/S0747-7171(89)80044-6}.

\bibitem[Tak95]{Takahashi1995}
Masako Takahashi,
``Parallel reductions in $\lambda$-calculus,''
\emph{Information and Computation} \textbf{118} (1995), no.~1, 120--127,
\texttt{doi:10.1006/inco.1995.1057}.

\bibitem[TeR03]{Terese2003}
{TeReSe},
\emph{Term Rewriting Systems},
Cambridge Tracts in Theoretical Computer Science, vol.~55,
Cambridge University Press, Cambridge, 2003.

\bibitem[Tur37]{Turing1937Computable}
Alan~M. Turing,
``On computable numbers, with an application to the Entscheidungsproblem,''
\emph{Proceedings of the London Mathematical Society} (2) \textbf{42} (1937), 230--265,
\texttt{doi:10.1112/plms/s2-42.1.230}; correction, ibid.\ \textbf{43} (1937), 544--546.

\bibitem[vEB90]{vanEmdeBoas1990Models}
Peter van~Emde~Boas,
``Machine models and simulations,''
in Jan van~Leeuwen (ed.), \emph{Handbook of Theoretical Computer Science, Volume~A: Algorithms and Complexity},
Elsevier / MIT Press, 1990, pp.~1--66.

\bibitem[Wal97]{Waldmann1997}
Johannes Waldmann,
\emph{Nimm Zwei},
Technical Report IR-432, Vrije Universiteit Amsterdam, 1997.

\bibitem[Wal98a]{Waldmann1998RTA}
Johannes Waldmann,
``Normalization of $S$-terms is decidable,''
in Tobias Nipkow (ed.), \emph{Rewriting Techniques and Applications (RTA '98)},
Lecture Notes in Computer Science, vol.~1379, Springer, Berlin, 1998, pp.~138--150,
\texttt{doi:10.1007/BFb0052367}.

\bibitem[Wal98b]{Waldmann1998Thesis}
Johannes Waldmann,
\emph{The Combinator S},
Ph.D.\ thesis, Fakult{\"a}t f{\"u}r Mathematik und Informatik, Friedrich-Schiller-Universit{\"a}t Jena, 1998.

\bibitem[Wal00]{Waldmann2000}
Johannes Waldmann,
``The combinator $S$,''
\emph{Information and Computation} \textbf{159} (2000), no.~1--2, 2--21,
\texttt{doi:10.1006/inco.2000.2874}.

\bibitem[Wat61]{Watanabe1961UTM}
Shigeru Watanabe,
``$5$-symbol $8$-state and $5$-symbol $6$-state universal Turing machines,''
\emph{Journal of the ACM} \textbf{8} (1961), no.~4, 476--483,
\texttt{doi:10.1145/321088.321090}.

\bibitem[Wol02]{Wolfram2002NKS}
Stephen Wolfram,
\emph{A New Kind of Science},
Wolfram Media, Champaign, IL, 2002.

\bibitem[Wol07]{Wolfram2007PrizeWon}
Stephen Wolfram,
``The prize is won; the simplest universal Turing machine is proved,''
\emph{Stephen Wolfram Writings}, October~24, 2007,
\texttt{https://writings.stephenwolfram.com/2007/10/the-prize-is-won-the-simplest-universal-turing-machine-is-proved/}.

\bibitem[Wol20]{Wolfram2020StoryOfComputation}
Stephen Wolfram,
``Combinators and the story of computation,''
\emph{Stephen Wolfram Writings}, December 2020, \texttt{arXiv:2102.09658}.

\bibitem[Wol21a]{Wolfram2021Challenge}
Stephen Wolfram,
``1920, 2020 and a \$20,000 prize: announcing the $S$ Combinator Challenge,''
\emph{Stephen Wolfram Writings}, June~8, 2021, \texttt{doi:10.31855/7709b85c-8e6},
official rules at \texttt{https://www.combinatorprize.org/}.

\bibitem[Wol21b]{Wolfram2021CombinatorsBook}
Stephen Wolfram,
\emph{Combinators: A Centennial View},
Wolfram Media, Champaign, IL, 2021, \texttt{arXiv:2103.12811}.

\bibitem[Woo98]{Woods1998}
Alan~R. Woods,
``Colouring rules for finite trees, and probabilities of monadic second order sentences,''
\emph{Random Structures \& Algorithms} \textbf{13} (1998), no.~3--4, 307--334.

\bibitem[WN06]{WoodsNeary2006TagFOCS}
Damien Woods and Turlough Neary,
``On the time complexity of $2$-tag systems and small universal Turing machines,''
in \emph{Proceedings of the 47th Annual IEEE Symposium on Foundations of Computer Science (FOCS 2006)},
IEEE, 2006, pp.~439--446, \texttt{doi:10.1109/FOCS.2006.57}.

\bibitem[WN09]{WoodsNeary2009Survey}
Damien Woods and Turlough Neary,
``The complexity of small universal Turing machines: a survey,''
\emph{Theoretical Computer Science} \textbf{410} (2009), no.~4--5, 443--450,
\texttt{doi:10.1016/j.tcs.2008.09.051}.

\bibitem[Zac78]{Zachos1978}
Efstathios Zachos,
\emph{Kombinatorische Logik und S-Terme},
Dissertation ETH Nr.~6214 (Berichte des Instituts f{\"u}r Informatik, Nr.~26),
Eidgen{\"o}ssische Technische Hochschule Z{\"u}rich, 1978,
\texttt{doi:10.3929/ethz-a-000136021}.

\bibitem[Zan95]{Zantema1995}
Hans Zantema,
``Termination of term rewriting by semantic labelling,''
\emph{Fundamenta Informaticae} \textbf{24} (1995), no.~1--2, 89--105,
\texttt{doi:10.3233/FI-1995-24124}.

\end{thebibliography}

\end{document}
